\documentclass[11pt, reqno]{amsart}

% GENUS workshop style: page geometry, fonts, hyperref, fancyhdr and the claim-status tags (genus-collab/meta)
\usepackage{genus}
% text width of the original note (170 mm), so that its displayed formulas still fit the line
\geometry{margin=2cm}
\usepackage[utf8]{inputenc}
\usepackage[T1]{fontenc}

%-------------------------------------------------------------------------------
% Packages (modelled on the preamble of [BS] and of the companion paper [LP2]; only what is needed)
\usepackage[dvipsnames]{xcolor}
\usepackage{amssymb}
\usepackage[mathscr]{eucal}
\usepackage{amsthm}
\usepackage{amsmath}
\usepackage{mathtools}
\definecolor{highlight}{HTML}{3465a4}
% hyperref is loaded by genus.sty, with the GENUS link colours
\hypersetup{hypertexnames=false, hyperfootnotes=false, linktocpage=true, bookmarksdepth=2}
\usepackage{cleveref}
\usepackage{url}
\usepackage{array}

%-------------------------------------------------------------------------------
% Typefaces. [BS] uses libertine with newtxmath; libertine needs fontaxes, which is not assumed
% here, so mathpazo (Palatino) is the substitute, as in the companion paper.
% (mathpazo and courier dropped: genus.sty sets Latin Modern)
\usepackage{microtype}
\linespread{1.06}
\setlength{\parskip}{4pt}
\setlength{\parindent}{0pt}
\setcounter{secnumdepth}{3}
% enumitem is not assumed: first-level enumerate is labelled (i), (ii), ... ; for named parts use
% description or \item[(F1)] with an explicit label.
\renewcommand{\labelenumi}{(\roman{enumi})}
\renewcommand{\theenumi}{(\roman{enumi})}
\numberwithin{equation}{section}

%-------------------------------------------------------------------------------
% Draft notes: visible, coloured, usable in text and in math. Set
% \showdraftnotesfalse to suppress them.
\newif\ifshowdraftnotes
\showdraftnotestrue
\newcommand{\draftnotetext}[1]{\textcolor{red}{\footnotesize[\textbf{Draft note:} #1]}}
\newcommand{\draftnote}[1]{\ifshowdraftnotes\ifmmode\text{\draftnotetext{#1}}\else\draftnotetext{#1}\fi\fi}

%-------------------------------------------------------------------------------
% Theorem environments (names as in [BS])
\theoremstyle{plain}
\newtheorem{thm}{Theorem}[section]
\newtheorem{lem}[thm]{Lemma}
\newtheorem{prop}[thm]{Proposition}
\newtheorem{cor}[thm]{Corollary}
\theoremstyle{definition}
\newtheorem{defn}[thm]{Definition}
\newtheorem{rmk}[thm]{Remark}
\newtheorem{notation}[thm]{Notation}
% genus.sty defines an example environment on its own counter; redefine it on the counter of thm
\let\example\undefined\let\endexample\undefined\expandafter\let\csname c@example\endcsname\undefined
\newtheorem{example}[thm]{Example}
\newtheorem{conv}[thm]{Convention}

\crefname{equation}{Eq.}{Eqs.}
\crefname{thm}{Theorem}{Theorems}
\crefname{lem}{Lemma}{Lemmas}
\crefname{prop}{Proposition}{Propositions}
\crefname{cor}{Corollary}{Corollaries}
\crefname{defn}{Definition}{Definitions}
\crefname{rmk}{Remark}{Remarks}
\crefname{notation}{Notation}{Notations}
\crefname{example}{Example}{Examples}
\crefname{conv}{Convention}{Conventions}
\crefname{section}{Section}{Sections}
\crefname{subsection}{Section}{Sections}
\crefname{appendix}{Appendix}{Appendices}

%-------------------------------------------------------------------------------
% Letters
\newcommand{\bbA}{\mathbb{A}}
\newcommand{\bbC}{\mathbb{C}}
\newcommand{\bbD}{\mathbb{D}}
\newcommand{\bbE}{\mathbb{E}}
\newcommand{\bbF}{\mathbb{F}}
\newcommand{\bbK}{\mathbb{K}}
\newcommand{\bbL}{\mathbb{L}}
\newcommand{\bbP}{\mathbb{P}}
\newcommand{\bbQ}{\mathbb{Q}}
\newcommand{\bbR}{\mathbb{R}}
\newcommand{\bbT}{\mathbb{T}}
\newcommand{\bbZ}{\mathbb{Z}}
\newcommand{\cA}{\mathcal{A}}
\newcommand{\cC}{\mathcal{C}}
\newcommand{\cE}{\mathcal{E}}
\newcommand{\cF}{\mathcal{F}}
\newcommand{\cG}{\mathcal{G}}
\newcommand{\cL}{\mathcal{L}}
\newcommand{\cN}{\mathcal{N}}
\newcommand{\cO}{\mathcal{O}}
\newcommand{\cP}{\mathcal{P}}
\newcommand{\cQ}{\mathcal{Q}}
\newcommand{\cW}{\mathcal{W}}
\newcommand{\cZ}{\mathcal{Z}}
\newcommand{\mgp}[1]{\mathsf{#1}}

%-------------------------------------------------------------------------------
% General
\newcommand{\dd}{\mathrm{d}}
\newcommand{\ri}{\mathrm{i}}
\newcommand{\re}{\mathrm{e}}
\newcommand{\de}{\partial}
% \llbracket: from newtxmath if available, otherwise a substitute (stmaryrd is not assumed)
\providecommand{\llbracket}{[\![}
\providecommand{\rrbracket}{]\!]}
\newcommand{\fps}[1]{\llbracket #1 \rrbracket}
\newcommand{\lser}[1]{(\!(#1)\!)}
\newcommand{\Frac}{\mathrm{Frac}}
\DeclareMathOperator{\Tr}{Tr}
\DeclareMathOperator{\End}{End}
\DeclareMathOperator{\Aut}{Aut}
\DeclareMathOperator{\Spec}{Spec}
\DeclareMathOperator*{\Res}{Res}
\DeclareMathOperator{\Disc}{Disc}
\DeclareMathOperator{\Li}{Li}
\DeclareMathOperator{\val}{val}
\DeclareMathOperator{\ev}{ev}
\newcommand{\tensor}{\otimes}
\newcommand{\coef}[1]{[#1]}                    % coefficient extraction, e.g. \coef{t^k}

%-------------------------------------------------------------------------------
% Geometry, tori, weights (weights as in [BS]: \alpha_i, \epsilon_i)
\newcommand{\Tmax}{\mgp{T}}                    % the five-torus
\newcommand{\TCY}{\mgp{T}_{\mathrm{CY}}}        % the Calabi--Yau four-torus E = \epsilon_1 + \epsilon_2
\newcommand{\TF}{\mgp{T}_{\mathrm{F}}}          % fibrewise torus \alpha_0 = 0, \alpha_i = \epsilon_i
\newcommand{\Tthree}{\mgp{T}_{3}}               % \alpha_0 = 0 inside \TCY
\newcommand{\wA}{\mathsf{A}}                    % \alpha_0 + \alpha_2 = E - \alpha_1
\newcommand{\wB}{\mathsf{B}}                    % \alpha_0 + \alpha_1 = E - \alpha_2
\newcommand{\uq}{\mathsf{u}}                    % (E - 2\alpha_1)(E - 2\alpha_2)
\newcommand{\Mbar}{\overline{M}}
\newcommand{\vir}{\mathrm{virt}}
\newcommand{\GW}{\mathsf{GW}}
\newcommand{\Hdg}{\bbE}                        % Hodge bundle
\newcommand{\Lamv}{\Lambda^{\!\vee}}           % \Lamv(u) = \sum_i (-1)^i \lambda_i u^{g-i}
\newcommand{\Hg}{\mathsf{Hg}}                  % \Hg_{\bw} = \prod_a \Lamv(w_a)/w_a^g (normalised Hodge class)
\newcommand{\Nb}{\cN_{\mathrm{cf}}}             % \cO(-1)_{\alpha_1} + \cO(-1)_{\alpha_2} + \cO_{-\epsilon_1} + \cO_{-\epsilon_2}
\newcommand{\bw}{\mathbf{w}}                   % (E,-\epsilon_1,-\epsilon_2)
\newcommand{\nw}{\mathsf{nw}}                  % (\alpha_1,\wA,-\epsilon_1,-\epsilon_2)
\newcommand{\wts}{\mathrm{wt}}                 % \wts_\pm = the five tangent weights at p_\pm
\newcommand{\eul}{\mathrm{e}}                  % \eul(\wts_\pm): product of the tangent weights
\newcommand{\sF}{\mathscr{F}}                  % \sF_g(\cQ) = \sum_{d\ge 0} \GW_{g,d}\cQ^d
\newcommand{\sG}{\mathscr{G}}                  % \sG_g = \sF_g - \sF_g(0)

%-------------------------------------------------------------------------------
% Fields and rings
\newcommand{\Kal}{\bbK_{\alpha}}               % Q(\alpha_0,\alpha_1,\alpha_2,\epsilon_1,\epsilon_2)
\newcommand{\KCY}{\bbK_{\mathrm{CY}}}          % Q(\alpha_0,\alpha_1,\alpha_2,\epsilon_1), \epsilon_2 = E - \epsilon_1
\newcommand{\Kred}{\bbK}                       % Q(\alpha_1,\alpha_2,\epsilon_1) (a_0-adic arguments)
\newcommand{\Kthree}{\bbK_{3}}                 % Q(\alpha_1,\alpha_2,\epsilon_1) as the function field of \Tthree
\newcommand{\Cz}{\bbC\lser{z}}                 % formal Laurent series in z
\newcommand{\Rz}{\mathsf{R}_{z}}               % \Cz((\tau))
\newcommand{\Rzo}{\mathsf{R}_{z}^{\circ}}      % \bbC[[z]]((\tau))
\newcommand{\Oan}{\mathscr{O}}                 % convergent power series in t (or \tau)
\newcommand{\Man}{\mathscr{M}}                 % meromorphic germs
\newcommand{\QA}{\mathsf{QA}}                  % quadratic algebra K(\cQ)[U]/(\Pq)
\newcommand{\RRing}{\mathbf{R}}                % \Kal[[t]] (Hodge side)
\newcommand{\Ccov}{\mathcal{C}}                % double cover \mathsf{y}^2 = \Disc(\cQ)
\newcommand{\ydc}{\mathsf{y}}                  % \ydc^2 = \Disc(\cQ)

%-------------------------------------------------------------------------------
% Base and genus-zero data
\newcommand{\thQ}{\theta_{\cQ}}                % \cQ d/d\cQ = d/dt
\newcommand{\thU}{\theta_{U}}                  % U d/dU
\newcommand{\Pq}{\mathsf{P}}                   % \Pq(U) = \wB\cQ U^2 - (\alpha_0 + \cQ(\alpha_1-\alpha_2))U - \alpha_2
\newcommand{\Ns}{\mathsf{N}}                   % \Ns(s) = \wB\tau s^2 - (E+\alpha_0\tau)s + \wA
\newcommand{\sast}{s^{\ast}}                   % singular root of \Ns
\newcommand{\stil}{\tilde{s}}                  % regular root of \Ns
\newcommand{\Lv}{\mathsf{L}}                   % eigenvalue function \Lv(U); \Lv_i = \Lv(U_i)
\newcommand{\qprod}{\circ}                     % quantum product
\newcommand{\eid}{\mathfrak{e}}                % idempotents
\newcommand{\fid}{\mathfrak{f}}                % normalised idempotents \Delta_i^{1/2}\eid_i
\newcommand{\Lsum}{\mathsf{S}}                 % \Lv_1 + \Lv_2 = E/\tau - \alpha_0
\newcommand{\Lprod}{\mathsf{\Pi}}              % \Lv_1\Lv_2 = \alpha_1\alpha_2/\tau
\newcommand{\lrem}{\ell}                       % \lrem(t) = \Lv_1 - E/t
\newcommand{\muq}{\mu}                         % E/\tau - \wA
\newcommand{\omq}{\omega}                      % (E+z)/\tau - (\wA+z)
\newcommand{\Aop}{\bbA}                        % matrix of the quantum differential equation on covectors
\newcommand{\chiv}{\boldsymbol{\chi}}          % \chiv_i = (1,\Lv_i)
\newcommand{\pim}{\varpi}                      % \pim(a) = \langle a,\chiv_2\rangle
\newcommand{\Thf}{\Theta}                      % \Delta_2/(\Lv_2-\Lv_1)
\newcommand{\Pg}{\mathsf{P}_{g}}               % [t^{2g}] \epsilon_1\epsilon_2 t^2/(4\sinh(\epsilon_1 t/2)\sinh(\epsilon_2 t/2))
\newcommand{\Ph}[1]{\mathsf{P}_{#1}}           % \Ph{h} for other genera
\newcommand{\gapg}{\mathrm{gap}_g}
\newcommand{\Ssh}{\mathsf{S}_{\circ}}          % \Ssh(y) = y/(2\sinh(y/2))
\newcommand{\Aeul}{\cA}                        % \Aeul_n(\cQ): \Li_{-n} = \Aeul_n/(1-\cQ)^{n+1}

%-------------------------------------------------------------------------------
% Euler integral, Laplace calculus, Gauss
\newcommand{\dphi}{\breve{\phi}}               % \phi_\tau(\sast) - \phi_0(\wA/E)
\newcommand{\alg}{\mathfrak{a}}                % the unit \alg(\tau) with \alg(0)=1
\newcommand{\Acal}{\mathscr{A}}                % ring of formal sums \sum K_{pq} v^p z^q
\newcommand{\gauss}[2]{\langle\!\langle #1\rangle\!\rangle_{#2}}   % Gaussian functional
\newcommand{\Lnorm}{\mathscr{N}}               % \Lnorm[f,a;u_0]: Gaussian-normalised formal Laplace expansion
\newcommand{\Lhat}{\widehat{\mathscr{N}}}      % normalised to leading coefficient 1
\newcommand{\Lfull}{\mathscr{N}^{\mathrm{full}}}
\newcommand{\Rr}{\mathsf{Rr}}                  % remainder of the phase, \Rr_f(v)
\newcommand{\Nlap}{\mathcal{N}}                % \Nlap(z,\tau) = \Lhat[\phi_\tau,(1-\tau s)^{-1};\sast]
\newcommand{\wick}{\mathsf{wick}}              % \wick_m: Wick polynomials
\newcommand{\Abeta}{\mathsf{A}_{\beta}}        % \Lnorm[\phi_0,1;\wA/E]
\newcommand{\ga}{\mathsf{a}}                   % 1 + \alpha_2/z
\newcommand{\gb}{\mathsf{b}}                   % 1 + \wA/z
\newcommand{\gc}{\mathsf{c}}                   % 2 + E/z
\newcommand{\Gauss}{\mathsf{Gauss}}            % Gauss's operator
\newcommand{\Fhyp}{\mathsf{F}}                 % {}_2F_1(\ga,\gb;\gc;\tau)
\newcommand{\Ephi}{\mathsf{X}}                 % \exp(-\dphi/z); coefficients \Ephi_n
\newcommand{\Hcoef}{\mathsf{H}}                % \Hcoef_n = [\tau^n](\Fhyp\Ephi)
\newcommand{\Kcoef}{\mathsf{K}}                % \Kcoef_n = [\tau^n](\Fhyp\Ephi\alg)
\newcommand{\Dz}{\mathsf{D}_{z}}               % z d/dt
\newcommand{\Dd}{\mathsf{D}}                   % abstract derivation symbol in the operator identity
\newcommand{\Dop}{\mathsf{D}}                  % \Dop_i = d/dt - \Delta_i'/(2\Delta_i)
\newcommand{\Eik}{\mathsf{Eik}}
\newcommand{\Ptr}{\mathsf{P}_{\mathrm{tr}}}
\newcommand{\Rtr}{\mathsf{R}_{\mathrm{tr}}}
\newcommand{\Ttr}{\mathsf{T}_{\mathrm{tr}}}
\newcommand{\Sph}{\mathsf{S}_{\phi}}           % eikonal phase
\newcommand{\Phia}{\Phi}                       % transport amplitude
\newcommand{\Ddisc}{\mathbb{D}}                % discs \Ddisc, \Ddisc', \Ddisc''
\newcommand{\Pol}{\mathrm{Pol}}                % \Pol_n(p)
\newcommand{\wrk}{\mathsf{w}}                  % \wrk_k = \tau^k \cW_k(U_1)  (\wr is the wreath-product symbol)
\newcommand{\zq}{\zeta}                        % \tau e^{\dphi/E}
\newcommand{\amp}{\mathsf{m}}                  % amplitude of a Laplace expansion (s03, s04)
\newcommand{\Dam}{\mathsf{D}_{\mathrm{am}}}     % \Dam q = d_tau q + q (d_tau f)/z (s03, s04)
\newcommand{\sQ}{\mathsf{s}}                    % 1/\cQ, the local coordinate at \cQ = \infty (s02, s07, appendix; not the Euler variable)

%-------------------------------------------------------------------------------
% Rows, factors, graph sums (one frame: frame B; point-frame objects carry a superscript pt)
\newcommand{\rpt}{\mathrm{r}^{\mathrm{pt}}}    % \rpt_{0,i}: point-frame identity row
\newcommand{\rZ}{\mathrm{r}^{Z}}               % \rZ_{0,i} = h_{\nw}\rpt_{0,i}: the row of R_Z
\newcommand{\rr}{\mathrm{r}}                   % \rr_{a,i}: frame-B rows
\newcommand{\vv}{\mathbf{r}}                   % \vv_i = (\rr_{0,i},\rr_{1,i})
\newcommand{\Rm}{\mathsf{R}}                   % R-matrix at generic weights
\newcommand{\RZ}{\mathsf{R}_{Z}}               % R-matrix with rows \rZ
\newcommand{\RB}{\mathsf{R}'}                  % frame-B R-matrix \RZ h_{\bw}^{-1}
\newcommand{\hw}{h_{\bw}}                      % h_{(E,-\epsilon_1,-\epsilon_2)}
\newcommand{\hnw}{h_{\nw}}                     % h_{(\alpha_1,\wA,-\epsilon_1,-\epsilon_2)}
\newcommand{\hreg}{h_{\mathrm{reg}}}           % h_{(\alpha_1,\alpha_2,\wA,\wB,-E,-E)}
\newcommand{\bk}{\mathsf{b}}                   % \bk_k = B_{2k}/(2k(2k-1))
\newcommand{\TFT}{\boldsymbol{\omega}}         % semisimple TFT
\newcommand{\Rinv}{\mathbf{B}}                 % \Rm^{-1}(\psi)
\newcommand{\Mdef}{\mathbf{M}}                 % defect matrix of the flatness proof
\newcommand{\Dil}{\mathsf{Dil}}                % \Dil_i(d) frame-B dilaton factors
\newcommand{\Edge}{\cE}                        % \Edge_{ij}(m,n) frame-B edge factors
\newcommand{\Ver}{\mathrm{Ver}}
\newcommand{\Edg}{\mathrm{Edg}}
\newcommand{\col}{\mathfrak{c}}                % colouring \col : \Ver(\Gamma) \to \{1,2\}
\newcommand{\GCont}{\mathrm{Cont}}             % \GCont(\Gamma,\col)
\newcommand{\Gsum}{\mathscr{G}}                % frame-B Hodge-vertex graph sum \Gsum_g
\newcommand{\Pij}{\cP}                         % \Pij_{ij}(z,w) = -\langle \vv_i(-z),\vv_j(-w)\rangle
\newcommand{\Kone}{\mathsf{K}_{1}}             % symplectic constant at idempotent 1
\newcommand{\Ktwo}{\mathsf{K}_{2}}
\newcommand{\brow}{\boldsymbol{\sigma}}        % \rho\vv_1(-\psi)
\newcommand{\Yv}{\mathbf{Y}}                   % t\brow
\newcommand{\Gk}{\mathsf{G}}                   % \Gk(z;t) = \rho\rr_{0,1}; \Gk_k = [t^k]\Gk
\newcommand{\gl}{\mathsf{g}}                   % \gl(\psi;t) = \Gk(-\psi;t)
\newcommand{\Dg}{\mathrm{D}\mathsf{g}}         % d/dt (t\gl)
\newcommand{\Rk}{\mathsf{V}}                   % \Rk_k = [t^k](z\Gk)
\newcommand{\Ebar}{\overline{\cE}}             % t^2\rho^2\Edge_{11}
\newcommand{\Nbar}{\overline{\mathsf{N}}}      % numerator of \Ebar
\newcommand{\Psh}{\mathfrak{p}}                % \Psh_k = \prod_{i\le k}((k+1)\alpha_2 - iE)((k+1)\wA - iE)
\newcommand{\Cj}{\mathsf{C}}                   % \Cj_j = \Psh_{j-1}/(j!(j-1)!E^{2j-4})
\newcommand{\Aleaf}{A}                         % A_j(t) leaf coefficient series
\newcommand{\Ahat}{\widehat{A}}                % t A_j

%-------------------------------------------------------------------------------
% Leaves and Hodge side
\newcommand{\lf}{\boldsymbol{\lambda}}         % \lf_j(\psi) = 1/(E - j\psi)
\newcommand{\Lsp}{\mathscr{L}}                 % \Lsp_m = span{\lf_1,...,\lf_m}
\newcommand{\qo}{\mathsf{q}_{\circ}}           % E - E^2\lf_1
\newcommand{\qt}{\mathsf{q}_{t}}               % dilaton series -\psi\rr_{0,1}(-\psi)
\newcommand{\Tdil}{T_{\mathrm{dil}}}           % \qo + \psi = -\psi^2/(E-\psi)
\newcommand{\Fpot}{\mathsf{F}^{\mathrm{pot}}}  % \Fpot_h
\newcommand{\dbase}{\mathfrak{d}}              % \dbase^s
\newcommand{\Dom}{\mathfrak{D}}                % neighbourhood of \qo
\newcommand{\Xiz}{\Xi}                         % \Xiz_h(\msM), scale E
\newcommand{\Xit}{\Xi^{\tau}}                  % scale-1 integral of [LP2]
\newcommand{\ZZ}{\mathrm{Z}}                   % \ZZ[h,n] = \Xiz_h(1^n)
\newcommand{\Mirr}{\mathsf{M}_{\mathrm{irr}}}  % \Mirr(h)
\newcommand{\dg}{\mathrm{dg}}                  % \dg(\msM) = \sum (j-1)
\newcommand{\msM}{\mathsf{M}}                  % multiset of integers >= 2
\newcommand{\msU}{\mathsf{U}}                  % multiset of positive integers
\newcommand{\tauw}{\tau_{\bw}}                 % -\epsilon_1/E, the scale-1 parameter of [LP2]
\newcommand{\st}{\sigma_{\!\bw}}               % -\epsilon_1\epsilon_2/E^2

%-------------------------------------------------------------------------------
% Three-torus and flop
\newcommand{\xc}{\mathsf{x}}                   % \alpha_1/\alpha_2 on \Tthree
\newcommand{\cc}{\mathsf{c}}                   % 1 + \xc = E/\alpha_2
\newcommand{\kap}{\varkappa}                   % \epsilon_1\epsilon_2\alpha_2^4/\alpha_1
\newcommand{\Zfl}{Z'}                          % the flop
\newcommand{\Xfl}{X'}

%-------------------------------------------------------------------------------
% Title

\title[Refined GW theory of the resolved conifold on the CY torus]{Refined Gromov--Witten theory of the resolved conifold on the Calabi--Yau torus: conifold regularity and rigidity}
\author{The GENUS collaboration}
\genusproject{solutions/resolved-conifold-cy-torus}

\begin{document}

\begin{abstract}
We study the equivariant Gromov--Witten theory of the resolved conifold times $\bbC^2$ on its Calabi--Yau four-torus, written as a Givental--Teleman graph sum with the normalised Hodge class of weights $(E,-\epsilon_1,-\epsilon_2)$ at every vertex. At the conifold point $\cQ=1$ one of the two idempotents is singular. We prove that every factor of the graph sum at the singular idempotent has a leaf structure: after one rescaling, the coefficient of $t^k$, $\cQ=\re^t$, is a combination of the functions $1/(E-j\psi)$ with $j\leq k+1$ (Theorem~A, at every non-degenerate weight point of the five-torus); that the graph sum equals the Gromov--Witten potential on the Calabi--Yau torus (Theorem~B); that every coloured stable graph with an edge is regular at the conifold point, the whole polar part being carried by the single-vertex graph with the value $\beta_{2g-2}(\epsilon_1,\epsilon_2)\,t^{2-2g}$ (Theorem~C); and that for $g\geq2$ the genus-$g$ potential equals $\Pg\Li_{3-2g}(\cQ)/(\epsilon_1\epsilon_2)$ up to its constant term, that is, $\epsilon_1\epsilon_2\GW_{g,d}=\Pg d^{2g-3}$ for $d\geq1$ (Theorem~D). With genus $0$ and $1$ this proves the rigidity conjecture of \cite{BS} for the resolved conifold, for every non-zero class and every group acting through the Calabi--Yau torus. The inputs are Teleman's classification, the $R$-matrix action of Pandharipande--Pixton--Zvonkine, virtual localisation, Mumford's formula, Stirling's series, and two theorems on Hodge integrals from the companion paper \cite{LP2}; in genus one and degree $d\geq2$ the formula of \cite{BS} on the subtori $\alpha_i=\epsilon_j$ is used as an input. The constant-map term is not determined.
\end{abstract}

\maketitle
\genusdisclaimer

\setcounter{tocdepth}{1}
\tableofcontents

% s01_intro. Section file; label prefix: int. The main theorems are stated here and proved in the body.
\section{Introduction}
\label{sec:int}

\draftnote{Working title; to be settled by the authors.}

%-------------------------------------------------------------------------------
\subsection{The setting}
\label{sec:int-setting}

Let $X$ be the resolved conifold, the total space of $\cO(-1)\oplus\cO(-1)\to\bbP^1$, and let $Z=X\times\bbC^2$. The five-torus $\Tmax=(\bbC^*)^5$ acts on $Z$ with two fixed points $p_\pm$ on the zero section, with tangent weights $(\alpha_0,\alpha_1,\alpha_2,-\epsilon_1,-\epsilon_2)$ at $p_+$ and $(-\alpha_0,\alpha_0+\alpha_1,\alpha_0+\alpha_2,-\epsilon_1,-\epsilon_2)$ at $p_-$: $\alpha_0$ is the weight of the tangent line of $\bbP^1$, $\alpha_1,\alpha_2$ are the weights of the two $\cO(-1)$ summands, and $-\epsilon_1,-\epsilon_2$ those of the two trivial directions (\cref{sec:set-target}). The \emph{Calabi--Yau torus} $\TCY\subset\Tmax$ is the four-torus on which the five tangent weights sum to zero,
\[
E\coloneqq\alpha_0+\alpha_1+\alpha_2=\epsilon_1+\epsilon_2 ,
\]
and we write $\wA\coloneqq\alpha_0+\alpha_2$, $\wB\coloneqq\alpha_0+\alpha_1$. The equivariant Gromov--Witten invariants $\GW_{g,d}$ of $Z$ in genus $g$ and degree $d[\bbP^1]$ are defined, as in \cite{BS}, by virtual localisation on the proper fixed loci of the moduli spaces of stable maps to $\bbP^1$ (\cref{def:set-inv}); on $\TCY$ they are elements of the function field $\KCY=\bbQ(\alpha_0,\alpha_1,\alpha_2,\epsilon_1)$, $\epsilon_2=E-\epsilon_1$. The genus-$g$ series is
\[
\sF_g(\cQ)=\sum_{d\geq0}\GW_{g,d}\,\cQ^d ,\qquad \cQ=\re^{t},
\]
with the constant-map term $d=0$ included; the series of \cite{BS} runs over $d\geq1$, and the difference is the constant $\sF_g(0)$. The \emph{conifold point} is $\cQ=1$, that is $t=0$, where the two critical points of the superpotential governing the genus-zero theory collide with a logarithmic singularity.

The tool is the Givental--Teleman expression of the higher-genus theory as a sum over stable graphs. On $\TCY$ the point $\cQ=0$ is semisimple (because $\alpha_0\neq0$), the equivariant classes of $Z$ form a semisimple cohomological field theory of rank two, and Teleman's classification writes it as the action of an $R$-matrix on the topological field theory with vertex weights $\Delta_i^{g_v-1}$, $i\in\{1,2\}$. The two idempotents behave very differently at the conifold point: one, the \emph{singular} idempotent (colour $1$), has $\Delta_1\sim-\epsilon_1\epsilon_2E^3/t^2$, the other, the \emph{regular} idempotent (colour $2$), has $\Delta_2$ analytic and non-zero at $t=0$. We write the graph sum in a particular frame, \emph{frame~B}: the $R$-matrix is multiplied on the right by the scalar series $\hw^{-1}$, $\hw=h_{(E,-\epsilon_1,-\epsilon_2)}$, which by Mumford's formula transfers the normalised Hodge class
\[
\Hg_{\bw}=\prod_{a=1}^{3}\frac{\Lamv(\bw_a)}{\bw_a^{g_v}},\qquad \bw=(E,-\epsilon_1,-\epsilon_2),\qquad \Lamv(u)=\sum_i(-1)^i\lambda_iu^{g_v-i},
\]
onto every vertex. The resulting \emph{frame-B Hodge-vertex graph sum} $\Gsum_g$ (\cref{def:set-graphsum}) is a sum over connected stable graphs of genus $g$ without legs, over colourings of their vertices by $\{1,2\}$, and over decorations, of products of vertex weights, dilaton factors and edge factors, all built from two rows $\rr_{0,i}(z;t)$, $\rr_{1,i}(z;t)$ per idempotent. The class $\bw$ is forced: it is the unique class, up to pairs $(v,-v)$, for which the singular row takes the value $E/(E+z)$ at the conifold point (\cref{lem:set-hmult}, \cref{cor:row-F2}). The rows themselves are normalised stationary-phase expansions of the Euler-type integrand $\re^{\cW/z}\dd U/U$ at the two critical points of the superpotential $\cW$, multiplied by one constant $\hnw=h_{(\alpha_1,\wA,-\epsilon_1,-\epsilon_2)}$ fixed at $\cQ=0$.

%-------------------------------------------------------------------------------
\subsection{Statements}
\label{sec:int-statements}

The notation is that of \cref{sec:set}: the non-degeneracy condition (ND) is $\alpha_1\alpha_2\wA\wB\epsilon_1\epsilon_2E\neq0$ (\cref{def:set-ND}), a \emph{weight point with (ND)} is a point of $\mathrm{Lie}(\TCY)$ at which (ND) holds, $\alpha_0=0$ being allowed, and ``at the generic point of $\TCY$'' means ``in $\KCY$''. The frame-B rows and factors are those of \cref{def:set-rows,def:set-factors}; $\rho(t)=(\Delta_1t^2/(-\epsilon_1\epsilon_2E^3))^{-1/2}$ with $\rho(0)=1$; the \emph{leaves} are $\lf_j(\psi)=1/(E-j\psi)$, $j\geq1$, the \emph{leaf spaces} are $\Lsp_m=\mathrm{span}(\lf_1,\dots,\lf_m)$, $\qo=E-E^2\lf_1$, $\qt(\psi)=\sum_{d\geq1}\Dil_1(d)\psi^d$ is the dilaton series of the singular idempotent, and $f_{j,b}(\psi)=\coef{w^b}\Edge_{1j}(\psi,w)$ are the legs (\cref{def:set-rho}); $\tau=1-1/\cQ$, and $\delta$, $\sast$, $\Lv_2$, $\alg$ are the algebraic functions of $\tau$ of \cref{def:set-alg}, \cref{lem:set-eig} and \cref{lem:row-hess}. The first theorem is the structural statement from which everything else follows.

\begin{thm}[Theorem~A: leaf structure at generic weights]\genusstatus{AI-proposed, unchecked}
\label{thm:int-leaf}
Let $(\alpha_0,\alpha_1,\alpha_2,\epsilon_1,\epsilon_2)$ be a weight point of $\TCY$ with (ND), or indeterminates with $E=\epsilon_1+\epsilon_2$ imposed. All factors are the frame-B factors at their true idempotents. Then:
\begin{description}
\item[(F1)] $\Delta_1=-\epsilon_1\epsilon_2E^3/(t^2\rho^2)$, where $\rho=\tau/(t\alg)$ is analytic at $t=0$ with $\rho(0)=1$, so that $\Delta_1=-(\epsilon_1\epsilon_2E^3/t^2)(1-\uq t/E^2+O(t^2))$, $\uq=(E-2\alpha_1)(E-2\alpha_2)$; $\Delta_2=\epsilon_1\epsilon_2\wB\delta\sast\Lv_2$ is analytic at $t=0$ with $\Delta_2(0)=\epsilon_1\epsilon_2\alpha_1\alpha_2\wA\wB/E\neq0$ and $\Delta_2'(0)=0$.
\item[(F2)] $\rr_{0,1}(z;0)=E/(E+z)$, and $\rho\,\rr_{1,1}(z;t)=E/t+\alpha_1\alpha_2/(E+z)+O(t)$.
\item[(F3)] Every frame-B dilaton factor $\Dil_i(d)$ ($i=1,2$, $d\geq1$) and every edge factor $\Edge_{ij}(m,n)$ ($i,j\in\{1,2\}$, $m,n\geq0$) is analytic at $t=0$.
\item[(L1)] $\coef{t^0}(\rho\qt)=\qo=E-E^2\lf_1$, and $\coef{t^k}(\rho\qt)\in\Lsp_{k+1}$ for every $k\geq1$.
\item[(L2)] For $j\in\{1,2\}$, every $b\geq0$ and every $k\geq0$, $\coef{t^k}(\rho f_{j,b})\in\Lsp_{k+1}$. Moreover $\Edge_{21}(b,a)=\Edge_{12}(a,b)$, and $\coef{t^k}f_{j,b}\in\Lsp_{k+1}$ as well.
\item[(L3)] For every $k\geq0$, $\coef{t^k}(\rho^2\Edge_{11}(\psi,\psi'))\in\mathrm{span}\{\lf_i(\psi)\lf_l(\psi'):i,l\geq1,\ (i-1)+(l-1)\leq k\}$, and $\Edge_{11}(a,b)=\Edge_{11}(b,a)$. The same support bound holds for $\coef{t^k}(\rho^e\Edge_{11})$ for every integer $e$.
\end{description}
Every membership is an identity of formal power series in $\psi$ (and $\psi'$), valid at every $\psi$-order and every $t$-order; the coefficients are rational functions of the weights, regular on (ND). The bounds in (L1) and (L3) are sharp at generic weights.
\end{thm}

In words: at the conifold point the vertex weight of the singular idempotent has a double pole with an exactly known unit, every factor is regular, and the $t$-expansion of every rescaled factor at the singular idempotent is \emph{leaf-structured}: the coefficient of $t^k$ is a combination of the simple-pole functions $1/(E-j\psi)$ with $j\leq k+1$, and of their products for loops, with no constant term. \Cref{thm:int-leaf} is proved as \cref{thm:lea-main}, where the sharpness statements and a closed form of every leaf coefficient are given (\cref{thm:lea-closed}). Its proof does not use the Calabi--Yau condition: the frame-B rows depend on $(\alpha_0,\alpha_1,\alpha_2)$ only, and $\epsilon_1\epsilon_2$ enters the factors as an overall constant, so the theorem holds at every weight point of the five-torus with (ND), with $\bw=(E,-\epsilon_1,-\epsilon_2)$ and $E=\alpha_0+\alpha_1+\alpha_2$ (\cref{rmk:lea-cy}). The Calabi--Yau condition enters the main chain once, in Theorem~C below, where it makes $\bw$ a sum-zero triple whose first weight is the leaf scale $E$.

\begin{thm}[Theorem~B: the graph sum is the potential]\genusstatus{AI-proposed, unchecked}
\label{thm:int-ident}
Let $g\geq2$. At the generic point of $\TCY$, $\sF_g$ lies in $\KCY(\cQ)$, is regular at $\cQ=0$ with expansion $\sum_{d\geq0}\GW_{g,d}\cQ^d$ there, and
\[
\sF_g=\Gsum_g ,
\]
the frame-B Hodge-vertex graph sum \eqref{eq:set-Gsum} with vertex class $\Hg_{\bw}$, $\bw=(E,-\epsilon_1,-\epsilon_2)$, and the frame-B factors of \cref{def:set-rows,def:set-factors}, in which $\RB=\RZ\hw^{-1}$ and $\RZ$ has the rows $\hnw\rpt_{0,i}$, $\nw=(\alpha_1,\wA,-\epsilon_1,-\epsilon_2)$. Each coloured term $\GCont(\Gamma,\col)$ is an element of the function field $\KCY(\Ccov)$ of the double cover $\ydc^2=\Disc(\cQ)$ of $\bbP^1_\cQ$, regular over every $\cQ\notin\{0,1,\infty\}$ with $\Disc(\cQ)\neq0$.
\end{thm}

\Cref{thm:int-ident} is proved as \cref{prop:ide-main}. It is an identity in $\KCY(\cQ)$, proved at the generic point of $\TCY$ without any specialisation: the field $\KCY$ embeds in $\bbC$, Teleman's classification applies there, and the $R$-matrix descends to $\KCY\lser{\cQ}$.

\begin{thm}[Theorem~C: regularity and the polar part]\genusstatus{AI-proposed, unchecked}
\label{thm:int-reg}
Fix $g\geq2$. In frame B with $\bw=(E,-\epsilon_1,-\epsilon_2)$, at the generic point of $\TCY$ or at any weight point with (ND):
\begin{enumerate}
\item[(a)] every coloured stable graph $(\Gamma,\col)$ of genus $g$ with at least one edge is regular at $t=0$ after summing over decorations;
\item[(b)] the single-vertex graph is regular at colour $2$, and at colour $1$ its polar part is exactly
\[
\gapg\,t^{2-2g},\qquad \gapg\coloneqq(-\epsilon_1\epsilon_2E^3)^{g-1}\ZZ[g,0]=\beta_{2g-2}(\epsilon_1,\epsilon_2)=\frac{(2g-3)!\,\Pg}{\epsilon_1\epsilon_2},
\]
with no polar term of order $t^{3-2g},\dots,t^{-1}$; here $\ZZ[g,0]$ is the Hodge integral of \cref{def:set-fpot}, and
\[
\Pg=\coef{t^{2g}}\frac{\epsilon_1\epsilon_2\,t^2}{4\sinh(\epsilon_1t/2)\sinh(\epsilon_2t/2)},\qquad
\beta_n(\epsilon_1,\epsilon_2)=(n-1)!\,\coef{t^n}\frac{1}{4\sinh(\epsilon_1t/2)\sinh(\epsilon_2t/2)} ;
\]
\item[(c)] hence the expansion of $\Gsum_g$ in $\KCY\lser t$ has polar part exactly $\gapg t^{2-2g}$, and, by \cref{thm:int-ident}, so has the image of $\sF_g$ under $\cQ=\re^t$.
\end{enumerate}
\end{thm}

\Cref{thm:int-reg} is proved as \cref{cor:reg-main}. Parts (a), (b) and the first half of (c) are statements about the decoration sums as defined in \cref{def:set-graphsum}; they rest on \cref{thm:int-leaf} and on the two theorems of the companion paper \cite{LP2} cited in \cref{sec:reg-cited}, and on the string and dilaton equations (\cref{sec:reg-hodge,sec:reg-string}). The whole polar part of the graph sum sits in the single-vertex graph at the singular colour.

\begin{thm}[Theorem~D: rigidity]\genusstatus{AI-proposed, unchecked}
\label{thm:int-rig}
For $g\geq2$, in $\KCY(\cQ)$,
\[
\sF_g(\cQ)-\sF_g(0)=\frac{\Pg\,\Li_{3-2g}(\cQ)}{\epsilon_1\epsilon_2},\qquad\text{that is}\qquad \epsilon_1\epsilon_2\,\GW_{g,d}=\Pg\,d^{2g-3}\quad\text{for every }d\geq1,
\]
with $\Li_{-n}(\cQ)=\sum_{d\geq1}d^n\cQ^d$. Both sides of the second identity are polynomials in the weights on $\mathrm{Lie}(\TCY)$, so it holds at every point of $\mathrm{Lie}(\TCY)$, including $\alpha_1=\alpha_2$, $E=0$, $\epsilon_1=0$ or $\epsilon_2=0$, and $\alpha_0=0$; in particular $\epsilon_1\epsilon_2\GW_{g,d}$ depends on $(\epsilon_1,\epsilon_2)$ only, symmetrically and through $\epsilon_1^2,\epsilon_2^2$.
\end{thm}

\Cref{thm:int-rig} is proved as \cref{thm:rig-main}. In genus $0$ one has $\epsilon_1\epsilon_2\GW_{0,d}=d^{-3}$ at every weight, by the Aspinwall--Morrison formula (\cref{prop:rig-g0}), and in genus $1$ one has $\epsilon_1\epsilon_2\GW_{1,d}=\Ph{1}/d=-(\epsilon_1^2+\epsilon_2^2)/(24d)$ on $\TCY$ (\cref{lem:rig-g1}); the genus-one statement is proved by direct localisation for $d=1$, and for $d\geq2$ it uses, as an input, the formula of \cite{BS} for the resolved conifold on the subtori on which $\alpha_i=\epsilon_j$. Together:

\begin{cor}[the rigidity conjecture of \cite{BS} for the resolved conifold]\genusstatus{AI-proposed, unchecked}
\label{cor:int-conj}
For every $d\geq1$ and every $g\geq0$, $\epsilon_1\epsilon_2\,\GW_g(Z,\TCY;d[\bbP^1])=\Pg\,d^{2g-3}$, so that
\[
\sum_{g\geq0}\GW_g(Z,\TCY;d[\bbP^1])=\frac1d\Big(2\sinh\frac{d\epsilon_1}2\Big)^{-1}\Big(2\sinh\frac{d\epsilon_2}2\Big)^{-1}.
\]
Hence the rigidity conjecture of \cite{BS} holds for $X$ the resolved conifold, every non-zero class $d[\bbP^1]$, and the group $\TCY$; by contravariance it holds for every subtorus of $\TCY$ on which the invariants are defined by localisation (the normal bundle $\Nb$ moving on the fixed loci) and for every connected group whose maximal torus acts on $Z$ through such a subtorus, in particular for every Calabi--Yau subtorus of the five-torus with this property. \draftnote{BS label: conj: rigidity.}
\end{cor}

\Cref{cor:int-conj} is \cref{cor:rig-conj}, where the conjecture is stated by content. The three-torus $\Tthree=\{\alpha_0=0\}\cap\TCY$, which contains the fibrewise torus $\TF$ ($\alpha_i=\epsilon_i$) on which the zero section is fixed pointwise, is not covered by the proof of \cref{thm:int-ident} directly, because $\cQ=0$ is not semisimple there; Theorems~A and~C hold on $\Tthree$ by a rescaling of the rows to the fibrewise ones and a ring map applied to the identity of \cref{thm:int-ident} (\cref{cor:rig-3T}), and rigidity on $\Tthree$ is the restriction of \cref{thm:int-rig} by polynomiality.

%-------------------------------------------------------------------------------
\subsection{Inputs}
\label{sec:int-inputs}

The results are conditional only on the following cited theorems, each named at the place where it enters.
\begin{enumerate}
\item Stirling's series for $\log\Gamma$ with remainder \cite[5.11.1 and Section~5.11(ii)]{DLMF}, together with the Beta integral \cite[5.12.1]{DLMF} and Laplace's method, which is proved in full (\cref{lem:row-laplace}). This is the only analytic input of Theorem~A; it enters through the normalised expansion of the Beta phase (\cref{lem:row-betaphase}, \cref{prop:row-beta}).
\item Teleman's classification of semisimple cohomological field theories \cite[Theorem~(2), (1.15), Prop.~8.7(ii)]{Tel12}, with the stability range of \cite[2.17--2.19, (2.15)]{Tel12} as reported there from \cite{Loo96,MW07,Har85,Iva90}; the $R$-matrix action and its group property \cite[Defs.~2.2, 2.13, Prop.~2.14]{PPZ15}; the virtual localisation formula \cite{GP99}; the twisted-theory axioms \cite[App.~1]{CG07}; Mumford's relation and formula \cite{Mum83}, \cite[Section~1.1]{FP00}; the non-vanishing of $\lambda_g\lambda_{g-1}\lambda_{g-2}$ \cite[Theorem~4]{FP00}; and the Aspinwall--Morrison formula \cite[p.~3]{FP00}. These enter Theorem~B; the localisation formula \cite{GP99} enters also Theorem~D (\cref{lem:rig-poly}) and genus $1$ (\cref{lem:rig-g1}), Mumford's relation enters genus $1$, and the Aspinwall--Morrison formula is the genus-zero statement (\cref{prop:rig-g0}). \draftnote{Teleman is cited with the numbering of the arXiv version; the journal numbering has not been compared.}
\item The two theorems of the companion paper \cite{LP2}: the vanishing theorem for the Hodge integrals $\Xiz_h(\msM)$ with the class $\Hg_{(1,\tau,-1-\tau)}$ at leaf degree $\dg(\msM)\leq2h-3$ (\cite[Theorem~8.1, Corollaries~8.2--8.4]{LP2}, quoted as \cref{thm:reg-L}), and the value of the single vertex, $(-\epsilon_1\epsilon_2E^3)^{g-1}\ZZ[g,0]=\beta_{2g-2}(\epsilon_1,\epsilon_2)$ (\cite[Theorem~9.1]{LP2}, quoted as \cref{thm:reg-VI}). Both rest on the Mari\~no--Vafa formula \cite[(1)--(3)]{LLZ03}. They enter Theorem~C only.
\item For genus $1$ and $d\geq2$ only, the formula of \cite{BS} for the resolved conifold on the subtori on which $\alpha_i=\epsilon_j$ (\cref{lem:rig-g1}). \draftnote{BS label: prop:ResConifold, eq:resconGWref.}
\end{enumerate}
Which result uses which: Theorem~A uses (i); Theorem~C uses Theorem~A and (iii); Theorem~B uses (i), through the normalisation constant at $\cQ=0$, and (ii); Theorem~D uses Theorems~B and~C and, in the polynomiality lemma, \cite{GP99}; genus $0$ uses the Aspinwall--Morrison formula; genus $1$ uses \cite{GP99}, Mumford's relation and, for $d\geq2$, (iv). None of the arguments of this note uses topological recursion, a mirror theorem, a degeneration argument or a finite-order computation; the one input proved by degeneration is (iv), used in genus $1$ for $d\geq2$ only; the finite-order checks that accompanied the work are collected in \cref{sec:evi}, and no proof step refers to them.

%-------------------------------------------------------------------------------
\subsection{What is not claimed}
\label{sec:int-notclaimed}

\begin{enumerate}
\item Nothing is claimed off $\TCY$. Theorem~A holds at every weight point of the five-torus with (ND), but Theorem~C does not: off $\TCY$ the class $\bw$ does not sum to zero, the vanishing theorem does not apply, and the coloured graphs with an edge acquire poles at $t=0$ (\cref{sec:evi-offCY}, as evidence only).
\item The constant-map term $\sF_g(0)=\GW_{g,0}$ is not determined. It is a rational function of all five weights (\cref{rmk:rig-d0}), it depends on $\alpha$ (\cref{rmk:rig-F0}), and it is not part of the conjecture, which concerns non-zero classes. The only statement about it proved here is the flop identity $\sF_g(\infty)=\sF_g(0)$ (\cref{prop:rig-flop}).
\item The conjecture for group actions on $Z$ that are not conjugate into the five-torus is not addressed.
\item The proofs of the cited theorems of \cref{sec:int-inputs} are not reproduced. In particular the two theorems of \cite{LP2} are quoted, with the scale conversion written out (\cref{lem:reg-scale}).
\end{enumerate}

%-------------------------------------------------------------------------------
\subsection{Strategy}
\label{sec:int-strategy}

\emph{The singular row (\cref{sec:row}).} In the Euler coordinate $U=1-\tau s$, $\tau=1-1/\cQ$, the superpotential becomes the Euler phase $\phi_\tau(s)=\wA\log s+\alpha_1\log(1-s)-\alpha_2\log(1-\tau s)$, whose critical points stay separated at $t=0$ although the critical points of $\cW$ collide. The point-frame row at the singular idempotent is the normalised formal Laplace expansion of $\re^{\phi_\tau/z}(1-\tau s)^{-1}\dd s$ at the critical point $\sast$ (\cref{prop:row-wickrow}); a formal Laplace calculus (coordinate invariance, parameter derivatives with a moving expansion point, and the vanishing of total derivatives) shows that this expansion, suitably conjugated, satisfies Gauss's hypergeometric equation in $\tau$ (\cref{prop:row-gaussode}). Uniqueness of the solution gives the closed form (\cref{thm:row-closed})
\[
\rr_{0,1}(z;t)=\frac{E}{E+z}\,{}_2F_1\Big(1+\frac{\alpha_2}z,1+\frac{\wA}z;2+\frac Ez;\tau\Big)\exp\Big(-\frac{\dphi(\tau)}z\Big)\alg(\tau),
\]
in which the constant $E/(E+z)$ comes from the Laplace expansion of a Beta integral, evaluated by Stirling's series (\cref{prop:row-beta}), and $\dphi$, $\alg$ are explicit algebraic functions of $\tau$. The regularity of the $t$-coefficients in $z$ follows (\cref{cor:row-reg}, and formally again by an eikonal and transport decomposition, \cref{prop:row-transport}).

\emph{Leaves (\cref{sec:lea}).} The key formula $\rho\,\rr_{0,1}=(E/(E+z))(\tau/t)\,{}_2F_1\exp(-\dphi/z)$ (\cref{prop:lea-key}) exhibits, order by order in $t$, a rational function of $z$ with simple poles at $z=-E/j$, $j\leq k+1$, vanishing at $\infty$ for $k\geq1$: this is (L1), by partial fractions. Row~2 of the quantum differential equation is the operator form of Gauss's equation (\cref{prop:lea-row2}); self-adjointness of the quantum product gives the symplectic constant $\Kone=1$ (\cref{prop:lea-K1}) and a first-order flow for the loop $t^2\rho^2\Edge_{11}$ whose right side is a rank-one product of leaf series (\cref{thm:lea-flow}); integrating from $t=0$ gives (F3) for $\Edge_{11}$ and (L3). At the regular idempotent, Row~2 is proved by the same Laplace calculus at the second critical point (\cref{thm:lea-row2reg}), which yields the orthogonality of the two rows, the regularity of the mixed edge and (L2) (\cref{prop:lea-L2}). The value of the regular row at $t=0$ is a second application of the Beta-phase lemma (\cref{prop:lea-regrows}).

\emph{Identification (\cref{sec:ide}).} Over $\KCY$, the point $\cQ=0$ is semisimple, so Teleman's theorem gives a unique symplectic $R$-matrix; a flatness argument of the kind used for the quantum differential equation, with a stability input to kill the defect, shows that it solves the quantum differential equation (\cref{prop:ide-flat}). The rows $\hnw\rpt_{0,i}$ form a second flat unitary solution (\cref{prop:ide-rows}); two such solutions differ by constant diagonal exponentials of odd series (\cref{lem:ide-ambig}), which are fixed at $\cQ=0$ by the Hodge integrals of constant maps and the non-vanishing of $\mathrm{ch}_{2k-1}(\Hdg)$ (\cref{prop:ide-const}). The passage to frame~B is the group property of the $R$-action for the scalar $\hw$ (\cref{prop:ide-frameB}).

\emph{Regularity (\cref{sec:reg}).} At a singular vertex of genus $h\geq2$ the dilaton homogeneity of the Hodge potential (\cref{lem:reg-H0}) rescales every half-edge by one factor $\rho$ and converts $\Delta_1^{h-1}$ into $(-\epsilon_1\epsilon_2E^3)^{h-1}t^{2-2h}$ exactly, by (F1). Expanding the rescaled factors in leaves by (L1)--(L3), the coefficient of $t^{2-2h+k}$ is a combination of Hodge integrals $\Xiz_h(\msU)$ with leaf degree $\dg(\msU)\leq k$; the parts equal to $1$ are stripped by the string equation (\cref{lem:reg-string}), and the vanishing theorem kills every term with $k\leq2h-3$. A count over the vertices then shows that every coloured graph with an edge is regular, while the single vertex keeps exactly the term $(-\epsilon_1\epsilon_2E^3)^{g-1}\ZZ[g,0]t^{2-2g}$ (\cref{thm:reg-main}), whose value is the cited \cref{thm:reg-VI}.

\emph{Rigidity (\cref{sec:rig}).} By Theorem~B, $\sF_g$ is a rational function of $\cQ$ whose poles can only lie over $\cQ=1$, $\cQ=\infty$ or the two roots of the discriminant $\Disc(\cQ)$ of the critical equation (\cref{lem:rig-poles}). There is no pole on $\Disc=0$: the invariants $\epsilon_1\epsilon_2\GW_{g,d}$, $d\geq1$, are polynomials in the weights by localisation on a proper space (\cref{lem:rig-poly}), the discriminant has a root of positive $\alpha_0$-adic valuation, and an $\alpha_0$-adic evaluation of the identity $N=(1-\cQ)^n\Disc^m\sG_g$ forces $m=0$ (\cref{lem:rig-disc}). There is no pole at $\cQ=\infty$, because every factor is a power series in $1/\cQ$ there (\cref{lem:rig-inf}). At $\cQ=1$ the polar part is given by Theorem~C and equals that of $\Pg\Li_{3-2g}(\cQ)/(\epsilon_1\epsilon_2)$ (\cref{lem:rig-li}); the difference is a rational function without poles, hence a constant, which vanishes at $\cQ=0$. The excluded loci of $\mathrm{Lie}(\TCY)$ are reached by polynomial continuation.

\emph{Dependencies.} \Cref{sec:set} is definitions and genus-zero data. \Cref{sec:row} uses \cref{sec:set}; \cref{sec:lea} uses \cref{sec:set,sec:row}; \cref{sec:ide} uses \cref{sec:set} and the Laplace calculus and Beta-phase lemma of \cref{sec:row}, not Theorem~A; \cref{sec:reg} uses \cref{sec:set}, Theorem~A and the cited theorems of \cite{LP2}, and uses Theorem~B only for the second half of \cref{cor:reg-main}(c); \cref{sec:rig} uses Theorems~B and~C, the genus-zero data of \cref{sec:set}, the localisation formula \cite{GP99} (\cref{lem:rig-poly}) and, for the flop identity and the three-torus only, the Beta-phase lemma \cref{lem:row-betaphase} and \cref{cor:lea-a0}. Teleman's theorem and the other inputs of item (ii) of \cref{sec:int-inputs} enter only through \cref{sec:ide}, that is, only through the identification of the graph sum with $\sF_g$.

%-------------------------------------------------------------------------------
\subsection{Relation to the results of \texorpdfstring{\cite{BS}}{[BS]}}
\label{sec:int-BS}

The paper \cite{BS} defines the refined invariants by localisation, formulates the rigidity conjecture for a rigid class on $X\times\bbC^2$ under a Calabi--Yau group action, proves the case in which the action on $X$ is itself Calabi--Yau (there the invariants are the unrefined ones times $(\epsilon_1\epsilon_2)^{g-1}$, by Mumford's relation), and proves for the resolved conifold the closed formula of \cref{cor:int-conj} on the subtori on which $\alpha_i=\epsilon_j$, by degeneration of the base curve and the relative invariants of local curves; it also verifies the formula on $\TCY$ by computer in low degree and genus. \draftnote{BS label: defn: refined GW invariant; conj: rigidity; prop: rigidity unref limit; prop:ResConifold.} What is new here is the statement on all of $\TCY$, for every $d\geq1$ and every genus $g\geq2$, genus by genus, and the method: for $g\geq2$ no degeneration is used, and the result comes from the structure of the Hodge-vertex graph sum at the conifold point, that is, from Theorems~A, B and~C. The genus-one input of \cite{BS} at $d\geq2$ is the one place where a result of \cite{BS} is used (\cref{rmk:rig-g1input}).

%-------------------------------------------------------------------------------
\subsection{Plan of the note}
\label{sec:int-plan}

\Cref{sec:set} fixes the target, the tori, the invariants, the genus-zero data (superpotential, Frobenius algebra, idempotents, discriminant), the rows and factors of the frame-B graph sum, the leaf spaces and the Hodge-side integrals. \Cref{sec:row} computes the singular row in closed form. \Cref{sec:lea} proves Theorem~A, with the closed forms of the leaf coefficients and the case $\alpha_0=0$. \Cref{sec:ide} proves Theorem~B. \Cref{sec:reg} quotes the two theorems of \cite{LP2} and proves Theorem~C. \Cref{sec:rig} proves Theorem~D, the genus-zero and genus-one statements, \cref{cor:int-conj}, the flop identity and the three-torus corollary. \Cref{sec:evi} records the finite-order checks; nothing in the body refers to it.


\subsection*{Acknowledgements}

This work was supported by the UKRI grant UKRI2394 and by the Leverhulme Trust grant LIP-2023-007.

% s02_setup. Section file; new macros at the top; label prefix: set.
\providecommand{\pc}{\mathsf{p}}            % pc_{k,i}: coefficients of the polynomials Pol_k (local)
\section{Set-up: the geometry, the torus and the genus-zero data}
\label{sec:set}

This section fixes the objects of the note: the target and its tori, the equivariant invariants, the genus-zero data of the resolved conifold at generic weights, the rows and factors of the frame-B Hodge-vertex graph sum, the leaf spaces and the Hodge-side integrals. Everything here is a definition or an elementary identity; nothing of the later sections is used, except where a forward reference says so explicitly. Throughout, idempotent (colour) $1$ is the singular one and idempotent (colour) $2$ the regular one, in the sense of \cref{def:set-alg}.

%-------------------------------------------------------------------------------
\subsection{The target, the tori and the weights}
\label{sec:set-target}

Let $X$ be the resolved conifold, the total space of $\cO(-1)\oplus\cO(-1)\to\bbP^1$, and let $Z\coloneqq X\times\bbC^2$, the total space of $\cO(-1)^{\oplus2}\oplus\cO^{\oplus2}\to\bbP^1$. The five-torus $\Tmax=(\bbC^*)^5$ acts on $Z$ with the two fixed points $p_+,p_-$ of the zero section $\bbP^1$, and with tangent weights
\begin{equation}
\label{eq:set-wts}
\wts_+=(\alpha_0,\alpha_1,\alpha_2,-\epsilon_1,-\epsilon_2),\qquad
\wts_-=(-\alpha_0,\alpha_0+\alpha_1,\alpha_0+\alpha_2,-\epsilon_1,-\epsilon_2)
\end{equation}
at $p_+$ and $p_-$: $\alpha_0$ is the weight of the tangent line of $\bbP^1$ at $p_+$, $\alpha_1,\alpha_2$ are the weights of the two $\cO(-1)$ summands at $p_+$, and $-\epsilon_1,-\epsilon_2$ are the weights of the two trivial directions. The $\Tmax$-linearised normal bundle of the zero section is $\Nb=\cO(-1)_{\alpha_1}\oplus\cO(-1)_{\alpha_2}\oplus\cO_{-\epsilon_1}\oplus\cO_{-\epsilon_2}$. \draftnote{compare the data of [BS, Appendix A] for the tangent weights.} We write throughout
\[
E\coloneqq\alpha_0+\alpha_1+\alpha_2,\qquad \wA\coloneqq\alpha_0+\alpha_2=E-\alpha_1,\qquad \wB\coloneqq\alpha_0+\alpha_1=E-\alpha_2,\qquad \uq\coloneqq(E-2\alpha_1)(E-2\alpha_2),
\]
so that $\wts_-=(-\alpha_0,\wB,\wA,-\epsilon_1,-\epsilon_2)$, $\wA\wB=\alpha_1\alpha_2+\alpha_0E$, $\uq=\alpha_0^2-(\alpha_1-\alpha_2)^2$ and $2\alpha_0E+4\alpha_1\alpha_2=E^2+\uq$ (indeed $\uq=E^2-2E(\alpha_1+\alpha_2)+4\alpha_1\alpha_2=E^2-2E(E-\alpha_0)+4\alpha_1\alpha_2$). The products of the tangent weights are
\[
\eul(\wts_+)=\alpha_0\alpha_1\alpha_2\epsilon_1\epsilon_2,\qquad \eul(\wts_-)=-\alpha_0\wA\wB\epsilon_1\epsilon_2 .
\]

\begin{defn}[The Calabi--Yau torus; non-degeneracy]
\label{def:set-ND}
The \emph{Calabi--Yau four-torus} $\TCY\subset\Tmax$ is the subtorus on which the sum of the five tangent weights at $p_+$ (equivalently at $p_-$) vanishes,
\[
\TCY:\qquad E=\alpha_0+\alpha_1+\alpha_2=\epsilon_1+\epsilon_2 .
\]
Inside $\TCY$ we use the three-torus $\Tthree:\ \alpha_0=0$ (so $E=\alpha_1+\alpha_2=\epsilon_1+\epsilon_2$) and the \emph{fibrewise torus} $\TF:\ \alpha_0=0,\ \alpha_i=\epsilon_i$, which fixes the zero section pointwise, with normal weights $(\epsilon_1,\epsilon_2,-\epsilon_1,-\epsilon_2)$. The \emph{non-degeneracy condition} is
\[
\text{(ND)}\qquad \alpha_1\alpha_2\wA\wB\epsilon_1\epsilon_2E\neq0 .
\]
A \emph{weight point with (ND)} is a point of $\mathrm{Lie}(\TCY)$, with complex (or rational) coordinates, at which (ND) holds; $\alpha_0=0$ is allowed. At the generic point of $\TCY$ all of $\alpha_0,\alpha_1,\alpha_2,\wA,\wB,\epsilon_1,\epsilon_2,E,\alpha_1-\alpha_2$ are non-zero. The point of $\TF$ in this notation has $\alpha_0=0$, $\wA=\alpha_2=\epsilon_2$, $\wB=\alpha_1=\epsilon_1$, and (ND) there reads $\epsilon_1\epsilon_2E\neq0$.
\end{defn}

The function fields in use are
\begin{equation}
\label{eq:set-fields}
\Kal\coloneqq\bbQ(\alpha_0,\alpha_1,\alpha_2,\epsilon_1,\epsilon_2),\qquad
\KCY\coloneqq\bbQ(\alpha_0,\alpha_1,\alpha_2,\epsilon_1)\ \text{ with }\ \epsilon_2\coloneqq E-\epsilon_1,\qquad
\Kred\coloneqq\bbQ(\alpha_1,\alpha_2,\epsilon_1),
\end{equation}
the five weights being indeterminates in $\Kal$; $\KCY$ is the function field of $\mathrm{Lie}(\TCY)$ and $\KCY=\Kred(\alpha_0)$. The same field $\Kred$, regarded as the function field of $\Tthree$, is written $\Kthree$. The statement ``at the generic point of $\TCY$'' means ``in $\KCY$'', where (ND) holds. The frame-B rows of \cref{def:set-rows} have coefficients in $\bbQ(\alpha_0,\alpha_1,\alpha_2)$ and never see $\epsilon_1,\epsilon_2$ separately, so nothing below depends on the choice of $\epsilon_1$ rather than $\epsilon_2$ as the fourth generator of $\KCY$.

\begin{lem}[Symmetries of the weights]\genusstatus{AI-proposed, unchecked}
\label{lem:set-sym}
The substitution $(\alpha_0,\alpha_1,\alpha_2)\mapsto(-\alpha_0,\alpha_0+\alpha_1,\alpha_0+\alpha_2)$ exchanges $\wts_+$ and $\wts_-$; it sends $(\alpha_0,\alpha_1,\alpha_2,\wA,\wB,E)$ to $(-\alpha_0,\wB,\wA,\alpha_2,\alpha_1,E)$, so it preserves $E$ and exchanges the pairs $(\alpha_1,\wA)\leftrightarrow(\wB,\alpha_2)$. It is induced by the automorphism of $Z$ lifting $z\mapsto1/z$ on $\bbP^1$ (the bundle $\cO(-1)^{\oplus2}\oplus\cO^{\oplus2}$ is $SL(2)$-homogeneous), so every equivariant invariant of \cref{def:set-inv}, as an element of $\KCY$, is invariant under it. Likewise $\alpha_1\leftrightarrow\alpha_2$ (exchange of the two $\cO(-1)$ summands) and $\epsilon_1\leftrightarrow\epsilon_2$ are symmetries.
\end{lem}

\begin{proof}
Direct from \eqref{eq:set-wts}: $-\alpha_0+(\alpha_0+\alpha_1)+(\alpha_0+\alpha_2)=E$, $(-\alpha_0)+(\alpha_0+\alpha_2)=\alpha_2$ and $(-\alpha_0)+(\alpha_0+\alpha_1)=\alpha_1$. The automorphism of $Z$ covering the inversion of $\bbP^1$ exchanges $p_+$ and $p_-$ and conjugates the action of $\Tmax$ into the action with the exchanged weights; the invariants of \cref{def:set-inv}, being defined from the fixed loci and the weights alone, are invariant.
\end{proof}

%-------------------------------------------------------------------------------
\subsection{The invariants and the series}
\label{sec:set-inv}

\begin{defn}[Invariants by localisation; the series]
\label{def:set-inv}
For $g\geq2$ and $d\geq0$ let $\GW_{g,d}\coloneqq\GW^{Z,\TCY}_{g,d[\bbP^1]}\in\KCY$ be the genus-$g$, degree $d[\bbP^1]$ $\TCY$-equivariant Gromov--Witten invariant of $Z$ in the definition by localisation: a $\TCY$-fixed stable map to $Z$ has image in the compact $\TCY$-invariant curves, that is, in the zero section, and is fixed by the rotation of weight $\alpha_0$ of $\bbP^1$ (as $\alpha_0\neq0$ in $\KCY$), so the $\TCY$-fixed loci of $\Mbar_{g,0}(Z,d)$ are the $\alpha_0$-fixed loci of the proper moduli space $\Mbar_{g,0}(\bbP^1,d)$, and the invariant is the sum over them of the residues of the virtual localisation formula \cite{GP99}, with $\Nb$ moving. \draftnote{BS label: defn: refined GW invariant, defn: refined GW series.} For the five-torus the same definition gives $\GW^{\Tmax}_{g,d}\in\Kal$. The genus-$g$ series is
\[
\sF_g(\cQ)\coloneqq\sum_{d\geq0}\GW_{g,d}\,\cQ^d,\qquad \sG_g\coloneqq\sF_g-\sF_g(0)=\sum_{d\geq1}\GW_{g,d}\,\cQ^d ,
\]
with $\cQ$ the degree variable of $\bbP^1$. We put $t\coloneqq\log\cQ$, $\cQ=\re^{t}$, and
\[
\begin{gathered}
\tau\coloneqq1-\frac1\cQ=1-\re^{-t},\qquad\text{so }\ \cQ=\frac1{1-\tau},\qquad \cQ\tau=\cQ-1,\qquad \frac{\dd\tau}{\dd t}=1-\tau,\\
\tau=t-\frac{t^2}2+\frac{t^3}6+O(t^4),\qquad t=-\log(1-\tau).
\end{gathered}
\]
The \emph{conifold point} is $\cQ=1$, that is, $t=\tau=0$. Finally, for $g\geq0$ and $n\geq1$,
\[
\begin{gathered}
\Pg\coloneqq\coef{t^{2g}}\,\frac{\epsilon_1\epsilon_2\,t^2}{4\sinh(\epsilon_1t/2)\sinh(\epsilon_2t/2)},\qquad
\beta_n(\epsilon_1,\epsilon_2)\coloneqq(n-1)!\,\coef{t^n}\,\frac{1}{4\sinh(\epsilon_1t/2)\sinh(\epsilon_2t/2)},\\
\Li_{-n}(\cQ)\coloneqq\sum_{d\geq1}d^n\cQ^d ,
\end{gathered}
\]
so that $\Ph{0}=1$; $\Pg$ is a polynomial in $\epsilon_1,\epsilon_2$ and $\beta_{2g-2}=(2g-3)!\,\Pg/(\epsilon_1\epsilon_2)$ for $g\geq2$ (\cref{lem:reg-beta}), and $\Li_{-n}$ is a rational function of $\cQ$ (\cref{lem:rig-li}).
\end{defn}

\begin{rmk}[The constant term]
\label{rmk:set-d0}
The series $\sF_g$ includes the term $d=0$, the constant-map term $\sF_g(0)=\GW_{g,0}$. The series of \cite{BS} sums over the non-zero classes, that is, over $d\geq1$; the difference is the constant $\sF_g(0)$, which the rigidity statement (\cref{thm:rig-main}) subtracts. The constant-map term is discussed in \cref{sec:rig} (\cref{rmk:rig-d0,prop:rig-flop,rmk:rig-F0}). The genus-$0$ and genus-$1$ invariants, of degree $d\geq1$, appear only in \cref{sec:rig}.
\end{rmk}

\emph{Derivatives and coefficients.} On functions of $t$ (at fixed auxiliary variables $z,w,\psi,\psi'$) a prime is $\dd/\dd t=\thQ=\cQ\,\dd/\dd\cQ$, and $f_\tau\coloneqq\dd f/\dd\tau$, so $f'=(1-\tau)f_\tau$. On functions of the Euler variable $s$ of \cref{def:set-alg} (phases, amplitudes) a prime or a superscript $(k)$ is $\dd/\dd s$, and $\de_\tau$ is the partial derivative in $\tau$ at fixed $s$. $\Dz\coloneqq z\,\dd/\dd t$. The symbols $\coef{t^k}$, $\coef{z^m}$, $\coef{\psi^a}$, $\coef{w^b}$ denote coefficients, taken coefficientwise in the other variables, and $X^{(k)}\coloneqq\coef{t^k}X$ ($k$ may be negative).

%-------------------------------------------------------------------------------
\subsection{The superpotential and the critical equation}
\label{sec:set-crit}

\begin{lem}[Superpotential, critical points]\genusstatus{AI-proposed, unchecked}
\label{lem:set-crit}
Let $\thU\coloneqq U\,\dd/\dd U$ and
\[
\cW(U)\coloneqq\alpha_1\log(1-\cQ U)-\alpha_2\log U+\wA\log(1-U),\qquad \cW_k\coloneqq\thU^k\cW\quad(k\geq1).
\]
The $\cW_k$ are single valued and rational in $U$ and $\cQ$; explicitly
\[
\cW_1=U\cW_U=-\frac{\alpha_1\cQ U}{1-\cQ U}-\alpha_2-\frac{\wA U}{1-U},\qquad
\cW_2=-\frac{\alpha_1\cQ U}{(1-\cQ U)^2}-\frac{\wA U}{(1-U)^2},\qquad
\de_t\cW=-\frac{\alpha_1\cQ U}{1-\cQ U}.
\]
\begin{enumerate}
\item The critical equation $U\cW_U=0$ is equivalent to
\begin{equation}
\label{eq:set-P}
\Pq(U)\coloneqq\wB\cQ U^2-\big(\alpha_0+\cQ(\alpha_1-\alpha_2)\big)U-\alpha_2=0 ,
\end{equation}
and $(1-\cQ U)(1-U)\,U\cW_U=\Pq(U)$ identically. The roots $U_1,U_2$ of $\Pq$ satisfy
\[
U_1+U_2=\frac{\alpha_0+\cQ(\alpha_1-\alpha_2)}{\wB\cQ},\qquad U_1U_2=-\frac{\alpha_2}{\wB\cQ},
\]
and the discriminant of $\Pq$ in $U$ is $(\alpha_0+\cQ(\alpha_1-\alpha_2))^2+4\alpha_2\wB\cQ=\Disc(\cQ)$ of \cref{lem:set-disc}, equal to $\cQ^2\delta^2$ with $\delta$ as in \cref{def:set-alg}.
\item $\Pq(1)=\wA(\cQ-1)$ and $\Pq(1/\cQ)=\alpha_1(1-\cQ)/\cQ=-\alpha_1\tau$. So, for $\cQ$ finite and non-zero, $U_i=1$ or $\cQ U_i=1$ happens only over $\cQ=1$.
\item Put $\Lv(U)\coloneqq-\alpha_1\cQ U/(1-\cQ U)=\de_t\cW$. At a critical point, $\Lv(U)=\alpha_2+\wA U/(1-U)$, equivalently $U/(1-U)=(\Lv(U)-\alpha_2)/\wA$ and $1/(1-U)=(\Lv(U)+\alpha_0)/\wA$. Moreover $1-\cQ U$ is a unit of $\KCY(\cQ)[U]/(\Pq)$ for $\tau\neq0$, and in that ring
\[
\Lv(U)=\frac{\wB U+\alpha_2}{\tau}.
\]
\item With $p_a\coloneqq(1-aU)^{-1}$ for $a\in\{1,\cQ\}$: $\thU\log(1-aU)=1-p_a$, $\thU p_a=p_a^2-p_a$, $\thU^k\log U=0$ for $k\geq2$, and, with $\Pol_0(p)\coloneqq p$, $\Pol_n(p)\coloneqq(p^2-p)\Pol_{n-1}'(p)$,
\begin{equation}
\label{eq:set-Wk}
\thU^k\log(1-aU)=-\Pol_{k-1}(p_a)\quad(k\geq2),\qquad \cW_k=-\alpha_1\Pol_{k-1}(p_\cQ)-\wA\Pol_{k-1}(p_1)\quad(k\geq2).
\end{equation}
$\Pol_n$ has degree $n+1$, leading coefficient $n!$, and $\Pol_n(0)=0$; we write $\Pol_{k-1}(p)=\sum_{i=1}^{k}\pc_{k-1,i}\,p^i$ with $\pc_{k-1,k}=(k-1)!$. In particular every $\cW_k$, $k\geq1$, is a polynomial in $p_1$ and $p_\cQ$ with coefficients in $\bbZ[\alpha_1,\alpha_2,\wA]$.
\end{enumerate}
\end{lem}

\begin{proof}
(i) Multiply $U\cW_U$ by $(1-\cQ U)(1-U)$:
\[
-\alpha_1\cQ U(1-U)-\alpha_2(1-\cQ U)(1-U)-\wA U(1-\cQ U)
=\cQ U^2(\alpha_1-\alpha_2+\wA)-U(\alpha_1\cQ-\alpha_2\cQ-\alpha_2+\wA)-\alpha_2 ,
\]
and $\alpha_1-\alpha_2+\wA=\wB$, $\alpha_1\cQ-\alpha_2\cQ-\alpha_2+\wA=\alpha_0+\cQ(\alpha_1-\alpha_2)$. \draftnote{The source writes the right side with an overall minus sign; the expansion above gives $+\Pq(U)$. The sign is immaterial for the critical equation.} The sum, product and discriminant of the roots are read off from \eqref{eq:set-P}; the discriminant is expanded in \cref{lem:set-disc}(i), and its identification with $\cQ^2\delta^2$ is \cref{lem:set-disc}(i) with \cref{def:set-alg}.
(ii) $\Pq(1)=\wB\cQ-\alpha_0-\cQ(\alpha_1-\alpha_2)-\alpha_2=\cQ(\alpha_0+\alpha_2)-(\alpha_0+\alpha_2)$, and $\Pq(1/\cQ)=\wB/\cQ-(\alpha_0+\cQ(\alpha_1-\alpha_2))/\cQ-\alpha_2=(\wB-\alpha_0)/\cQ-\alpha_1=\alpha_1(1-\cQ)/\cQ$.
(iii) $U\cW_U=0$ gives $-\alpha_1\cQ U/(1-\cQ U)=\alpha_2+\wA U/(1-U)$; then $1/(1-U)=1+U/(1-U)=(\wA+\Lv-\alpha_2)/\wA=(\Lv+\alpha_0)/\wA$. $1-\cQ U$ is a unit modulo $\Pq$ because $\Pq(1/\cQ)=-\alpha_1\tau\neq0$ by (ii). Finally
\begin{align*}
(\wB U+\alpha_2)(1-\cQ U)+\alpha_1\cQ U\tau&=\wB U-\wB\cQ U^2+\alpha_2-\alpha_2\cQ U+\alpha_1U(\cQ-1)\\
&=-\wB\cQ U^2+U\big(\wB-\alpha_1+\cQ(\alpha_1-\alpha_2)\big)+\alpha_2=-\Pq(U),
\end{align*}
using $\alpha_1\cQ U\tau=\alpha_1U(\cQ-1)$ and $\wB-\alpha_1=\alpha_0$; so $(\wB U+\alpha_2)/\tau=-\alpha_1\cQ U/(1-\cQ U)$ modulo $\Pq$.
(iv) $\thU\log(1-aU)=-aU/(1-aU)=1-p_a$ and $\thU p_a=aU/(1-aU)^2=p_a^2-p_a$; by induction $\thU^k\log(1-aU)=-\thU^{k-1}p_a=-\Pol_{k-1}(p_a)$, since $\thU\,P(p_a)=(p_a^2-p_a)P'(p_a)$ for a polynomial $P$. The degree, leading coefficient and vanishing at $0$ follow by induction from the recursion. Then \eqref{eq:set-Wk} is $\cW=\alpha_1\log(1-\cQ U)-\alpha_2\log U+\wA\log(1-U)$ differentiated $k\geq2$ times.
\end{proof}

The two roots are labelled by their behaviour at the conifold point: $U_1\to1$ and $U_2\to-\alpha_2/\wB$ as $t\to0$ (\cref{def:set-alg}, \cref{lem:set-eig}). Different labellings, by the behaviour at $\cQ=0$ and at $\cQ=\infty$, are used in \cref{lem:set-limits}(ii), (iii) and in \cref{sec:ide}; the text says which one is in force.

%-------------------------------------------------------------------------------
\subsection{Frobenius data}
\label{sec:set-frob}

\begin{prop}[Frobenius data in the flat basis]\genusstatus{AI-proposed, unchecked}
\label{prop:set-frob}
Let $\{1,H\}$ be the flat basis, $H$ the equivariant hyperplane class of $\bbP^1$ with $H|_{p_+}=0$, $H|_{p_-}=-\alpha_0$. The genus-zero potential of the $\TCY$-equivariant theory of $Z$ in the coordinates $t_0,t_1$ dual to $1,H$ is
\[
\mathsf{Prep}=\mathsf{Prep}^{\mathrm{cl}}+\frac{\Li_3(\re^{t_1})}{\epsilon_1\epsilon_2},\qquad
\mathsf{Prep}^{\mathrm{cl}}=\frac16\Big[\frac{t_0^3}{\alpha_0\alpha_1\alpha_2\epsilon_1\epsilon_2}-\frac{(t_0-\alpha_0t_1)^3}{\alpha_0\wA\wB\epsilon_1\epsilon_2}\Big],
\]
with $\cQ=\re^{t_1}$ (\cref{lem:ide-g0}; only its third derivatives are used here). Then:
\begin{enumerate}
\item the metric $\eta_{ij}\coloneqq\mathsf{Prep}_{0ij}$ and its inverse are
\[
\eta=\frac{1}{\wA\wB\epsilon_1\epsilon_2}\begin{pmatrix}E/(\alpha_1\alpha_2)&1\\1&-\alpha_0\end{pmatrix},\qquad
\eta^{-1}=\epsilon_1\epsilon_2\begin{pmatrix}\alpha_0\alpha_1\alpha_2&\alpha_1\alpha_2\\\alpha_1\alpha_2&-E\end{pmatrix},
\]
independent of $t$;
\item $1\qprod a=a$ for all $a$, and
\[
H\qprod H=-\frac{\alpha_1\alpha_2}{\tau}\,1+\Big(\frac E\tau-\alpha_0\Big)H ;
\]
\item the minimal polynomial of $H\qprod$ is $\mathsf{Leq}(\Lv)\coloneqq\Lv^2-(E/\tau-\alpha_0)\Lv+\alpha_1\alpha_2/\tau$; its roots $\Lv_1,\Lv_2$ (the eigenvalues of $H\qprod$) satisfy $\Lv_1+\Lv_2=E/\tau-\alpha_0$ and $\Lv_1\Lv_2=\alpha_1\alpha_2/\tau$;
\item the product is commutative, associative and Frobenius: $\eta(a\qprod b,d)=\eta(a,b\qprod d)$.
\end{enumerate}
Covectors $a=(a_0,a_1)$ (components paired with $1$ and $H$) are paired by $\eta^{-1}$,
\begin{equation}
\label{eq:set-pairing}
\langle a,b\rangle\coloneqq\sum\eta^{a'b'}a_{a'}b_{b'}=\epsilon_1\epsilon_2\big[\alpha_0\alpha_1\alpha_2\,a_0b_0+\alpha_1\alpha_2(a_0b_1+a_1b_0)-E\,a_1b_1\big],
\end{equation}
symmetric with constant coefficients, so that $\frac{\dd}{\dd t}\langle a,b\rangle=\langle a',b\rangle+\langle a,b'\rangle$.
\end{prop}

\begin{proof}
The classical part is the localisation formula at the two fixed points: $\mathsf{Prep}^{\mathrm{cl}}=\frac16\sum_{p=p_\pm}(t_0+t_1H|_p)^3/\eul(\wts_p)$ with $H|_{p_+}=0$, $H|_{p_-}=-\alpha_0$ and the products of the tangent weights $\eul(\wts_+)=\alpha_0\alpha_1\alpha_2\epsilon_1\epsilon_2$, $\eul(\wts_-)=-\alpha_0\wA\wB\epsilon_1\epsilon_2$ (equivalently, $H$ satisfies $H(H+\alpha_0)=0$, whose roots $0,-\alpha_0$ are the fixed points). Its third derivatives are
\[
\begin{gathered}
\mathsf{Prep}^{\mathrm{cl}}_{000}=\frac{1}{\alpha_0\alpha_1\alpha_2\epsilon_1\epsilon_2}-\frac{1}{\alpha_0\wA\wB\epsilon_1\epsilon_2}=\frac{E}{\alpha_1\alpha_2\wA\wB\epsilon_1\epsilon_2},\qquad
\mathsf{Prep}^{\mathrm{cl}}_{001}=\frac{1}{\wA\wB\epsilon_1\epsilon_2},\\
\mathsf{Prep}^{\mathrm{cl}}_{011}=\frac{-\alpha_0}{\wA\wB\epsilon_1\epsilon_2},\qquad
\mathsf{Prep}^{\mathrm{cl}}_{111}=\frac{\alpha_0^2}{\wA\wB\epsilon_1\epsilon_2},
\end{gathered}
\]
using $\wA\wB-\alpha_1\alpha_2=\alpha_0E$; none has $\alpha_0$ in a denominator. (i) $\eta_{ij}=\mathsf{Prep}^{\mathrm{cl}}_{0ij}$ gives $\eta$. Its determinant is $(-\alpha_0E-\alpha_1\alpha_2)/(\alpha_1\alpha_2(\wA\wB\epsilon_1\epsilon_2)^2)=-1/(\alpha_1\alpha_2\wA\wB\epsilon_1^2\epsilon_2^2)$, since $\alpha_0E+\alpha_1\alpha_2=\wA\wB$; inverting,
\[
\eta^{-1}=-\alpha_1\alpha_2\wA\wB\epsilon_1^2\epsilon_2^2\cdot\frac{1}{\wA\wB\epsilon_1\epsilon_2}\begin{pmatrix}-\alpha_0&-1\\-1&E/(\alpha_1\alpha_2)\end{pmatrix}=\epsilon_1\epsilon_2\begin{pmatrix}\alpha_0\alpha_1\alpha_2&\alpha_1\alpha_2\\\alpha_1\alpha_2&-E\end{pmatrix}.
\]
(ii) $(\dd/\dd t_1)^3\Li_3(\re^{t_1})=\Li_0(\cQ)=\cQ/(1-\cQ)=-1/\tau$ (as $1-\cQ=-\cQ\tau$), so $\mathsf{Prep}_{111}=\mathsf{Prep}^{\mathrm{cl}}_{111}-1/(\epsilon_1\epsilon_2\tau)$, and the other third derivatives are classical. The structure constants are $c^i_{jk}=\sum_m\eta^{im}\mathsf{Prep}_{mjk}$, $e_j\qprod e_k=\sum_ic^i_{jk}e_i$ with $e_0=1$, $e_1=H$. Since $\mathsf{Prep}_{0jk}=\eta_{jk}$, $c^i_{0k}=\delta_{ik}$, which is $1\qprod a=a$. Further
\begin{align*}
c^0_{11}&=\epsilon_1\epsilon_2\Big[\alpha_0\alpha_1\alpha_2\cdot\frac{-\alpha_0}{\wA\wB\epsilon_1\epsilon_2}+\alpha_1\alpha_2\Big(\frac{\alpha_0^2}{\wA\wB\epsilon_1\epsilon_2}-\frac{1}{\epsilon_1\epsilon_2\tau}\Big)\Big]=-\frac{\alpha_1\alpha_2}{\tau},\\
c^1_{11}&=\epsilon_1\epsilon_2\Big[\alpha_1\alpha_2\cdot\frac{-\alpha_0}{\wA\wB\epsilon_1\epsilon_2}-E\Big(\frac{\alpha_0^2}{\wA\wB\epsilon_1\epsilon_2}-\frac{1}{\epsilon_1\epsilon_2\tau}\Big)\Big]=-\frac{\alpha_0(\alpha_1\alpha_2+\alpha_0E)}{\wA\wB}+\frac E\tau=\frac E\tau-\alpha_0 .
\end{align*}
(iii) By (ii), $(H\qprod)^2-(E/\tau-\alpha_0)(H\qprod)+(\alpha_1\alpha_2/\tau)=0$ on the two-dimensional algebra generated by $H$, and $H\qprod$ is not scalar, so $\mathsf{Leq}$ is its minimal polynomial; the relations between the roots are Vieta's. (iv) The algebra is generated by $H$, and $\eta(e_j\qprod e_k,e_l)=\mathsf{Prep}_{ljk}$ is totally symmetric. The pairing \eqref{eq:set-pairing} is $\eta^{-1}$ written out.
\end{proof}

\draftnote{compare BS App. A for the Frobenius data; the note derives them from the genus-0 potential, whose derivation is \cref{lem:ide-g0}.} On $\TF$ these data read $\eta^{-1}=\begin{psmallmatrix}0&\xc^2\\\xc^2&-\cc\xc\end{psmallmatrix}$ and $H\qprod H=-(\xc/\tau)1+(\cc/\tau)H$ with $\xc=\alpha_1=\epsilon_1$, $\alpha_2=\epsilon_2=1$, $\cc=1+\xc=E$.

%-------------------------------------------------------------------------------
\subsection{Idempotents, the vertex weights, the quadratic algebra}
\label{sec:set-idem}

\begin{defn}[The quadratic algebra and its involution]
\label{def:set-QA}
Let $\QA\coloneqq\KCY(\cQ)[U]/(\Pq(U))$. It has the two characters $\ev_i:U\mapsto U_i$ ($i=1,2$), ring homomorphisms into the algebraic functions of $\cQ$, jointly injective because $\Pq$ has distinct roots ($\Disc\neq0$, \cref{lem:set-disc}). The \emph{involution} of $\QA$ is the ring automorphism $U\mapsto U_1+U_2-U$ (the Galois conjugation), which exchanges the two characters: $\ev_1\circ(\text{inv})=\ev_2$ and $\ev_2\circ(\text{inv})=\ev_1$. By \cref{lem:set-disc}(v), $\QA$ is the function field $\KCY(\Ccov)$ of the double cover $\Ccov\to\bbP^1_\cQ$, $\ydc^2=\Disc(\cQ)$, and the involution is its deck involution. By \cref{lem:set-crit}(iii), $\Lv(U)=(\wB U+\alpha_2)/\tau$ lies in $\QA$, and $\ev_i\Lv(U)=\Lv_i$.
\end{defn}

Every factor of idempotent $2$ defined below is obtained from the corresponding factor of idempotent $1$ by applying the involution, that is, by evaluating the same expression in $U$ at the other root; this is because the two scalar series $\hnw$ and $\hw^{-1}$ of \cref{def:set-rows} are the same at both idempotents (by \cref{def:set-rows}; that this is the correct normalisation at both idempotents is \cref{lem:ide-norm}). This is made precise in \cref{lem:lea-conj}.

\begin{prop}[Idempotents and vertex weights]\genusstatus{AI-proposed, unchecked}
\label{prop:set-idem}
Let $\Lv_1\neq\Lv_2$ be the roots of $\mathsf{Leq}$.
\begin{enumerate}
\item $\eid_1=(H-\Lv_2)/(\Lv_1-\Lv_2)$ and $\eid_2=(H-\Lv_1)/(\Lv_2-\Lv_1)$ are the idempotents: $\eid_i\qprod\eid_j=\delta_{ij}\eid_i$, $H\qprod\eid_i=\Lv_i\eid_i$, $\eid_1+\eid_2=1$.
\item $\eta(\eid_i,\eid_j)=0$ for $i\neq j$. Put $\Delta_i\coloneqq1/\eta(\eid_i,\eid_i)$; the normalised idempotents are $\fid_i\coloneqq\Delta_i^{1/2}\eid_i$. The \emph{vertex weight} of a vertex of genus $g_v$ and colour $i$ is $\Delta_i^{g_v-1}$.
\item For $\{i,j\}=\{1,2\}$,
\begin{equation}
\label{eq:set-Delta}
\begin{split}
\Delta_i&=\frac{\alpha_1\alpha_2\wA\wB\epsilon_1\epsilon_2(\Lv_i-\Lv_j)}{\alpha_1\alpha_2-E\Lv_j}
=\epsilon_1\epsilon_2(\Lv_i-\Lv_j)(\alpha_1\alpha_2-E\Lv_i)\\
&=\epsilon_1\epsilon_2(\alpha_0\alpha_1\alpha_2+2\alpha_1\alpha_2\Lv_i-E\Lv_i^2)
=\alpha_1\wA\epsilon_1\epsilon_2\,\cW_2(U_i),
\end{split}
\end{equation}
and $(\alpha_1\alpha_2-E\Lv_1)(\alpha_1\alpha_2-E\Lv_2)=\alpha_1\alpha_2\wA\wB$. The four expressions of \eqref{eq:set-Delta}, with $\Lv_i=\Lv(U)$, $\Lv_j=E/\tau-\alpha_0-\Lv(U)$ and $\cW_2=\cW_2(U)$, are equal as elements of $\QA$.
\item $\Delta_1\Delta_2=-\epsilon_1^2\epsilon_2^2\,\alpha_1\alpha_2\wA\wB(\Lv_1-\Lv_2)^2$ and $\dfrac1{\Delta_1}+\dfrac1{\Delta_2}=\dfrac{E}{\alpha_1\alpha_2\wA\wB\epsilon_1\epsilon_2}=\eta(1,1)$.
\end{enumerate}
\end{prop}

\begin{proof}
(i) By \cref{prop:set-frob}(ii), (iii), $(H-\Lv_1)\qprod(H-\Lv_2)=H\qprod H-(\Lv_1+\Lv_2)H+\Lv_1\Lv_2=0$. Hence $(H-\Lv_1)\qprod\eid_1=0$, that is $H\qprod\eid_1=\Lv_1\eid_1$, and $\eid_1\qprod\eid_2=0$; with $\eid_1+\eid_2=1$ this gives $\eid_i\qprod\eid_i=\eid_i$. (ii) $\eta(\eid_i,\eid_j)=\eta(\eid_i\qprod\eid_j,1)=0$ for $i\neq j$ by the Frobenius property.
(iii) First form: $\eta(\eid_i,\eid_i)=\eta(\eid_i\qprod\eid_i,1)=\eta(\eid_i,1)=\big(\eta(H,1)-\Lv_j\eta(1,1)\big)/(\Lv_i-\Lv_j)$, and by \cref{prop:set-frob}(i) this is
\[
\frac{1}{\Lv_i-\Lv_j}\Big[\frac{1}{\wA\wB\epsilon_1\epsilon_2}-\frac{\Lv_jE}{\alpha_1\alpha_2\wA\wB\epsilon_1\epsilon_2}\Big]=\frac{\alpha_1\alpha_2-E\Lv_j}{\alpha_1\alpha_2\wA\wB\epsilon_1\epsilon_2(\Lv_i-\Lv_j)}.
\]
The product identity: $(\alpha_1\alpha_2-E\Lv_1)(\alpha_1\alpha_2-E\Lv_2)=\alpha_1^2\alpha_2^2-E\alpha_1\alpha_2(E/\tau-\alpha_0)+E^2\alpha_1\alpha_2/\tau=\alpha_1\alpha_2(\alpha_1\alpha_2+\alpha_0E)=\alpha_1\alpha_2\wA\wB$; so $\alpha_1\alpha_2\wA\wB/(\alpha_1\alpha_2-E\Lv_j)=\alpha_1\alpha_2-E\Lv_i$, which is the second form. Third form: $(\Lv_i-\Lv_j)(\alpha_1\alpha_2-E\Lv_i)=\alpha_1\alpha_2\Lv_i-E\Lv_i^2-\alpha_1\alpha_2\Lv_j+E\Lv_i\Lv_j$, with $-\alpha_1\alpha_2\Lv_j=\alpha_1\alpha_2\Lv_i-\alpha_1\alpha_2(E/\tau-\alpha_0)$ and $E\Lv_i\Lv_j=E\alpha_1\alpha_2/\tau$; the terms in $\tau^{-1}$ cancel and $\alpha_0\alpha_1\alpha_2+2\alpha_1\alpha_2\Lv_i-E\Lv_i^2$ remains. Fourth form: at a critical point $\Lv=-\alpha_1\cQ U/(1-\cQ U)$, so $1/(1-\cQ U)=1+\cQ U/(1-\cQ U)=(\alpha_1-\Lv)/\alpha_1$ and $-\alpha_1\cQ U/(1-\cQ U)^2=\Lv(\alpha_1-\Lv)/\alpha_1$; and $U/(1-U)=(\Lv-\alpha_2)/\wA$, $1/(1-U)=(\Lv+\alpha_0)/\wA$ (\cref{lem:set-crit}(iii)), so $\wA U/(1-U)^2=(\Lv-\alpha_2)(\Lv+\alpha_0)/\wA$. Hence
\begin{align*}
\alpha_1\wA\,\cW_2(U)&=\wA\Lv(\alpha_1-\Lv)-\alpha_1(\Lv-\alpha_2)(\Lv+\alpha_0)\\
&=-(\wA+\alpha_1)\Lv^2+\alpha_1(\wA-\alpha_0+\alpha_2)\Lv+\alpha_0\alpha_1\alpha_2=-E\Lv^2+2\alpha_1\alpha_2\Lv+\alpha_0\alpha_1\alpha_2 .
\end{align*}
Each of the four expressions is an element of $\QA$ (with $\Lv=\Lv(U)\in\QA$ and $1-U$, $1-\cQ U$ units of $\QA$ for $\tau\neq0$, by \cref{lem:set-crit}(ii)), and the computations above show that they have equal values under $\ev_1$ and under $\ev_2$; as the two characters are jointly injective, they are equal in $\QA$.
(iv) The product follows from the second form of \eqref{eq:set-Delta} and the product identity; the sum is $\eta(\eid_1,\eid_1)+\eta(\eid_2,\eid_2)=\eta(\eid_1+\eid_2,1)=\eta(1,1)$.
\end{proof}

%-------------------------------------------------------------------------------
\subsection{Algebraic functions of \texorpdfstring{$\tau$}{tau}; the critical polynomial}
\label{sec:set-alg}

\begin{defn}[Algebraic functions of $\tau$]
\label{def:set-alg}
\begin{description}
\item[$\delta$.] Put
\[
\delta^2\coloneqq(E+\alpha_0\tau)^2-4\wA\wB\tau=E^2-(E^2+\uq)\tau+\alpha_0^2\tau^2=(E-\alpha_0\tau)^2-4\alpha_1\alpha_2\tau ;
\]
the three forms agree since $4\wA\wB=4\alpha_1\alpha_2+4\alpha_0E$ and $2\alpha_0E+4\alpha_1\alpha_2=E^2+\uq$. Let $\delta$ be the root analytic near $\tau=0$ with $\delta(0)=E$. It is analytic on the disc $\Ddisc\coloneqq\{|\tau|<r_\Ddisc\}$, $r_\Ddisc$ the smallest modulus of a root of $\delta^2$; $r_\Ddisc$ is finite under (ND): for $\alpha_0\neq0$, $\delta^2$ is quadratic in $\tau$, and at $\alpha_0=0$, where $E^2+\uq=4\alpha_1\alpha_2\neq0$, $\delta^2=E^2-4\alpha_1\alpha_2\tau$. Further $\Ddisc'\coloneqq\Ddisc\cap\{|\tau|<1\}$ and $\Ddisc''\coloneqq\Ddisc'\setminus\{0\}$. (On $\TF$: $\delta^2=\cc^2-4\xc\tau$.)
\item[The critical polynomial.] $\Ns(s)\coloneqq\wB\tau s^2-(E+\alpha_0\tau)s+\wA$, with roots
\[
\sast\coloneqq\frac{E+\alpha_0\tau-\delta}{2\wB\tau}=\frac{2\wA}{E+\alpha_0\tau+\delta},\qquad
\stil\coloneqq\frac{E+\alpha_0\tau+\delta}{2\wB\tau}=\frac{\wA}{\wB\tau\sast},
\]
so that $\sast+\stil=(E+\alpha_0\tau)/(\wB\tau)$ and $\sast\stil=\wA/(\wB\tau)$. $\sast$ is analytic on $\Ddisc$ with $\sast(0)=\wA/E$: the denominator $E+\alpha_0\tau+\delta$ vanishes only where $(E+\alpha_0\tau)^2=\delta^2$, that is $4\wA\wB\tau=0$, that is $\tau=0$, where it equals $2E\neq0$. $\tau\stil$ is analytic with value $E/\wB$ at $0$. One has $\Ns'(\sast)=2\wB\tau\sast-(E+\alpha_0\tau)=-\delta$, $\Ns'(\stil)=+\delta$, and $\delta=E+\alpha_0\tau-2\wB\tau\sast$.
\item[The critical points.] $U_1\coloneqq1-\tau\sast$ and $U_2\coloneqq1-\tau\stil$; these are the two roots of $\Pq$ (\cref{lem:set-eig}), with $U_1\to1$ and $U_2\to1-E/\wB=-\alpha_2/\wB$ as $t\to0$. $U_1$ is the root with the singular vertex weight (\cref{lem:set-limits}); idempotent $1$ is the \emph{singular} idempotent and idempotent $2$ the \emph{regular} one.
\item[The Euler phase.] $\phi_\tau(s)\coloneqq\wA\log s+\alpha_1\log(1-s)-\alpha_2\log(1-\tau s)$, with principal branches near $(\tau,s)=(0,\wA/E)$; $\phi_0(\wA/E)=\wA\log(\wA/E)+\alpha_1\log(\alpha_1/E)$. Put $\dphi(\tau)\coloneqq\phi_\tau(\sast(\tau))-\phi_0(\wA/E)$, an analytic germ at $\tau=0$ with $\dphi(0)=0$.
\end{description}
\end{defn}

\begin{lem}[The critical polynomial]\genusstatus{AI-proposed, unchecked}
\label{lem:set-N}
Let $\tau\neq0$.
\begin{enumerate}
\item[(a)] $\phi_\tau'(s)=\dfrac{\wA}{s}-\dfrac{\alpha_1}{1-s}+\dfrac{\alpha_2\tau}{1-\tau s}=\dfrac{\Ns(s)}{s(1-s)(1-\tau s)}$. The critical points of $\phi_\tau$ in $\bbC\setminus\{0,1,1/\tau\}$ are $\sast$ and $\stil$. On $\Ddisc'$, $\sast\notin\{0,1,1/\tau\}$; on $\Ddisc''$, $\stil\notin\{0,1,1/\tau\}$ and $U_2\notin\{0,1,1/\cQ\}$.
\item[(b)] At $s=\sast$ (and, for (ii) and (iv), also at $s=\stil$; (i) and (iii) hold at $\stil$ with $\delta$ replaced by $-\delta$ and $\wB\sast-\alpha_0$ by $\wB\stil-\alpha_0$):
\begin{description}
\item[(i)] $2\wB\tau\sast-(E+\alpha_0\tau)=-\delta$;
\item[(ii)] $(1-\tau)\alpha_2s=(\wA-\wB s)(1-\tau s)$, equivalently $(E-\wB\tau s)(1-s)=\alpha_1(1-\tau s)$; moreover $E-\wB\tau\sast=\delta+\tau\Lv_2$ with $\Lv_2\coloneqq\wB\sast-\alpha_0$;
\item[(iii)] $\dfrac{\dd\sast}{\dd\tau}=\dfrac{\sast\Lv_2}{\delta}$;
\item[(iv)] $\wA-Es=-\tau s(\wB s-\alpha_0)$.
\end{description}
Consequently (envelope formula)
\begin{equation}
\label{eq:set-envelope}
\dphi_\tau=\frac{\alpha_2\sast}{1-\tau\sast}=\frac{\wA-\wB\sast}{1-\tau},\qquad \dphi'=\wA-\wB\sast ,
\end{equation}
and $\dphi$ lies in $\tau\Oan$, with $\dphi=(\alpha_2\wA/E)\tau+O(\tau^2)$.
\end{enumerate}
\end{lem}

\begin{proof}
(a) The derivative is immediate from the definition of $\phi_\tau$. Over the common denominator the numerator is
\[
\wA(1-s)(1-\tau s)-\alpha_1s(1-\tau s)+\alpha_2\tau s(1-s)=\wA-(\wA+\alpha_1)s-\wA\tau s+\alpha_2\tau s+(\wA+\alpha_1-\alpha_2)\tau s^2=\Ns(s),
\]
using $\wA+\alpha_1=E$, $\alpha_2-\wA=-\alpha_0$, $\wA+\alpha_1-\alpha_2=\wB$. The roots of $\Ns$ are $\sast,\stil$ (\cref{def:set-alg}). $s=0$ is not a root ($\Ns(0)=\wA\neq0$); $\Ns(1)=\wB\tau-E-\alpha_0\tau+\wA=(\wB-\alpha_0)\tau-\alpha_1=\alpha_1(\tau-1)$, so a root equals $1$ only if $\tau=1$ ($\alpha_1\neq0$); $\Ns(1/\tau)=\wB/\tau-E/\tau-\alpha_0+\wA=(\wB-E)/\tau+\alpha_2=\alpha_2(1-1/\tau)$, so a root equals $1/\tau$ only if $\tau=1$ ($\alpha_2\neq0$). On $\Ddisc''$ both roots are finite ($\wB\tau\neq0$) and $\tau\neq1$; on $\Ddisc'$ the root $\sast$ is analytic with $\sast(0)=\wA/E\notin\{0,1\}$. Finally $U_2=0$ means $\stil=1/\tau$, and $U_2\in\{1,1/\cQ\}$ means $\stil\in\{0,1\}$, because $1-\cQ U_2=1-\cQ+\cQ\tau\stil=-\cQ\tau(1-\stil)$.

(b) (i) is $\Ns'(\sast)=-\delta$. For (ii), first form:
\begin{align*}
(1-\tau)\alpha_2s-(\wA-\wB s)(1-\tau s)&=\alpha_2s-\alpha_2\tau s-\wA+\wA\tau s+\wB s-\wB\tau s^2\\
&=-\big[\wB\tau s^2-(\wB+\alpha_2)s+(\alpha_2-\wA)\tau s+\wA\big]=-\Ns(s),
\end{align*}
since $\wB+\alpha_2=E$ and $\alpha_2-\wA=-\alpha_0$. Second form:
\[
(E-\wB\tau s)(1-s)-\alpha_1(1-\tau s)=E-Es-\wB\tau s+\wB\tau s^2-\alpha_1+\alpha_1\tau s=\wA-Es+\wB\tau s^2-(\wB-\alpha_1)\tau s=\Ns(s).
\]
Both vanish at both roots. By (i), $\delta=E+\alpha_0\tau-2\wB\tau\sast$, and $\tau\Lv_2=\wB\tau\sast-\alpha_0\tau$, so $E-\wB\tau\sast=\delta+\tau\Lv_2$. (iii) Differentiate $\Ns(\sast(\tau),\tau)=0$ in $\tau$: $\wB\sast{}^2+2\wB\tau\sast\sast_\tau-\alpha_0\sast-(E+\alpha_0\tau)\sast_\tau=0$, that is $\Ns'(\sast)\sast_\tau=-(\wB\sast{}^2-\alpha_0\sast)=-\sast\Lv_2$; with (i), $\sast_\tau=\sast\Lv_2/\delta$. At $\stil$ the same computation gives $\stil_\tau=-\stil(\wB\stil-\alpha_0)/\delta$. (iv) is $\Ns(s)=0$ rearranged: $\wA-Es=\alpha_0\tau s-\wB\tau s^2$.

Envelope: $\frac{\dd}{\dd\tau}\phi_\tau(\sast(\tau))=(\de_\tau\phi_\tau)(\sast)+\phi_\tau'(\sast)\sast_\tau=(\de_\tau\phi_\tau)(\sast)=\alpha_2\sast/(1-\tau\sast)$, since $\phi'_\tau(\sast)=0$ and $\de_\tau\phi_\tau(s)=\alpha_2s/(1-\tau s)$; then the first form of (ii) gives $(\wA-\wB\sast)/(1-\tau)$, and $\dd/\dd t=(1-\tau)\dd/\dd\tau$. Since $\sast$, $1-\tau\sast$ are analytic on $\Ddisc$ with $\sast(0)=\wA/E$, $\dphi_\tau$ is analytic at $0$ with value $\alpha_2\wA/E$, and $\dphi(0)=0$.
\end{proof}

On $\TF$, (ii) is the factorisation $1-s=\xc s(1-\tau s)$; at generic weights $\Ns$ has no factorisation of this form, and the two forms of (ii) are what replaces it in \cref{sec:row,sec:lea}.

\begin{lem}[Eigenvalues in terms of $\sast$]\genusstatus{AI-proposed, unchecked}
\label{lem:set-eig}
\begin{enumerate}
\item With $U=1-\tau s$ ($\tau\neq0$), $U$ is a root of $\Pq$ if and only if $\Ns(s)=0$; indeed $\Pq(1-\tau s)=\cQ\tau\Ns(s)$ as polynomials in $s$. The roots of $\Pq$ are $U_1=1-\tau\sast$ and $U_2=1-\tau\stil$.
\item At a root $U_i$, $(\de_t\cW)(U_i)=\Lv(U_i)=\alpha_2+\wA U_i/(1-U_i)$, and
\[
\Lv_1=\frac{\wA}{\tau\sast}-\alpha_0,\qquad \Lv_2=\wB\sast-\alpha_0=\frac{\wA}{\tau\stil}-\alpha_0,\qquad
\Lv_1+\Lv_2=\frac E\tau-\alpha_0,\qquad \Lv_1\Lv_2=\frac{\alpha_1\alpha_2}{\tau},\qquad \Lv_1-\Lv_2=\frac\delta\tau .
\]
So $\Lv_1,\Lv_2$ are the roots of $\mathsf{Leq}$, the eigenvalues of $H\qprod$, and the $\Lv_2$ of \cref{lem:set-N} is this one.
\item $\Lv_1=E/t+\lrem(t)$ with $\lrem\coloneqq\wA/(\tau\sast)-\alpha_0-E/t\in\Oan$, $\lrem(0)=E/2-\alpha_0-\alpha_1\alpha_2/E$. $\Lv_2\in\Oan$ with $\Lv_2(0)=\alpha_1\alpha_2/E$. Also $\sast=\wA/E+s_1\tau+O(\tau^2)$ with $s_1\coloneqq\sast_\tau(0)=\wA\alpha_1\alpha_2/E^3$, and $\delta_\tau(0)=-(E^2+\uq)/(2E)$.
\end{enumerate}
\end{lem}

\begin{proof}
(i) Compare the two sides of $\Pq(1-\tau s)=\cQ\tau\Ns(s)$: both are quadratic in $s$; the coefficients of $s^2$ are $\wB\cQ\tau^2$ and $\cQ\tau\cdot\wB\tau$; the constant terms are $\Pq(1)=\wA(\cQ-1)=\wA\cQ\tau$ and $\cQ\tau\wA$; the coefficients of $s$ are $-\tau\Pq'(1)=-\tau\big(2\wB\cQ-\alpha_0-\cQ(\alpha_1-\alpha_2)\big)=-\tau\big(\cQ E+\alpha_0(\cQ-1)\big)=-\cQ\tau E-\alpha_0\cQ\tau^2$ and $-\cQ\tau(E+\alpha_0\tau)$. The roots correspond. (ii) The first statement is \cref{lem:set-crit}(iii). With $1-U_i=\tau s_i$ ($s_1=\sast$, $s_2=\stil$): $\Lv_i=\alpha_2+\wA(1-\tau s_i)/(\tau s_i)=\wA/(\tau s_i)-\alpha_0$. From $\sast\stil=\wA/(\wB\tau)$, $\wA/(\tau\stil)=\wB\sast$. Sum: $\wA(\sast+\stil)/(\tau\sast\stil)-2\alpha_0=(E+\alpha_0\tau)/\tau-2\alpha_0$. Product: $(\wA/(\tau\sast)-\alpha_0)(\wB\sast-\alpha_0)=\wA\wB/\tau-\alpha_0(\Lv_1+\alpha_0)-\alpha_0(\Lv_2+\alpha_0)+\alpha_0^2=(\wA\wB-\alpha_0E)/\tau=\alpha_1\alpha_2/\tau$. Difference: $\wA/(\tau\sast)-\wB\sast=(\wA-\wB\tau\sast{}^2)/(\tau\sast)=(E+\alpha_0\tau-2\wB\tau\sast)/\tau=\delta/\tau$, by $\Ns(\sast)=0$ and \cref{lem:set-N}(b)(i). By \cref{prop:set-frob}(iii) the roots of $\mathsf{Leq}$ are determined by their sum and product. (iii) $1/\tau=1/t+1/2+O(t)$. From \cref{lem:set-N}(b)(iii), $s_1=\sast(0)\Lv_2(0)/\delta(0)$ with $\Lv_2(0)=\wB\wA/E-\alpha_0=(\wA\wB-\alpha_0E)/E=\alpha_1\alpha_2/E$, so $s_1=(\wA/E)(\alpha_1\alpha_2/E)/E$. Then $\wA/(\tau\sast)=(E/\tau)\big(1-(E/\wA)s_1\tau+O(\tau^2)\big)=E/\tau-(E^2/\wA)s_1+O(\tau)=E/t+E/2-\alpha_1\alpha_2/E+O(t)$, which gives $\lrem\in\Oan$ and $\lrem(0)$. $\Lv_2=\wB\sast-\alpha_0$ is analytic with the stated value. Finally $2\delta\delta_\tau=-(E^2+\uq)+2\alpha_0^2\tau$ gives $\delta_\tau(0)$.
\end{proof}

On $\TF$: $\Lv_1=1/(\tau\sast)$, $\Lv_2=\xc\sast$, $\lrem(0)=\cc/2-\xc/\cc$.

\begin{lem}[$\Delta_2$ in terms of $\sast$]\genusstatus{AI-proposed, unchecked}
\label{lem:set-Delta2}
\begin{enumerate}
\item $\alpha_1\alpha_2-E\Lv_2=-\wB\tau\sast\Lv_2$, and
\[
\Delta_2=\epsilon_1\epsilon_2\wB\,\delta\,\sast\Lv_2,\qquad
\Delta_1\Delta_2=-\epsilon_1^2\epsilon_2^2\,\alpha_1\alpha_2\wA\wB\,\frac{\delta^2}{\tau^2}.
\]
Consequently $\Delta_2/(\Lv_2-\Lv_1)=\epsilon_1\epsilon_2(\alpha_1\alpha_2-E\Lv_2)=-\epsilon_1\epsilon_2\wB\tau\sast\Lv_2$ (this quantity is named $\Thf$ in \cref{sec:lea}).
\item $\Delta_2\in\Oan$ with $\Delta_2(0)=\epsilon_1\epsilon_2\alpha_1\alpha_2\wA\wB/E\neq0$ and $\Delta_2'(0)=0$.
\end{enumerate}
\end{lem}

\begin{proof}
(i) By \cref{lem:set-eig}(ii), $\alpha_1\alpha_2-E\Lv_2=\alpha_1\alpha_2-E(\wB\sast-\alpha_0)=\wA\wB-E\wB\sast=\wB(\wA-E\sast)=-\wB\tau\sast(\wB\sast-\alpha_0)$ by \cref{lem:set-N}(b)(iv). The second form of \eqref{eq:set-Delta} with $\Lv_2-\Lv_1=-\delta/\tau$ gives $\Delta_2=\epsilon_1\epsilon_2(-\delta/\tau)(-\wB\tau\sast\Lv_2)$. The product is \cref{prop:set-idem}(iv) with $(\Lv_1-\Lv_2)^2=\delta^2/\tau^2$, and the last formula is the second form of \eqref{eq:set-Delta} divided by $\Lv_2-\Lv_1$.
(ii) $\Delta_2=\epsilon_1\epsilon_2\wB\delta\sast\Lv_2$ is a product of elements of $\Oan$, with value $\epsilon_1\epsilon_2\wB E(\wA/E)(\alpha_1\alpha_2/E)$ at $0$. Derivative in $\tau$ at $0$, with $\delta_\tau(0)=-(E^2+\uq)/(2E)$, $\sast_\tau(0)=s_1=\wA\alpha_1\alpha_2/E^3$ and $(\Lv_2)_\tau(0)=\wB s_1$ (\cref{lem:set-eig}(iii)):
\[
(\delta\sast\Lv_2)_\tau(0)=-\frac{E^2+\uq}{2E}\cdot\frac{\wA}{E}\cdot\frac{\alpha_1\alpha_2}{E}+E\,s_1\frac{\alpha_1\alpha_2}{E}+E\frac{\wA}{E}\wB s_1
=\frac{\wA\alpha_1\alpha_2}{E^3}\Big[-\frac{E^2+\uq}{2}+\alpha_1\alpha_2+\wA\wB\Big]=0,
\]
because $\wA\wB+\alpha_1\alpha_2=2\alpha_1\alpha_2+\alpha_0E=(E^2+\uq)/2$. Since $\dd/\dd t=(1-\tau)\dd/\dd\tau$, also $\Delta_2'(0)=0$.
\end{proof}

The companion statements for $\Delta_1$ (the form $\Delta_1=-\alpha_1\wA\epsilon_1\epsilon_2\delta U_1/(\tau^2\sast(1-\sast))=-\epsilon_1\epsilon_2E^3\alg^2/\tau^2$, which is (F1)) need the Hessian of the Euler phase and are \cref{lem:row-hess} and \cref{prop:lea-key}; the leading term is \cref{lem:set-limits}(i).

%-------------------------------------------------------------------------------
\subsection{The discriminant and the double cover; behaviour at \texorpdfstring{$\cQ=1,0,\infty$}{Q = 1, 0, infinity}}
\label{sec:set-disc}

\begin{lem}[The discriminant]\genusstatus{AI-proposed, unchecked}
\label{lem:set-disc}
Let $\Disc(\cQ)\coloneqq\alpha_0^2+\big(2\alpha_0(\alpha_1+\alpha_2)+4\alpha_1\alpha_2\big)\cQ+(\alpha_1-\alpha_2)^2\cQ^2$.
\begin{enumerate}
\item $\Disc=\cQ^2\tau^2(\Lv_1-\Lv_2)^2=\cQ^2\delta^2=(\alpha_0+\cQ(\alpha_1-\alpha_2))^2+4\alpha_2\wB\cQ$, the discriminant of $\Pq(U)$ in $U$.
\item The discriminant of $\Disc$ as a polynomial in $\cQ$ is $\big(2\alpha_0(\alpha_1+\alpha_2)+4\alpha_1\alpha_2\big)^2-4(\alpha_1-\alpha_2)^2\alpha_0^2=16\alpha_1\alpha_2\wA\wB$.
\item $\Disc(1)=E^2$, $\Disc(0)=\alpha_0^2$, and the leading coefficient is $(\alpha_1-\alpha_2)^2$. So, at the generic point of $\TCY$, $\Disc$ has no root over $\cQ\in\{0,1\}$ and has two distinct finite roots.
\item $\Disc$ is irreducible over $\KCY=\Kred(\alpha_0)$.
\item The double cover $\Ccov\to\bbP^1_\cQ$ defined by $\ydc^2=\Disc(\cQ)$ has function field $\KCY(\Ccov)=\QA$, where $\QA=\KCY(\cQ)[U]/(\Pq(U))$ is the quadratic algebra of \cref{def:set-QA}, the two descriptions being identified by
\[
2\wB\cQ U-\big(\alpha_0+\cQ(\alpha_1-\alpha_2)\big)=\pm\ydc ;
\]
it is ramified exactly over the two roots of $\Disc$ and unramified over $\cQ=0,1,\infty$.
\end{enumerate}
\end{lem}

\begin{proof}
(i) By \cref{lem:set-eig}(ii), $\tau^2(\Lv_1-\Lv_2)^2=\delta^2=(E-\alpha_0\tau)^2-4\alpha_1\alpha_2\tau$. Multiply by $\cQ^2$ and expand with $\cQ\tau=\cQ-1$: $\cQ^2(E-\alpha_0\tau)^2=(\cQ(E-\alpha_0)+\alpha_0)^2=\cQ^2(E-\alpha_0)^2+2\alpha_0\cQ(E-\alpha_0)+\alpha_0^2$ and $-4\alpha_1\alpha_2\tau\cQ^2=-4\alpha_1\alpha_2\cQ^2+4\alpha_1\alpha_2\cQ$; with $(E-\alpha_0)^2-4\alpha_1\alpha_2=(\alpha_1-\alpha_2)^2$ and $2\alpha_0(E-\alpha_0)=2\alpha_0(\alpha_1+\alpha_2)$ this is $\Disc$. The last form is a direct expansion: $(\alpha_0+\cQ(\alpha_1-\alpha_2))^2+4\alpha_2(\alpha_0+\alpha_1)\cQ=\alpha_0^2+\cQ\big(2\alpha_0(\alpha_1-\alpha_2)+4\alpha_0\alpha_2+4\alpha_1\alpha_2\big)+\cQ^2(\alpha_1-\alpha_2)^2$.
(ii) With $p\coloneqq\alpha_1+\alpha_2$, $q\coloneqq\alpha_1\alpha_2$: $(2\alpha_0p+4q)^2-4\alpha_0^2(p^2-4q)=16\alpha_0pq+16q^2+16\alpha_0^2q=16q(\alpha_0^2+\alpha_0p+q)=16\alpha_1\alpha_2(\alpha_0+\alpha_1)(\alpha_0+\alpha_2)$.
(iii) $\Disc(1)=\alpha_0^2+2\alpha_0(\alpha_1+\alpha_2)+(\alpha_1+\alpha_2)^2=E^2$; the rest is read off. Since $\alpha_1\neq\alpha_2$ and, by (ii), the discriminant is non-zero, the two roots are finite and distinct.
(iv) A quadratic over a field is reducible if and only if its discriminant is a square. $16\alpha_1\alpha_2(\alpha_0+\alpha_1)(\alpha_0+\alpha_2)$, as a polynomial in $\alpha_0$ over $\Kred$, has the two simple roots $-\alpha_1\neq-\alpha_2$, so it is not a square in $\Kred[\alpha_0]$, hence not in $\Kred(\alpha_0)$ (if $f=(p/q)^2$ with $p,q$ coprime in the unique factorisation domain $\Kred[\alpha_0]$, then $q$ is a unit).
(v) Solving $\Pq(U)=0$ gives $U=\big(\alpha_0+\cQ(\alpha_1-\alpha_2)\pm\sqrt{\Disc}\big)/(2\wB\cQ)$, so $\KCY(\cQ)[U]/(\Pq)=\KCY(\cQ)[\ydc]/(\ydc^2-\Disc)$. The cover $\ydc^2=\Disc$ is ramified where $\Disc$ has a simple zero, that is at its two roots (simple by (ii)); over $\cQ=0$ and $\cQ=1$ it is unramified by (iii); over $\cQ=\infty$, in the coordinate $\sQ=1/\cQ$, $\ydc^2=\cQ^2\big((\alpha_1-\alpha_2)^2+O(\sQ)\big)$ and $\ydc/\cQ\to\pm(\alpha_1-\alpha_2)\neq0$, so it is unramified there as well.
\end{proof}

\begin{lem}[Behaviour at $\cQ=1$, $\cQ=0$, $\cQ=\infty$]\genusstatus{AI-proposed, unchecked}
\label{lem:set-limits}
\begin{enumerate}
\item At $\cQ=1$ ($\tau\to0$): $\Lv_1=E/\tau+O(1)$ and $\Lv_2\to\alpha_1\alpha_2/E$;
\[
t^2\Delta_1\in\Oan\ \text{ with }\ t^2\Delta_1\big|_{t=0}=-\epsilon_1\epsilon_2E^3,\qquad \Delta_2(0)=\epsilon_1\epsilon_2\alpha_1\alpha_2\wA\wB/E\neq0 .
\]
Idempotent $1$ is therefore the singular one. The exact unit $t^2\Delta_1/(-\epsilon_1\epsilon_2E^3)=\rho^{-2}$ is (F1) of \cref{thm:lea-main} (\cref{prop:lea-key}).
\item At $\cQ=0$, label the roots of $\Pq$ by $U_+\to-\alpha_2/\alpha_0$ and $\cQ U_-\to\alpha_0/\wB$. Then $\Lv_+\to0$, $\Lv_-\to-\alpha_0$, and
\[
\Delta_+(0)=\eul(\wts_+)=\alpha_0\alpha_1\alpha_2\epsilon_1\epsilon_2,\qquad \Delta_-(0)=\eul(\wts_-)=-\alpha_0\wA\wB\epsilon_1\epsilon_2,
\]
both non-zero in $\KCY$ because $\alpha_0\neq0$: the point $\cQ=0$ is semisimple at $\alpha_0\neq0$. (On $\Tthree$ and on $\TF$ it is not; see \cref{lem:rig-3Tid}.)
\item At $\cQ=\infty$ ($\sQ\coloneqq1/\cQ\to0$): $U_1(\sQ)$ and $V_2(\sQ)\coloneqq\cQ U_2$ are power series in $\sQ$ over $\KCY$ with $U_1\to(\alpha_1-\alpha_2)/\wB$ and $V_2\to-\alpha_2/(\alpha_1-\alpha_2)$; $\Lv_1\to\alpha_1$, $\Lv_2\to\alpha_2$, and
\begin{equation}
\label{eq:set-Dinf}
\begin{gathered}
\Delta_1(\infty)=\epsilon_1\epsilon_2(\alpha_1-\alpha_2)(\alpha_1\alpha_2-E\alpha_1)=-\alpha_1\wB(\alpha_1-\alpha_2)\epsilon_1\epsilon_2,\\
\Delta_2(\infty)=\epsilon_1\epsilon_2(\alpha_2-\alpha_1)(\alpha_1\alpha_2-E\alpha_2)=\alpha_2\wA(\alpha_1-\alpha_2)\epsilon_1\epsilon_2,
\end{gathered}
\end{equation}
both non-zero in $\KCY$. Here the labels $1,2$ are those of the two points of $\Ccov$ over $\cQ=\infty$ at which $U\to(\alpha_1-\alpha_2)/\wB$, respectively $U\to0$.
\end{enumerate}
\end{lem}

\begin{proof}
(i) By \cref{lem:set-eig}(iii), $\Lv_1=E/t+\lrem$ with $\lrem\in\Oan$ and $\Lv_2\in\Oan$ with $\Lv_2(0)=\alpha_1\alpha_2/E$; and $E/t=E/\tau+O(1)$. By the second form of \eqref{eq:set-Delta},
\[
t^2\Delta_1=\epsilon_1\epsilon_2\,t(\Lv_1-\Lv_2)\,t(\alpha_1\alpha_2-E\Lv_1)=\epsilon_1\epsilon_2\big(E+t(\lrem-\Lv_2)\big)\big(\alpha_1\alpha_2t-E^2-Et\lrem\big),
\]
a product of elements of $\Oan$ with value $\epsilon_1\epsilon_2E(-E^2)$ at $t=0$. $\Delta_2(0)$ is \cref{lem:set-Delta2}(ii).
(ii) At $\cQ=0$, $\Pq(U)=-\alpha_0U-\alpha_2$ has the root $U_+(0)=-\alpha_2/\alpha_0$; for the other root put $V\coloneqq\cQ U$, so that $\cQ\Pq(U)=\wB V^2-(\alpha_0+\cQ(\alpha_1-\alpha_2))V-\alpha_2\cQ$, which at $\cQ=0$ is $V(\wB V-\alpha_0)$, with the simple root $V=\alpha_0/\wB$. Both are simple roots at $\cQ=0$, so $U_+$ and $V_-\coloneqq\cQ U_-$ are power series in $\cQ$. Then $\Lv_+=-\alpha_1\cQ U_+/(1-\cQ U_+)\to0$ and $\Lv_-=-\alpha_1V_-/(1-V_-)\to-\alpha_1(\alpha_0/\wB)/(1-\alpha_0/\wB)=-\alpha_0\alpha_1/(\wB-\alpha_0)=-\alpha_0$. By the second form of \eqref{eq:set-Delta}: $\Delta_+(0)=\epsilon_1\epsilon_2(0+\alpha_0)(\alpha_1\alpha_2)$ and $\Delta_-(0)=\epsilon_1\epsilon_2(-\alpha_0)(\alpha_1\alpha_2+\alpha_0E)=-\alpha_0\wA\wB\epsilon_1\epsilon_2$.
(iii) Dividing \eqref{eq:set-P} by $\cQ$: $\wB U^2-(\alpha_0\sQ+\alpha_1-\alpha_2)U-\alpha_2\sQ=0$, which at $\sQ=0$ has the simple roots $U=(\alpha_1-\alpha_2)/\wB$ and $U=0$ (simple because $\alpha_1\neq\alpha_2$). For the second root put $V\coloneqq\cQ U$: $\wB\sQ V^2-(\alpha_0\sQ+\alpha_1-\alpha_2)V-\alpha_2=0$, with the simple root $V\to-\alpha_2/(\alpha_1-\alpha_2)$ at $\sQ=0$. Hence $U_1(\sQ)$ and $V_2(\sQ)$ are power series in $\sQ$. At $\tau=1$, $\Lv_1+\Lv_2\to E-\alpha_0=\alpha_1+\alpha_2$ and $\Lv_1\Lv_2\to\alpha_1\alpha_2$, so $\{\Lv_1,\Lv_2\}\to\{\alpha_1,\alpha_2\}$; at the point where $U\to(\alpha_1-\alpha_2)/\wB\neq0$, $\Lv=-\alpha_1\cQ U/(1-\cQ U)\to\alpha_1$. The values \eqref{eq:set-Dinf} follow from the second form of \eqref{eq:set-Delta}, using $\alpha_2-E=-\wB$, $\alpha_1-E=-\wA$. The same values follow from the fourth form: $\cW_2(U_1)\to-\wA U/(1-U)^2$ at $U=(\alpha_1-\alpha_2)/\wB$, $1-U=\wA/\wB$, giving $-(\alpha_1-\alpha_2)\wB/\wA$ and $\Delta_1(\infty)=\alpha_1\wA\epsilon_1\epsilon_2\cdot(-(\alpha_1-\alpha_2)\wB/\wA)$; and $\cW_2(U_2)\to-\alpha_1V/(1-V)^2$ at $V=-\alpha_2/(\alpha_1-\alpha_2)$, $1-V=\alpha_1/(\alpha_1-\alpha_2)$, giving $\alpha_2(\alpha_1-\alpha_2)/\alpha_1$ and $\Delta_2(\infty)=\alpha_1\wA\epsilon_1\epsilon_2\cdot\alpha_2(\alpha_1-\alpha_2)/\alpha_1$.
\end{proof}

%-------------------------------------------------------------------------------
\subsection{The rows}
\label{sec:set-rows}

For a finite multiset $\mathsf{v}$ of non-zero weights put
\begin{equation}
\label{eq:set-h}
h_{\mathsf{v}}(z)\coloneqq\exp\Big[-\sum_{k\geq1}\bk_k\,z^{2k-1}\sum_{a}\mathsf{v}_a^{1-2k}\Big],\qquad \bk_k\coloneqq\frac{B_{2k}}{2k(2k-1)},
\end{equation}
with $B_{2k}$ the Bernoulli numbers. Then $h_{\mathsf{u}+\mathsf{v}}=h_{\mathsf{u}}h_{\mathsf{v}}$ (union of multisets), $h_{-\mathsf{v}}=h_{\mathsf{v}}^{-1}$, $h_{(v,-v)}=1$, and $h_{\mathsf{v}}(z)h_{\mathsf{v}}(-z)=1$ (the exponent is odd in $z$); $h_{\mathsf{v}}\in1+z\bbQ(\mathsf{v})\fps z$ is independent of $t$. The series in use are
\[
\hw=h_{(E,-\epsilon_1,-\epsilon_2)},\qquad \hnw=h_{(\alpha_1,\wA,-\epsilon_1,-\epsilon_2)},\qquad \hreg=h_{(\alpha_1,\alpha_2,\wA,\wB,-E,-E)},\qquad h_{\wts_\pm},
\]
with $\bw\coloneqq(E,-\epsilon_1,-\epsilon_2)$ and $\nw\coloneqq(\alpha_1,\wA,-\epsilon_1,-\epsilon_2)$.

\begin{defn}[Rows]
\label{def:set-rows}
\begin{description}
\item[Point-frame identity row.] For $i=1,2$,
\begin{equation}
\label{eq:set-rpt}
\rpt_{0,i}(z;t)\coloneqq\sum_{m\geq0}z^m\,\wick_m\big(\cW_2(U_i),\cW_3(U_i),\ldots\big),
\end{equation}
where $\wick_m$ is the universal polynomial
\[
\wick_m\coloneqq\sum_{n=m}^{3m}\frac{(2n)!\,(-1)^n}{n!\,2^n\,\cW_2^{\,n}}\sum_{p}\ \prod_k\frac{\cW_k^{m_k}}{(k!)^{m_k}\,m_k!},
\]
$p$ running over the partitions of $2n$ into exactly $n-m$ parts $\geq3$, $m_k$ the multiplicity of $k$ in $p$. Equivalently, $\wick_m$ is the coefficient of $z^m$ in
\[
\sum_{n\geq0}\frac{(2n)!}{n!\,2^n}\Big(\frac{-z}{\cW_2}\Big)^{n}\,\coef{v^{2n}}\exp\Big[\sum_{k\geq3}\frac{\cW_kv^k}{k!\,z}\Big],
\]
a term with $j$ factors $\cW_k$ and $v$-degree $2n$ having $z$-order $n-j$ (so $j=n-m$, and $n\leq3m$ because the parts are $\geq3$). Thus $\wick_0=1$, and $\wick_m$ is a polynomial in $\cW_2^{-1},\cW_3,\ldots,\cW_{2m+2}$ with rational coefficients, every monomial $\cW_2^{-n}\prod_i\cW_{k_i}$ having $\sum_ik_i=2n$. The row \eqref{eq:set-rpt} is the normalised formal Laplace expansion of $\re^{\cW/z}\dd U/U$ at the critical point $U_i$ (\cref{prop:row-wickrow} for $i=1$, \cref{thm:lea-row2reg} for $i=2$).
\item[The row with the $\cQ=0$ constant.] $\rZ_{0,i}\coloneqq\hnw(z)\,\rpt_{0,i}(z;t)$. The matrix with these rows is $\RZ$; its meaning (the $R$-matrix of the CohFT, normalised at $\cQ=0$) is \cref{prop:ide-rows,lem:ide-norm,prop:ide-const}. On $\TF$, $\nw=(\xc,1,-\xc,-1)$ and $\hnw=1$.
\item[Frame-B rows.]
\begin{equation}
\label{eq:set-rowB}
\rr_{0,i}(z;t)\coloneqq h_{(\alpha_1,\wA,-E)}(z)\,\rpt_{0,i}(z;t)=\hw(z)^{-1}\hnw(z)\,\rpt_{0,i}(z;t),
\end{equation}
the second equality by \cref{lem:set-hmult}; the matrix with these rows is $\RB=\RZ\hw^{-1}$, and \emph{frame B} means this frame (\cref{prop:ide-frameB}). The frame-B rows depend on $(\alpha_0,\alpha_1,\alpha_2)$ only. In the sequel $\rr_{a,i}$ denotes the frame-B row unless ``point frame'' or $\RZ$ is said.
\item[The $H$-row.]
\begin{equation}
\label{eq:set-Hrow}
\rr_{1,i}(z)\coloneqq(\Lv_i+z\Dop_i)\,\rr_{0,i}(z),\qquad \Dop_i\coloneqq\frac{\dd}{\dd t}-\frac{\Delta_i'}{2\Delta_i},
\end{equation}
that is $\rr_{1,i}[0]=\Lv_i\rr_{0,i}[0]$ and $\rr_{1,i}[m]=\Lv_i\rr_{0,i}[m]+\Dop_i\rr_{0,i}[m-1]$ for $m\geq1$, where $\rr[m]\coloneqq\coef{z^m}\rr$. The rule \eqref{eq:set-Hrow} is linear over $\bbC\fps z$ with $t$-independent coefficients, so it holds in every frame (point frame, $\RZ$, frame B), the scalar series $h$ being independent of $t$. The \emph{covector} of idempotent $i$ is $\vv_i(z)\coloneqq(\rr_{0,i}(z),\rr_{1,i}(z))$, and $\vv_i[p]\coloneqq\coef{z^p}\vv_i$.
\end{description}
\end{defn}

\begin{lem}[The multiset identity]\genusstatus{AI-proposed, unchecked}
\label{lem:set-hmult}
\begin{equation}
\label{eq:set-hmult}
\hnw(z)\,h_{(E,-\alpha_1,-\wA)}(z)=\hw(z),\qquad\text{equivalently}\qquad \hw^{-1}\hnw\,h_{(E,-\alpha_1,-\wA)}=1,\qquad \hw^{-1}\hnw=h_{(\alpha_1,\wA,-E)} .
\end{equation}
Moreover, for two multisets $\bw',\bw$ of non-zero weights, $h_{\bw'}=h_{\bw}$ if and only if $\bw'$ and $\bw$ agree up to pairs $(v,-v)$. In particular $\bw=(E,-\epsilon_1,-\epsilon_2)$ is, up to such pairs, the unique class with $h_{\bw}=\hnw h_{(E,-\alpha_1,-\wA)}$.
\end{lem}

\begin{proof}
As multisets, $\nw+(E,-\alpha_1,-\wA)=(\alpha_1,\wA,-\epsilon_1,-\epsilon_2,E,-\alpha_1,-\wA)=(E,-\epsilon_1,-\epsilon_2)+(\alpha_1,-\alpha_1)+(\wA,-\wA)$, and $h_{(v,-v)}=1$; this is \eqref{eq:set-hmult}. For the uniqueness: $h_{\mathsf{v}}$ depends on $\mathsf{v}$ only through the odd power sums $\sum_a\mathsf{v}_a^{1-2k}$, $k\geq1$, of the reciprocals $1/\mathsf{v}_a$, and $h_{\mathsf{v}}=1$ if and only if all of them vanish (the $\bk_k$ are non-zero). The even power sums of the multiset $\{1/\mathsf{v}_a\}$ and of its negative $\{-1/\mathsf{v}_a\}$ agree, and the odd ones are negatives of each other; so all odd power sums vanish if and only if the two multisets have the same power sums, if and only if (Newton's identities, in characteristic $0$: the power sums of a finite multiset determine it) $\{1/\mathsf{v}_a\}$ equals its negative, if and only if $\mathsf{v}$ is a union of pairs $(v,-v)$. Applied to $\mathsf{v}=\bw'+(-\bw)$: $h_{\bw'}=h_{\bw}$ if and only if $h_{\bw'}h_{\bw}^{-1}=h_{\bw'+(-\bw)}=1$, if and only if $\bw'+(-\bw)$ is a union of pairs, that is, $\bw'$ and $\bw$ agree up to pairs.
\end{proof}

The significance of the class $\bw$ is that it makes the frame-B constant at $t=0$ the geometric series: $\rr_{0,1}(z;0)=E/(E+z)$ (\cref{cor:row-F2}), because the point-frame row at $t=0$ is $h_{(E,-\alpha_1,-\wA)}E/(E+z)$ (\cref{prop:row-beta}). The same scalar $\hnw$ multiplies both idempotents (\cref{lem:ide-norm}), so the frame-B rows of the regular idempotent carry the non-trivial constant $h_{(\alpha_1,\wA,-E)}\rpt_{0,2}(z;0)$ (\cref{prop:lea-regrows}). On $\TF$, $\hnw=1$ and \eqref{eq:set-hmult} is $\hw=h_{(\cc,-\xc,-1)}=h_{(\xc,1,-\cc)}^{-1}$.

%-------------------------------------------------------------------------------
\subsection{The factors}
\label{sec:set-factors}

\begin{defn}[Dilaton and edge factors, legs]
\label{def:set-factors}
All rows are the frame-B rows of \cref{def:set-rows}.
\begin{description}
\item[Dilaton factors.] $\Dil_i(d)\coloneqq(-1)^d\,\rr_{0,i}[d-1]$ for $d\geq1$. The \emph{dilaton series} at idempotent $i$ is
\[
\qt^{(i)}(\psi)\coloneqq\sum_{d\geq1}\Dil_i(d;t)\,\psi^d=\sum_{m\geq0}(-1)^{m+1}\rr_{0,i}[m]\,\psi^{m+1}=-\psi\,\rr_{0,i}(-\psi;t),\qquad \qt\coloneqq\qt^{(1)} .
\]
Since $\rr_{0,i}[0]=1$ ($\wick_0=1$ and $h[0]=1$), $\Dil_i(1)=-1$ at every $t$: the term $d=1$ is the dilaton shift, $\Dil_i(1)\psi=-\psi$.
\item[The pairing of the rows.]
\begin{equation}
\label{eq:set-Pij}
\Pij_{ij}(z,w)\coloneqq-\langle\vv_i(-z),\vv_j(-w)\rangle=\sum_{p,q}\Pij^{ij}_{pq}z^pw^q,\qquad \Pij^{ij}_{pq}=-(-1)^{p+q}\langle\vv_i[p],\vv_j[q]\rangle,
\end{equation}
with the pairing \eqref{eq:set-pairing}.
\item[Edge factors.]
\begin{equation}
\label{eq:set-edge}
\Edge_{ij}(z,w)\coloneqq\frac{\Pij_{ij}(z,w)-\Pij_{ij}(-w,w)}{z+w}=\sum_{m,n\geq0}\Edge_{ij}(m,n)\,z^mw^n ;
\end{equation}
the numerator vanishes on $z=-w$, so the quotient lies in $\bbC\fps{z,w}$ coefficientwise in $t$ (\cref{lem:set-edgegen}). In leaf language $\Edge_{11}(\psi,\psi')$ means $\Edge_{11}(z=\psi,w=\psi')$, without change of sign.
\item[Legs.] $f_{j,b}(\psi)\coloneqq\sum_a\Edge_{1j}(a,b)\psi^a=\coef{w^b}\Edge_{1j}(\psi,w)$, for $j\in\{1,2\}$ and $b\geq0$.
\end{description}
\end{defn}

\begin{lem}[Edge factors as generating functions]\genusstatus{AI-proposed, unchecked}
\label{lem:set-edgegen}
For any two rows $\vv_i,\vv_j$ (no symplecticity is assumed), the right side of \eqref{eq:set-edge} lies in $\bbC\fps{z,w}$ coefficientwise in $t$, and its coefficients are
\begin{equation}
\label{eq:set-EdN}
\Edge_{ij}(m,n)=\sum_{k=0}^{n}(-1)^k\,\Pij^{ij}_{m+1+k,\,n-k} .
\end{equation}
\end{lem}

\begin{proof}
For $p\geq1$, $(z^p-(-w)^p)/(z+w)=\sum_{l=0}^{p-1}z^l(-w)^{p-1-l}$; for $p=0$ the quotient is $0$. So the right side of \eqref{eq:set-edge} equals $\sum_{p,q}\Pij^{ij}_{pq}w^q\sum_{l=0}^{p-1}z^l(-w)^{p-1-l}$, a power series in $z,w$. Its coefficient of $z^mw^n$ collects the terms with $l=m$, $p=m+1+k$ and $q=n-k$, where $k=p-1-m\in[0,n]$, each with the sign $(-1)^k$; this is \eqref{eq:set-EdN}.
\end{proof}

\begin{rmk}
\label{rmk:set-edgesym}
Once the symplectic relations $\langle\vv_i(z),\vv_j(-z)\rangle=\delta_{ij}\Delta_i$ are known (\cref{thm:lea-qde}(ii)), $\Pij_{ij}(-w,w)=-\delta_{ij}\Delta_i$ and \eqref{eq:set-edge} reads $\Edge_{ij}(z,w)=\big[\delta_{ij}\Delta_i-\langle\vv_i(-z),\vv_j(-w)\rangle\big]/(z+w)$, which is symmetric under $(i,z)\leftrightarrow(j,w)$: $\Edge_{ij}(m,n)=\Edge_{ji}(n,m)$. Until then the order of the two ends of an edge is part of the data in \cref{def:set-graphsum}.
\end{rmk}

\begin{defn}[$\rho$, the leaves, the leaf spaces]
\label{def:set-rho}
By \cref{lem:set-limits}(i), $t^2\Delta_1/(-\epsilon_1\epsilon_2E^3)$ is a unit of $\Oan$ with value $1$ at $t=0$; let
\[
\rho(t)\coloneqq\Big(\frac{\Delta_1t^2}{-\epsilon_1\epsilon_2E^3}\Big)^{-1/2},\qquad\text{the branch with }\rho(0)=1 .
\]
The \emph{leaves} are $\lf_j(\psi)\coloneqq1/(E-j\psi)$ for $j\geq1$ (and $\lf_0=1/E$, which appears only in \cref{rmk:reg-L0}); in the row variable $z=-\psi$, $\lf_j=1/(E+jz)$. The \emph{leaf spaces} are $\Lsp_m\coloneqq\mathrm{span}(\lf_1,\ldots,\lf_m)$ and $\Lsp_m^+\coloneqq(\text{constants})+\Lsp_m$. Further
\[
\qo\coloneqq E-E^2\lf_1=-\frac{E\psi}{E-\psi},\qquad \Tdil(\psi)\coloneqq\qo+\psi=-\frac{\psi^2}{E-\psi}.
\]
\end{defn}

\begin{rmk}[Scope]
\label{rmk:set-scope}
The objects of the theorems of \cref{sec:row,sec:lea} are the formal series in $z$ (and $w$) defined by \eqref{eq:set-rpt}--\eqref{eq:set-edge} for all $m,n$, with coefficients that are functions of $t$ (algebraic functions of $\tau$, by \cref{lem:set-rings}). Each statement is an identity of formal power series in $\psi$ (and $\psi'$, $w$), valid at every $\psi$-order and every $t$-order.
\end{rmk}

%-------------------------------------------------------------------------------
\subsection{Coloured stable graphs and the graph sum}
\label{sec:set-graphsum}

For a stable curve of genus $h$, with $\lambda_i=c_i(\Hdg)$ the Chern classes of the Hodge bundle, let $\Lamv(u)\coloneqq\sum_i(-1)^i\lambda_iu^{h-i}$. For a multiset $\mathsf{v}$ of non-zero weights the \emph{normalised Hodge class} is $\Hg_{\mathsf{v}}\coloneqq\prod_{a\in\mathsf{v}}\Lamv(\mathsf{v}_a)/\mathsf{v}_a^{h}$. The vertex class of the note is
\begin{equation}
\label{eq:set-Hg}
\Hg_{\bw}=\prod_{a=1}^{3}\frac{\Lamv(\bw_a)}{\bw_a^{h}},\qquad \bw=(E,-\epsilon_1,-\epsilon_2),
\end{equation}
the same at both colours and independent of $t$. It is a pull-back under the forgetful maps, so the string and dilaton equations hold for it when $h\geq2$.

\begin{defn}[Coloured stable graphs; the decoration sum; the graph sum]
\label{def:set-graphsum}
Fix $g\geq2$. Let $\Gamma$ run over the connected stable graphs of genus $g$ without legs, with vertex set $\Ver(\Gamma)$, edge set $\Edg(\Gamma)$, vertex genera $g_v$ and $n_v$ half-edges at $v$ (stability: $2g_v-2+n_v>0$; $g=h^1(\Gamma)+\sum_vg_v$). A \emph{colouring} is a map $\col:\Ver(\Gamma)\to\{1,2\}$, $v\mapsto\col_v$; colour $1$ is singular, colour $2$ regular. A \emph{decoration} assigns an exponent $a_h\geq0$ to each half-edge $h$ (we write $a_e,b_e$ for the exponents at the two half-edges of an edge $e$, in a chosen order of its ends $v(e),v'(e)$) and to each vertex $v$ an integer $m_v\geq0$ with an ordered tuple $(d_{v,1},\ldots,d_{v,m_v})$ of \emph{dilaton degrees} $d_{v,k}\geq2$, with weight $1/m_v!$. The \emph{decoration sum} of $(\Gamma,\col)$ is
\begin{multline}
\label{eq:set-graphsum}
\GCont(\Gamma,\col)\coloneqq\sum_{\mathrm{dec}}\ \prod_{v\in\Ver(\Gamma)}\Big[\Delta_{\col_v}^{\,g_v-1}\,\frac{1}{m_v!}\int_{\Mbar_{g_v,n_v+m_v}}\Hg_{\bw}\prod_{h\text{ at }v}\psi_h^{a_h}\prod_{k=1}^{m_v}\psi_{n_v+k}^{d_{v,k}}\ \prod_{k=1}^{m_v}\Dil_{\col_v}(d_{v,k})\Big]\\
\times\prod_{e\in\Edg(\Gamma)}\Edge_{\col_{v(e)}\col_{v'(e)}}(a_e,b_e),
\end{multline}
and the \emph{frame-B Hodge-vertex graph sum} is
\begin{equation}
\label{eq:set-Gsum}
\Gsum_g\coloneqq\sum_{\Gamma}\frac{1}{|\Aut\Gamma|}\sum_{\col:\Ver(\Gamma)\to\{1,2\}}\GCont(\Gamma,\col).
\end{equation}
\end{defn}

The decoration sum is finite: the integral vanishes unless $\sum_ha_h+\sum_kd_{v,k}\leq3g_v-3+n_v+m_v$ at every vertex, and as $d_{v,k}\geq2$ this forces $m_v\leq3g_v-3+n_v$; so in genus $g$ only finitely many $z$-coefficients of the rows enter, and $\GCont(\Gamma,\col)$ is a finite sum of products of the factors of \cref{prop:set-idem,def:set-factors}. This is what gives ``regular at the conifold point after summing over decorations'' its meaning in \cref{sec:reg}: the Laurent expansion of $\GCont(\Gamma,\col)$ in $t$ (\cref{lem:set-rings}) has no negative powers. The independence of \eqref{eq:set-graphsum} of the chosen order of the ends of each edge is \cref{rmk:set-edgesym}. Equation \eqref{eq:set-graphsum} is a definition; the regularity statements of \cref{sec:reg} are statements about it, and its identification with the Gromov--Witten series $\sF_g$ is \cref{prop:ide-main}, used only in \cref{sec:rig}.

\emph{The dilaton shift.} The decorations carry dilaton degrees $d\geq2$ only. Since $\Dil_i(1)\psi=-\psi$, the sum over dilaton decorations at a vertex of genus $h$ and colour $i$ with half-edge series $f_1,\ldots,f_n$ equals $\Delta_i^{h-1}\sum_m\frac{1}{m!}\int\Hg_{\bw}\prod_rf_r(\psi_r)\prod_{k\leq m}(\qt^{(i)}+\psi)(\psi_{n+k})$, that is $\Delta_i^{h-1}\dd^n\Fpot_h(\qt^{(i)})[f_1,\ldots,f_n]$ in the dilaton-shifted form of \cref{def:set-fpot}; this identification is what makes the homogeneity step of \cref{thm:reg-main} apply.

%-------------------------------------------------------------------------------
\subsection{Rings, general principles, leaf identities}
\label{sec:set-rings}

\begin{lem}[Rings and principles]\genusstatus{AI-proposed, unchecked}
\label{lem:set-rings}
\begin{description}
\item[Rings.] $\Oan\coloneqq\bbC\{t\}=\bbC\{\tau\}$ (convergent power series; $\tau(t)$ is biholomorphic near $0$ and $\tau/t$, $t/\tau$ are units of $\Oan$); $\Man\coloneqq\bbC(\{t\})=\bbC(\{\tau\})$, the field of meromorphic germs; $\Oan(\Ddisc)$, $\Oan(\Ddisc')$ the functions analytic on $\Ddisc,\Ddisc'$. $\Cz=\bbC\lser z$ (formal Laurent series), in which $\bbC(z)$ is embedded by Laurent expansion at $z=0$. $\Rz\coloneqq\Cz\lser\tau$ with the $\Cz$-linear derivation $\dd/\dd t=(1-\tau)\dd/\dd\tau$; $\Rzo\coloneqq\bbC\fps z\lser\tau=\bbC\fps{z,\tau}[1/\tau]$, a subring of $\Rz$ and (expanding in $z$) of $\bbC\lser\tau\fps z$, compatibly with products and $\dd/\dd t$. $\iota:\Oan\fps z\to\bbC\fps{z,\tau}$, ``Taylor-expand every $z$-coefficient'', is an injective ring homomorphism commuting with $\dd/\dd\tau$ and with multiplication by elements of $\Oan$. $\mathrm{refl}:z\mapsto-z$ (coefficientwise) commutes with $\dd/\dd t$; $f(-z)\coloneqq\mathrm{refl}(f)$. $\Man\lser z$ denotes formal Laurent series in $z$, bounded below, with coefficients in $\Man$ ($\dd/\dd t$ coefficientwise).
\item[Transfer.] An identity in $\Rz$ between elements of $\Rzo$ is an identity of formal double series: reading off $z^m$ gives identities of formal Laurent series in $\tau$, and, when these are expansions of meromorphic germs (as for all rows), identities of meromorphic functions of $t$ near $0$. ``To all orders in $z$'' means this; in $\Man\lser z$ it means ``coefficientwise in $z$, as identities in $\Man$''.
\item[(P) Complex and indeterminate weights.] Every object of this note built from the genus-zero data ($\tau$, $\sast$, $\delta$, $\dphi$, $\alg$, $\rho$, $\Lv_i$, $\Delta_i$, $\Thf$, the rows, $h_{(\alpha_1,\wA,-E)}$, the dilaton and edge factors) is obtained from the weights and $\tau$ by field operations, $\dd/\dd t$, and the expansions of the square roots $\delta$, $\alg$ (units at $\tau=0$) and of logarithms of units. At $\tau=0$ the only denominators are $\alpha_1,\alpha_2,\wA,\wB,E,\epsilon_1,\epsilon_2$: $\sast(0)=\wA/E$, $1-\sast(0)=\alpha_1/E$, $U_2(0)=-\alpha_2/\wB$, $1-U_2(0)=E/\wB$, $\phi_0''(\wA/E)=-E^3/(\alpha_1\wA)$, $\Delta_2(0)=\epsilon_1\epsilon_2\alpha_1\alpha_2\wA\wB/E$, and the Gauss denominators $E+jz$ of \cref{sec:row}. Hence the $t$-, $z$- and $\psi$-coefficients are rational functions of the weights with rational coefficients, regular on (ND); an identity between such coefficients proved at every complex weight point of a non-empty open subset of the weight space (more generally, of a set on which no non-zero rational function of the weights vanishes identically) holds in $\Kal$ (or $\KCY$), and conversely an identity or a membership proved over $\Kal$ specialises to every weight point with (ND). \Cref{sec:lea} is written over $\Kal$ where it says so.
\item[Leaf spaces.] $\Lsp_m$ has dimension $m$ (the poles $E/j$, $1\leq j\leq m$, are distinct because $E\neq0$). A rational function $f$ of $\psi$ lies in $\Lsp_m^+$ if and only if its poles are simple, lie in $\{E/j:1\leq j\leq m\}$, and $f$ is bounded at $\infty$; it lies in $\Lsp_m$ if and only if moreover $f(\psi)\to0$ as $\psi\to\infty$. A rational function regular at $0$ is identified with its Taylor series at $0$ (an injective linear map, compatible with products and with the substitution $z\mapsto-\psi$); thus $\lf_j(\psi)=\sum_{n\geq0}j^n\psi^n/E^{n+1}$.
\item[(T) $\tau$-orders versus $t$-orders.] For a vector space $V$ and $f=\sum_nf_n\tau^n\in V\fps\tau$, $V\fps\tau=V\fps t$ and
\[
\coef{t^k}f=\sum_{n=0}^{k}u_{k,n}f_n,\qquad u_{k,n}\coloneqq\coef{t^k}\tau^n\in\bbQ,\quad u_{k,k}=1,\quad u_{k,n}=0\ \text{for }n>k,
\]
with the inverse unitriangular relation from $t=-\log(1-\tau)$. So for every increasing chain $V_0\subset V_1\subset\cdots$ of subspaces, ($\coef{\tau^n}f\in V_n$ for $n\leq k$) if and only if ($\coef{t^n}f\in V_n$ for $n\leq k$). Multiplication by a series in $t$ that does not involve the auxiliary variables also preserves such a filtration: $\coef{t^k}(af)=\sum_ma_m\coef{t^{k-m}}f$.
\end{description}
\end{lem}

\begin{proof}
Everything is a definition or standard, except (T) and the characterisation of $\Lsp_m^+$. (T): $\tau=t-t^2/2+\cdots$, so $\tau^n=t^n+O(t^{n+1})$, which gives $u_{k,k}=1$ and $u_{k,n}=0$ for $n>k$; the relation is therefore unitriangular and invertible. The filtration statement follows: $\coef{t^k}f=f_k+\sum_{n<k}u_{k,n}f_n$ with $f_n\in V_n\subset V_k$, and conversely by the inverse relation. Leaf spaces: an element of $\Lsp_m^+$ clearly has the stated properties; conversely, by partial fractions (PF) below, a rational function with simple poles in $\{E/j\}$ and bounded at $\infty$ is a constant plus a combination of the $1/(\psi-E/j)=-j\lf_j$, and the constant is the value at $\infty$.
\end{proof}

\begin{lem}[Leaf identities, partial fractions]\genusstatus{AI-proposed, unchecked}
\label{lem:set-leafid}
For $j,m\geq1$:
\begin{description}
\item[(LAM1)] $\coef{\psi^n}\lf_j=j^n/E^{n+1}$;
\item[(LAM2)] $\psi\lf_j=(E\lf_j-1)/j$, that is $-\psi\lf_j=(1-E\lf_j)/j$;
\item[(LAM3)] $(E-m\psi)\lf_j=m/j+E(1-m/j)\lf_j$; in particular $(E-j\psi)\lf_j=1$;
\item[(LAM4)] for $f=\sum_{j=1}^{M}f_j\lf_j$ one has $f=-\frac1\psi\sum_jf_j/j+O(\psi^{-2})$ at $\psi=\infty$, so $f=O(\psi^{-2})$ if and only if $\sum_jf_j/j=0$;
\item[(LAM5)] $\lf_i\lf_j=\dfrac{i\lf_i-j\lf_j}{E(i-j)}$ for $i\neq j$;
\item[(LAM6)] $\lf_j(\psi)-\lf_j(-w)=j(\psi+w)\lf_j(\psi)\lf_j(-w)$.
\item[(PF)] Let $S$ be a finite set of points and $f$ a rational function whose poles all lie in $S$, are of order $\leq1$, and which is bounded at $\infty$. Then $f=f(\infty)+\sum_{p\in S}\Res_{z=p}(f)/(z-p)$. For $S=\{-E/j:1\leq j\leq M\}$ (distinct and non-zero since $E\neq0$),
\[
f=f(\infty)+\sum_{j=1}^{M}\frac{\beta_j}{E+jz},\qquad \beta_j=j\Res_{z=-E/j}f=\lim_{z\to-E/j}(E+jz)f(z).
\]
\end{description}
\end{lem}

\begin{proof}
(LAM1) is the geometric series $1/(E-j\psi)=E^{-1}\sum_n(j\psi/E)^n$. $(E-j\psi)\lf_j=1$ gives (LAM2). (LAM3): $(E-m\psi)\lf_j=1+(j-m)\psi\lf_j=1+\frac{j-m}{j}(E\lf_j-1)$. (LAM4): $\lf_j=-1/(j\psi)-E/(j^2\psi^2)-\cdots$ at $\psi=\infty$. (LAM5): $i\lf_i-j\lf_j=\dfrac{i(E-j\psi)-j(E-i\psi)}{(E-i\psi)(E-j\psi)}=(i-j)E\lf_i\lf_j$. (LAM6): $\dfrac{1}{E-j\psi}-\dfrac{1}{E+jw}=\dfrac{j(\psi+w)}{(E-j\psi)(E+jw)}$. (PF): the difference $h\coloneqq f-f(\infty)-\sum_p\Res_p(f)/(z-p)$ has no pole and $h(\infty)=0$; a rational function without poles is a polynomial, and a polynomial with limit $0$ at $\infty$ is $0$. For the second form, $1/(z+E/j)=j/(E+jz)$.
\end{proof}

%-------------------------------------------------------------------------------
\subsection{The Hodge side}
\label{sec:set-hodge}

Fix $h\geq2$. Let $\RRing\coloneqq\RRing_0\fps t$ with $\RRing_0$ a $\bbQ$-algebra in which $\epsilon_1,\epsilon_2,E$ are invertible (for instance $\bbC$ or $\KCY$). The class $\Hg_{\bw}$ is that of \eqref{eq:set-Hg}, and $\Tdil$, $\qo$ are those of \cref{def:set-rho}, so that $\qo=\Tdil-\psi$.

\begin{defn}
\label{def:set-fpot}
\begin{description}
\item[Base-point derivatives.] For $\mathsf{b}_1,\ldots,\mathsf{b}_s\in\RRing\fps\psi$, put
\[
\dbase^s(\mathsf{b}_1,\ldots,\mathsf{b}_s)\coloneqq\sum_{m\geq0}\frac{1}{m!}\int_{\Mbar_{h,s+m}}\Hg_{\bw}\prod_{i=1}^{s}\mathsf{b}_i(\psi_i)\prod_{k=1}^{m}\Tdil(\psi_{s+k}).
\]
This is a finite sum, since $\Tdil=O(\psi^2)$ forces $m\leq3h-3+s$, and each power series is cut off by the dimension.
\item[The potential on the neighbourhood $\Dom$ of $\qo$.] Let $\Dom\coloneqq\qo+t\RRing\fps\psi$. For $\mathbf{q}=\qo+\upsilon\in\Dom$ and $f_1,\ldots,f_n\in\RRing\fps\psi$, define
\[
\dd^n\Fpot_h(\mathbf{q})[f_1,\ldots,f_n]\coloneqq\sum_{r\geq0}\frac{1}{r!}\,\dbase^{n+r}(f_1,\ldots,f_n,\underbrace{\upsilon,\ldots,\upsilon}_{r}),\qquad \Fpot_h(\mathbf{q})\coloneqq\dd^0\Fpot_h(\mathbf{q}),
\]
which converges $t$-adically because $\upsilon=O(t)$. Here $\upsilon$ may have a term linear in $\psi$: $\Dom$ contains $\rho\qt$ (below), whose $\psi$-coefficient is $-\rho$; so $\mathbf{q}$ need not be of the form $-\psi+O(\psi^2)$. This is why the domain has to be stated.
\item[The dilaton-shifted form.] If $\mathbf{q}=-\psi+O(\psi^2)$, put
\[
\Fpot_h(\mathbf{q})\coloneqq\sum_{m\geq0}\frac{1}{m!}\int_{\Mbar_{h,m}}\Hg_{\bw}\prod_{i=1}^{m}(\mathbf{q}+\psi)(\psi_i),
\]
a finite sum, and similarly for $\dd^n\Fpot_h(\mathbf{q})[f_1,\ldots,f_n]$. On the overlap with $\Dom$ the two definitions agree: expand $\prod(\Tdil+\upsilon)$ binomially, and the factor $\frac{1}{m!}\binom mr$ reindexes to $\frac{1}{r!}\frac{1}{(m-r)!}$. This is the form in which the decorations of \eqref{eq:set-graphsum} appear (the dilaton shift paragraph of \cref{sec:set-graphsum}): the vertex factor at a vertex of genus $h\geq2$ and colour $i$ with half-edge series $f_1,\ldots,f_n$ is $\Delta_i^{h-1}\dd^n\Fpot_h(\qt^{(i)})[f_1,\ldots,f_n]$, multilinear in the slots, applied term by term to loops and to edges between two vertices.
\item[$\qt$ and $\Dom$.] At idempotent $1$, $\qt\in\Dom$: every $\Dil_1(d)$ lies in $\RRing$ and $\qt|_{t=0}=\qo$ (\cref{cor:row-F2}, \cref{thm:lea-L1}(iv)); and $\rho\in1+t\RRing$ by (F1) of \cref{thm:lea-main}, so $\rho\qt\in\Dom$ as well. At idempotent $2$, $\qt^{(2)}|_{t=0}=-\psi\,\hreg(-\psi)\neq\qo$ (\cref{prop:lea-regrows}), so $\qt^{(2)}\notin\Dom$; only the dilaton-shifted form is used there, and nothing is rescaled.
\item[The numbers $\Xiz_h$.] For a multiset $\msU=\{j_1,\ldots,j_r\}$ of positive integers,
\[
\Xiz_h(\msU)\coloneqq\dd^r\Fpot_h(\qo)[\lf_{j_1},\ldots,\lf_{j_r}]=\dbase^r(\lf_{j_1},\ldots,\lf_{j_r})=\sum_{m\geq0}\frac{1}{m!}\int_{\Mbar_{h,r+m}}\Hg_{\bw}\prod_{i=1}^{r}\lf_{j_i}(\psi_i)\prod_{k=1}^{m}\Tdil(\psi_{r+k}),
\]
the leaf-scale-$E$ integrals with the class $\Hg_{\bw}$. Further $J_{\msU}\coloneqq\sum_ij_i$, $\dg(\msU)\coloneqq\sum_i(j_i-1)$, and $\ZZ[h,n]\coloneqq\Xiz_h(1^n)$, so that $\ZZ[h,0]=\Fpot_h(\qo)$.
\end{description}
\end{defn}

\begin{defn}[The statement $\Mirr(h)$]
\label{def:set-mirr}
$\Mirr(h)$ is the statement: $\Xiz_h(\msM)=0$ for every \emph{non-empty} multiset $\msM$ of integers $\geq2$ with $\dg(\msM)\leq2h-3$.
\end{defn}

The empty multiset is excluded because $\Xiz_h(\emptyset)=\ZZ[h,0]$ is the value of the pole (\cref{thm:reg-VI}). Parts equal to $1$ are not in $\Mirr$: they are removed by the string identity (\cref{lem:reg-string}), and part $0$ never occurs (\cref{rmk:reg-nol0}). The statement $\Mirr(h)$ for the class $\Hg_{\bw}$, at the generic point of $\TCY$ and at every weight point with $\epsilon_1\epsilon_2E\neq0$, is the cited \cref{thm:reg-L}.

%-------------------------------------------------------------------------------
\subsection{The main statements}
\label{sec:set-statements}

The main theorems of the note are stated in \cref{sec:int-statements}, in the notation fixed in this section, and are proved later: Theorem~A, the leaf structure at generic weights (\cref{thm:int-leaf}), is \cref{thm:lea-main}; Theorem~B, the identification of the graph sum with the potential (\cref{thm:int-ident}), is \cref{prop:ide-main}; Theorem~C, the regularity of the coloured graphs with an edge and the polar part of the vertex graph (\cref{thm:int-reg}), is \cref{cor:reg-main}; Theorem~D, rigidity (\cref{thm:int-rig}), is \cref{thm:rig-main}; and the rigidity conjecture of \cite{BS} for the resolved conifold (\cref{cor:int-conj}) is \cref{cor:rig-conj}. Every object occurring in their statements has now been defined: the tori and the condition (ND) (\cref{def:set-ND}), the invariants and the series $\sF_g$, $\Pg$, $\beta_n$, $\Li_{-n}$ (\cref{def:set-inv}), the vertex weights $\Delta_i$ (\cref{prop:set-idem}), the frame-B rows and factors (\cref{def:set-rows,def:set-factors}), $\rho$, the leaves $\lf_j$ and the leaf spaces $\Lsp_m$ (\cref{def:set-rho}), the coloured graph sum $\Gsum_g$ (\cref{def:set-graphsum}) and the Hodge integral $\ZZ[g,0]$ (\cref{def:set-fpot}).

% s03_rows. Section file; new macros at the top; label prefix: row.
% Local macros (used in this section only):
\providecommand{\kE}{k_E}                              % s(1-s)(1-tau s)^{-2}
\providecommand{\kph}{\kappa_{\mathrm{ph}}}            % constant of the phase, used nowhere
\providecommand{\Lt}{\Lnorm_{\tau}}                    % \Lnorm[\phi_\tau,(1-\tau s)^{-1};s^*]
\providecommand{\Ngauss}{\Nlap_{\mathrm{G}}}            % the conjugated expansion that solves Gauss's equation
\providecommand{\Spht}{\mathsf{S}_{\phi,\tau}}              % d/dtau of the eikonal phase
\providecommand{\Sphtt}{\mathsf{S}_{\phi,\tau\tau}}         % second tau-derivative of the eikonal phase
\providecommand{\Abetam}[1]{\mathsf{A}_{\beta,#1}}         % z-coefficients of \Abeta
\providecommand{\Bfn}{\mathrm{B}}                      % Euler's Beta function
\section{The singular row: Euler integral, Gauss equation and the Beta constant}
\label{sec:row}

This section computes the frame-B identity row $\rr_{0,1}(z;t)$ at the singular idempotent in closed form, as a product of Gauss's hypergeometric series, an exponential of the critical value of an Euler phase and an algebraic unit of $\tau$, and derives from it the regularity of its $t$-coefficients and the value $\rr_{0,1}(z;0)=E/(E+z)$. Throughout, the weights satisfy the non-degeneracy condition (ND) of \cref{def:set-ND}, either as complex numbers or as indeterminates; by the passage (P) of \cref{lem:set-rings} an identity of the kind proved here holds at indeterminate weights if and only if it holds at every complex weight point with (ND). We use the Euler variable $s$, the critical polynomial $\Ns(s)=\wB\tau s^2-(E+\alpha_0\tau)s+\wA$ with its roots $\sast$, $\stil$, the algebraic functions $\delta$, $\sast$, $\stil$ and the discs $\Ddisc$, $\Ddisc'=\Ddisc\cap\{|\tau|<1\}$, $\Ddisc''=\Ddisc'\smallsetminus\{0\}$ of \cref{def:set-alg}, the identities of \cref{lem:set-N}, the eigenvalues $\Lv_1,\Lv_2$ of \cref{lem:set-eig}, and the point-frame row $\rpt_{0,1}$ of \cref{def:set-rows}. Each statement is either \emph{formal} (an identity of formal series, valid at every weight point with (ND)) or \emph{analytic} (a statement about convergent integrals, with explicit hypotheses); we say which. The only analytic statements are \cref{lem:row-laplace,lem:row-betaphase,prop:row-beta}, and the only analytic inputs are Laplace's method, which is proved in full in \cref{lem:row-laplace}, and Stirling's series with remainder \cite[5.11.1 and Section~5.11(ii)]{DLMF}.

On functions of $t$ a prime denotes $\dd/\dd t=\thQ$, and $f_\tau$ denotes $\dd f/\dd\tau$, so that $f'=(1-\tau)f_\tau$ (recall $\tau=1-\re^{-t}$, $\dd\tau/\dd t=1-\tau$). On functions written in the Euler variable $s$ (phases, amplitudes) a prime or a superscript $(k)$ means $\dd/\dd s$, and $\de_\tau$ is the partial derivative in $\tau$ at fixed $s$. $\coef{t^k}$, $\coef{z^m}$, $\coef{\tau^n}$ denote coefficients, taken coefficientwise in the other variables.

%-------------------------------------------------------------------------------
\subsection{The Euler coordinate}
\label{sec:row-euler}

Recall the superpotential $\cW(U)=\alpha_1\log(1-\cQ U)-\alpha_2\log U+\wA\log(1-U)$ and $\cW_k=\thU^k\cW$ (\cref{lem:set-crit}). The \emph{Euler phase} is
\begin{equation}
\label{eq:row-phase}
\phi_\tau(s)\coloneqq\wA\log s+\alpha_1\log(1-s)-\alpha_2\log(1-\tau s),
\end{equation}
so that $\phi_0(s)=\wA\log s+\alpha_1\log(1-s)$, and, by \cref{lem:set-N},
\begin{equation}
\label{eq:row-phaseder}
\phi_\tau'(s)=\frac{\wA}{s}-\frac{\alpha_1}{1-s}+\frac{\alpha_2\tau}{1-\tau s}=\frac{\Ns(s)}{s(1-s)(1-\tau s)},
\end{equation}
because $\wA(1-s)(1-\tau s)-\alpha_1s(1-\tau s)+\alpha_2\tau s(1-s)=\wA-(\wA+\alpha_1)s-\wA\tau s+\alpha_2\tau s+(\wA+\alpha_1-\alpha_2)\tau s^2=\Ns(s)$, using $\wA+\alpha_1=E$, $\alpha_2-\wA=-\alpha_0$ and $\wA+\alpha_1-\alpha_2=\wB$. The critical points of $\phi_\tau$ in $\bbC\smallsetminus\{0,1,1/\tau\}$ are therefore the roots $\sast,\stil$ of $\Ns$; on $\Ddisc'$ the root $\sast$ is not in $\{0,1,1/\tau\}$, and on $\Ddisc''$ neither is $\stil$ (\cref{lem:set-N}). The critical value at the singular root, normalised at $\tau=0$, is the germ $\dphi$ of \cref{def:set-alg},
\begin{equation}
\label{eq:row-dphi}
\dphi(\tau)=\phi_\tau(\sast(\tau))-\phi_0(\wA/E),\qquad \phi_0(\wA/E)=\wA\log\frac{\wA}{E}+\alpha_1\log\frac{\alpha_1}{E},
\end{equation}
with principal branches near $(\tau,s)=(0,\wA/E)$; it lies in $\tau\Oan$, and by the envelope formula of \cref{lem:set-N},
\begin{equation}
\label{eq:row-dphider}
\dphi_\tau=\frac{\alpha_2\sast}{1-\tau\sast}=\frac{\wA-\wB\sast}{1-\tau},\qquad \dphi'=\wA-\wB\sast=\alpha_2-\Lv_2 .
\end{equation}
(The last equality is $\Lv_2=\wB\sast-\alpha_0$ of \cref{lem:set-eig} with $\wA-\alpha_2=\alpha_0$.)

\begin{lem}[The Euler coordinate; formal]\genusstatus{AI-proposed, unchecked}
\label{lem:row-euler}
For $\tau\neq0$ put $U=1-\tau s$. Then
\[
1-U=\tau s,\qquad 1-\cQ U=(1-\cQ)(1-s)=-\cQ\tau(1-s),\qquad \frac{\dd U}{U}=-\frac{\tau\,\dd s}{1-\tau s},
\]
and on every simply connected domain of the $s$-plane avoiding $s=0,1,1/\tau$, with compatible branches of the logarithms,
\begin{align}
\label{eq:row-euler1}
\cW(1-\tau s)&=\alpha_1\log\big(-\cQ\tau(1-s)\big)-\alpha_2\log(1-\tau s)+\wA\log(\tau s)\notag\\
&=\phi_\tau(s)+E\log\tau+\alpha_1\log\cQ+\alpha_1\log(-1),\\
\label{eq:row-euler2}
\re^{\cW/z}\,\frac{\dd U}{U}&=\kph(z)\,\cQ^{\alpha_1/z}\,\tau^{E/z+1}\,\re^{\phi_\tau(s)/z}\,(1-\tau s)^{-1}\,\dd s,
\end{align}
where $\kph(z)\coloneqq-\re^{\alpha_1\log(-1)/z}$.
\end{lem}

\begin{proof}
$1-\cQ U=1-\cQ+\cQ\tau s=(1-\cQ)(1-s)$ because $\cQ\tau=\cQ-1$. In \eqref{eq:row-euler1} the term $\alpha_1\log(-\cQ\tau(1-s))$ contributes $\alpha_1\log(1-s)+\alpha_1\log\tau+\alpha_1\log\cQ+\alpha_1\log(-1)$ and $\wA\log(\tau s)$ contributes $\wA\log s+\wA\log\tau$; the two multiples of $\log\tau$ add up to $(\alpha_1+\wA)\log\tau=E\log\tau$. Exponentiating \eqref{eq:row-euler1} divided by $z$ and multiplying by $\dd U/U=-\tau\,\dd s/(1-\tau s)$ gives \eqref{eq:row-euler2}.
\end{proof}

The constant $\kph$ depends on $z$ and on the branches only, not on $s$ or $t$; it plays no role below. In Pochhammer form, $\re^{\phi_\tau(s)/z}(1-\tau s)^{-1}=s^{\wA/z}(1-s)^{\alpha_1/z}(1-\tau s)^{-\alpha_2/z-1}$. The logarithmic points $U=1$, $U=1/\cQ$, $U=0$ of $\cW$ are $s=0$, $s=1$, $s=1/\tau$. At $t=0$ the first two collide ($1=1/\cQ$) and the singular critical point $U_1$ runs into them; in the coordinate $s$ they stay at $0$ and $1$, the critical point stays at $\sast(0)=\wA/E$ (strictly between them when $\alpha_1,\wA>0$), and $U=0$ escapes to $s=\infty$.

%-------------------------------------------------------------------------------
\subsection{Hessians and the tensor rule}
\label{sec:row-hess}

Define
\begin{equation}
\label{eq:row-algdef}
\alg(\tau)\coloneqq(1-\tau\sast)\Big(\frac{\phi_\tau''(\sast)}{\phi_0''(\wA/E)}\Big)^{1/2},
\end{equation}
the square root with $\alg(0)=1$; that it is a well-defined unit of $\Oan$ is part of the next lemma.

\begin{lem}[Hessians, $\alg^2$, the tensor rule; formal]\genusstatus{AI-proposed, unchecked}
\label{lem:row-hess}
Let $\tau\neq0$. (The parts are labelled (c) and (d); parts (a) and (b) of the same statement, the critical points and the identities (i)--(iv), are \cref{lem:set-N}.)
\begin{enumerate}
\item[(c)] At a root $s_c$ of $\Ns$, $\phi_\tau''(s_c)=\Ns'(s_c)/\big(s_c(1-s_c)(1-\tau s_c)\big)$; hence
\begin{equation}
\label{eq:row-hess}
\phi_\tau''(\sast)=-\frac{\delta}{\sast(1-\sast)U_1},\qquad
\phi_\tau''(\stil)=\frac{\delta}{\stil(1-\stil)U_2},\qquad
\phi_0''(\wA/E)=-\frac{E^3}{\alpha_1\wA}.
\end{equation}
Both critical points are non-degenerate on $\Ddisc''$, and
\begin{equation}
\label{eq:row-alg}
\alg^2=\frac{U_1^2\,\phi_\tau''(\sast)}{\phi_0''(\wA/E)}=\frac{\alpha_1\wA\,\delta\,U_1}{E^3\,\sast(1-\sast)}
\end{equation}
is a unit of $\Oan$ with value $1$ at $\tau=0$, so $\alg$ is well defined. Equivalently,
\begin{equation}
\label{eq:row-algp}
\alg^2=\frac{\wA\,\delta\,(E-\wB\tau\sast)}{E^3\,\sast},\qquad E-\wB\tau\sast=\delta+\tau\Lv_2,
\end{equation}
and $E-\wB\tau\sast$ has no zero on $\Ddisc$. Since $\delta$ and $\sast$ are units of $\Oan(\Ddisc)$ as well, $\alg^2$ is a unit of $\Oan(\Ddisc)$, and $\alg$, $1/\alg$ are analytic on $\Ddisc$.
\item[(d)] (Tensor rule.) For a function $G$ of $U$ and the composite $G(1-\tau s)$,
\[
\frac{\dd}{\dd s}G(1-\tau s)=-\frac{\tau}{1-\tau s}(\thU G)(1-\tau s),\qquad
\frac{\dd^2}{\dd s^2}G(1-\tau s)=\Big(\frac{\tau}{1-\tau s}\Big)^2\big[\thU^2G-\thU G\big](1-\tau s).
\]
Since $\phi_\tau=\cW(1-\tau s)+\mathrm{const}$ by \eqref{eq:row-euler1} and $\cW_1(U_i)=0$,
\begin{equation}
\label{eq:row-tensor}
\cW_2(U_i)=\frac{U_i^2\,\phi_\tau''(s_i)}{\tau^2}\quad((s_1,s_2)\coloneqq(\sast,\stil)),\qquad
\cW_2(U_1)=-\frac{\delta\,U_1}{\tau^2\sast(1-\sast)},\qquad
\cW_2(U_2)=\frac{\delta\,U_2}{\tau^2\stil(1-\stil)}.
\end{equation}
\end{enumerate}
\end{lem}

\begin{proof}
(c) Since $\phi_\tau'=\Ns/(s(1-s)(1-\tau s))$ and $\Ns(s_c)=0$, differentiating gives $\phi_\tau''(s_c)=\Ns'(s_c)/(s_c(1-s_c)(1-\tau s_c))$: the derivative of the denominator multiplies $\Ns(s_c)=0$. By \cref{lem:set-N}, $\Ns'(\sast)=2\wB\tau\sast-(E+\alpha_0\tau)=-\delta$ and $\Ns'(\stil)=\delta$, which gives the first two formulas of \eqref{eq:row-hess} with $1-\tau\sast=U_1$, $1-\tau\stil=U_2$. At $\tau=0$, $\phi_0''(s)=-\wA/s^2-\alpha_1/(1-s)^2$, and at $s=\wA/E$, $1-s=\alpha_1/E$, so $\phi_0''(\wA/E)=-\wA(E/\wA)^2-\alpha_1(E/\alpha_1)^2=-E^2(1/\wA+1/\alpha_1)=-E^3/(\alpha_1\wA)$. On $\Ddisc$ the function $\delta$ has no zero, and $\sast$, $1-\sast$, $U_1$ are units of $\Oan$ (values $\wA/E$, $\alpha_1/E$, $1$ at $\tau=0$), so $\phi_\tau''(\sast)$ is a unit of $\Oan$; likewise $\phi_\tau''(\stil)\neq0$ on $\Ddisc''$. Then \eqref{eq:row-alg} follows from \eqref{eq:row-hess}:
\[
\alg^2=U_1^2\cdot\Big(-\frac{\delta}{\sast(1-\sast)U_1}\Big)\cdot\Big(-\frac{\alpha_1\wA}{E^3}\Big)=\frac{\alpha_1\wA\,\delta\,U_1}{E^3\sast(1-\sast)},
\]
with value $E\cdot(\alpha_1\wA/E^3)\cdot(E/\wA)\cdot(E/\alpha_1)=1$ at $\tau=0$. So $\alg^2$ is a unit of $\Oan$ with value $1$, and $\alg$ is its square root with value $1$ (two square roots of a unit of $\Oan$ with the same value at $0$ coincide). For \eqref{eq:row-algp}, the second form of identity (ii) of \cref{lem:set-N}, $(E-\wB\tau\sast)(1-\sast)=\alpha_1(1-\tau\sast)$, reads $U_1/(1-\sast)=(E-\wB\tau\sast)/\alpha_1$; inserting it into \eqref{eq:row-alg} gives the first formula of \eqref{eq:row-algp}, and $E-\wB\tau\sast=\delta+\tau\Lv_2$ is identity (i) of \cref{lem:set-N}, $\delta=E+\alpha_0\tau-2\wB\tau\sast$, combined with $\tau\Lv_2=\wB\tau\sast-\alpha_0\tau$. If $\wB\tau\sast=E$ then $\Ns(\sast)=\sast(\wB\tau\sast-E-\alpha_0\tau)+\wA=-\alpha_0E/\wB+\wA=(\wA\wB-\alpha_0E)/\wB=\alpha_1\alpha_2/\wB\neq0$ (we used $\wA\wB=\alpha_1\alpha_2+\alpha_0E$), a contradiction; so $E-\wB\tau\sast$ has no zero on $\Ddisc$. As $\delta$ has no zero on $\Ddisc$ and $\sast=2\wA/(E+\alpha_0\tau+\delta)$ is analytic and non-vanishing on $\Ddisc$ (\cref{def:set-alg}), $\alg^2$ is a unit of $\Oan(\Ddisc)$; the disc is simply connected, so its square root $\alg$ with $\alg(0)=1$ and the inverse $1/\alg$ are analytic on $\Ddisc$.

(d) $\dd U/\dd s=-\tau$ and $\thU=U\,\dd/\dd U$ give $\dd/\dd s=-(\tau/U)\thU$ on composites. Differentiating once more, $\dd/\dd s\,(-\tau/(1-\tau s))=-\tau^2/(1-\tau s)^2$, which gives the second formula. At a critical point $\thU\cW=\cW_1=0$, so $\dd^2/\dd s^2[\cW(1-\tau s)]=(\tau/U)^2\cW_2$ there; with \eqref{eq:row-euler1} and \eqref{eq:row-hess} this is \eqref{eq:row-tensor}.
\end{proof}

\begin{rmk}
\label{rmk:row-hessfib}
$\cW_2(U_1)$ has a double pole at $\tau=0$: the coordinate $y=\log U$ degenerates there ($\dd\log(1-\tau s)/\dd s=O(\tau)$), the coordinate $s$ does not. At $\alpha_0=0$ (where $\wA=\alpha_2$, $\wB=\alpha_1$) the critical equation $\Ns(s)=0$ factorises as $\alpha_2(1-s)=\alpha_1s(1-\tau s)$, and \eqref{eq:row-alg} becomes $\alg^2=\alpha_2^2\delta/(E^3\sast{}^2)$; at generic weights there is no factorisation of $\Ns(s)$ of this form, and the two forms of identity (ii) of \cref{lem:set-N} replace it in \cref{sec:row-eik} and in \cref{sec:lea}. The forms \eqref{eq:row-alg} and \eqref{eq:row-algp} of $\alg^2$ are also the two forms of $\Delta_1$ in \cref{prop:lea-key}.
\end{rmk}

%-------------------------------------------------------------------------------
\subsection{The rescaled derivatives \texorpdfstring{$\tau^k\cW_k(U_1)$}{tau\textasciicircum k W\_k(U\_1)}}
\label{sec:row-wr}

For $k\geq1$,
\begin{equation}
\label{eq:row-phik}
\begin{gathered}
\phi_\tau^{(k)}(s)=(k-1)!\Big[(-1)^{k-1}\wA s^{-k}-\alpha_1(1-s)^{-k}+\alpha_2\tau^k(1-\tau s)^{-k}\Big],\\
\phi_0^{(k)}(\wA/E)=(k-1)!\,E^k\Big[(-1)^{k-1}\wA^{1-k}-\alpha_1^{1-k}\Big],
\end{gathered}
\end{equation}
since $(\wA/E)^{-k}=E^k\wA^{-k}$ and $(1-\wA/E)^{-k}=(\alpha_1/E)^{-k}=E^k\alpha_1^{-k}$. For $a\in\{1,\cQ\}$ put $p_a\coloneqq(1-aU)^{-1}$. Then $\thU\log(1-aU)=-aU/(1-aU)=1-p_a$ and $\thU p_a=aU/(1-aU)^2=p_a^2-p_a$. Define polynomials $\Pol_0(p)\coloneqq p$ and $\Pol_n(p)\coloneqq(p^2-p)\Pol_{n-1}'(p)$ (so $\Pol_1=p^2-p$, $\Pol_2=2p^3-3p^2+p$). By induction on $k$,
\[
\thU^k\log(1-aU)=-\Pol_{k-1}(p_a)\quad(k\geq2):
\]
for $k=2$, $\thU(1-p_a)=-(p_a^2-p_a)=-\Pol_1(p_a)$, and $\thU\Pol_{n-1}(p_a)=\Pol_{n-1}'(p_a)\thU p_a=(p_a^2-p_a)\Pol_{n-1}'(p_a)=\Pol_n(p_a)$. Also $\thU^k\log U=0$ for $k\geq2$ since $\thU\log U=1$. Hence
\begin{equation}
\label{eq:row-Wk}
\cW_k=-\alpha_1\Pol_{k-1}(p_\cQ)-\wA\Pol_{k-1}(p_1)\qquad(k\geq2).
\end{equation}
$\Pol_n$ has degree $n+1$ and leading coefficient $n!$ (the degree rises by one and the leading coefficient is multiplied by $n$ at each step), and $\Pol_n(0)=0$; write $\Pol_{k-1}(p)=\sum_{i=1}^{k}\pc_{k-1,i}\,p^i$ with $\pc_{k-1,k}=(k-1)!$.

\begin{lem}[Rescaled derivatives; formal]\genusstatus{AI-proposed, unchecked}
\label{lem:row-wr}
For $k\geq2$ and $\tau\in\Ddisc''$ put $\wrk_k(\tau)\coloneqq\tau^k\cW_k(U_1)$. Then
\begin{equation}
\label{eq:row-wr}
\wrk_k=-\sum_{i=1}^{k}\pc_{k-1,i}\,\tau^{k-i}\Big[\alpha_1\Big(-\frac{1-\tau}{1-\sast}\Big)^{i}+\wA\,\sast{}^{-i}\Big],
\end{equation}
whose right side is analytic on $\Ddisc'$, $\tau=0$ included, with
\[
\wrk_k(0)=-(k-1)!\,E^k\Big[(-1)^k\alpha_1^{1-k}+\wA^{1-k}\Big]=(-1)^k\phi_0^{(k)}(\wA/E),\qquad
\wrk_2=-\frac{\delta\,U_1}{\sast(1-\sast)}\ \text{(no zero on $\Ddisc'$)}.
\]
\end{lem}

\begin{proof}
For $\tau\in\Ddisc''$, $U_1=1-\tau\sast$ is not in $\{1,1/\cQ\}$ ($\tau\sast\neq0$, and $\cQ U_1=1$ would force $\sast=1$), so $\cW_k(U_1)$ is defined: by \eqref{eq:row-Wk} $\cW_k$ has poles only at $U=1$ and $U=1/\cQ$. At $U_1$, $1-U_1=\tau\sast$ and $1-\cQ U_1=-\cQ\tau(1-\sast)$ (\cref{lem:row-euler}), so, with $1/\cQ=1-\tau$,
\[
\tau\,p_1(U_1)=\frac{1}{\sast},\qquad \tau\,p_\cQ(U_1)=-\frac{1-\tau}{1-\sast}.
\]
Multiplying \eqref{eq:row-Wk} by $\tau^k$ and writing $\tau^k\Pol_{k-1}(p)=\sum_i\pc_{k-1,i}\tau^{k-i}(\tau p)^i$ gives \eqref{eq:row-wr}. Only powers $\tau^{k-i}$ with $0\leq k-i\leq k-1$ occur, and $\sast$, $1/\sast$, $1/(1-\sast)$ are analytic on $\Ddisc'$ (\cref{lem:set-N}: $\sast\notin\{0,1\}$ there), so the right side is analytic on $\Ddisc'$. At $\tau=0$ only $i=k$ survives; with $\sast(0)=\wA/E$, $1-\sast(0)=\alpha_1/E$ and $\pc_{k-1,k}=(k-1)!$,
\[
\wrk_k(0)=-(k-1)!\Big[\alpha_1\Big(-\frac{E}{\alpha_1}\Big)^k+\wA\Big(\frac{E}{\wA}\Big)^k\Big]=-(k-1)!\,E^k\Big[(-1)^k\alpha_1^{1-k}+\wA^{1-k}\Big],
\]
which equals $(-1)^k\phi_0^{(k)}(\wA/E)$ by \eqref{eq:row-phik}. Finally $\wrk_2=\tau^2\cW_2(U_1)$ is \eqref{eq:row-tensor}, and $\delta$, $U_1$, $\sast$, $1-\sast$ have no zero on $\Ddisc'$.
\end{proof}

So $\cW_k(U_1)=\tau^{-k}\wrk_k=O(\tau^{-k})$ and $\cW_2(U_1)^{-1}=O(\tau^2)$: the coordinate $y=\log U$ behaves at $t=0$ like $-\tau$ times a good coordinate (\cref{rmk:row-unif}).

%-------------------------------------------------------------------------------
\subsection{The formal Laplace calculus}
\label{sec:row-laplace}

This subsection contains no weights; it is the formal counterpart of Laplace's method for integrals $\int\re^{f/z}\amp\,\dd u$ through a non-degenerate critical point, together with the one analytic statement (\cref{lem:row-laplace}) that connects it with actual integrals. Only \cref{lem:row-laplace} is analytic; everything else is formal.

\begin{defn}[The ring $\Acal$ and the Gaussian functional]
\label{def:row-Acal}
Let $\Acal$ be the set of formal sums $K=\sum_{p\geq0,\,q\in\bbZ}K_{pq}v^pz^q$ (coefficients complex numbers, or analytic functions of parameters) such that, with the weight $\mathrm{wt}(p,q)\coloneqq p/2+q$, (A1) $\mathrm{wt}$ is bounded below on the support of $K$, and (A2) for every $N$ only finitely many $(p,q)$ with $\mathrm{wt}\leq N$ occur. $\Acal$ is a ring (weights add; every coefficient of a product is a finite sum). It contains $\bbC\fps{v}[z,1/z]$ and $\mathrm{Exp}[\Rr(v)/z]\coloneqq\sum_j\Rr^j/(j!z^j)$ for every $\Rr\in v^3\bbC\fps{v}$ (the term $(\Rr/z)^j$ has weight $\geq j/2$). $\dd/\dd v$ is a derivation of $\Acal$ lowering weights by $1/2$, multiplication by $v$ raises them by $1/2$, and $\dd/\dd v\,\mathrm{Exp}[\Rr/z]=(\Rr'/z)\mathrm{Exp}[\Rr/z]$. For $a\neq0$ the \emph{Gaussian functional} (the formal moments of a centred Gaussian of variance $-z/a$) is
\begin{equation}
\label{eq:row-gaussfun}
\gauss{K}{a}\coloneqq\sum_{p\ \mathrm{even},\,q}K_{pq}\,(p-1)!!\,\Big(-\frac za\Big)^{p/2}z^q\quad\in\Cz\qquad((-1)!!\coloneqq1).
\end{equation}
Each power $z^w$ receives contributions only from monomials of weight $w$, finitely many by (A2), and $w$ is bounded below by (A1); $\gauss{\cdot}{a}$ is linear over $\Cz$ and over functions of parameters, and maps elements of weight $\geq N$ into $z^N\bbC\fps{z}$.
\end{defn}

\begin{lem}[Gaussian calculus; formal]\genusstatus{AI-proposed, unchecked}
\label{lem:row-gauss}
For $K\in\Acal$:
\begin{enumerate}
\item[(G0)] scaling: $\gauss{K(\lambda v)}{\lambda^2a}=\gauss{K(v)}{a}$ for $\lambda\neq0$;
\item[(G1)] integration by parts: $\gauss{\dd K/\dd v}{a}=-(a/z)\gauss{vK}{a}$;
\item[(G2)] variance derivative: if $a=a(\tau)$ and $K$ depends analytically on $\tau$ (coefficientwise, with $\tau$-independent support), then
\[
\frac{\dd}{\dd\tau}\gauss{K}{a(\tau)}=\gauss{\dd K/\dd\tau}{a}+\frac{a_\tau}{2z}\gauss{v^2K}{a}+\frac{a_\tau}{2a}\gauss{K}{a}.
\]
\end{enumerate}
\end{lem}

\begin{proof}
Each side is $\gauss{\mathsf{L}K}{\cdot}$ for an operator $\mathsf{L}$ on $\Acal$ that shifts weights by a bounded amount, so every $z$-coefficient involves finitely many monomials of $K$, and it suffices to take $K=v^p$ (for (G2), together with the product rule for the coefficients). Put $m_p\coloneqq\gauss{v^p}{a}$, so $m_{2n}=(2n-1)!!(-z/a)^n$ and $m_{\mathrm{odd}}=0$. (G0): $\lambda^{2n}(2n-1)!!(-z/(\lambda^2a))^n=m_{2n}$. (G1): for $p=2n+1$, $\gauss{p\,v^{2n}}{a}=(2n+1)!!(-z/a)^n=-(a/z)(2n+1)!!(-z/a)^{n+1}=-(a/z)m_{2n+2}$; for $p$ even both sides vanish. (G2): $\dd/\dd\tau\,(-z/a)^n=-n(a_\tau/a)(-z/a)^n$, and $(a_\tau/(2z))m_{2n+2}+(a_\tau/(2a))m_{2n}=m_{2n}(a_\tau/(2a))[-(2n+1)+1]=-n(a_\tau/a)m_{2n}$.
\end{proof}

\begin{defn}[Normalised formal Laplace expansions]
\label{def:row-Lnorm}
Data: a point $u_0$, a \emph{phase} $f$ analytic near $u_0$ (only its Taylor series matters) with $f'(u_0)=0$ and $f_2\coloneqq f''(u_0)\neq0$, and an \emph{amplitude} $\amp\in\Oan_{u_0}[z,1/z]$ (a Laurent polynomial in $z$ with coefficients analytic near $u_0$). Write
\[
f(u_0+v)=f(u_0)+\tfrac12f_2v^2+\Rr_f(v),\qquad \Rr_f(v)\coloneqq\sum_{k\geq3}f^{(k)}(u_0)\frac{v^k}{k!},
\]
and define
\begin{align*}
\Lnorm[f,\amp;u_0]&\coloneqq\gauss{\amp(u_0+v)\,\mathrm{Exp}[\Rr_f(v)/z]}{f_2}\quad\in\Cz,\\
\Lhat[f,\amp;u_0]&\coloneqq\Lnorm[f,\amp;u_0]/\amp(u_0)\quad\text{for }\amp\in\Oan_{u_0},\ \amp(u_0)\neq0,\\
\Lfull[f,\amp;u_0]&\coloneqq\re^{f(u_0)/z}(2\pi z)^{1/2}(-f_2)^{-1/2}\,\Lnorm[f,\amp;u_0],
\end{align*}
the Gaussian-normalised formal Laplace expansion, the normalised expansion (leading coefficient $1$), and the unnormalised expansion (a formal symbol, with a chosen root). $\Lfull$ is the formal counterpart of $\int\re^{f/z}\amp\,\dd u$ over a path through $u_0$ (\cref{lem:row-laplace}). Immediate properties:
\begin{enumerate}
\item[(D1)] (Wick form.) $\Lnorm[f,\amp;u_0]=\sum_{n\geq0}\frac{(2n)!}{n!\,2^n}\big(-\frac{z}{f_2}\big)^n\coef{v^{2n}}\big(\amp(u_0+v)\mathrm{Exp}[\Rr_f(v)/z]\big)$. For $\amp=1$ the coefficient of $z^m$ is the Wick polynomial $\wick_m$ of \cref{def:set-rows} with $\cW_k$ replaced by $f^{(k)}(u_0)$.
\item[(D2)] For $\amp\in\Oan_{u_0}$, $\Lnorm[f,\amp;u_0]$ lies in $\bbC\fps{z}$ with constant term $\amp(u_0)$, and its $z^m$ coefficient is a polynomial with rational coefficients in $1/f_2$, $f^{(3)}(u_0),\dots,f^{(2m+2)}(u_0)$, $\amp(u_0),\dots,\amp^{(2m)}(u_0)$. (Both factors have weight $\geq0$; the weight-$m$ part involves monomials $v^a$ of $\amp$ with $a/2\leq m$ and terms $f^{(k)}v^k/(k!z)$ of weight $k/2-1\leq m$.)
\item[(D3)] $\Lnorm$ depends only on the Taylor coefficients of $f$ and $\amp$ at $u_0$ (so a translation $u=u_0+v$ changes nothing), is unchanged when a constant is added to $f$, and is linear in $\amp$ over $\bbC[z,1/z]$ and over functions independent of $u$.
\item[(D4)] (Rescaling of the phase.) For a constant $\lambda\neq0$ and an amplitude $\amp$ independent of $z$, $\Lnorm[\lambda f,\amp;u_0](z)=\Lnorm[f,\amp;u_0](z/\lambda)$: indeed $\Rr_{\lambda f}/z=\Rr_f/(z/\lambda)$, so every $z$-dependence of $\amp(u_0+v)\mathrm{Exp}[\Rr_{\lambda f}/z]$ is through $z/\lambda$, and $(-z/(\lambda f_2))^{p/2}=(-(z/\lambda)/f_2)^{p/2}$ in \eqref{eq:row-gaussfun}.
\item[(D5)] (The parameters of the Beta phase.) By (D2) and \eqref{eq:row-phik}, the $z^m$ coefficient of $\Lnorm[\phi_0,1;\wA/E]$ is a polynomial with rational coefficients in $1/\phi_0''(\wA/E)=-\alpha_1\wA/E^3$ and in the $\phi_0^{(k)}(\wA/E)$ of \eqref{eq:row-phik}, hence an element of $\bbQ(\alpha_1,\wA)=\bbQ(\alpha_1,\wA,E)$, homogeneous of degree $-m$ in the weights (in (D1) a term with $j$ factors $f^{(k_i)}v^{k_i}/(k_i!z)$ in the $v^{2n}$ coefficient has $z$-order $n-j$ and degree $j-n$ in the weights, as each $f^{(k)}(\wA/E)$ has degree $1$ and $1/f_2$ has degree $-1$).
\end{enumerate}
\end{defn}

\begin{prop}[(P3): total derivatives; formal]\genusstatus{AI-proposed, unchecked}
\label{prop:row-P3}
For every $\mathsf{k}\in\Oan_{u_0}[z,1/z]$: $\Lnorm[f,\mathsf{k}'+f'\mathsf{k}/z;u_0]=0$. (The amplitude $\mathsf{k}'+f'\mathsf{k}/z$ is that of a total derivative: $\re^{f/z}(\mathsf{k}'+f'\mathsf{k}/z)=\dd/\dd u\,(\re^{f/z}\mathsf{k})$.)
\end{prop}

\begin{proof}
Put $K(v)\coloneqq\mathsf{k}(u_0+v)\mathrm{Exp}[\Rr_f(v)/z]\in\Acal$. Since $f'(u_0+v)=f_2v+\Rr_f'(v)$, we have $(\mathsf{k}'+f'\mathsf{k}/z)(u_0+v)\mathrm{Exp}[\Rr_f(v)/z]=\dd K/\dd v+(f_2/z)vK$, and $\gauss{\dd K/\dd v}{f_2}+(f_2/z)\gauss{vK}{f_2}=0$ by (G1).
\end{proof}

\begin{prop}[(P2): derivative in a parameter with a moving expansion point; formal]\genusstatus{AI-proposed, unchecked}
\label{prop:row-P2}
Let $f(\tau,s)$ and $\amp(\tau,s)$ be analytic near $(\tau_0,s_0)$ ($\amp$ may be a Laurent polynomial in $z$ with such coefficients), let $s_c(\tau)$ be analytic with $s_c(\tau_0)=s_0$, $\de_sf(\tau,s_c(\tau))=0$ and $f_2(\tau)\coloneqq\de_s^2f(\tau,s_c(\tau))\neq0$. Put $f_0(\tau)\coloneqq f(\tau,s_c(\tau))$ and $\mathsf{n}(\tau)\coloneqq\Lnorm[f(\tau,\cdot),\amp(\tau,\cdot);s_c(\tau)]$. Then every $z$-coefficient of $\mathsf{n}$ is analytic in $\tau$, and
\begin{equation}
\label{eq:row-P2}
\de_\tau\mathsf{n}+\Big(\frac{f_{0,\tau}}{z}-\frac{f_{2,\tau}}{2f_2}\Big)\mathsf{n}=\Lnorm[f(\tau,\cdot),\Dam\amp;s_c(\tau)],\qquad \Dam\amp\coloneqq\de_\tau\amp+\amp\,\frac{\de_\tau f}{z},
\end{equation}
where $f_{0,\tau}=(\de_\tau f)(\tau,s_c(\tau))$ (by the critical-point condition) and $f_{2,\tau}$ is the total derivative of $f_2$. Equivalently, with $(-f_2(\tau))^{-1/2}$ chosen continuously, $\de_\tau\Lfull[f,\amp;s_c(\tau)]=\Lfull[f,\Dam\amp;s_c(\tau)]$: the expansion of the integrand $\re^{f/z}\amp$ commutes with $\dd/\dd\tau$ although the expansion point moves.
\end{prop}

\begin{proof}
Analyticity follows from (D2), as $\de_s^kf(\tau,s_c(\tau))$, $\de_s^k\amp(\tau,s_c(\tau))$ and $1/f_2(\tau)$ are analytic.

\emph{Step 1 (shifted variable).} Let $\tilde f(\tau,v)\coloneqq f(\tau,s_c(\tau)+v)$ and $\tilde\amp(\tau,v)\coloneqq\amp(\tau,s_c(\tau)+v)$. By (D3), $\mathsf{n}=\Lnorm[\tilde f,\tilde\amp;0]$, and $\de_v\tilde f(\tau,0)=0$ for all $\tau$. Write $\tilde f=f_0(\tau)+\frac12f_2(\tau)v^2+\tilde\Rr(\tau,v)$ and $\tilde K\coloneqq\tilde\amp\,\mathrm{Exp}[\tilde\Rr/z]\in\Acal$. By (G2), with $\de_\tau\tilde K=(\de_\tau\tilde\amp+\tilde\amp\,\de_\tau\tilde\Rr/z)\mathrm{Exp}[\tilde\Rr/z]$,
\[
\de_\tau\mathsf{n}=\gauss{(\de_\tau\tilde\amp+\tilde\amp\,\de_\tau\tilde\Rr/z)\mathrm{Exp}[\tilde\Rr/z]}{f_2}+\frac{f_{2,\tau}}{2z}\gauss{v^2\tilde K}{f_2}+\frac{f_{2,\tau}}{2f_2}\mathsf{n}.
\]
Since $\de_\tau\tilde f=f_{0,\tau}+\frac12f_{2,\tau}v^2+\de_\tau\tilde\Rr$,
\begin{align*}
\Lnorm[\tilde f,\de_\tau\tilde\amp+\tilde\amp\,\de_\tau\tilde f/z;0]
&=\gauss{(\de_\tau\tilde\amp+\tilde\amp\,\de_\tau\tilde\Rr/z)\mathrm{Exp}[\tilde\Rr/z]}{f_2}+\frac{f_{0,\tau}}{z}\mathsf{n}+\frac{f_{2,\tau}}{2z}\gauss{v^2\tilde K}{f_2}\\
&=\de_\tau\mathsf{n}+\Big(\frac{f_{0,\tau}}{z}-\frac{f_{2,\tau}}{2f_2}\Big)\mathsf{n}.
\end{align*}

\emph{Step 2 (back to $s$).} $\de_\tau\tilde\amp=(\de_\tau\amp+s_{c,\tau}\de_s\amp)(\tau,s_c+v)$ and $\de_\tau\tilde f=(\de_\tau f+s_{c,\tau}\de_sf)(\tau,s_c+v)$, so
\[
\de_\tau\tilde\amp+\tilde\amp\,\de_\tau\tilde f/z=\big[\Dam\amp+s_{c,\tau}\big(\de_s\amp+(\de_sf)\amp/z\big)\big](\tau,s_c+v).
\]
The last term is $s_{c,\tau}$ times an amplitude of the form $\mathsf{k}'+f'\mathsf{k}/z$ (with $\mathsf{k}=\amp(\tau,\cdot)$), whose $\Lnorm$ vanishes by \cref{prop:row-P3}. Translating back by (D3) gives \eqref{eq:row-P2}. The unnormalised form follows from $\de_\tau[\re^{f_0/z}(-f_2)^{-1/2}\mathsf{n}]=\re^{f_0/z}(-f_2)^{-1/2}[\de_\tau\mathsf{n}+(f_{0,\tau}/z-f_{2,\tau}/(2f_2))\mathsf{n}]$.
\end{proof}

In \eqref{eq:row-P2}, $f_{0,\tau}/z$ comes from the moving critical value $\re^{f_0/z}$, $-f_{2,\tau}/(2f_2)$ from the Gaussian normalisation $(-f_2)^{-1/2}$, and the motion of the expansion point is a total derivative.

\begin{prop}[(P1): coordinate invariance; formal]\genusstatus{AI-proposed, unchecked}
\label{prop:row-P1}
Let $h$ be analytic near $v_0$ with $h(v_0)=u_0$ and $h_1\coloneqq h'(v_0)\neq0$, and let $f$, $\amp$ be as in \cref{def:row-Lnorm} with $\amp\in\Oan_{u_0}$. Then $f\circ h$ has a non-degenerate critical point at $v_0$, $(f\circ h)(v_0)=f(u_0)$, $(f\circ h)''(v_0)=f_2h_1^2$, and
\begin{equation}
\label{eq:row-P1}
\Lnorm[f\circ h,(\amp\circ h)\,h';v_0]=h_1\,\Lnorm[f,\amp;u_0].
\end{equation}
Exactly what is invariant: (i) the critical value $\re^{f(u_0)/z}$; (ii) the unnormalised expansion, $\Lfull[f\circ h,(\amp\circ h)h';v_0]=\Lfull[f,\amp;u_0]$, when the roots are related by $(-(f\circ h)''(v_0))^{1/2}=h_1(-f_2)^{1/2}$ (the Jacobian $h'$ in the amplitude is compensated by $h_1$ in the Gaussian normalisation; for real data and $h_1>0$ both roots are positive, and $h_1<0$ is the orientation reversal of the path); (iii) the normalised expansion, with no convention at all: $\Lhat[f\circ h,(\amp\circ h)h';v_0]=\Lhat[f,\amp;u_0]$, because $((\amp\circ h)h')(v_0)=h_1\amp(u_0)$.
\end{prop}

\begin{proof}
For $\varepsilon\in\bbC$ let $h_\varepsilon(v)\coloneqq u_0+h_1(v-v_0)+\varepsilon\big(h(v)-u_0-h_1(v-v_0)\big)$; then $h_\varepsilon(v_0)=u_0$ and $h_\varepsilon'(v_0)=h_1$ for every $\varepsilon$, $h_0$ is affine and $h_\varepsilon=h$ at $\varepsilon=1$. On a neighbourhood $V$ of $v_0$ (uniform for $|\varepsilon|\leq2$) $h_\varepsilon'$ has no zero and $h_\varepsilon(V)$ lies in the domain of $f$, $\amp$. Put $f_\varepsilon\coloneqq f\circ h_\varepsilon$, $\amp_\varepsilon\coloneqq(\amp\circ h_\varepsilon)h_\varepsilon'$ and $n(\varepsilon)\coloneqq\Lnorm[f_\varepsilon,\amp_\varepsilon;v_0]$. Each $f_\varepsilon$ has a critical point at $v_0$ with critical value $f(u_0)$ and $f_\varepsilon''(v_0)=f_2h_1^2$, independent of $\varepsilon$; the Taylor coefficients of $f_\varepsilon$, $\amp_\varepsilon$ at $v_0$ are polynomials in $\varepsilon$, so by (D2) each $z$-coefficient of $n(\varepsilon)$ is a polynomial in $\varepsilon$. Now \eqref{eq:row-P2} with $\tau=\varepsilon$, $s_c(\varepsilon)=v_0$, $f_{0,\varepsilon}=0$, $f_{2,\varepsilon}=0$ gives $n'(\varepsilon)=\Lnorm[f_\varepsilon,\de_\varepsilon\amp_\varepsilon+\amp_\varepsilon\,\de_\varepsilon f_\varepsilon/z;v_0]$. Let $\eta_\varepsilon\coloneqq\de_\varepsilon h_\varepsilon=h(v)-u_0-h_1(v-v_0)$ and $\xi\coloneqq\eta_\varepsilon/h_\varepsilon'$. Then $\de_\varepsilon f_\varepsilon=(f'\circ h_\varepsilon)\eta_\varepsilon=\xi f_\varepsilon'$, and $\de_\varepsilon\amp_\varepsilon=(\amp'\circ h_\varepsilon)\eta_\varepsilon h_\varepsilon'+(\amp\circ h_\varepsilon)\eta_\varepsilon'=(\xi\amp_\varepsilon)'$ (because $\xi\amp_\varepsilon=\eta_\varepsilon(\amp\circ h_\varepsilon)$). So the amplitude is $(\xi\amp_\varepsilon)'+f_\varepsilon'(\xi\amp_\varepsilon)/z$, a total derivative, and $n'(\varepsilon)=0$ by \cref{prop:row-P3}; hence $n(1)=n(0)$. Finally, with $w=v-v_0$, $f_0(v_0+w)=f(u_0+h_1w)$ has $\Rr_{f_0}(w)=\Rr_f(h_1w)$ and $f_0''(v_0)=h_1^2f_2$, so by (G0) with $\lambda=h_1$, $n(0)=h_1\gauss{\amp(u_0+h_1w)\mathrm{Exp}[\Rr_f(h_1w)/z]}{h_1^2f_2}=h_1\Lnorm[f,\amp;u_0]$. (ii) and (iii) restate \eqref{eq:row-P1}. (Here $\eta_\varepsilon$ is an auxiliary function of this proof, not the pairing.)
\end{proof}

\begin{lem}[Laplace's method; analytic]\genusstatus{AI-proposed, unchecked}
\label{lem:row-laplace}
Let $I$ be an open real interval, $f$ and $\amp$ real-analytic on $I$, $u_0\in I$ with $f'(u_0)=0$, $f''(u_0)<0$, and assume $\sup_{I\smallsetminus(u_0-e,u_0+e)}f<f(u_0)$ for every $e>0$ and $\int_I\re^{f/z_0}|\amp|\,\dd u<\infty$ for some $z_0>0$. Then for every $M$, as $z\to0^+$,
\begin{equation}
\label{eq:row-laplace}
\int_I\re^{f(u)/z}\amp(u)\,\dd u=\re^{f(u_0)/z}\Big(\frac{2\pi z}{-f''(u_0)}\Big)^{1/2}\Big[\sum_{m<M}\big(\coef{z^m}\Lnorm[f,\amp;u_0]\big)z^m+O(z^M)\Big].
\end{equation}
\end{lem}

\begin{proof}
\emph{Localisation.} Fix $e>0$ small and $\eta>0$ with $f\leq f(u_0)-\eta$ outside $(u_0-e,u_0+e)$. There, for $0<z<z_0$, $\re^{f/z}=\re^{f/z_0}\re^{f(1/z-1/z_0)}\leq\re^{f/z_0}\re^{(f(u_0)-\eta)(1/z-1/z_0)}$, so the outer part of the integral is $O(\re^{(f(u_0)-\eta)/z})=\re^{f(u_0)/z}O(z^M)$ for every $M$.

\emph{Morse coordinate.} $f(u)-f(u_0)=(u-u_0)^2q(u)$ with $q$ real-analytic and $q(u_0)=f''(u_0)/2<0$; shrinking $e$, $q<0$ on $(u_0-e,u_0+e)$, and $w(u)\coloneqq(u-u_0)\big(2q(u)/f''(u_0)\big)^{1/2}$ is real-analytic and strictly increasing there, with $w(u_0)=0$, $w'(u_0)=1$ and $f(u)=f(u_0)+\frac12f''(u_0)w^2$. Let $u(w)$ be the inverse map, $[w_-,w_+]$ ($w_-<0<w_+$) the image interval and $G(w)\coloneqq\amp(u(w))\,u'(w)$.

\emph{Watson's lemma.} Put $\mathsf{a}\coloneqq-f''(u_0)>0$ and expand $G(w)=\sum_{p<2M}G_pw^p+O(|w|^{2M})$ on $[w_-,w_+]$. Each $\int_{w_-}^{w_+}\re^{-\mathsf{a}w^2/(2z)}w^p\,\dd w$ equals the integral over $\bbR$, which is $(2\pi z/\mathsf{a})^{1/2}(p-1)!!(z/\mathsf{a})^{p/2}$ for $p$ even and $0$ for $p$ odd, up to $O(\re^{-b/z})$ with $b>0$, and $\int_\bbR\re^{-\mathsf{a}w^2/(2z)}|w|^{2M}\,\dd w=O(z^{M+1/2})$. So the inner integral is
\[
\re^{f(u_0)/z}\Big(\frac{2\pi z}{\mathsf{a}}\Big)^{1/2}\Big[\sum_{n<M}G_{2n}(2n-1)!!\Big(\frac z{\mathsf{a}}\Big)^n+O(z^M)\Big],
\]
and by \eqref{eq:row-gaussfun} and (D3)
\[
\sum_nG_{2n}(2n-1)!!\Big(\frac z{\mathsf{a}}\Big)^n=\gauss{G}{-\mathsf{a}}=\Lnorm[\tfrac12f''(u_0)w^2,G;0].
\]
Finally \eqref{eq:row-P1} with $h=u(w)$, $h_1=1$, and (D3) for the constant $f(u_0)$ give $\Lnorm[f\circ u,(\amp\circ u)u';0]=\Lnorm[f,\amp;u_0]$.
\end{proof}

\Cref{lem:row-laplace} is used only for Beta integrals, in \cref{lem:row-betaphase,prop:row-beta}.

\begin{lem}[Homogeneity of the Wick coefficients; formal]\genusstatus{AI-proposed, unchecked}
\label{lem:row-homog}
Every monomial of $\wick_m$ is a rational multiple of $\cW_2^{-n}\prod_{i=1}^{n-m}\cW_{k_i}$ with all $k_i\geq3$ and $\sum_ik_i=2n$. Consequently, for every $\lambda\neq0$, as an identity of rational functions of the $\cW_k$,
\begin{equation}
\label{eq:row-homog}
\wick_m(\lambda^2\cW_2,\lambda^3\cW_3,\dots,\lambda^{2m+2}\cW_{2m+2})=\wick_m(\cW_2,\cW_3,\dots,\cW_{2m+2}),
\end{equation}
and $\Lnorm[f(u_0+\lambda\,\cdot),1;0]=\Lnorm[f,1;u_0]$ for $f$ as in \cref{def:row-Lnorm}.
\end{lem}

\begin{proof}
By the Wick formula of \cref{def:set-rows}, the monomials of the $z^m$ coefficient are $\cW_2^{-n}\prod_{k\in p}\cW_k$ with $p$ a partition of $2n$ into $n-m$ parts $\geq3$; under $\cW_k\mapsto\lambda^k\cW_k$ such a monomial is multiplied by $\lambda^{-2n}\lambda^{2n}=1$. For the second statement put $f_\lambda(v)\coloneqq f(u_0+\lambda v)$; then $f_\lambda'(0)=0$ and $f_\lambda^{(k)}(0)=\lambda^kf^{(k)}(u_0)$, so by (D1) the two $z^m$ coefficients are $\wick_m$ at the arguments $\lambda^kf^{(k)}(u_0)$ and $f^{(k)}(u_0)$, which agree by \eqref{eq:row-homog}. (Equivalently: (G0) for $K=\mathrm{Exp}[\Rr_f/z]$.)
\end{proof}

In words: giving $\cW_k$ the weight $k$, every Wick monomial has weight $-2n+\sum k_i=0$, so the Wick formula is blind to linear rescalings of the coordinate. This is why the degenerate coordinate $y=\log U$ is harmless.

%-------------------------------------------------------------------------------
\subsection{The point-frame row is the normalised expansion of the Euler integrand}
\label{sec:row-wickrow}

Recall from \cref{def:set-rows} that the point-frame identity row at the singular idempotent is
\[
\rpt_{0,1}(z;t)=\sum_{m\geq0}z^m\wick_m\big(\cW_2(U_1),\cW_3(U_1),\dots\big).
\]
Put
\begin{equation}
\label{eq:row-Ndef}
\begin{gathered}
\Nlap(z,\tau)\coloneqq\Lhat[\phi_\tau,(1-\tau s)^{-1};\sast(\tau)]=U_1\,\Lnorm[\phi_\tau,(1-\tau s)^{-1};\sast(\tau)],\\
\Abeta(z)\coloneqq\Lnorm[\phi_0,1;\wA/E],\qquad \Abetam{m}\coloneqq\coef{z^m}\Abeta .
\end{gathered}
\end{equation}

\begin{prop}[formal]\genusstatus{AI-proposed, unchecked}
\label{prop:row-wickrow}
\begin{enumerate}
\item For every $\tau\in\Ddisc''$ (every small $t\neq0$), coefficientwise in $z$,
\begin{equation}
\label{eq:row-wickrow}
\rpt_{0,1}(z;t)=\Nlap(z,\tau).
\end{equation}
\item For every $m\geq0$ and $\tau\in\Ddisc''$,
\begin{equation}
\label{eq:row-wick13}
\coef{z^m}\rpt_{0,1}(z;t)=\wick_m\big(\cW_2(U_1),\cW_3(U_1),\dots\big)=\wick_m\big(\wrk_2(\tau),\wrk_3(\tau),\dots,\wrk_{2m+2}(\tau)\big),
\end{equation}
a polynomial in $1/\wrk_2,\wrk_3,\dots,\wrk_{2m+2}$ with rational coefficients. Hence the $z^m$ coefficient of $\rpt_{0,1}$, an algebraic function of $\cQ$, extends analytically to all of $\Ddisc'$, $\tau=0$ included, and its value at $\tau=0$ is $\Abetam{m}$.
\item The $z^m$ coefficient $\Nlap_m$ of $\Nlap$ is analytic on $\Ddisc'$ and $\Nlap(z,0)=\Abeta(z)$. So \eqref{eq:row-wickrow} holds on all of $\Ddisc'$, $\Nlap$ lies in $\Oan(\Ddisc')\fps{z}$, and the Taylor series at $t=0$ of the $z^m$ coefficient of $\rpt_{0,1}$ is that of $\Nlap_m(\tau(t))$: \eqref{eq:row-wickrow} holds in $\bbC\fps{z}\fps{t}$.
\end{enumerate}
\end{prop}

\begin{proof}
\emph{(1) The point-frame row is a normalised expansion in $y=\log U$.} Fix $\tau\in\Ddisc''$. Then $U_1\neq0$ (as $\sast\neq1/\tau$), so $y_1\coloneqq\log U_1$ (any branch) is defined and $\cW_y(y)\coloneqq\cW(\re^y)$ is analytic near $y_1$, with $\cW_y^{(k)}(y_1)=\cW_k(U_1)$ (since $\dd/\dd y=\thU$), $\cW_y'(y_1)=0$ and $\cW_y''(y_1)=\cW_2(U_1)=-\delta U_1/(\tau^2\sast(1-\sast))\neq0$ by \eqref{eq:row-tensor} ($\delta U_1\neq0$ on $\Ddisc''$). In this coordinate the integrand is $\re^{\cW/z}\,\dd U/U=\re^{\cW_y(y)/z}\,\dd y$, with amplitude $1$. By (D1) and the Wick formula of \cref{def:set-rows}, $\rpt_{0,1}(z;t)=\Lnorm[\cW_y,1;y_1]=\Lhat[\cW_y,1;y_1]$, to every order in $z$.

\emph{(2) Change of coordinate; proof of (i).} $h(s)\coloneqq\log(1-\tau s)$, with the branch for which $h(\sast)=y_1$, is analytic near $\sast$, and $h'(\sast)=-\tau/U_1\neq0$. By \eqref{eq:row-euler1}, $\cW_y\circ h=\phi_\tau+\mathrm{const}$ near $\sast$, and the amplitude $1$ becomes $(1\circ h)h'=-\tau(1-\tau s)^{-1}$. By \cref{prop:row-P1}(iii), (D3) and linearity (the factor $-\tau$),
\[
\Lhat[\cW_y,1;y_1]=\Lhat[\phi_\tau,-\tau(1-\tau s)^{-1};\sast]=\Lhat[\phi_\tau,(1-\tau s)^{-1};\sast]=\Nlap(z,\tau),
\]
and $\Lhat=\Lnorm/(1-\tau\sast)^{-1}=U_1\Lnorm$.

\emph{(3) Proof of (ii).} The first equality in \eqref{eq:row-wick13} is the definition. By \eqref{eq:row-homog} with $\lambda=\tau$, $\wick_m(\cW_k(U_1))=\wick_m(\tau^k\cW_k(U_1))=\wick_m(\wrk_k)$. By \cref{lem:row-wr} every $\wrk_k$ is analytic on $\Ddisc'$ and $\wrk_2$ has no zero there, so the right side of \eqref{eq:row-wick13} is analytic on $\Ddisc'$. At $\tau=0$, $\wrk_k(0)=(-1)^k\phi_0^{(k)}(\wA/E)$, and \eqref{eq:row-homog} with $\lambda=-1$ followed by (D1) ($\phi_0'(\wA/E)=E-E=0$, $\phi_0''(\wA/E)=-E^3/(\alpha_1\wA)\neq0$) gives $\wick_m(\wrk_k(0))=\wick_m(\phi_0^{(k)}(\wA/E))=\coef{z^m}\Lnorm[\phi_0,1;\wA/E]=\Abetam{m}$. No coordinate change, i.e.\ no use of (P1), occurs here.

\emph{(4) Proof of (iii).} By (D2), $\Nlap_m$ is a polynomial in $1/\phi_\tau''(\sast)$, in the $\phi_\tau^{(k)}(\sast)$ ($k\leq2m+2$), in the derivatives $\de_s^a(1-\tau s)^{-1}$ at $\sast$, which equal $a!\,\tau^aU_1^{-a-1}$ ($a\leq2m$), and in $U_1$. By \eqref{eq:row-phik} all of these are analytic on $\Ddisc'$, since $\sast$ is analytic there, $\sast\notin\{0,1,1/\tau\}$, and $\phi_\tau''(\sast)\neq0$ (\cref{lem:row-hess}). At $\tau=0$, $\sast(0)=\wA/E$ and $(1-0\cdot s)^{-1}=1$, so $\Nlap(z,0)=\Abeta(z)$. By (i) the analytic functions of (ii) and $\Nlap_m$ agree on $\Ddisc''$, and by (ii) and $\Nlap(z,0)=\Abeta(z)$ also at $\tau=0$. Finally $t\mapsto\tau=1-\re^{-t}$ is analytic with $\tau(0)=0$.
\end{proof}

\begin{rmk}[The degenerate coordinate; a uniform version without continuation]
\label{rmk:row-unif}
In $y$ the inputs of the Wick formula blow up at $t=0$ ($\cW_k(U_1)=\tau^{-k}\wrk_k$, $\cW_2(U_1)^{-1}=O(\tau^2)$, $h'(\sast)=O(\tau)$), yet every single Wick monomial is analytic there: a monomial $\cW_2^{-n}\prod\cW_{k_i}$ has $\sum k_i=2n$, so it equals $\wrk_2^{-n}\prod\wrk_{k_i}$; no cancellation between monomials is needed. Geometrically $y$ is $-\tau$ times a coordinate that is good uniformly near $\tau=0$. Let
\[
\tilde y\coloneqq-\frac{\log(1-\tau s)}{\tau}=\sum_{n\geq1}\frac{\tau^{n-1}s^n}{n},\qquad
s=\varsigma_\tau(\tilde y)\coloneqq\frac{1-\re^{-\tau\tilde y}}{\tau}=\sum_{n\geq1}\frac{(-\tau)^{n-1}\tilde y^n}{n!}
\]
($\tilde y=s$ at $\tau=0$). $\varsigma$ is entire in $(\tau,\tilde y)$, $\varsigma_0=\mathrm{id}$, $\dd\varsigma_\tau/\dd\tilde y=\re^{-\tau\tilde y}=1-\tau\varsigma_\tau(\tilde y)$, $U=\re^{-\tau\tilde y}$, i.e.\ $y=-\tau\tilde y$ and $\dd/\dd\tilde y=-\tau\thU$. By \eqref{eq:row-euler1},
\begin{equation}
\label{eq:row-unif}
\cW(\re^{-\tau\tilde y})=\phi_\tau(\varsigma_\tau(\tilde y))+E\log\tau+\alpha_1\log\cQ+\alpha_1\log(-1),
\end{equation}
and the first term on the right is analytic in $(\tau,\tilde y)$ near $(0,\wA/E)$. With $\tilde y_1(\tau)\coloneqq-\log(U_1)/\tau$ (principal branch; $\tilde y_1(0)=\wA/E$, $\varsigma_\tau(\tilde y_1)=\sast$), differentiating \eqref{eq:row-unif} $k\geq2$ times gives
\begin{equation}
\label{eq:row-unif2}
(-\tau)^k\cW_k(U_1)=(-1)^k\wrk_k(\tau)=\frac{\dd^k}{\dd\tilde y^k}\big[\phi_\tau\circ\varsigma_\tau\big](\tilde y_1(\tau)),
\end{equation}
which re-proves \cref{lem:row-wr} near $\tau=0$. For $\tau\neq0$, \cref{lem:row-homog} ($\lambda=-\tau$), \eqref{eq:row-unif} and (D3) give $\rpt_{0,1}(z;t)=\Lnorm[\cW_y,1;y_1]=\Lnorm[\phi_\tau\circ\varsigma_\tau,1;\tilde y_1]$. The right side is defined for every small $\tau$, $\tau=0$ included (second derivative $\wrk_2\neq0$, equal to $-E^3/(\alpha_1\wA)$ at $\tau=0$; coefficients analytic in $\tau$ by (D2)). \Cref{prop:row-P1} with $h=\varsigma_\tau$, $v_0=\tilde y_1$, $u_0=\sast$, $h_1=U_1$ and $\amp=(1-\tau s)^{-1}$, for which $(\amp\circ h)h'=\re^{-\tau\tilde y}/(1-\tau\varsigma_\tau(\tilde y))=1$, gives $\Lnorm[\phi_\tau\circ\varsigma_\tau,1;\tilde y_1]=U_1\Lnorm[\phi_\tau,(1-\tau s)^{-1};\sast]=\Nlap(z,\tau)$ at every small $\tau$, $\tau=0$ included.
\end{rmk}

%-------------------------------------------------------------------------------
\subsection{Gauss's equation for the formal expansion}
\label{sec:row-gauss}

Gauss's parameters and series are
\begin{equation}
\label{eq:row-gaussdef}
\begin{gathered}
\ga\coloneqq1+\frac{\alpha_2}{z},\qquad \gb\coloneqq1+\frac{\wA}{z},\qquad \gc\coloneqq2+\frac Ez,\qquad
\Fhyp(z,\tau)\coloneqq{}_2F_1(\ga,\gb;\gc;\tau)=\sum_{n\geq0}\Fhyp_n(z)\tau^n,\\
\Fhyp_n=\frac{(\ga)_n(\gb)_n}{(\gc)_n\,n!}=\frac{\prod_{j=1}^{n}(\alpha_2+jz)(\wA+jz)}{n!\,z^n\prod_{j=2}^{n+1}(E+jz)}\in\Kal(z),
\end{gathered}
\end{equation}
using $(1+v/z)_n=z^{-n}\prod_{j=1}^n(v+jz)$ and $(2+E/z)_n=z^{-n}\prod_{j=2}^{n+1}(E+jz)$. Gauss's operator, acting on $\Cz\fps{\tau}$ with $\dd/\dd\tau$-derivatives, is
\begin{equation}
\label{eq:row-gaussop}
\begin{gathered}
\Gauss[Y]\coloneqq\tau(1-\tau)Y_{\tau\tau}+\big(\gc-(\ga+\gb+1)\tau\big)Y_\tau-\ga\gb\,Y,\\
\ga+\gb+1=3+\frac{\alpha_2+\wA}{z},\qquad \ga\gb=1+\frac{\alpha_2+\wA}{z}+\frac{\alpha_2\wA}{z^2}.
\end{gathered}
\end{equation}
Further, $\Ephi\coloneqq\mathrm{Exp}[-\dphi/z]=\sum_l(-\dphi)^l/(l!z^l)$, an element of $\bbC[1/z]\fps{\tau}$ ($\tau$-adically convergent since $\dphi=O(\tau)$), with $\Ephi_n\coloneqq\coef{\tau^n}\Ephi$ a polynomial of degree $\leq n$ in $1/z$; $\Ephi^{-1}=\mathrm{Exp}[\dphi/z]=\Ephi(-z)$ and $z\Ephi'=-\dphi'\Ephi$. Finally $\Hcoef_n\coloneqq\coef{\tau^n}(\Fhyp\Ephi)=\sum_{n_1+n_2=n}\Fhyp_{n_1}\Ephi_{n_2}$ and $\Kcoef_n\coloneqq\coef{\tau^n}(\Fhyp\Ephi\alg)$, elements of $\Kal(z)$ (the Taylor coefficients of $\alg$ and $\dphi$ lie in $\Kal$ by (P) of \cref{lem:set-rings}).

\begin{lem}[The amplitude identity; formal]\genusstatus{AI-proposed, unchecked}
\label{lem:row-amp}
Put $f\coloneqq\phi_\tau$, $\amp_0\coloneqq(1-\tau s)^{-1}$ and $\Dam q\coloneqq\de_\tau q+q\,(\de_\tau f)/z$ as in \eqref{eq:row-P2}, so that $\re^{f/z}\Dam q=\de_\tau(\re^{f/z}q)$. Since $\de_\tau f=\alpha_2s/(1-\tau s)$,
\[
\Dam\amp_0=\ga\,s(1-\tau s)^{-2},\qquad \Dam^2\amp_0=\ga(\ga+1)\,s^2(1-\tau s)^{-3},
\]
and, with $\kE\coloneqq s(1-s)(1-\tau s)^{-2}$,
\begin{equation}
\label{eq:row-amp}
\tau(1-\tau)\Dam^2\amp_0+\big(\gc-(\ga+\gb+1)\tau\big)\Dam\amp_0-\ga\gb\,\amp_0=-\ga\,\big(\kE'+f'\kE/z\big).
\end{equation}
This is an identity of Laurent polynomials in $z$ with coefficients rational in $(s,\tau)$, valid identically in $(s,\tau)$, not only at a critical point.
\end{lem}

\begin{proof}
$\de_\tau\amp_0=s(1-\tau s)^{-2}$, so $\Dam\amp_0=s(1-\tau s)^{-2}(1+\alpha_2/z)=\ga s(1-\tau s)^{-2}$; and $\Dam[s(1-\tau s)^{-2}]=2s^2(1-\tau s)^{-3}+(\alpha_2/z)s^2(1-\tau s)^{-3}=(\ga+1)s^2(1-\tau s)^{-3}$. For \eqref{eq:row-amp}, multiply by $(1-\tau s)^3/\ga$. Since $f'\kE=\Ns(s)(1-\tau s)^{-3}$ by \eqref{eq:row-phaseder} and $\kE'=[(1-2s)(1-\tau s)+2\tau s(1-s)](1-\tau s)^{-3}$, the claim becomes
\[
\tau(1-\tau)(\ga+1)s^2+\big(\gc-(\ga+\gb+1)\tau\big)s(1-\tau s)-\gb(1-\tau s)^2=-(1-2s)(1-\tau s)-2\tau s(1-s)-\frac{\Ns(s)}{z},
\]
an identity between Laurent polynomials in $z$ of order $\geq-1$. Coefficient of $1/z$ on the left: $\tau(1-\tau)\alpha_2s^2+\big(E-(\alpha_2+\wA)\tau\big)s(1-\tau s)-\wA(1-\tau s)^2$; expanding,
\begin{align*}
&\tau\alpha_2s^2-\tau^2\alpha_2s^2+Es-E\tau s^2-(\alpha_2+\wA)\tau s+(\alpha_2+\wA)\tau^2s^2-\wA+2\wA\tau s-\wA\tau^2s^2\\
&\quad=-\wA+\big[E+(\wA-\alpha_2)\tau\big]s+(\alpha_2-E)\tau s^2+(-\alpha_2+\alpha_2+\wA-\wA)\tau^2s^2\\
&\quad=-\wA+(E+\alpha_0\tau)s-\wB\tau s^2=-\Ns(s),
\end{align*}
using $\wA-\alpha_2=\alpha_0$ and $E-\alpha_2=\wB$. Coefficient of $z^0$ on the left: $2\tau(1-\tau)s^2+(2-3\tau)s(1-\tau s)-(1-\tau s)^2=2\tau s^2-2\tau^2s^2+2s-3\tau s-2\tau s^2+3\tau^2s^2-1+2\tau s-\tau^2s^2=2s-\tau s-1$; on the right, $-(1-2s)(1-\tau s)-2\tau s(1-s)=-1+\tau s+2s-2\tau s^2-2\tau s+2\tau s^2=2s-\tau s-1$.
\end{proof}

Identity \eqref{eq:row-amp} is Euler's integration by parts for the integrand $s^{\gb-1}(1-s)^{\gc-\gb-1}(1-\tau s)^{-\ga}$, written for amplitudes (\cref{rmk:row-eulerint}).

\begin{prop}[Gauss's equation for the formal expansion; formal]\genusstatus{AI-proposed, unchecked}
\label{prop:row-gaussode}
Let $\Lt\coloneqq\Lnorm[\phi_\tau,(1-\tau s)^{-1};\sast(\tau)]$, an element of $\Oan(\Ddisc')\fps{z}$ by \cref{prop:row-wickrow}(iii), and
\[
\Ngauss(z,\tau)\coloneqq\mathrm{Exp}[\dphi(\tau)/z]\,\Big(\frac{\phi_\tau''(\sast)}{\phi_0''(\wA/E)}\Big)^{-1/2}\iota(\Lt)\quad\in\Cz\fps{\tau}
\]
(root with value $1$ at $\tau=0$; $\iota$ is the Taylor expansion of \cref{lem:set-rings}). Then $\Gauss[\Ngauss]=0$ in $\Cz\fps{\tau}$, $\Ngauss(z,0)=\Abeta(z)$, which lies in $1+z\Kal\fps{z}$, and $\Ngauss=\mathrm{Exp}[\dphi/z]\,\Nlap/\alg$.
\end{prop}

\begin{proof}
Here, in the notation of \cref{prop:row-P2}, $f=\phi_\tau$, $s_c=\sast$, $f_0(\tau)=\phi_\tau(\sast(\tau))$, so $f_{0,\tau}=\dphi_\tau$, and $f_2(\tau)=\phi_\tau''(\sast(\tau))$. Let $\nabla\coloneqq\dd/\dd\tau+\dphi_\tau/z-f_{2,\tau}/(2f_2)$, acting on $\Oan((z))$-valued expressions. \Cref{prop:row-P2}, applied to each $z$-power of an amplitude $q$, says $\nabla\Lnorm[q]=\Lnorm[\Dam q]$ (expansion point $\sast(\tau)$ throughout); applying it twice and using (D3) (linearity over $\bbC[\tau][z,1/z]$),
\[
\big[\tau(1-\tau)\nabla^2+(\gc-(\ga+\gb+1)\tau)\nabla-\ga\gb\big]\Lt=\Lnorm\big[\text{left side of \eqref{eq:row-amp}}\big]=-\ga\,\Lnorm[\kE'+f'\kE/z]=0
\]
by \cref{prop:row-P3}. Let $\mathsf{J}\coloneqq\mathrm{Exp}[\dphi/z](f_2(\tau)/f_2(0))^{-1/2}\in\bbC[1/z]\fps{\tau}$. In $\Cz\fps{\tau}$, $\dd/\dd\tau\,(\mathsf{J}Y)=\mathsf{J}\nabla Y$, because $\mathsf{J}_\tau=\mathsf{J}(\dphi_\tau/z-f_{2,\tau}/(2f_2))$ ($\iota$ commutes with $\dd/\dd\tau$ and with multiplication by convergent functions of $\tau$, and $(\mathrm{Exp}[\dphi/z])_\tau=(\dphi_\tau/z)\mathrm{Exp}[\dphi/z]$). Hence $\Gauss[\mathsf{J}\iota(\Lt)]=\mathsf{J}\,\iota\big([\tau(1-\tau)\nabla^2+\dots]\Lt\big)=0$. At $\tau=0$: $\sast(0)=\wA/E$, $(1-0\cdot s)^{-1}=1$ and $\dphi(0)=0$, so $\Ngauss(z,0)=\Lnorm[\phi_0,1;\wA/E]=\Abeta(z)$, which lies in $1+z\Kal\fps{z}$ by (D5). Finally $\Nlap=U_1\Lt$ and $\alg=U_1(f_2(\tau)/f_2(0))^{1/2}$ by \eqref{eq:row-algdef}, so $\Ngauss=\mathrm{Exp}[\dphi/z]\Nlap/\alg$.
\end{proof}

No boundary term occurs at any stage: the formal functional kills total derivatives exactly (\cref{prop:row-P3}), which is the formal shadow of the vanishing boundary terms of Euler's integration by parts, valid for all weights and all $z$.

%-------------------------------------------------------------------------------
\subsection{Uniqueness and the closed form}
\label{sec:row-closed}

\begin{lem}[Uniqueness; formal]\genusstatus{AI-proposed, unchecked}
\label{lem:row-unique}
If $Y=\sum_kY_k\tau^k\in\Cz\fps{\tau}$ satisfies $\Gauss[Y]=0$, then $Y=Y_0\Fhyp$.
\end{lem}

\begin{proof}
The coefficient of $\tau^k$ in $\Gauss[Y]$ is
\[
(k+1)kY_{k+1}-k(k-1)Y_k+\gc(k+1)Y_{k+1}-(\ga+\gb+1)kY_k-\ga\gb Y_k=(k+1)(k+\gc)Y_{k+1}-(k+\ga)(k+\gb)Y_k .
\]
Now $k+\gc=(E+(k+2)z)/z$ is a unit of $\Cz$ because $E\neq0$ (this is the only place where the parameters enter), and $k+1\neq0$, so $Y_{k+1}=(k+\ga)(k+\gb)Y_k/((k+1)(k+\gc))$, whence $Y_k=Y_0\Fhyp_k$ by induction.
\end{proof}

So the solutions of Gauss's equation in $\Cz\fps{\tau}$ form the line $\Cz\Fhyp$; the second local exponent at $\tau=0$, $1-\gc=-1-E/z$, is not in this ring. (The same recursion shows $\Gauss[\Fhyp]=0$: $\Fhyp_{n+1}/\Fhyp_n=(n+\ga)(n+\gb)/((n+1)(n+\gc))$. The exponents of Gauss's equation at $\tau=0$ are $0$ and $1-\gc$; those of the Picard--Fuchs operator $z^2(\cQ-1)\thQ^2-z(\alpha_0+(\alpha_1+\alpha_2)\cQ)\thQ+\alpha_1\alpha_2\cQ$ at $\cQ=1$ are $0$ and $1+E/z$: only $E$ enters at $\cQ=1$. This is not used.)

\begin{thm}[The closed form; formal]\genusstatus{AI-proposed, unchecked}
\label{thm:row-closed}
In $\Cz\fps{\tau}$,
\begin{equation}
\label{eq:row-closed}
\Nlap(z,\tau)=\Abeta(z)\,\Fhyp(z,\tau)\,\Ephi(z,\tau)\,\alg(\tau).
\end{equation}
\end{thm}

\begin{proof}
By \cref{prop:row-gaussode,lem:row-unique}, $\Ngauss=\Ngauss(z,0)\Fhyp=\Abeta\Fhyp$. Multiply by $\alg\,\Ephi$ ($\alg$ is in $\Oan$, $\Ephi$ is the inverse of $\mathrm{Exp}[\dphi/z]$ in $\bbC[1/z]\fps{\tau}$) and use $\Ngauss=\mathrm{Exp}[\dphi/z]\Nlap/\alg$.
\end{proof}

%-------------------------------------------------------------------------------
\subsection{The Beta constant}
\label{sec:row-beta}

The constant $\Abeta(z)$ of \eqref{eq:row-closed} is the normalised Laplace expansion of a Beta integral. Its evaluation is the one analytic input of this section and of \cref{sec:lea}; it enters through the following lemma and proposition, whose proofs use \cref{lem:row-laplace} and Stirling's series with remainder, and nothing else that is not formal. The lemma is also used in \cref{prop:lea-regrows}, \cref{lem:ide-norm} and \cref{prop:rig-flop}. Recall from \eqref{eq:set-h} that for a finite multiset $\mathsf{v}=(v_1,\dots,v_r)$ of non-zero elements, $h_{\mathsf{v}}(z)\coloneqq\mathrm{Exp}\big[-\sum_{k\geq1}\bk_kz^{2k-1}\sum_av_a^{1-2k}\big]$ with $\bk_k=B_{2k}/(2k(2k-1))$, the $B_{2k}$ being the Bernoulli numbers.

\begin{lem}[The Beta phase; analytic input]\genusstatus{AI-proposed, unchecked}
\label{lem:row-betaphase}
Let $a$, $b$ be indeterminates, or complex numbers with $ab(a+b)\neq0$, and let $f(y)\coloneqq ay+b\log(1-\re^y)$, the phase $a\log U+b\log(1-U)$ in the coordinate $y=\log U$. Then $f$ has a critical point at $y_0\coloneqq\log(a/(a+b))$, i.e.\ at $U=a/(a+b)$, with $f''(y_0)=-a(a+b)/b\neq0$, and
\begin{equation}
\label{eq:row-betaphase}
\Lnorm[f,1;y_0]=h_{(-a,-b,a+b)}(z)\qquad\text{in }\bbQ(a,b)\fps{z}.
\end{equation}
\end{lem}

\begin{proof}
\emph{Critical point and rationality.} $f'(y)=a-b\,\re^y/(1-\re^y)$ vanishes if and only if $\re^y/(1-\re^y)=a/b$, i.e.\ $\re^y=a/(a+b)$; and $f''(y)=-b\,\re^y/(1-\re^y)^2$, which at $y_0$ equals $-b\,(a/(a+b))\,((a+b)/b)^2=-a(a+b)/b$. For $k\geq2$, by the computation of \cref{sec:row-wr} with $\thU=\dd/\dd y$, $f^{(k)}(y_0)=-b\,\Pol_{k-1}(p)$ with $p=1/(1-\re^{y_0})=(a+b)/b$; so $1/f''(y_0)$ and all $f^{(k)}(y_0)$ lie in $\bbQ(a,b)$, and by (D2) every $z$-coefficient of the left side of \eqref{eq:row-betaphase} lies in $\bbQ(a,b)$. The $z$-coefficients of the right side are polynomials in $1/a$, $1/b$, $1/(a+b)$ with rational coefficients. A rational function of $(a,b)$ that vanishes on the open quadrant $a,b>0$ vanishes identically, so it suffices to prove \eqref{eq:row-betaphase} coefficientwise for all real $a,b>0$. Fix such $a,b$.

\emph{Laplace.} On $I\coloneqq(-\infty,0)$ the function $f$ is real-analytic and strictly concave ($f''<0$), and it tends to $-\infty$ at both ends ($ay\to-\infty$ as $y\to-\infty$, $b\log(1-\re^y)\to-\infty$ as $y\to0^-$); so $y_0$ is its unique maximum, the supremum of $f$ off $(y_0-e,y_0+e)$ is $\max(f(y_0-e),f(y_0+e))<f(y_0)$, and $\int_I\re^{f}\,\dd y=\int_0^1U^{a}(1-U)^{b}\,\dd U/U=\Bfn(a,b+1)<\infty$ (substitute $U=\re^y$, $\dd y=\dd U/U$). For real $z>0$ the same substitution and \cite[5.12.1]{DLMF} give
\[
\int_I\re^{f(y)/z}\,\dd y=\int_0^1U^{a/z}(1-U)^{b/z}\,\frac{\dd U}{U}=\Bfn\Big(\frac az,\frac bz+1\Big).
\]
With $f^\ast\coloneqq f(y_0)=a\log a+b\log b-(a+b)\log(a+b)$ and $-f''(y_0)=a(a+b)/b$, \cref{lem:row-laplace} (amplitude $1$, $z_0=1$) gives, for every $M$, as $z\to0^+$,
\begin{equation}
\label{eq:row-betalap}
\Bfn\Big(\frac az,\frac bz+1\Big)=\re^{f^\ast/z}\Big(\frac{2\pi z\,b}{a(a+b)}\Big)^{1/2}\Big[\sum_{m<M}\big(\coef{z^m}\Lnorm[f,1;y_0]\big)z^m+O(z^M)\Big].
\end{equation}

\emph{Stirling.} $\Gamma(u+1)=u\Gamma(u)$ gives $\Bfn(p,q+1)=\Bfn(p,q)\,q/(p+q)$; for $p=a/z$, $q=b/z$ the factor is $b/(a+b)$, independent of $z$. Stirling's series with remainder \cite[5.11.1]{DLMF} states that for $u\to+\infty$ and every $K\geq1$
\begin{equation}
\label{eq:row-stirling}
\log\Gamma(u)=\Big(u-\frac12\Big)\log u-u+\frac12\log(2\pi)+\sum_{k=1}^{K}\bk_k\,u^{1-2k}+O(u^{-2K-1}),\qquad \bk_k=\frac{B_{2k}}{2k(2k-1)};
\end{equation}
for real positive $u$ the remainder after the term $k=K$ is bounded in absolute value by the first neglected term $\bk_{K+1}u^{-2K-1}$ and has its sign \cite[Section~5.11(ii)]{DLMF}. Apply \eqref{eq:row-stirling} to
\[
\log\Bfn\Big(\frac az,\frac bz\Big)=\log\Gamma\Big(\frac az\Big)+\log\Gamma\Big(\frac bz\Big)-\log\Gamma\Big(\frac{a+b}z\Big),
\]
all three arguments tending to $+\infty$. The terms $-u$ cancel: $-a/z-b/z+(a+b)/z=0$. In the terms $(u-\frac12)\log u$, written as $(u-\frac12)(\log w-\log z)$ for $u=w/z$, the coefficient of $\log z$ is $-(a/z-\frac12)-(b/z-\frac12)+((a+b)/z-\frac12)=\frac12$, and the remaining terms give
\[
\frac{a\log a+b\log b-(a+b)\log(a+b)}{z}-\frac12\big(\log a+\log b-\log(a+b)\big)=\frac{f^\ast}{z}+\frac12\log\frac{a+b}{ab}.
\]
The three constants $\frac12\log(2\pi)$ leave one $\frac12\log(2\pi)$. In the tails, $u^{1-2k}=z^{2k-1}w^{1-2k}$ for $u=w/z$, so they give $\sum_{k\leq K}\bk_kz^{2k-1}\big(a^{1-2k}+b^{1-2k}-(a+b)^{1-2k}\big)+O(z^{2K+1})$. Since $1-2k$ is odd,
\[
a^{1-2k}+b^{1-2k}-(a+b)^{1-2k}=-\big[(-a)^{1-2k}+(-b)^{1-2k}+(a+b)^{1-2k}\big],
\]
so the tail is the truncation to $k\leq K$ of $-\sum_k\bk_kz^{2k-1}\sum_vv^{1-2k}$ over the multiset $(-a,-b,a+b)$, that is, of $\log h_{(-a,-b,a+b)}(z)$. Hence
\[
\Bfn\Big(\frac az,\frac bz\Big)=\re^{f^\ast/z}\Big(\frac{2\pi z\,(a+b)}{ab}\Big)^{1/2}\mathrm{Exp}\Big[\sum_{k=1}^{K}\bk_kz^{2k-1}\big(a^{1-2k}+b^{1-2k}-(a+b)^{1-2k}\big)+O(z^{2K+1})\Big],
\]
and the exponential of (a polynomial in $z$ without constant term) $+\,O(z^{2K+1})$ is the truncation of the formal exponential plus $O(z^{2K+1})$. Multiplying by $b/(a+b)$ and using $\big(2\pi z(a+b)/(ab)\big)^{1/2}\,b/(a+b)=\big(2\pi z\,b/(a(a+b))\big)^{1/2}$,
\[
\Bfn\Big(\frac az,\frac bz+1\Big)=\re^{f^\ast/z}\Big(\frac{2\pi z\,b}{a(a+b)}\Big)^{1/2}\Big[\big(h_{(-a,-b,a+b)}(z)\big)_{<M}+O(z^M)\Big]
\]
for every $M$ (take $2K+1\geq M$).

\emph{Uniqueness.} The prefactors in this display and in \eqref{eq:row-betalap} are the same positive function of $z$, so the two power series are asymptotic expansions as $z\to0^+$ of the same function; asymptotic power series are unique, so they coincide coefficientwise.
\end{proof}

\draftnote{The Stirling step, which the source gives in condensed form (``as in the proof of the Beta constant''), is written out here: the series \eqref{eq:row-stirling} for $\log\Gamma(a/z)+\log\Gamma(b/z)-\log\Gamma((a+b)/z)$, its tail $\sum_k\bk_kz^{2k-1}(a^{1-2k}+b^{1-2k}-(a+b)^{1-2k})$, and the identification of this tail with the exponent of $h_{(-a,-b,a+b)}$ through the oddness of $1-2k$. To be checked by the authors.}

\begin{prop}[The Beta constant; analytic input]\genusstatus{AI-proposed, unchecked}
\label{prop:row-beta}
\begin{equation}
\label{eq:row-beta}
\Abeta(z)=\Lnorm[\phi_0,1;\wA/E]=h_{(E,-\alpha_1,-\wA)}(z)\,\frac{E}{E+z}\qquad\text{in }\bbQ(\alpha_1,\wA)\fps{z}.
\end{equation}
In words: the normalised Laplace expansion of the Beta integral $\Bfn(\wA/z+1,\alpha_1/z+1)=\int_0^1s^{\wA/z}(1-s)^{\alpha_1/z}\,\dd s=\int_0^1\re^{\phi_0(s)/z}\,\dd s$ is the Stirling tail $h_{(E,-\alpha_1,-\wA)}$ of $\Bfn(\wA/z,\alpha_1/z)$ times the elementary factor $E/(E+z)$ coming from $\Bfn(p+1,q+1)=\Bfn(p,q)\,pq/((p+q)(p+q+1))$.
\end{prop}

\begin{proof}
The proof is the one of \cref{lem:row-betaphase}, for the phase $\phi_0$ in the variable $s$ with amplitude $1$.

\emph{(1) Reduction to $\alpha_1,\wA>0$.} Both sides have $z$-coefficients in $\bbQ(\alpha_1,\wA)$: the left by (D5), the right by definition ($E/(E+z)=\sum_n(-z/E)^n$ with $E=\alpha_1+\wA$). Two rational functions of $(\alpha_1,\wA)$ that agree on the open quadrant are equal. Fix real $\alpha_1,\wA>0$.

\emph{(2) Laplace.} For real $z>0$, $\Bfn(\wA/z+1,\alpha_1/z+1)=\int_0^1s^{\wA/z}(1-s)^{\alpha_1/z}\,\dd s=\int_0^1\re^{\phi_0(s)/z}\,\dd s$. On $I=(0,1)$, $\phi_0=\wA\log s+\alpha_1\log(1-s)$ is real-analytic, $\phi_0''=-\wA/s^2-\alpha_1/(1-s)^2<0$, $\phi_0'=(\wA-Es)/(s(1-s))$ vanishes only at $\wA/E$, and $\phi_0\to-\infty$ at both ends; so $\phi_0$ is strictly concave with its unique maximum at $\wA/E$, its supremum off $(\wA/E-e,\wA/E+e)$ is $\max(\phi_0(\wA/E\pm e))<\phi_0(\wA/E)$, and $\int_0^1\re^{\phi_0}\,\dd s=\Bfn(\wA+1,\alpha_1+1)<\infty$. \Cref{lem:row-laplace} (amplitude $1$, $z_0=1$), with $-\phi_0''(\wA/E)=E^3/(\alpha_1\wA)$ by \eqref{eq:row-hess} and $\phi_0(\wA/E)=\wA\log\wA+\alpha_1\log\alpha_1-E\log E=:f^\ast$, gives, as $z\to0^+$,
\[
\Bfn\Big(\frac{\wA}z+1,\frac{\alpha_1}z+1\Big)=\re^{f^\ast/z}\Big(\frac{2\pi z\,\alpha_1\wA}{E^3}\Big)^{1/2}\Big[\sum_{m<M}\Abetam{m}z^m+O(z^M)\Big].
\]

\emph{(3) Stirling.} $\Gamma(u+1)=u\Gamma(u)$ gives $\Bfn(p+1,q+1)=\Bfn(p,q)\,pq/((p+q)(p+q+1))$; for $p=\wA/z$, $q=\alpha_1/z$ this factor is $(\wA\alpha_1/z^2)/\big((E/z)((E+z)/z)\big)=\wA\alpha_1/(E(E+z))$. The Stirling computation in the proof of \cref{lem:row-betaphase}, with $(a,b)=(\wA,\alpha_1)$ and $a+b=E$, gives
\[
\Bfn\Big(\frac{\wA}z,\frac{\alpha_1}z\Big)=\re^{f^\ast/z}\Big(\frac{2\pi z\,E}{\alpha_1\wA}\Big)^{1/2}\mathrm{Exp}\Big[\sum_{k=1}^{K}\bk_kz^{2k-1}\big(\wA^{1-2k}+\alpha_1^{1-2k}-E^{1-2k}\big)+O(z^{2K+1})\Big],
\]
whose exponent is the truncation of $\log h_{(E,-\alpha_1,-\wA)}(z)=-\sum_k\bk_kz^{2k-1}\big(E^{1-2k}-\alpha_1^{1-2k}-\wA^{1-2k}\big)$. Multiplying by $\wA\alpha_1/(E(E+z))$, which is analytic at $z=0$, and using $\big(2\pi zE/(\alpha_1\wA)\big)^{1/2}\,\wA\alpha_1/E=\big(2\pi z\,\alpha_1\wA/E^3\big)^{1/2}\,E$,
\[
\Bfn\Big(\frac{\wA}z+1,\frac{\alpha_1}z+1\Big)=\re^{f^\ast/z}\Big(\frac{2\pi z\,\alpha_1\wA}{E^3}\Big)^{1/2}\Big[\Big(h_{(E,-\alpha_1,-\wA)}(z)\frac{E}{E+z}\Big)_{<M}+O(z^M)\Big].
\]

\emph{(4)} The prefactors in (2) and (3) are the same positive function of $z$, so the two power series are asymptotic expansions of the same function as $z\to0^+$; asymptotic power series are unique, so they coincide.
\end{proof}

\begin{rmk}
\label{rmk:row-betaroutes}
(i) \Cref{prop:row-beta} is not an instance of \cref{lem:row-betaphase}: the lemma expands the phase $\wA\log U+\alpha_1\log(1-U)$ in $y=\log U$ with amplitude $1$ in $y$, i.e.\ with amplitude $1/U$ in $U$ (the measure $\dd U/U$), whereas $\Abeta$ is the expansion of $\phi_0$ in $s$ with amplitude $1$ (the measure $\dd s$); by \cref{prop:row-P1}(iii) the two normalised expansions differ by this amplitude. Comparing the two results, $\Lnorm[\phi_0,1;\wA/E]=\Lnorm[\phi_0\circ\exp,1;\log(\wA/E)]\cdot E/(E+z)$, the factor being $\Bfn(\wA/z+1,\alpha_1/z+1)/\Bfn(\wA/z,\alpha_1/z+1)=(\wA/z)/((E+z)/z)$ times $E/\wA$. \draftnote{This relation between the two expansions is a consequence of the two statements as proved, not a step of either proof; the sources do not state it directly, which is why the proposition is proved separately rather than derived from the lemma.}

(ii) Equivalently, the proposition at the pair $(\wA,\alpha_1)$ is the rescaling, by (D4), of its own case $\wA=1$: with $\xc\coloneqq\alpha_1/\wA$ ($\xc\neq0,-1$ by (ND)) and $\mathsf{ph}(s)\coloneqq\log s+\xc\log(1-s)$, whose critical point is $1/(1+\xc)=\wA/E$, one has $\phi_0=\wA\,\mathsf{ph}$, so $\Abeta(z)=\Lnorm[\mathsf{ph},1;1/(1+\xc)](z/\wA)$. The case $\wA=1$ of \eqref{eq:row-beta} reads $\Lnorm[\mathsf{ph},1;1/(1+\xc)](z')=h_{(1+\xc,-\xc,-1)}(z')\,(1+\xc)/(1+\xc+z')$, and at $z'=z/\wA$ the exponent of $h$ becomes $-\sum_k\bk_k(z/\wA)^{2k-1}\sum_aw_a^{1-2k}=-\sum_k\bk_kz^{2k-1}\sum_a(\wA w_a)^{1-2k}$ with $\wA\cdot(1+\xc,-\xc,-1)=(E,-\alpha_1,-\wA)$, while $(1+\xc)/(1+\xc+z/\wA)=E/(E+z)$. Rationality in the weights (step (1)) is what replaces, at generic weights, the passage from positive to arbitrary parameters.
\end{rmk}

\begin{example}[The coefficient of $z$, by hand]
\label{ex:row-beta1}
This is an illustration of \eqref{eq:row-beta}, not part of its proof. By (D1), $\Abetam{1}=f_4/(8f_2^2)-5f_3^2/(24f_2^3)$ with $f_k\coloneqq\phi_0^{(k)}(\wA/E)$ (the terms $n=2$ with one factor $f_4v^4/(4!z)$ and $n=3$ with two factors $f_3v^3/(3!z)$). From \eqref{eq:row-phik}, $f_2=-E^3/(\alpha_1\wA)$, $f_3=2E^3(\alpha_1^2-\wA^2)/(\alpha_1\wA)^2$, $f_4=-6E^4(\alpha_1^3+\wA^3)/(\alpha_1\wA)^3$; so $f_4/(8f_2^2)=-3(\alpha_1^3+\wA^3)/(4\alpha_1\wA E^2)$, $5f_3^2/(24f_2^3)=-5(\alpha_1^2-\wA^2)^2/(6\alpha_1\wA E^3)$, and, with $\alpha_1^3+\wA^3=E(\alpha_1^2-\alpha_1\wA+\wA^2)$ and $(\alpha_1^2-\wA^2)^2=E^2(\alpha_1-\wA)^2$,
\[
\Abetam{1}=\frac{-9(\alpha_1^2-\alpha_1\wA+\wA^2)+10(\alpha_1-\wA)^2}{12\,\alpha_1\wA E}=\frac{\alpha_1^2-11\alpha_1\wA+\wA^2}{12\,\alpha_1\wA E}.
\]
On the right of \eqref{eq:row-beta}, with $\bk_1=1/12$, $\coef{z^1}\big(h_{(E,-\alpha_1,-\wA)}E/(E+z)\big)=\frac1{12}\big(\frac1{\alpha_1}+\frac1{\wA}-\frac1E\big)-\frac1E=\frac{E^2-13\alpha_1\wA}{12\,\alpha_1\wA E}=\frac{\alpha_1^2-11\alpha_1\wA+\wA^2}{12\,\alpha_1\wA E}$.
\end{example}

%-------------------------------------------------------------------------------
\subsection{Regularity, the frame-B row, and (F2)}
\label{sec:row-main}

\begin{cor}[Regularity of the $\tau$-coefficients; formal given \eqref{eq:row-beta}]\genusstatus{AI-proposed, unchecked}
\label{cor:row-reg}
For every $k\geq0$, $\Kcoef_k=\coef{\tau^k}(\Fhyp\Ephi\alg)$ and $\Hcoef_k=\coef{\tau^k}(\Fhyp\Ephi)$ are rational functions of $z$ with coefficients in $\Kal$, without pole at $z=0$.
\end{cor}

\begin{proof}
$\Kcoef_k=\sum_{i+j+l=k}\Fhyp_i\Ephi_j\coef{\tau^l}\alg$ is a finite sum of products of elements of $\Kal(z)$, $\Kal[1/z]$ and $\Kal$ ($\alg$ and $\dphi$ have Taylor coefficients in $\Kal$), so it lies in $\Kal(z)$. By \eqref{eq:row-closed}, $\Kcoef_k=\coef{\tau^k}(\Nlap/\Abeta)$ in $\Cz$: $\Abeta$ is a unit of $\bbC\fps{z}$ (constant term $1$) and $\Nlap$ lies in $\bbC\fps{z}\fps{\tau}$ (\cref{prop:row-wickrow}(iii)), so the Laurent expansion of $\Kcoef_k$ at $z=0$ has no negative powers. Only $\Abeta\in1+z\bbC\fps{z}$ is used here, not its value. Since $\alg\in1+\tau\Kal\fps{\tau}$ does not depend on $z$, $\Hcoef_k=\sum_{j\leq k}\Kcoef_j\coef{\tau^{k-j}}(1/\alg)$ is a $\Kal$-linear combination of $\Kcoef_0,\dots,\Kcoef_k$.
\end{proof}

\paragraph{The constant line.} By \cref{def:set-rows} the frame-B identity row is $\rr_{0,1}=h_{(\alpha_1,\wA,-E)}\rpt_{0,1}=\hw^{-1}\hnw\rpt_{0,1}$, where the equality of the two scalars is the multiset identity $\hnw\,h_{(E,-\alpha_1,-\wA)}=\hw$ of \cref{lem:set-hmult}: as multisets, $\nw+(E,-\alpha_1,-\wA)=(\alpha_1,\wA,-\epsilon_1,-\epsilon_2,E,-\alpha_1,-\wA)=(E,-\epsilon_1,-\epsilon_2)+(\alpha_1,-\alpha_1)+(\wA,-\wA)$, and pairs $v,-v$ contribute nothing to the odd power sums in the exponent. Equivalently $\hw^{-1}\hnw\,h_{(E,-\alpha_1,-\wA)}=1$. This is where the class $\bw=(E,-\epsilon_1,-\epsilon_2)$ is forced: by \cref{lem:set-hmult} it is the unique class, up to pairs $v,-v$, for which the frame-B constant of the singular row is the geometric series $E/(E+z)$. The same scalar $\hnw$ multiplies both idempotents (\cref{lem:ide-norm}), so the frame-B rows of the regular idempotent carry the non-trivial constant $h_{(\alpha_1,\wA,-E)}\rpt_{0,2}(z;0)$ (\cref{prop:lea-regrows}).

\begin{thm}[The frame-B singular row; formal given \eqref{eq:row-beta}]\genusstatus{AI-proposed, unchecked}
\label{thm:row-main}
In $\bbC\fps{z}\fps{t}$, and coefficientwise in $z$ as analytic functions of $t$ near $0$,
\begin{equation}
\label{eq:row-main}
\rr_{0,1}(z;t)=\frac{E}{E+z}\,\Fhyp(z,\tau)\,\Ephi(z,\tau)\,\alg(\tau),\qquad \tau=1-\re^{-t}.
\end{equation}
Precisely: for every $k$, $\coef{t^k}\rr_{0,1}(z;t)=\sum_{j\leq k}\big(\coef{t^k}\tau^j\big)\frac{E}{E+z}\Kcoef_j(z)$, a rational function of $z$ regular at $z=0$, whose Taylor coefficients in $z$ are the $t^k$ Taylor coefficients of the $z$-coefficients of $\rr_{0,1}$.
\end{thm}

\begin{proof}
By \cref{def:set-rows} and the constant line, $\rr_{0,1}=\hw^{-1}\hnw\rpt_{0,1}$, with a $t$-independent scalar; by \cref{prop:row-wickrow} this equals $\hw^{-1}\hnw\Nlap$; by \eqref{eq:row-closed} it equals $\hw^{-1}\hnw\Abeta\Fhyp\Ephi\alg$; and by \cref{prop:row-beta} and the constant line, $\hw^{-1}\hnw\Abeta=\hw^{-1}\hnw h_{(E,-\alpha_1,-\wA)}\,E/(E+z)=E/(E+z)$. The passage from $\tau$ to $t$ is the unitriangular substitution (T) of \cref{lem:set-rings} and preserves regularity. Coefficientwise meaning: the right side of \eqref{eq:row-main} is a priori an element of $\Cz\fps{\tau}$; by \cref{cor:row-reg} it lies in $\bbC\fps{z}\fps{\tau}$, and the chain above says that it equals $\iota(\hw^{-1}\hnw\Nlap)$ with $\hw^{-1}\hnw\Nlap\in\Oan(\Ddisc')\fps{z}$ (\cref{prop:row-wickrow}(iii)). So each $z$-coefficient of the right side is the Taylor series of a function analytic on $\Ddisc'$, and its sum is the corresponding $z$-coefficient of $\rr_{0,1}$. By \eqref{eq:row-transport} below the right side also equals $(E/(E+z))\,\Phia\,\alg$ with $\Phia\in\Oan(\Ddisc)\fps{z}$.
\end{proof}

\begin{cor}[(F2) at $t=0$]\genusstatus{AI-proposed, unchecked}
\label{cor:row-F2}
$\rr_{0,1}(z;0)=E/(E+z)$. Also $\rpt_{0,1}(z;0)=\Abeta(z)=h_{(E,-\alpha_1,-\wA)}(z)\,E/(E+z)$, and for the row $\rZ_{0,1}=\hnw\rpt_{0,1}$ with the $\cQ=0$ constant, $\rZ_{0,1}(z;0)=h_{(E,-\epsilon_1,-\epsilon_2)}(z)\,E/(E+z)=\hw(z)\,E/(E+z)$.
\end{cor}

\begin{proof}
$t=0$ means $\tau=0$, and $\Fhyp(z,0)=1$, $\dphi(0)=0$, $\alg(0)=1$. The value of $\rpt_{0,1}(z;0)$ is \cref{prop:row-wickrow}(ii) with \cref{prop:row-beta}, and the last statement is the constant line.
\end{proof}

%-------------------------------------------------------------------------------
\subsection{A purely formal proof of the regularity: eikonal and transport equations}
\label{sec:row-eik}

This subsection uses neither the Laplace identification (\cref{sec:row-laplace,sec:row-wickrow,sec:row-beta}) nor the value of $\Abeta$; it shows directly that the $z$-expansion of $\Ephi\Fhyp$ exists and has coefficients analytic in $\tau$, on the whole disc $\Ddisc$. Everything here is formal apart from the analyticity of the explicit algebraic functions $\delta$, $\sast$, $\Lv_2$ on $\Ddisc$.

\begin{lem}[Eikonal and transport decomposition; formal]\genusstatus{AI-proposed, unchecked}
\label{lem:row-eik}
Write $\Fhyp=\mathrm{Exp}[\Sph/z]\,\Phia$ with $\Sph\in\tau\Oan$. With $\gc=E/z+2$, $\ga+\gb+1=(\alpha_2+\wA)/z+3$ and $\ga\gb=\alpha_2\wA/z^2+(\alpha_2+\wA)/z+1$, as operators on $\Cz\fps{\tau}$,
\begin{equation}
\label{eq:row-eik}
\mathrm{Exp}[-\Sph/z]\,\Gauss\big[\mathrm{Exp}[\Sph/z]\,\Phia\big]=z^{-2}\,\Eik(\Spht)\,\Phia+z^{-1}\big(\Ptr\,\Phia_\tau+\Rtr\,\Phia\big)+\Ttr\,\Phia,
\end{equation}
where
\begin{align*}
\Eik(p)&\coloneqq\tau(1-\tau)p^2+\big(E-(\alpha_2+\wA)\tau\big)p-\alpha_2\wA, &
\Ptr&\coloneqq2\tau(1-\tau)\Spht+E-(\alpha_2+\wA)\tau,\\
\Rtr&\coloneqq\tau(1-\tau)\Sphtt+(2-3\tau)\Spht-(\alpha_2+\wA), &
\Ttr&\coloneqq\tau(1-\tau)\frac{\dd^2}{\dd\tau^2}+(2-3\tau)\frac{\dd}{\dd\tau}-1 .
\end{align*}
Moreover:
\begin{enumerate}
\item The eikonal equation $\Eik(p)=0$ is quadratic in $p$ with discriminant $\delta^2$; its root regular at $\tau=0$ has $p(0)=\alpha_2\wA/E$, the other root is $-E/\tau+O(1)$ (exponent $1-\gc$).
\item $\Sph\coloneqq\dphi$ solves the eikonal equation: $\Eik(\dphi_\tau)=0$, with $\dphi_\tau(0)=\alpha_2\wA/E$.
\item For $\Sph=\dphi$: $\Ptr=\delta$, and
\[
\Rtr=\alpha_0-2\wB\sast-\frac{\wB\tau\sast\Lv_2}{\delta},\qquad \Rtr(0)=-\frac{2\alpha_1\alpha_2+\alpha_0E}{E},
\]
analytic on $\Ddisc$. With $\Eik=0$, $z\,\mathrm{Exp}[-\Sph/z]\Gauss[\mathrm{Exp}[\Sph/z]\Phia]=\delta\,\Phia_\tau+\Rtr\,\Phia+z\,\Ttr\,\Phia$.
\end{enumerate}
\end{lem}

\begin{proof}
\emph{Derivation of \eqref{eq:row-eik}.} $\Fhyp_\tau=\re^{\Sph/z}(\Phia_\tau+\Spht\Phia/z)$ and $\Fhyp_{\tau\tau}=\re^{\Sph/z}(\Phia_{\tau\tau}+2\Spht\Phia_\tau/z+\Sphtt\Phia/z+\Spht^2\Phia/z^2)$. Collect the powers of $z$ in $\re^{-\Sph/z}\Gauss[\re^{\Sph/z}\Phia]$ with \eqref{eq:row-gaussop}: the coefficient of $z^{-2}$ is $\tau(1-\tau)\Spht^2+(E-(\alpha_2+\wA)\tau)\Spht-\alpha_2\wA$ times $\Phia$; that of $z^{-1}$ is $2\tau(1-\tau)\Spht\Phia_\tau+\tau(1-\tau)\Sphtt\Phia+(E-(\alpha_2+\wA)\tau)\Phia_\tau+(2-3\tau)\Spht\Phia-(\alpha_2+\wA)\Phia$; that of $z^0$ is $\tau(1-\tau)\Phia_{\tau\tau}+(2-3\tau)\Phia_\tau-\Phia$.

\emph{(i)} The discriminant of $\Eik$ in $p$ is
\[
\big(E-(\alpha_2+\wA)\tau\big)^2+4\alpha_2\wA\tau(1-\tau)=E^2-\big[2E(\alpha_2+\wA)-4\alpha_2\wA\big]\tau+(\wA-\alpha_2)^2\tau^2=\delta^2,
\]
because $\wA-\alpha_2=\alpha_0$ and $2E(\alpha_2+\wA)-4\alpha_2\wA=2E(E-\alpha_1+\alpha_2)-4\alpha_2(E-\alpha_1)=2E(E-\alpha_1-\alpha_2)+4\alpha_1\alpha_2=2\alpha_0E+4\alpha_1\alpha_2=E^2+\uq$, and $\delta^2=E^2-(E^2+\uq)\tau+\alpha_0^2\tau^2$ (\cref{def:set-alg}). At $\tau=0$ the equation reads $Ep-\alpha_2\wA=0$, so the regular root has $p(0)=\alpha_2\wA/E$; the other root is $-E/\tau+O(1)$.

\emph{(ii)} With $s=\sast$ and $p\coloneqq\dphi_\tau=(\wA-\wB s)/(1-\tau)$ (\eqref{eq:row-dphider}),
\begin{align*}
(1-\tau)\Eik(p)&=\tau(\wA-\wB s)^2+\big(E-(\alpha_2+\wA)\tau\big)(\wA-\wB s)-\alpha_2\wA(1-\tau)\\
&=\wB^2\tau s^2+\wB\tau s(\alpha_2+\wA-2\wA)+\wA\tau(\wA-\alpha_2-\wA+\alpha_2)+\wA(E-\alpha_2)-E\wB s\\
&=\wB\big[\wB\tau s^2-\alpha_0\tau s-Es+\wA\big]=\wB\,\Ns(s)=0,
\end{align*}
using $\alpha_2-\wA=-\alpha_0$ and $E-\alpha_2=\wB$. And $p(0)=\wA-\wB\wA/E=\wA(E-\wB)/E=\wA\alpha_2/E$.

\emph{(iii)} $\Ptr=(\dd\Eik/\dd p)(\Spht)$, so $\Ptr^2$ is the discriminant $\delta^2$; directly, $\Ptr=2\tau(\wA-\wB\sast)+E-(\alpha_2+\wA)\tau=E+(\wA-\alpha_2)\tau-2\wB\tau\sast=E+\alpha_0\tau-2\wB\tau\sast=\delta$ by identity (i) of \cref{lem:set-N}. For $\Rtr$: differentiating $(1-\tau)\Spht=\wA-\wB\sast$ gives $-\Spht+(1-\tau)\Sphtt=-\wB\sast_\tau$, i.e.\ $\tau(1-\tau)\Sphtt=\tau\Spht-\wB\tau\sast_\tau$, so
\[
\Rtr=2(1-\tau)\Spht-\wB\tau\sast_\tau-(\alpha_2+\wA)=2(\wA-\wB\sast)-\alpha_2-\wA-\wB\tau\sast_\tau=\alpha_0-2\wB\sast-\frac{\wB\tau\sast\Lv_2}{\delta},
\]
using $\sast_\tau=\sast\Lv_2/\delta$ (identity (iii) of \cref{lem:set-N}). At $\tau=0$, $\Rtr(0)=\alpha_0-2\wB\wA/E=(\alpha_0E-2\wA\wB)/E=-(2\alpha_1\alpha_2+\alpha_0E)/E$, since $\wA\wB=\alpha_1\alpha_2+\alpha_0E$. On $\Ddisc$ the functions $\sast$, $\Lv_2=\wB\sast-\alpha_0$ and $1/\delta$ are analytic, so $\Rtr$ is analytic on $\Ddisc$. The last display is \eqref{eq:row-eik} multiplied by $z$ with $\Eik(\Spht)=0$ and $\Ptr=\delta$.
\end{proof}

\begin{prop}[The transport equations; formal]\genusstatus{AI-proposed, unchecked}
\label{prop:row-transport}
For $\Phia=\sum_{m\geq0}z^m\Phia_m$, the transport equations
\[
\delta\,(\Phia_m)_\tau+\Rtr\,\Phia_m+\Ttr\,\Phia_{m-1}=0\qquad(m\geq0,\ \Phia_{-1}\coloneqq0)
\]
have unique solutions analytic on $\Ddisc$ with $\Phia_0(0)=1$ and $\Phia_m(0)=0$ for $m\geq1$, namely
\[
\Phia_0=\mathrm{Exp}\Big[-\int_0^\tau\frac{\Rtr}{\delta}\Big]=\frac1{\alg},\qquad
\Phia_m(\tau)=-\Phia_0(\tau)\int_0^\tau\frac{\Ttr[\Phia_{m-1}](u)}{\delta(u)\Phia_0(u)}\,\dd u,
\]
and in $\Cz\fps{\tau}$
\begin{equation}
\label{eq:row-transport}
\Ephi\,\Fhyp=\Phia,\qquad \Phia\in\Oan(\Ddisc)\fps{z},\qquad \coef{z^0}\Phia=\frac1{\alg}.
\end{equation}
\end{prop}

\begin{proof}
On $\Ddisc$ the function $\delta$ has no zero and $\Rtr$ is analytic (\cref{lem:row-eik}(iii)); so the first-order linear equations have, inductively in $m$, unique analytic solutions on the (simply connected) disc with the prescribed values at $0$, given by the displayed integrals. By \eqref{eq:row-eik} with $\Eik(\dphi_\tau)=0$ and $\Ptr=\delta$,
\[
z\,\mathrm{Exp}[-\dphi/z]\,\Gauss\big[\mathrm{Exp}[\dphi/z]\iota(\Phia)\big]=\sum_mz^m\big(\delta(\Phia_m)_\tau+\Rtr\Phia_m+\Ttr\Phia_{m-1}\big)=0,
\]
so $\tilde Y\coloneqq\mathrm{Exp}[\dphi/z]\iota(\Phia)$ satisfies $\Gauss[\tilde Y]=0$, and $\tilde Y(z,0)=\Phia(z,0)=1$; hence $\tilde Y=\Fhyp$ by \cref{lem:row-unique}, i.e.\ $\Ephi\Fhyp=\Phia$. It remains to identify $\Phia_0$ with $1/\alg$, i.e.\ to show $\Rtr=\delta(\log\alg)_\tau$: both $\Phia_0$ and $1/\alg$ are analytic on $\Ddisc$ with value $1$ at $0$ (for $1/\alg$ by \cref{lem:row-hess}(c) through \eqref{eq:row-algp}), and $\Phia_0$ is the unique such solution of $\delta\Phia_0'+\Rtr\Phia_0=0$, so this identity suffices. By \eqref{eq:row-alg},
\[
2(\log\alg)_\tau=\frac{\delta_\tau}{\delta}+\frac{U_{1,\tau}}{U_1}-\frac{\sast_\tau}{\sast}+\frac{\sast_\tau}{1-\sast}.
\]
From $\delta=E+\alpha_0\tau-2\wB\tau\sast$ we get $\delta_\tau=\alpha_0-2\wB\sast-2\wB\tau\sast_\tau=\Rtr-\wB\tau\sast_\tau$. So the claim $2\Rtr/\delta=2(\log\alg)_\tau$ becomes, after multiplying by $\delta$ and inserting $\sast_\tau=\sast\Lv_2/\delta$ and $U_{1,\tau}=-\sast-\tau\sast_\tau=-\sast(\delta+\tau\Lv_2)/\delta$,
\[
-\frac{\sast(\delta+\tau\Lv_2)}{U_1}-\Lv_2+\frac{\Lv_2\sast}{1-\sast}-\delta_\tau=\frac{2\wB\tau\sast\Lv_2}{\delta}
\]
(the right side is $2\Rtr-2\delta_\tau=2\wB\tau\sast_\tau$). Inserting $\delta_\tau=\alpha_0-2\wB\sast-2\wB\tau\sast\Lv_2/\delta$, whose last term cancels the right side, the claim reads
\[
-\frac{\sast(\delta+\tau\Lv_2)}{U_1}-\Lv_2+\frac{\Lv_2\sast}{1-\sast}-\alpha_0+2\wB\sast=0 .
\]
Since $-\Lv_2-\alpha_0+2\wB\sast=\wB\sast$, dividing by $\sast$ this is
\[
\frac{\delta+\tau\Lv_2}{U_1}=\frac{\Lv_2}{1-\sast}+\wB=\frac{\Lv_2+\wB(1-\sast)}{1-\sast}=\frac{\wB-\alpha_0}{1-\sast}=\frac{\alpha_1}{1-\sast},
\]
i.e.\ $(\delta+\tau\Lv_2)(1-\sast)=\alpha_1U_1$, which is the second form of identity (ii) of \cref{lem:set-N}, $(E-\wB\tau\sast)(1-\sast)=\alpha_1(1-\tau\sast)$, with $E-\wB\tau\sast=\delta+\tau\Lv_2$ (\eqref{eq:row-algp}).
\end{proof}

Hence $\Kcoef_k=\coef{\tau^k}(\Phia\alg)$ lies in $\bbC\fps{z}$, which is \cref{cor:row-reg} again, and \eqref{eq:row-transport} says more: the $z$-coefficients of $\Ephi\Fhyp$ are analytic on the whole disc $\Ddisc$, up to the nearest zero of $\delta^2$, where $\sast$ and $\stil$ coalesce. Comparison with \eqref{eq:row-closed}: $\Nlap=\Abeta\,(\Phia\alg)$; the Laplace route fixes the constant $\Abeta$, the transport route the $\tau$-dependence. (The identification $\Phia_0=1/\alg$ also follows formally from \cref{prop:row-gaussode}, whose $z^0$ coefficient is $\coef{z^0}(\Nlap/\Abeta)/\alg=1/\alg$.) Both proofs of the absence of poles at $z=0$ are kept: \cref{cor:row-reg} needs only $\Abeta\in1+z\bbC\fps{z}$, \cref{prop:row-transport} needs no Laplace expansion at all.

%-------------------------------------------------------------------------------
\subsection{First orders}
\label{sec:row-first}

\begin{rmk}[First orders]
\label{rmk:row-first}
From \eqref{eq:row-gaussdef}, $\Fhyp_1=\ga\gb/\gc=(\alpha_2+z)(\wA+z)/(z(E+2z))$; and $\Ephi_1=-\dphi_1/z$ with $\dphi_1\coloneqq\coef{\tau^1}\dphi=\dphi_\tau(0)=\alpha_2\wA/E$ (\eqref{eq:row-dphider} at $\tau=0$). So
\[
\Hcoef_1=\Fhyp_1+\Ephi_1=\frac{E(\alpha_2+z)(\wA+z)-\alpha_2\wA(E+2z)}{Ez(E+2z)}=\frac{Ez+E(\alpha_2+\wA)-2\alpha_2\wA}{E(E+2z)}:
\]
the poles at $z=0$ cancel. Next, $\coef{\tau^1}\alg=-(E^2+\uq)/(2E^2)$: by \eqref{eq:row-alg}, $\log\alg^2=\log(\delta/E)+\log U_1-\log(E\sast/\wA)-\log(E(1-\sast)/\alpha_1)$, whose first-order terms in $\tau$ are $-(E^2+\uq)/(2E^2)$, $-\wA/E$, $-\alpha_1\alpha_2/E^2$ and $+\wA\alpha_2/E^2$ respectively, using $\delta=E-(E^2+\uq)\tau/(2E)+O(\tau^2)$ and $\sast_\tau(0)=\sast(0)\Lv_2(0)/E=\wA\alpha_1\alpha_2/E^3$ (identity (iii) of \cref{lem:set-N} with $\Lv_2(0)=\alpha_1\alpha_2/E$, \cref{lem:set-eig}); their sum is $-(E^2+\uq)/(2E^2)-(\wA\wB+\alpha_1\alpha_2)/E^2=-(E^2+\uq)/E^2$, since $\wA\wB+\alpha_1\alpha_2=2\alpha_1\alpha_2+\alpha_0E=(E^2+\uq)/2$. Therefore
\[
\coef{\tau^1}\rr_{0,1}=\frac{E}{E+z}\Big(\Hcoef_1-\frac{E^2+\uq}{2E^2}\Big)=\frac{E}{E+z}\cdot\frac{2E\big(Ez+E(\alpha_2+\wA)-2\alpha_2\wA\big)-(E^2+\uq)(E+2z)}{2E^2(E+2z)},
\]
and therefore
\[
\coef{\tau^1}\rr_{0,1}=-\frac{\uq\,z}{E(E+z)(E+2z)},
\]
where in the numerator the constant terms cancel because $2E(\alpha_2+\wA)-4\alpha_2\wA=E^2+\uq$ (proof of \cref{lem:row-eik}(i)), and the coefficient of $z$ is $2E^2-2E^2-2\uq=-2\uq$. Since $\tau=t+O(t^2)$, $\coef{t^1}\rr_{0,1}=\coef{\tau^1}\rr_{0,1}$.
\end{rmk}

%-------------------------------------------------------------------------------
\subsection{Canonical coordinates}
\label{sec:row-canon}

\begin{lem}[Canonical coordinates; formal]\genusstatus{AI-proposed, unchecked}
\label{lem:row-canon}
Put $u_i\coloneqq\cW(U_i)$ for $i=1,2$ and $\muq\coloneqq E/\tau-\wA$. Then $\dd u_i/\dd t=\Lv_i$. Explicitly, $u_1=\alpha_1t+E\log\tau+\phi_\tau(\sast)+\alpha_1\log(-1)$, and
\[
\Lv_1=\muq+\dphi',\qquad \Lv_1-\dphi'=\frac E\tau-\wA,\qquad \Lv_2=\muq+(\alpha_2-\Lv_1),\qquad \frac{\dd}{\dd t}\phi_\tau(\stil)=\wA-\wB\stil=\alpha_2-\Lv_1 .
\]
\end{lem}

\begin{proof}
At a critical point, $(\de_t\cW)(U)=\Lv(U)$ (\cref{lem:set-crit}), and
\[
\frac{\dd u_i}{\dd t}=(\de_t\cW)(U_i)+\cW_1(U_i)\,\frac{\dd\log U_i}{\dd t}=(\de_t\cW)(U_i)=\Lv_i .
\]
The explicit form of $u_1$ is \eqref{eq:row-euler1} at $s=\sast$ with $\log\cQ=t$. Hence $u_1'=\alpha_1+E(1-\tau)/\tau+\dd\phi_\tau(\sast)/\dd t$, and $\dd\phi_\tau(\sast)/\dd t=\dphi'=\wA-\wB\sast$ by \eqref{eq:row-dphider}; since $\alpha_1+E(1-\tau)/\tau=E/\tau-\wA+(\alpha_1+\wA-E)=\muq$, this gives $\Lv_1=\muq+\dphi'$. Directly,
\[
\muq+\dphi'=E/\tau-\wB\sast=\wA/(\tau\sast)-\alpha_0=\Lv_1
\]
(\cref{lem:set-eig}), because $\Ns(\sast)=0$ divided by $\tau\sast$ reads $\wB\sast-(E+\alpha_0\tau)/\tau+\wA/(\tau\sast)=0$. At $\stil$, identity (ii) of \cref{lem:set-N} at the second root gives $\dd\phi_\tau(\stil)/\dd t=(1-\tau)\,\alpha_2\stil/(1-\tau\stil)$, which equals $\wA-\wB\stil$ (the envelope formula, as for $\sast$), and $\wB\stil=\wA/(\tau\sast)=\Lv_1+\alpha_0$ by $\sast\stil=\wA/(\wB\tau)$, so $\wA-\wB\stil=\alpha_2-\Lv_1$ and $\muq+\alpha_2-\Lv_1=E/\tau-\alpha_0-\Lv_1=\Lv_2$ by $\Lv_1+\Lv_2=E/\tau-\alpha_0$ (\cref{lem:set-eig}).
\end{proof}

\begin{rmk}[Euler's integral; motivation, used nowhere]
\label{rmk:row-eulerint}
Under $\mathrm{Re}(\wA/z)>-1$, $\mathrm{Re}(\alpha_1/z)>-1$, $|\tau|<1$, the integral $J(z,\tau)\coloneqq\int_0^1s^{\wA/z}(1-s)^{\alpha_1/z}(1-\tau s)^{-\alpha_2/z-1}\,\dd s$ equals $\Bfn(\wA/z+1,\alpha_1/z+1)\,\Fhyp(z,\tau)$: this is Euler's integral representation with $(a,b,g)=(\ga,\gb,\gc)$ (expand $(1-\tau s)^{-\ga}$ and use $\int_0^1s^{\gb-1+n}(1-s)^{\gc-\gb-1}\,\dd s=\Bfn(\gb,\gc-\gb)(\gb)_n/(\gc)_n$). For general $a,b,g$ and $\mathsf{K}\coloneqq s^{b-1}(1-s)^{g-b-1}(1-\tau s)^{-a}$ the exact identity
\[
\big[\tau(1-\tau)\de_\tau^2+(g-(a+b+1)\tau)\de_\tau-ab\big]\mathsf{K}=-a\,\de_s\big[s^b(1-s)^{g-b}(1-\tau s)^{-a-1}\big]
\]
holds: after division by $a\mathsf{K}(1-\tau s)^{-2}$ both sides equal $-b+gs+(b-a-1)\tau s+(a+1-g)\tau s^2$. Under the stated hypotheses the boundary terms vanish and $\Gauss[J]=0$; the recursion of \cref{lem:row-unique} at numerical $z$ ($k+\gc\neq0$ since $\mathrm{Re}\,\gc>0$) re-proves $J=J(z,0)\Fhyp$. The formal expansion obeys Gauss's equation because of the pointwise identity \eqref{eq:row-amp} alone, with no boundary terms and no hypothesis on $z$; \eqref{eq:row-amp} is this identity with $(a,b,g)=(\ga,\gb,\gc)$, for which $\mathsf{K}=\re^{\phi_\tau/z}(1-\tau s)^{-1}$.
\end{rmk}

% s04_leaves. Section file; new macros at the top; label prefix: lea.
% Local macros (used in this section only).
\providecommand{\leaVn}{\mathbf{n}}                       % normalised vector t rho v_1
\providecommand{\leaPr}{\mathrm{Pr}}                      % projectors Pr_1, Pr_2 onto the eigenlines
\providecommand{\leaPhiF}{\boldsymbol{\Phi}_{\Fhyp}}       % (F, (D_z + omega) F) in the Wronskian remark
\providecommand{\leaLt}{\Lnorm_{2}}                        % Lnorm[phi_tau, (1 - tau s)^{-1}; s~]
\providecommand{\leaDt}{\mathsf{D}_{\stil}}                % d/dt + f_0'/z - f_2'/(2 f_2)
\providecommand{\leaDtau}{\mathsf{D}_{\stil,\tau}}         % (1 - tau)^{-1} D_{s~}
\providecommand{\leaNtw}{\mathsf{N}_{12}}                  % rho P_{12}
\providecommand{\leaDisc}{\Ddisc_{1}}                      % a disc around tau_0 with closure in D''
\providecommand{\TFsup}{\mathrm{F}}                        % superscript for the fibrewise point (0, x, 1, x, 1)

\section{The leaf structure at generic weights}
\label{sec:lea}

This section proves Theorem~A. Throughout, the weights $(\alpha_0,\alpha_1,\alpha_2,\epsilon_1,\epsilon_2)$ are complex numbers satisfying the non-degeneracy condition (ND) of \cref{def:set-ND}, or indeterminates; all factors are the frame-B factors of \cref{def:set-rows,def:set-factors} at their true idempotents, idempotent~$1$ being the critical point $U_1$ of the superpotential that runs into the logarithmic singularity at $t=0$ (the singular idempotent, colour~$1$) and idempotent~$2$ the one that stays away from it (the regular idempotent, colour~$2$). The objects $\rho$, $\lf_j$, $\Lsp_m$, $\qt$, $f_{j,b}$, $\qo$ are those of \cref{def:set-rho,def:set-factors}:
\begin{gather*}
\rho(t) = \Bigl(\frac{\Delta_1 t^2}{-\epsilon_1\epsilon_2E^3}\Bigr)^{-1/2}\ (\rho(0)=1),\qquad
\lf_j(\psi) = \frac{1}{E-j\psi},\qquad \Lsp_m = \operatorname{span}(\lf_1,\dots,\lf_m),\\
\qt(\psi) = \sum_{d\geq 1}\Dil_1(d;t)\psi^d,\qquad f_{j,b}(\psi) = \coef{w^b}\Edge_{1j}(\psi,w),\qquad \qo = E - E^2\lf_1 .
\end{gather*}
The results of \cref{sec:row} are used in the following form: the frame-B singular identity row is $\rr_{0,1}(z;t) = (E/(E+z))\,\Fhyp(z,\tau)\,\Ephi(z,\tau)\,\alg(\tau)$ (\cref{thm:row-main}), with $\Fhyp = {}_2F_1(\ga,\gb;\gc;\tau)$, $\ga = 1+\alpha_2/z$, $\gb = 1+\wA/z$, $\gc = 2+E/z$, $\Ephi = \exp(-\dphi/z)$ and $\alg$ the unit of \cref{lem:row-hess}; its $\tau$-coefficients are rational functions of $z$ without pole at $z=0$ (\cref{cor:row-reg}); $\rr_{0,1}(z;0) = E/(E+z)$ (\cref{cor:row-F2}); and $\Lv_1 = \muq + \dphi'$ with $\muq = E/\tau - \wA$ (\cref{lem:row-canon}). A prime denotes $d/dt = (1-\tau)\,d/d\tau$.

\subsection{Statements}
\label{sec:lea-statements}

\begin{thm}[Theorem A: leaf structure at generic weights]\genusstatus{AI-proposed, unchecked}
\label{thm:lea-main}
Let $(\alpha_0,\alpha_1,\alpha_2,\epsilon_1,\epsilon_2)$ be complex numbers on $\TCY$ with (ND), or indeterminates with $E = \epsilon_1+\epsilon_2$ imposed. All factors are the frame-B factors of \cref{sec:set} (rows $\RB = \RZ\hw^{-1}$ with the class $\bw = (E,-\epsilon_1,-\epsilon_2)$, $\RZ$ the rows with the $\cQ = 0$ constant $\hnw$), at their true idempotents. Then:
\begin{description}
\item[(F1)] $\Delta_1 = -\epsilon_1\epsilon_2E^3/(t^2\rho^2)$, where $\rho = \tau/(t\alg)$ is analytic at $t = 0$ with $\rho(0) = 1$; so $\Delta_1 = -(\epsilon_1\epsilon_2E^3/t^2)(1 - \uq t/E^2 + O(t^2))$. The function $\Delta_2 = \epsilon_1\epsilon_2\wB\delta\sast\Lv_2$ is analytic at $t = 0$, with $\Delta_2(0) = \epsilon_1\epsilon_2\alpha_1\alpha_2\wA\wB/E \neq 0$ and $\Delta_2'(0) = 0$.
\item[(F2)] $\rr_{0,1}(z;0) = E/(E+z)$, and $\rho\,\rr_{1,1}(z;t) = E/t + \alpha_1\alpha_2/(E+z) + O(t)$.
\item[(F3)] Every frame-B dilaton factor $\Dil_i(d)$ ($i = 1,2$, $d \geq 1$) and every edge factor $\Edge_{ij}(m,n)$ ($i,j\in\{1,2\}$, $m,n\geq 0$) is analytic at $t = 0$.
\item[(L1)] $\coef{t^0}(\rho\qt) = \qo = E - E^2\lf_1$, and $\coef{t^k}(\rho\qt)$ lies in $\Lsp_{k+1}$ for every $k\geq 1$.
\item[(L2)] For $j\in\{1,2\}$, every $b\geq 0$ and every $k\geq 0$, $\coef{t^k}(\rho f_{j,b})$ lies in $\Lsp_{k+1}$. Moreover $\Edge_{21}(b,a) = \Edge_{12}(a,b)$, and $\coef{t^k}f_{j,b}$ lies in $\Lsp_{k+1}$ as well.
\item[(L3)] For every $k\geq 0$, $\coef{t^k}(\rho^2\Edge_{11}(\psi,\psi'))$ lies in $\operatorname{span}\{\lf_i(\psi)\lf_l(\psi') : i,l\geq 1,\ (i-1)+(l-1)\leq k\}$, and $\Edge_{11}(a,b) = \Edge_{11}(b,a)$. The same support bound holds for $\coef{t^k}(\rho^e\Edge_{11})$ for every integer $e$.
\end{description}
Every membership is an identity of formal power series in $\psi$ (and $\psi'$), valid at every $\psi$-order and every $t$-order; the coefficients are rational functions of the weights, regular on (ND). Sharpness: with
\[
\Psh_k := \prod_{i=1}^{k}\bigl((k+1)\alpha_2 - iE\bigr)\bigl((k+1)\wA - iE\bigr)\qquad(\Psh_0 = 1),
\]
the top leaf $\lf_{k+1}$ occurs in $\coef{t^k}(\rho\qt)$ if and only if $\Psh_k\neq 0$, i.e.\ if and only if $E\notin\{(k+1)\alpha_2/i,\,(k+1)\wA/i : 1\leq i\leq k\}$; every entry of the anti-diagonal $(i-1)+(l-1) = k$ of (L3) is non-zero if and only if $\Psh_{i-1}\Psh_{l-1}\neq 0$ (\cref{cor:lea-top,cor:lea-anti}). Large-$z$ form: coefficientwise in $t$, $z\,\rr_{0,1}(z;t)\to E/\rho(t)$ and $\Edge_{1j}(z,w)\to 0$ as $z\to\infty$ at fixed order in $w$, in the sense of the rational functions of $z$ of \cref{lem:lea-tcoef} and \cref{cor:lea-alllegs}.
\end{thm}

\begin{rmk}[where the Calabi--Yau condition enters]
\label{rmk:lea-cy}
No step of the proof below uses $E = \epsilon_1+\epsilon_2$: the frame-B rows are $h_{(\alpha_1,\wA,-E)}\rpt_{0,i}$ (\cref{lem:set-hmult}) and do not involve $\epsilon_1,\epsilon_2$ at all; $\Delta_i$, $\Edge_{ij}$ and the pairing carry $\epsilon_1\epsilon_2$ as an overall factor. \Cref{thm:lea-main} therefore holds at every weight point with (ND), with $\bw := (E,-\epsilon_1,-\epsilon_2)$ and $E := \alpha_0+\alpha_1+\alpha_2$. The Calabi--Yau condition is what makes $\bw$ a sum-zero triple whose first weight is the leaf scale $E$, which is what the two theorems of \cite{LP2} cited in \cref{sec:reg} (\cref{thm:reg-L,thm:reg-VI}) need; so it is needed for the regularity theorem built on \cref{thm:lea-main}, not for \cref{thm:lea-main} itself.
\end{rmk}

\begin{thm}[quantum differential equation and symplecticity]\genusstatus{AI-proposed, unchecked}
\label{thm:lea-qde}
\begin{enumerate}
\item For $i = 1,2$ the second row of the quantum differential equation holds to all orders in $z$: $z\Dop_i\vv_i = (\Aop - \Lv_i)\vv_i$, with $\vv_i = (\rr_{0,i},\rr_{1,i})$, $\Dop_i = d/dt - \Delta_i'/(2\Delta_i)$ and $\Aop = \begin{bmatrix}0&1\\-\alpha_1\alpha_2/\tau & E/\tau-\alpha_0\end{bmatrix}$ acting on covectors (\cref{prop:lea-row2,thm:lea-row2reg}).
\item $\langle\vv_i(w),\vv_j(-w)\rangle = \delta_{ij}\Delta_i$ for all $i,j$, to all orders in $w$ (\cref{prop:lea-K1,cor:lea-orth,cor:lea-K2}).
\item $\vv_2(z;0) = \hreg(z)\,(1,\alpha_1\alpha_2/E)$ with $\hreg = h_{(\alpha_1,\alpha_2,\wA,\wB,-E,-E)}$ (\cref{prop:lea-regrows}; the scalar factor is the analytic input of \cref{prop:lea-regrows}(a), the direction $(1,\alpha_1\alpha_2/E)$ is formal).
\end{enumerate}
\end{thm}

\begin{thm}[closed forms]\genusstatus{AI-proposed, unchecked}
\label{thm:lea-closed}
In frame B,
\begin{align*}
\rho\,\rr_{0,1}(z;t) &= \frac{E}{E+z}\,\frac{\tau}{t}\,{}_2F_1\Bigl(1+\frac{\alpha_2}{z},\,1+\frac{\wA}{z};\,2+\frac{E}{z};\,\tau\Bigr)\exp\Bigl(-\frac{\dphi(\tau)}{z}\Bigr),\\
\rho\,\qt(\psi) &= E + \sum_{j\geq 1}\Aleaf_j(t)\lf_j(\psi),\qquad \Aleaf_j(t) = (-1)^j\Cj_j\,t^{-1}\zq(t)^j\,{}_2F_1\Bigl(\frac{j\wB}{E},\frac{j\alpha_1}{E};\,j;\,\tau\Bigr),\\
&\qquad \Cj_j := \frac{\Psh_{j-1}}{j!\,(j-1)!\,E^{2j-4}},\qquad \zq := \tau\,e^{\dphi/E},\\
\frac{d}{dt}\bigl[t^2\rho^2\Edge_{11}(\psi,\psi')\bigr] &= \Thf(t)\,\Dg(\psi)\,\Dg(\psi'),\qquad \gl = -\frac{1}{E}\sum_j j\Aleaf_j\lf_j,\\
\Thf &= \epsilon_1\epsilon_2(\alpha_1\alpha_2 - E\Lv_2) = -\epsilon_1\epsilon_2\wB\tau\sast\Lv_2,
\end{align*}
with closed forms of every leaf and loop coefficient and of the $t = 0$ values of $\Edge_{12}$, $\Edge_{22}$ (\cref{sec:lea-closed}). The first line is the key formula.
\end{thm}

The key formula is proved in \cref{prop:lea-key}, the flow in \cref{thm:lea-flow}, and the remaining closed forms in \cref{sec:lea-closed}. The proof of \cref{thm:lea-main,thm:lea-qde} is assembled in \cref{sec:lea-assembly}.

\subsection{The metric in the Euler coordinate, the key formula, the quantum differential equation and the symplectic constant at the singular idempotent}
\label{sec:lea-qde}

\begin{prop}[the Hessian in the Euler coordinate; $\rho$; the key formula]\genusstatus{AI-proposed, unchecked}
\label{prop:lea-key}
\begin{enumerate}
\item The vertex weight of the singular idempotent is
\[
\Delta_1 = \alpha_1\wA\epsilon_1\epsilon_2\,\cW_2(U_1) = \alpha_1\wA\epsilon_1\epsilon_2\,U_1^2\,\phi_\tau''(\sast)\,\tau^{-2} = -\epsilon_1\epsilon_2E^3\,\alg^2\,\tau^{-2}.
\]
\item $\rho = \tau/(t\alg)$, a unit of $\Oan$ with $\rho(0) = 1$, and
\begin{equation}
\label{eq:lea-rho}
t^2\rho^2\Delta_1 = -\epsilon_1\epsilon_2E^3\quad\text{exactly};
\end{equation}
$\Delta_1t^2/(-\epsilon_1\epsilon_2E^3) = (t/\tau)^2\alg^2$ is a unit of $\Oan$ with value $1$ at $t = 0$. This is (F1) for $\Delta_1$.
\item (Key formula.)
\begin{equation}
\label{eq:lea-key}
\Gk(z;t) := \rho\,\rr_{0,1}(z;t) = \frac{E}{E+z}\,\frac{\tau}{t}\,\Fhyp(z,\tau)\,\Ephi(z,\tau);
\end{equation}
$\Gk$ lies in $\bbC\fps{z,t}$, $\Gk(z;0) = E/(E+z)$, and every $\Gk_k := \coef{t^k}\Gk$ is a rational function of $z$ without pole at $z = 0$ whose Taylor series is $\coef{t^k}(\rho\,\rr_{0,1})$.
\item $\rho = 1 + \uq t/(2E^2) + O(t^2)$ and $\Delta_1 = -(\epsilon_1\epsilon_2E^3/t^2)(1 - \uq t/E^2 + O(t^2))$.
\end{enumerate}
\end{prop}

\begin{proof}
(i) The first equality is the fourth form of $\Delta_1$ in \cref{prop:set-idem}; the second is the tensor rule $\cW_2(U_i) = U_i^2\phi_\tau''(s_i)/\tau^2$ of \cref{lem:row-hess}; the third is the definition $\alg^2 = U_1^2\phi_\tau''(\sast)/\phi_0''(\wA/E)$ with $\phi_0''(\wA/E) = -E^3/(\alpha_1\wA)$ (\cref{lem:row-hess}).

(ii) By (i), $\rho^{-2} = \Delta_1t^2/(-\epsilon_1\epsilon_2E^3) = (t\alg/\tau)^2$, a unit of $\Oan$ with value $1$ at $t = 0$ since $\alg(0) = 1$ and $\tau/t = 1 - t/2 + O(t^2)$. Its square root with value $1$ at $t = 0$ is $t\alg/\tau$ (two square roots of a unit of $\Oan$ with the same value at $0$ coincide), so $\rho = \tau/(t\alg)$ and \eqref{eq:lea-rho} holds exactly.

(iii) By \cref{thm:row-main} and (ii), $\rho\,\rr_{0,1} = \frac{\tau}{t\alg}\,\frac{E}{E+z}\,\Fhyp\,\Ephi\,\alg$, which is \eqref{eq:lea-key}. The rational form of $\Gk_k$: write $\Kcoef_j := \coef{\tau^j}(\Fhyp\Ephi\alg)$, so that $\rr_{0,1} = \sum_j(E/(E+z))\Kcoef_j\tau^j$ and $\Gk = \sum_j\rho\,(E/(E+z))\,\Kcoef_j\tau^j$; hence
\[
\Gk_k = \sum_{m\leq k}\bigl(\coef{t^m}\rho\bigr)\sum_{j\leq k-m}\bigl(\coef{t^{k-m}}\tau^j\bigr)\frac{E}{E+z}\Kcoef_j ,
\]
a finite combination of rational functions regular at $z = 0$ (\cref{cor:row-reg}) whose Taylor series is $\sum_m(\coef{t^m}\rho)\coef{t^{k-m}}\rr_{0,1} = \coef{t^k}(\rho\,\rr_{0,1})$ by the coefficientwise meaning of \cref{thm:row-main}. $\Gk_0 = E/(E+z)$ since $\Kcoef_0 = 1$.

(iv) By \cref{rmk:row-first}, $\alg = 1 - \frac{E^2+\uq}{2E^2}\tau + O(\tau^2)$, and $\tau/t = 1 - t/2 + O(t^2)$; so $\rho = (\tau/t)/\alg = 1 + \bigl[\frac{E^2+\uq}{2E^2} - \frac12\bigr]t + O(t^2) = 1 + \frac{\uq}{2E^2}t + O(t^2)$, and $\rho^{-2} = 1 - \uq t/E^2 + O(t^2)$.
\end{proof}

\begin{lem}[the operator identity]\genusstatus{AI-proposed, unchecked}
\label{lem:lea-op}
Let $\mathsf{B} := \bbC[z,1/z][\tau,1/\tau,1/(1-\tau)]$, and consider the ring of differential operators generated by $\mathsf{B}$ and one symbol $\Dd$ subject to $\Dd a = a\Dd + z a'$ for $a\in\mathsf{B}$ (so $\Dd(\tau) = z(1-\tau)$); put $d/d\tau := (z(1-\tau))^{-1}\Dd$. With $\omq = (E+z)/\tau - (\wA+z)$ and Gauss's operator of \eqref{eq:row-gaussop}, written with this $d/d\tau$,
\[
\Gauss[Y] = \tau(1-\tau)Y_{\tau\tau} + (\gc - (\ga+\gb+1)\tau)Y_\tau - \ga\gb\,Y,
\]
one has
\begin{equation}
\label{eq:lea-op}
(\Dd+\omq)^2 - \Bigl(\frac{E}{\tau}-\alpha_0\Bigr)(\Dd+\omq) + \frac{\alpha_1\alpha_2}{\tau} = \frac{(1-\tau)z^2}{\tau}\,\Gauss .
\end{equation}
The same holds for any symbol $\Dd'$ with the same Leibniz rule $\Dd'(aY) = z a'Y + a\Dd'Y$. In particular, for $\Dd = \Dz = z\,d/dt$ and every $f\in\Rz$,
\[
\Bigl[(\Dz+\omq)^2 - \Bigl(\frac{E}{\tau}-\alpha_0\Bigr)(\Dz+\omq) + \frac{\alpha_1\alpha_2}{\tau}\Bigr]f = \frac{(1-\tau)z^2}{\tau}\,\Gauss[f].
\]
The case $\Dd = \Dz$ is $\Dz$ conjugated by $\mathsf{Pf} := (1-\tau)^{-\alpha_1/z}\tau^{E/z+1}$: formally $\Dz(\mathsf{Pf}\,f) = \mathsf{Pf}\,(\Dz+\omq)f$, since $z(\log\mathsf{Pf})' = \alpha_1 + E(1-\tau)/\tau + z(1-\tau)/\tau = \muq + z(1-\tau)/\tau = \omq$.
\end{lem}

\begin{proof}
Write $\Lsum := E/\tau - \alpha_0$. For $f$ in the ring, $(\Dd+\omq)^2f - \Lsum(\Dd+\omq)f + (\alpha_1\alpha_2/\tau)f = \Dd^2f + (2\omq - \Lsum)\Dd f + [\Dd(\omq) + \omq(\omq-\Lsum) + \alpha_1\alpha_2/\tau]f$, with $\Dd f = z(1-\tau)f_\tau$ and $\Dd^2f = z^2(1-\tau)^2f_{\tau\tau} - z^2(1-\tau)f_\tau$ (because $\Dd(1-\tau) = -z(1-\tau)$).

Coefficient of $f$. $\Dd(\omq) = z\omq' = z(1-\tau)\,d\omq/d\tau = -z(1-\tau)(E+z)/\tau^2$. Next, $\omq - \Lsum = z/\tau - (\alpha_2+z)$, since $(E+z)-E = z$ and $-(\wA+z)+\alpha_0 = -(\alpha_2+z)$; so
\[
\omq(\omq-\Lsum) = \frac{z(E+z)}{\tau^2} - \frac{(E+z)(\alpha_2+z) + z(\wA+z)}{\tau} + (\wA+z)(\alpha_2+z).
\]
The $\tau^{-2}$ terms of $\Dd(\omq) + \omq(\omq-\Lsum)$ cancel, the $\tau^{-1}$ terms add up to $[z(E+z) - (E+z)(\alpha_2+z) - z(\wA+z) + \alpha_1\alpha_2]/\tau$, and
\begin{align*}
z(E+z) - (E+z)(\alpha_2+z) - z(\wA+z) + \alpha_1\alpha_2 &= -[\alpha_2E + (\alpha_2+\wA)z + z^2 - \alpha_1\alpha_2]\\
&= -[\alpha_2\wA + (\alpha_2+\wA)z + z^2] = -(\wA+z)(\alpha_2+z),
\end{align*}
using $\alpha_2E - \alpha_1\alpha_2 = \alpha_2\wA$. So the coefficient of $f$ is $-(\wA+z)(\alpha_2+z)/\tau + (\wA+z)(\alpha_2+z) = -(1-\tau)(\wA+z)(\alpha_2+z)/\tau$.

Coefficient of $f_\tau$: $-z^2(1-\tau) + z(1-\tau)(2\omq-\Lsum) = z(1-\tau)[2\omq - \Lsum - z] = z(1-\tau)\bigl[(E+2z)/\tau - (2\wA - \alpha_0 + 3z)\bigr] = z(1-\tau)\bigl[(E+2z)/\tau - (\alpha_2+\wA+3z)\bigr]$, since $2\wA-\alpha_0 = \wA + \alpha_2$.

Coefficient of $f_{\tau\tau}$: $z^2(1-\tau)^2$.

On the right, $((1-\tau)z^2/\tau)\Gauss[f]$ has $f_{\tau\tau}$-coefficient $z^2(1-\tau)^2$, $f_\tau$-coefficient
\[
(1-\tau)z^2\bigl[(2+E/z)/\tau - (3+(\alpha_2+\wA)/z)\bigr] = z(1-\tau)\bigl[(E+2z)/\tau - (\alpha_2+\wA+3z)\bigr],
\]
and $f$-coefficient
\[
-\frac{(1-\tau)z^2}{\tau}\Bigl(1+\frac{\alpha_2}{z}\Bigr)\Bigl(1+\frac{\wA}{z}\Bigr) = -\frac{(1-\tau)(\alpha_2+z)(\wA+z)}{\tau}.
\]
The three coefficients agree. Only the Leibniz rule and $\Dd(\tau) = z(1-\tau)$ were used, which gives the statement for $\Dd'$. For $\mathsf{Pf}$: $z(\log\mathsf{Pf})' = z(1-\tau)\bigl[\alpha_1/(z(1-\tau)) + (E/z+1)/\tau\bigr] = \alpha_1 + E(1-\tau)/\tau + z(1-\tau)/\tau$, and $\alpha_1 + E(1-\tau)/\tau = E/\tau - \wA = \muq$ because $\alpha_1 - E = -\wA$.
\end{proof}

\begin{prop}[the second row of the quantum differential equation at the singular idempotent]\genusstatus{AI-proposed, unchecked}
\label{prop:lea-row2}
In $\Rzo$ ($\rr_{0,1}\in\bbC\fps{z,\tau}$, $\rr_{1,1}\in\tau^{-1}\bbC\fps{z,\tau}$), hence to all orders in $z$ as identities of meromorphic functions of $t$ near $0$:
\begin{align*}
\text{Row 1:}\quad & \rr_{1,1} = (\Lv_1 + z\Dop_1)\rr_{0,1};\\
\text{Row 2:}\quad & (\Lv_1 + z\Dop_1)\rr_{1,1} = -\frac{\alpha_1\alpha_2}{\tau}\rr_{0,1} + \Bigl(\frac{E}{\tau}-\alpha_0\Bigr)\rr_{1,1},\quad\text{equivalently}\quad z\Dop_1\rr_{1,1} = \Lv_2(\rr_{1,1} - \Lv_1\rr_{0,1}),
\end{align*}
i.e.\ $(\Dop_1\rr_{1,1})[m-1] = \Lv_2(\rr_{1,1}[m] - \Lv_1\rr_{0,1}[m])$ for every $m\geq 1$. In matrix form $z\Dop_1\vv_1 = (\Aop-\Lv_1)\vv_1$, and for $\leaVn := t\rho\,\vv_1$, an element of $\bbC\fps{z,t}^2$,
\begin{equation}
\label{eq:lea-qde}
z\,\leaVn' = (\Aop-\Lv_1)\leaVn .
\end{equation}
\end{prop}

\begin{proof}
Row 1 is the definition of the H-row (\cref{def:set-rows}). By \eqref{eq:lea-rho}, $\Delta_1'/(2\Delta_1) = -(t\rho)'/(t\rho)$, so $\Dop_1f = (t\rho)^{-1}(t\rho f)'$, i.e.
\begin{equation}
\label{eq:lea-conj}
\Lv_1 + z\Dop_1 = (t\rho)^{-1}\bigl(z\,\tfrac{d}{dt} + \Lv_1\bigr)(t\rho)\qquad\text{as operators on }\Rz .
\end{equation}
(i) Conjugation. For $f\in\Rz$, using $z\Ephi' = -\dphi'\,\Ephi$, $\tau' = 1-\tau$ and $\Lv_1 - \dphi' = \muq$ (\cref{lem:row-canon}),
\[
\bigl(z\tfrac{d}{dt} + \Lv_1\bigr)(\tau\Ephi f) = \tau\Ephi\bigl[zf' + (\Lv_1 - \dphi' + z(1-\tau)/\tau)f\bigr] = \tau\Ephi\,(\Dz+\omq)f,
\]
because $\muq + z(1-\tau)/\tau = (E+z)/\tau - (\wA+z) = \omq$.

(ii) By \eqref{eq:lea-key}, $t\rho\,\rr_{0,1} = t\Gk = (E/(E+z))\,\tau\Ephi\Fhyp$. By \eqref{eq:lea-conj} and (i),
\begin{align*}
t\rho\,\rr_{1,1} &= \bigl(z\tfrac{d}{dt}+\Lv_1\bigr)(t\Gk) = \frac{E}{E+z}\,\tau\Ephi\,(\Dz+\omq)\Fhyp,\\
t\rho\,(\Lv_1+z\Dop_1)\rr_{1,1} &= \bigl(z\tfrac{d}{dt}+\Lv_1\bigr)(t\rho\,\rr_{1,1}) = \frac{E}{E+z}\,\tau\Ephi\,(\Dz+\omq)^2\Fhyp .
\end{align*}
(iii) Hence, by \eqref{eq:lea-op} and $\Gauss[\Fhyp] = 0$ (\cref{lem:row-unique,thm:row-closed}),
\begin{align*}
t\rho\Bigl[(\Lv_1+z\Dop_1)\rr_{1,1} + \frac{\alpha_1\alpha_2}{\tau}\rr_{0,1} - \Bigl(\frac{E}{\tau}-\alpha_0\Bigr)\rr_{1,1}\Bigr]
&= \frac{E}{E+z}\,\tau\Ephi\Bigl[(\Dz+\omq)^2 - \Bigl(\frac{E}{\tau}-\alpha_0\Bigr)(\Dz+\omq) + \frac{\alpha_1\alpha_2}{\tau}\Bigr]\Fhyp\\
&= \frac{E}{E+z}\,\Ephi\,(1-\tau)z^2\,\Gauss[\Fhyp] = 0,
\end{align*}
and $t\rho$ is a unit of $\Rz$. The second form of Row 2 uses $E/\tau - \alpha_0 - \Lv_1 = \Lv_2$ and $\alpha_1\alpha_2/\tau = \Lv_1\Lv_2$ (\cref{lem:set-eig}). The components of $(\Aop-\Lv_1)\vv_1$ are $\rr_{1,1} - \Lv_1\rr_{0,1} = z\Dop_1\rr_{0,1}$ (Row 1) and $-(\alpha_1\alpha_2/\tau)\rr_{0,1} + (E/\tau-\alpha_0-\Lv_1)\rr_{1,1} = z\Dop_1\rr_{1,1}$ (Row 2), and $z(t\rho f)' = t\rho\,z\Dop_1f$, which gives \eqref{eq:lea-qde}.

Membership: $\Lv_1$ and $\Delta_1'/\Delta_1$ lie in $\tau^{-1}\Oan$, so $\Lv_1 + z\Dop_1$ maps $\bbC\fps{z,\tau}$ into $\tau^{-1}\bbC\fps{z,\tau}$ and $\tau^{-1}\bbC\fps{z,\tau}$ into $\tau^{-2}\bbC\fps{z,\tau}$; with \cref{thm:row-main} every term lies in $\Rzo$, and the identity in $\Rz$ between elements of $\Rzo$ is an identity of formal double series, hence (TRANSFER, \cref{lem:set-rings}) of meromorphic functions of $t$ coefficientwise in $z$. $\leaVn = (t\Gk,\,(t\Lv_1+z)\Gk + zt\Gk')$ lies in $\bbC\fps{z,t}^2$, since $t\Lv_1\in\Oan$ (\cref{lem:set-eig}).
\end{proof}

Given Row 1 and \eqref{eq:lea-key}, Row 2 is equivalent to $\Gauss[\Fhyp] = 0$, since $(E/(E+z))\Ephi(1-\tau)z^2$ is a unit of $\Rz$; nothing about the regular idempotent is used. From step (ii), since $\Fhyp$ is a unit of $\Cz\fps{\tau}$,
\begin{equation}
\label{eq:lea-Hrow}
\rr_{1,1} = \rr_{0,1}\Bigl[\frac{E+z}{\tau} - (\wA+z) + z(1-\tau)\frac{\Fhyp_\tau}{\Fhyp}\Bigr],\qquad\text{equivalently}\qquad \rho\,\rr_{1,1} = \Bigl(\Lv_1 + \frac{z}{t}\Bigr)\Gk + z\Gk' .
\end{equation}
In terms of $\Fhyp$ alone, with $\Psi_{\Fhyp} := \tau(\Dz+\omq)\Fhyp = z\tau(1-\tau)\Fhyp_\tau + (E+z-(\wA+z)\tau)\Fhyp$, Row 2 reads $z(1-\tau)(\Psi_{\Fhyp})_\tau = \alpha_2\Psi_{\Fhyp} - \alpha_1\alpha_2\Fhyp$, because $z(1-\tau)(\Psi_{\Fhyp})_\tau - \alpha_2\Psi_{\Fhyp} + \alpha_1\alpha_2\Fhyp = (1-\tau)z^2\Gauss[\Fhyp]$ (the coefficient $\alpha_2$ is $-[\omq - z(1-\tau)/\tau - (E/\tau-\alpha_0)] = \wA - \alpha_0$, i.e.\ $\Lv_2 + \dphi'$).

\begin{prop}[the operator $\Aop$]\genusstatus{AI-proposed, unchecked}
\label{prop:lea-A}
\begin{enumerate}
\item $\Aop$ is the transpose of the matrix $\begin{bmatrix}0 & -\alpha_1\alpha_2/\tau\\ 1 & E/\tau-\alpha_0\end{bmatrix}$ of $H\qprod$ in the basis $(1,H)$ (columns: the coordinates of $H\qprod 1$, $H\qprod H$); $\Aop$ acts on covectors by $s\mapsto s(H\qprod\,\cdot\,)$.
\item $\chiv_i = (1,\Lv_i)$ (the character $a_0\,1 + a_1H\mapsto a_0 + a_1\Lv_i$; $\chiv_i(\eid_j) = \delta_{ij}$) satisfies $\Aop\chiv_i = \Lv_i\chiv_i$.
\item $\Aop$ is self-adjoint for $\langle\,,\rangle$: the bilinear form $(a,b)\mapsto\langle\Aop a,b\rangle$ has the symmetric matrix
\[
\epsilon_1\epsilon_2\begin{bmatrix}-(\alpha_1\alpha_2)^2/\tau & \alpha_1\alpha_2E/\tau\\ \alpha_1\alpha_2E/\tau & \alpha_1\alpha_2+\alpha_0E - E^2/\tau\end{bmatrix},
\]
so $\langle\Aop a,b\rangle = \langle a,\Aop b\rangle$.
\item $\langle\chiv_i,\chiv_j\rangle = \delta_{ij}\Delta_i$.
\item $\Aop - \Lv_1 = (\Lv_2-\Lv_1)\leaPr_2$ with $\leaPr_2a := (\pim(a)/\Delta_2)\chiv_2$ the $\langle\,,\rangle$-orthogonal projector onto $\bbC\chiv_2$ along $\bbC\chiv_1$, where
\[
\pim(a) := \langle a,\chiv_2\rangle = \epsilon_1\epsilon_2\bigl[\alpha_1\alpha_2(\alpha_0+\Lv_2)a_0 + (\alpha_1\alpha_2 - E\Lv_2)a_1\bigr] = \epsilon_1\epsilon_2\wB\sast\bigl(\alpha_1\alpha_2a_0 - \tau\Lv_2a_1\bigr),
\]
and, for all covectors $a,b$,
\begin{equation}
\label{eq:lea-proj}
\langle(\Aop-\Lv_1)a,b\rangle = \frac{\Lv_2-\Lv_1}{\Delta_2}\pim(a)\pim(b) = \frac{\pim(a)\pim(b)}{\Thf},\qquad
\Thf := \frac{\Delta_2}{\Lv_2-\Lv_1} = \epsilon_1\epsilon_2(\alpha_1\alpha_2 - E\Lv_2) = -\epsilon_1\epsilon_2\wB\tau\sast\Lv_2 .
\end{equation}
\end{enumerate}
\end{prop}

\begin{proof}
(i) $H\qprod 1 = H$ and $H\qprod H = -(\alpha_1\alpha_2/\tau)1 + (E/\tau-\alpha_0)H$ (\cref{prop:set-frob}); $(\Aop s)_0 = s_1 = s(H\qprod 1)$ and $(\Aop s)_1 = -(\alpha_1\alpha_2/\tau)s_0 + (E/\tau-\alpha_0)s_1 = s(H\qprod H)$.
(ii) $\Aop(1,\Lv) = (\Lv,\,-\alpha_1\alpha_2/\tau + (E/\tau-\alpha_0)\Lv)$, which is $\Lv(1,\Lv)$ if and only if $\Lv^2 - (E/\tau-\alpha_0)\Lv + \alpha_1\alpha_2/\tau = 0$, the equation of $\Lv_1,\Lv_2$ (\cref{lem:set-eig}).
(iii) Write $\eta^{-1} = \epsilon_1\epsilon_2\begin{bmatrix}p&q\\q&r\end{bmatrix}$ with $p = \alpha_0\alpha_1\alpha_2$, $q = \alpha_1\alpha_2$, $r = -E$ (\cref{prop:set-frob}). Then $\langle\Aop a,b\rangle = \epsilon_1\epsilon_2[p(\Aop a)_0b_0 + q((\Aop a)_0b_1 + (\Aop a)_1b_0) + r(\Aop a)_1b_1]$ with $(\Aop a)_0 = a_1$, $(\Aop a)_1 = -(\alpha_1\alpha_2/\tau)a_0 + (E/\tau-\alpha_0)a_1$. The coefficient of $a_0b_0$ is $-q\alpha_1\alpha_2/\tau = -(\alpha_1\alpha_2)^2/\tau$; of $a_0b_1$: $-r\alpha_1\alpha_2/\tau = \alpha_1\alpha_2E/\tau$; of $a_1b_0$: $p + q(E/\tau-\alpha_0) = \alpha_0\alpha_1\alpha_2 + \alpha_1\alpha_2E/\tau - \alpha_0\alpha_1\alpha_2 = \alpha_1\alpha_2E/\tau$; of $a_1b_1$: $q + r(E/\tau-\alpha_0) = \alpha_1\alpha_2 + \alpha_0E - E^2/\tau$. The matrix is symmetric. (Conceptually: $H\qprod$ is $\eta$-self-adjoint.)
(iv) $\langle\chiv_i,\chiv_j\rangle = \epsilon_1\epsilon_2[\alpha_0\alpha_1\alpha_2 + \alpha_1\alpha_2(\Lv_i+\Lv_j) - E\Lv_i\Lv_j]$. For $i\neq j$ this is $\epsilon_1\epsilon_2[\alpha_0\alpha_1\alpha_2 + \alpha_1\alpha_2(E/\tau-\alpha_0) - E\alpha_1\alpha_2/\tau] = 0$; for $i = j$ it is $\epsilon_1\epsilon_2[\alpha_0\alpha_1\alpha_2 + 2\alpha_1\alpha_2\Lv_i - E\Lv_i^2] = \Delta_i$ by the third form of \cref{prop:set-idem}.
(v) By (iv), $\chiv_1,\chiv_2$ is an orthogonal basis with non-zero norms, so $a = \sum_i(\langle a,\chiv_i\rangle/\Delta_i)\chiv_i$, $1 = \leaPr_1 + \leaPr_2$ and, by (ii), $\Aop = \Lv_1\leaPr_1 + \Lv_2\leaPr_2$, whence $\Aop - \Lv_1 = (\Lv_2-\Lv_1)\leaPr_2$ and $\langle\leaPr_2a,b\rangle = \pim(a)\pim(b)/\Delta_2$. For $\pim$: $\alpha_0\alpha_1\alpha_2 + \alpha_1\alpha_2\Lv_2 = \alpha_1\alpha_2(\alpha_0+\Lv_2) = \alpha_1\alpha_2\wB\sast$ and $\alpha_1\alpha_2 - E\Lv_2 = -\wB\tau\sast\Lv_2$ (\cref{lem:set-eig,lem:set-Delta2}). $(\Lv_2-\Lv_1)/\Delta_2 = 1/\Thf$ by definition, and $\Delta_2 = \epsilon_1\epsilon_2(\Lv_2-\Lv_1)(\alpha_1\alpha_2 - E\Lv_2)$ (second form of \cref{prop:set-idem}) gives the two other forms of $\Thf$.
\end{proof}

Note that $\pim(\chiv_1) = \epsilon_1\epsilon_2\wB\sast(\alpha_1\alpha_2 - \tau\Lv_1\Lv_2) = 0$ is literally $\Lv_1\Lv_2 = \alpha_1\alpha_2/\tau$. Self-adjointness is linear over $\bbC((t))$ and \eqref{eq:lea-proj} is bilinear; both extend coefficientwise to covectors with entries in series rings in $\psi$, $\psi'$.

\begin{prop}[the symplectic constant $\Kone = 1$]\genusstatus{AI-proposed, unchecked}
\label{prop:lea-K1}
\begin{enumerate}
\item The series
\[
\Kone(z) := \frac{\langle\vv_1(z),\vv_1(-z)\rangle}{\Delta_1} = -\frac{\langle\leaVn(z),\leaVn(-z)\rangle}{\epsilon_1\epsilon_2E^3}
\]
lies in $\bbC\fps{z,t}$ and is independent of~$t$: $\Kone\in\bbC\fps{z}$.
\item Put $\brow(\psi) := \rho\,\vv_1(-\psi)$. Then, with $\gl(\psi;t) = \Gk(-\psi;t)$ and $'$ the $t$-derivative at fixed $\psi$,
\[
\brow(\psi) = \bigl(\gl,\ (\Lv_1 - \psi/t)\gl - \psi\gl'\bigr);
\]
the series $t\brow(\psi)$ lies in $\bbC\fps{\psi,t}^2$,
\[
\brow^{(-1)} := \coef{t^{-1}}\brow(\psi) = (0,E),\qquad \gl(\psi)|_{t=0} = \frac{E}{E-\psi} = E\lf_1(\psi)\quad\text{(F2)}.
\]
\item $\Kone = 1$, i.e.\ $\langle\vv_1(w),\vv_1(-w)\rangle = \Delta_1$ to all orders in $w$.
\item Consequences: $\langle\brow(\psi),\brow(-\psi)\rangle = \rho^2\Delta_1 = -\epsilon_1\epsilon_2E^3/t^2$; the subtracted term in the edge formula of \cref{lem:set-edgegen} is $\Pij_{11}(-w,w) = -\Delta_1$, so
\begin{equation}
\label{eq:lea-E11sigma}
\rho^2\Edge_{11}(\psi,\psi') = \frac{-\epsilon_1\epsilon_2E^3/t^2 - \langle\brow(\psi),\brow(\psi')\rangle}{\psi+\psi'};
\end{equation}
and $-\psi\,\frac{d}{dt}(t\brow(\psi)) = (\Aop-\Lv_1)(t\brow(\psi))$, i.e.\ $\frac{d\brow}{dt} = -\frac{\brow}{t} - \frac{1}{\psi}(\Aop-\Lv_1)\brow$.
\end{enumerate}
\end{prop}

\begin{proof}
(i) $(t\rho)^2 = -\epsilon_1\epsilon_2E^3/\Delta_1$ by \eqref{eq:lea-rho} gives the second expression, which lies in $\bbC\fps{z,t}$ by \cref{prop:lea-row2}. Applying $z\mapsto -z$ to \eqref{eq:lea-qde} gives $-z\,\leaVn(-z)' = (\Aop-\Lv_1)\leaVn(-z)$ ($\Aop$ and $\Lv_1$ do not involve $z$). Since $\eta^{-1}$ is constant in $t$,
\[
z\,\langle\leaVn(z),\leaVn(-z)\rangle' = \langle(\Aop-\Lv_1)\leaVn(z),\leaVn(-z)\rangle - \langle\leaVn(z),(\Aop-\Lv_1)\leaVn(-z)\rangle = 0
\]
by \cref{prop:lea-A}(iii). Thus $(1-\tau)\,d\Kone/d\tau = 0$, and the kernel of $d/d\tau$ on $\bbC\fps{z,\tau}$ is $\bbC\fps{z}$.

(ii) By Row 1, \eqref{eq:lea-conj} and \eqref{eq:lea-key}, $\rho\,\rr_{1,1} = \rho(t\rho)^{-1}(z\frac{d}{dt}+\Lv_1)(t\Gk) = (\Lv_1 + z/t)\Gk + z\Gk'$, which is \eqref{eq:lea-Hrow}. By \cref{lem:set-eig}, $t\Lv_1\in\Oan$ with value $E$ at $t = 0$, so $t\rho\,\rr_{1,1} = (t\Lv_1+z)\Gk + zt\Gk'$ lies in $\bbC\fps{z,t}$, and at $t = 0$ it equals $(E+z)\cdot E/(E+z) = E$: the $z$-dependence cancels exactly because of the geometric form $E/(E+z)$ of (F2). Substituting $z = -\psi$ gives the two components of $\brow$, $t\brow\in\bbC\fps{\psi,t}^2$, $\brow^{(-1)} = (0,E)$, and $\gl|_{t=0} = \Gk(-\psi;0) = E/(E-\psi)$.

(iii) $\leaVn(z) = t\rho\,\vv_1(z) = t\brow(-z)$, so $\Kone(z) = -(\epsilon_1\epsilon_2E^3)^{-1}\langle t\brow(-z),t\brow(z)\rangle$. By (ii) this lies in $\bbC\fps{z,t}$ with value at $t = 0$ equal to $-(\epsilon_1\epsilon_2E^3)^{-1}\langle(0,E),(0,E)\rangle = -(\epsilon_1\epsilon_2E^3)^{-1}\cdot\epsilon_1\epsilon_2(-E)E^2 = 1$. By (i), $\Kone = 1$.

(iv) $\langle\brow(\psi),\brow(-\psi)\rangle = \rho^2\langle\vv_1(-\psi),\vv_1(\psi)\rangle = \rho^2\Delta_1\Kone(-\psi) = \rho^2\Delta_1$, and $\Pij_{11}(-w,w) = -\langle\vv_1(w),\vv_1(-w)\rangle = -\Delta_1\Kone(w) = -\Delta_1$; with $z = \psi$, $w = \psi'$ and $\rho^2\Pij_{11}(\psi,\psi') = -\langle\brow(\psi),\brow(\psi')\rangle$, the formula $\Edge_{11} = [\Pij_{11}(z,w) - \Pij_{11}(-w,w)]/(z+w)$ of \cref{lem:set-edgegen} gives \eqref{eq:lea-E11sigma}. Finally \eqref{eq:lea-qde} at $z = -\psi$, with $\leaVn(-\psi) = t\brow(\psi)$, is $-\psi(t\brow(\psi))' = (\Aop-\Lv_1)(t\brow(\psi))$; divide by $-\psi t$. The division by $\psi$ is legitimate: $\vv_1(0) = \rr_{0,1}[0]\chiv_1$ by Row 1, and $(\Aop-\Lv_1)\chiv_1 = 0$, so $(\Aop-\Lv_1)\brow$ vanishes at $\psi = 0$.
\end{proof}

Only the $t\to 0$ data (F2) and $\brow^{(-1)} = (0,E)$ enter the limit in (iii), and they come from Row 1 and \eqref{eq:lea-key} alone; the $t$-independence (i) uses Row 2, i.e.\ the full closed form.

\begin{rmk}[$\Kone = 1$ as a Wronskian identity for Gauss's equation]
\label{rmk:lea-wronsk}
By step (ii) of the proof of \cref{prop:lea-row2}, $\leaVn = (E/(E+z))\,\tau\Ephi\,\leaPhiF$ with $\leaPhiF := (\Fhyp,(\Dz+\omq)\Fhyp)$. Since $\Ephi(z)\Ephi(-z) = 1$ and $(E/(E+z))(E/(E-z)) = E^2/(E^2-z^2)$,
\[
\Kone(z) = -\frac{\tau^2\langle\leaPhiF(z),\leaPhiF(-z)\rangle}{\epsilon_1\epsilon_2E(E^2-z^2)},\qquad\text{so}\qquad \Kone = 1\iff \tau^2\langle\leaPhiF(z),\leaPhiF(-z)\rangle = -\epsilon_1\epsilon_2E(E^2-z^2).
\]
The latter follows from $\Gauss[\Fhyp] = 0$ alone: $(\Dz+\omq)\leaPhiF = \Aop\leaPhiF$ (Row 2 in $\Fhyp$-form, by \eqref{eq:lea-op}), and $z\mapsto -z$ gives $(-\Dz+\omq(-z))\leaPhiF(-z) = \Aop\leaPhiF(-z)$; by self-adjointness
\[
z\,\frac{d}{dt}\langle\leaPhiF(z),\leaPhiF(-z)\rangle = \bigl(\omq(-z)-\omq(z)\bigr)\langle\leaPhiF(z),\leaPhiF(-z)\rangle = -2z\,\frac{1-\tau}{\tau}\langle\leaPhiF(z),\leaPhiF(-z)\rangle,
\]
since $\omq(-z)-\omq(z) = -2z/\tau + 2z$; hence $\frac{d}{dt}\bigl(\tau^2\langle\leaPhiF(z),\leaPhiF(-z)\rangle\bigr) = 0$. At $\tau = 0$, $\tau\leaPhiF(z)\to(0,E+z)$, which gives $\epsilon_1\epsilon_2(-E)(E+z)(E-z)$. Given \eqref{eq:lea-key} this is a second proof of \cref{prop:lea-K1}(iii).
\end{rmk}

\begin{rmk}[flat sections; not used]
\label{rmk:lea-flat}
Fix $\kappa_0$ with $\kappa_0^2 = -\epsilon_1\epsilon_2E^3$ and put $\nu_1 := t\rho/\kappa_0$, so that $\nu_1^2 = 1/\Delta_1$. \Cref{prop:lea-row2} says that $\nu_1\vv_1$, multiplied by the formal symbol $e^{u_1/z}$ with $u_1' = \Lv_1$ (\cref{lem:row-canon}), is a flat section of $z\,d/dt - \Aop$: ``$z\,ds/dt = \Aop s$''. Since $\Lv_1 = \muq + \dphi'$, the identity-component of this section is $\kappa_1(z)\,\cQ^{\alpha_1/z}\tau^{E/z+1}\Fhyp$ with $\kappa_1(z) := E/(\kappa_0(E+z))$, $t$-independent, the only genuinely formal factor being $\tau^{E/z}$. Replacing $\rr_{0,1}$ by $\kappa\,\rr_{0,1}$ with $\kappa$ a unit of $\Cz$ replaces $\rr_{1,1}$ by $\kappa\,\rr_{1,1}$, preserves Row 2 and replaces $\Kone(z)$ by $\kappa(z)\kappa(-z)\Kone(z)$; for the frame change $\kappa = \hw$, $\hw(z)\hw(-z) = 1$, so $\Kone$ and the symplectic relations are the same in both frames. The indicial exponents of $\nabla^2 - (E/\tau-\alpha_0)\nabla + \alpha_1\alpha_2/\tau$ at $\tau = 0$ are $0$ and $1 + E/z$ (the $\alpha_0$ and $\alpha_1\alpha_2$ terms are of lower $\tau$-order); the singular section carries the exponent $1 + E/z$, the regular one (\cref{rmk:lea-E22}) the exponent $0$. Nothing below uses this remark.
\end{rmk}

\subsection{The dilaton leaves: (L1)}
\label{sec:lea-L1}

This subsection and the next two are written over the field $\Kal$ (the passage (P) of \cref{lem:set-rings}): every function of $t$ is an element of $\Kal((t))$, the expansion at $t = 0$ of the corresponding germ; products of series in $\Kal((t))\fps{\psi}$ or $\Kal((t))\fps{\psi,\psi'}$ are expanded order by order with finite sums (all of them lie in $t^{-2}\Kal\fps{t}\fps{\psi,\psi'}$); and elements of $\Kal(z)\fps{t}$, such as $\Gk$ and $z\Gk$, are compared with series in $z$ or $\psi$ only through \cref{prop:lea-key}(iii): $\Gk_k$ is regular at $z = 0$ with Taylor series $\coef{t^k}(\rho\,\rr_{0,1})$, and $\gl_k := \coef{t^k}\gl$ is the Taylor series of $\Gk_k(-\psi)$. Recall $\Fhyp = \sum_n\Fhyp_n\tau^n$ with
\[
\Fhyp_n = \frac{(\ga)_n(\gb)_n}{(\gc)_n\,n!} = \frac{\prod_{j=1}^{n}(\alpha_2+jz)(\wA+jz)}{n!\,z^n\prod_{j=2}^{n+1}(E+jz)}\in\Kal(z),
\]
$\Ephi_m := \coef{\tau^m}\Ephi$ and $\Hcoef_n := \coef{\tau^n}(\Fhyp\Ephi) = \sum_{n_1+n_2 = n}\Fhyp_{n_1}\Ephi_{n_2}$.

\begin{lem}[the $\tau$-coefficients]\genusstatus{AI-proposed, unchecked}
\label{lem:lea-taucoef}
Let $n,m\geq 0$.
\begin{enumerate}
\item $\Fhyp_n$ has poles only at $z = 0$ (of order exactly $n$) and at $z = -E/j$, $2\leq j\leq n+1$ (of order $\leq 1$), and $\Fhyp_n(\infty) = 1/(n+1)$.
\item $\Ephi_m$ is a polynomial in $1/z$ of degree $\leq m$ with $\Ephi_m(\infty) = \delta_{m0}$.
\item $\Hcoef_n$ has poles only at $z = 0$ (of order $\leq n$) and at $z = -E/j$, $2\leq j\leq n+1$ (of order $\leq 1$), and $\Hcoef_n(\infty) = 1/(n+1)$.
\end{enumerate}
\end{lem}

\begin{proof}
(i) The numerator $\prod_{j=1}^{n}(\alpha_2+jz)(\wA+jz)$ has degree $2n$, value $(\alpha_2\wA)^n\neq 0$ at $z = 0$ (so the pole at $0$ has order exactly $n$) and leading coefficient $(n!)^2$. The denominator $n!\,z^n\prod_{j=2}^{n+1}(E+jz)$ has degree $2n$, leading coefficient $n!\,(n+1)!$, a zero of order $n$ at $0$ and simple zeros at the $n$ distinct non-zero points $-E/j$ (distinct and non-zero because $E\neq 0$). A zero of the numerator (at $-\alpha_2/i$ or $-\wA/i$) can only lower a pole order; the poles at $-E/j$ are therefore of order $\leq 1$ (exactly $1$ over $\Kal$, and at every weight point with $E\notin\{j\alpha_2/i,\,j\wA/i : 1\leq i\leq n,\ 2\leq j\leq n+1\}$). $\Fhyp_n(\infty) = (n!)^2/(n!\,(n+1)!) = 1/(n+1)$, independent of the weights. Only the degree and the leading coefficient of the numerator are used, not the position or multiplicity of its zeros.
(ii) $\dphi\in\tau\Kal\fps{\tau}$ gives $\dphi^l\in\tau^l\Kal\fps{\tau}$, so $\Ephi_m = \sum_{l=0}^{m}\frac{(-1)^l}{l!}z^{-l}\coef{\tau^m}(\dphi^l)$, a polynomial in $1/z$ of degree $\leq m$ with constant term $\coef{\tau^m}\dphi^0 = \delta_{m0}$.
(iii) $\Ephi_{n_2}$ is regular away from $0$ with a pole of order $\leq n_2$ at $0$; the poles of $\Fhyp_{n_1}\Ephi_{n_2}$ are therefore at $0$ (order $\leq n_1+n_2 = n$) and at $-E/j$, $2\leq j\leq n_1+1\leq n+1$ (simple); at $\infty$, $\Fhyp_{n_1}\Ephi_{n_2}\to\delta_{n_2,0}/(n_1+1)$, so $\Hcoef_n(\infty) = 1/(n+1)$.
\end{proof}

\begin{lem}[the $t$-coefficients of $z\Gk$]\genusstatus{AI-proposed, unchecked}
\label{lem:lea-tcoef}
For $k\geq 0$ put $\Rk_k(z) := \coef{t^k}(z\Gk) = z\Gk_k(z)$. Then
\begin{equation}
\label{eq:lea-Rk}
\Rk_k = \frac{Ez}{E+z}\sum_{n=0}^{k}d_{k,n}\Hcoef_n,\qquad d_{k,n} := \coef{t^k}\Bigl(\frac{\tau}{t}\tau^n\Bigr)\in\bbQ,\qquad d_{k,k} = 1 .
\end{equation}
So $\Rk_k$ is a rational function with poles only at $z = 0$ and at $z = -E/j$, $1\leq j\leq k+1$, of order $\leq 1$ at each $-E/j$; it is bounded at $\infty$ and
\begin{equation}
\label{eq:lea-Rinf}
\Rk_k(\infty) = E\,\delta_{k0}.
\end{equation}
\end{lem}

\begin{proof}
$z\Gk = (Ez/(E+z))(\tau/t)\sum_n\Hcoef_n\tau^n$ with $\tau/t\in 1 + t\bbQ\fps{t}$ and $\tau^n\in t^n(1 + t\bbQ\fps{t})$; this is the principle (T) of \cref{lem:set-rings} and gives the finite sum \eqref{eq:lea-Rk}. Poles: by \cref{lem:lea-taucoef}(iii) each $\Hcoef_n$ with $n\leq k$ has poles only at $0$ and at $-E/j$ with $2\leq j\leq n+1\leq k+1$, simple at the latter; the factor $Ez/(E+z)$ adds a simple pole at $-E$ (this is $j = 1$), which is not one of the $-E/j$ with $j\geq 2$. Limit: $Ez/(E+z)\to E$ and $\Hcoef_n\to 1/(n+1)$, so
\[
\Rk_k(\infty) = E\sum_{n\leq k}\frac{d_{k,n}}{n+1} = E\coef{t^k}\Bigl(\frac{\tau}{t}\sum_{n\geq 0}\frac{\tau^n}{n+1}\Bigr) = E\coef{t^k}\Bigl(\frac{-\log(1-\tau)}{t}\Bigr) = E\coef{t^k}1 = E\delta_{k0},
\]
since the terms $n > k$ do not contribute to $\coef{t^k}$ and $-\log(1-\tau) = t$.
\end{proof}

The limit is coefficientwise in $\tau$: $\Fhyp(\infty,\tau) := \sum_n\Fhyp_n(\infty)\tau^n = -\log(1-\tau)/\tau$ is the limit of $\Fhyp$ coefficient by coefficient in $\tau$, for every weight point, because $\Fhyp_n(\infty) = 1/(n+1)$ does not depend on the weights; only this is used, and nothing is claimed about ${}_2F_1(\ga,\gb;\gc;\tau)$ as $z\to\infty$ at fixed $\tau$. For $k\geq 1$ the limit vanishes because $(\tau/t)\Fhyp(\infty,\tau) = 1$ is $t$-independent: this is the precise sense in which the normalisation by $\rho$ removes $\lf_0$. Equivalently $\lim_{z\to\infty}z\,\rr_{0,1}(z;t) = E/\rho(t)$, coefficientwise in $t$.

\begin{rmk}[no pole at $z = 0$]
\label{rmk:lea-nopole}
By \cref{prop:lea-key}(iii), $\Gk_k$ is regular at $z = 0$, so $\Rk_k = z\Gk_k$ is regular at $0$ with $\Rk_k(0) = 0$. In \eqref{eq:lea-Rk} the separate terms do have poles at $0$ ($z\Hcoef_n$ can have a pole of order $n-1$); that they cancel is the regularity theorem of \cref{sec:row}: every $\Hcoef_n$ is regular at $z = 0$, for either of two reasons: (a) \cref{cor:row-reg}, which rests on the closed form \cref{thm:row-closed} and on the fact that $\Abeta$ is a unit of $\bbC\fps{z}$ (its value is not used): $\Hcoef_n = \sum_{m\leq n}(\coef{\tau^{n-m}}(1/\alg))\Kcoef_m$ with every $\Kcoef_m$ regular at $0$; (b) \cref{prop:row-transport}, purely formal: $\Fhyp\Ephi = \Phia$ has $z$-coefficients analytic in $\tau$, so the Laurent expansion at $z = 0$ of its $\tau^n$-coefficient $\Hcoef_n$ has no negative powers. Conversely the regularity of all $\Hcoef_n$ gives that of all $\Gk_k$ by (T), since $E/(E+z)$ is a unit at $z = 0$ and $\tau/t$ is a $z$-independent unit of $\Kal\fps{t}$. By hand (\cref{rmk:row-first}): $\Hcoef_1 = \Fhyp_1 + \Ephi_1 = [Ez + E(\alpha_2+\wA) - 2\alpha_2\wA]/(E(E+2z))$, in which the poles of $\Fhyp_1 = (\alpha_2+z)(\wA+z)/(z(E+2z))$ and $\Ephi_1 = -\alpha_2\wA/(Ez)$ at $z = 0$ cancel.
\end{rmk}

\begin{prop}[(L1) in the variable $z$: partial fractions]\genusstatus{AI-proposed, unchecked}
\label{prop:lea-pf}
$\Rk_0 = Ez/(E+z) = E - E^2/(E+z)$. For $k\geq 1$ there are $a_j^{(k)}\in\Kal$, $1\leq j\leq k+1$, with
\begin{equation}
\label{eq:lea-pf}
\Rk_k(z) = \coef{t^k}(z\Gk) = \sum_{j=1}^{k+1}\frac{a_j^{(k)}}{E+jz},\qquad a_j^{(k)} = j\Res_{z=-E/j}\Rk_k,\qquad \sum_{j=1}^{k+1}a_j^{(k)} = 0 .
\end{equation}
\end{prop}

\begin{proof}
For $k = 0$, $(\tau/t)\Fhyp\Ephi = \Hcoef_0 = 1$ at $t = 0$, so $\Gk_0 = E/(E+z)$. For $k\geq 1$, \cref{lem:lea-tcoef} and \cref{rmk:lea-nopole} show that the poles of $\Rk_k$ lie in $\{-E/j : 1\leq j\leq k+1\}$ and are simple, and $\Rk_k(\infty) = 0$ by \eqref{eq:lea-Rinf}; so the partial-fraction principle (PF) of \cref{lem:set-leafid}, with $S = \{-E/j : 1\leq j\leq k+1\}$ (distinct and non-zero since $E\neq 0$), gives \eqref{eq:lea-pf} without constant term ($f(\infty) = 0$: a rational function with simple poles that tends to $0$ at $\infty$ is the sum of its principal parts, and $\Res_{z=-E/j}$ of $1/(E+jz)$ is $1/j$). Evaluating at $z = 0$ gives $0 = \Rk_k(0) = E^{-1}\sum_ja_j^{(k)}$, and $E\neq 0$.
\end{proof}

With $a_1^{(0)} := -E^2$ and $a_j^{(k)} := 0$ for $j > k+1$, the sum rule reads $\sum_ja_j^{(k)} = -E^2\delta_{k0}$ for every $k\geq 0$.

\begin{thm}[(L1), all $k$]\genusstatus{AI-proposed, unchecked}
\label{thm:lea-L1}
Write $q^{(k)} := \coef{t^k}(\rho\qt)$. Then $q^{(0)} = \qo = E - E^2\lf_1$ and, for $k\geq 1$,
\[
q^{(k)} = \sum_{j=1}^{k+1}a_j^{(k)}\lf_j\in\Lsp_{k+1},
\]
with the $a_j^{(k)}$ of \eqref{eq:lea-pf}. Moreover $\gl_0 = E\lf_1$, and for $k\geq 1$
\begin{equation}
\label{eq:lea-gk}
\gl_k = \sum_{j=1}^{k+1}g_j^{(k)}\lf_j = -\frac{1}{E}\sum_{j=1}^{k+1}j\,a_j^{(k)}\lf_j\in\Lsp_{k+1},\qquad \gl_k = O(\psi^{-2})\ \text{at}\ \psi = \infty .
\end{equation}
\end{thm}

\begin{proof}
By \cref{def:set-factors}, $\qt(\psi) = -\psi\,\rr_{0,1}(-\psi;t)$. Multiply by $\rho$, which does not involve $\psi$: $\rho\qt(\psi) = -\psi\,\gl(\psi;t)$, and $\coef{t^k}(\rho\qt)(\psi) = \Rk_k(-\psi)$ as formal series (Taylor expansion at $0$ commutes with $z\mapsto -\psi$). Since $1/(E+jz) = \lf_j(\psi)$ at $z = -\psi$, \cref{prop:lea-pf} gives the first two statements. For \eqref{eq:lea-gk}: $\Gk_k = \Rk_k/z$ and $1/(z(E+jz)) = E^{-1}(1/z - j/(E+jz))$, so by the sum rule
\[
\Gk_k(z) = \frac{1}{Ez}\sum_ja_j^{(k)} - \frac{1}{E}\sum_j\frac{j\,a_j^{(k)}}{E+jz} = -\frac{1}{E}\sum_j\frac{j\,a_j^{(k)}}{E+jz};
\]
substitute $z = -\psi$. Since $\Rk_k = O(1/z)$ at $\infty$ for $k\geq 1$, $\Gk_k = O(z^{-2})$, i.e.\ $\gl_k = O(\psi^{-2})$; equivalently, by (LAM4), $\sum_jg_j^{(k)}/j = -E^{-1}\sum_ja_j^{(k)} = 0$.
\end{proof}

Remarks. (i) $g_j^{(k)} = -(j/E)a_j^{(k)}$ for all $k\geq 0$ (at $k = 0$, $g_1^{(0)} = E$). (ii) Conversely, by (LAM2), (L1) at order $k\geq 1$ is equivalent to ``$\gl_k\in\Lsp_{k+1}$ and $\gl_k = O(\psi^{-2})$'': $-\psi\gl_k = \sum_jg_j^{(k)}/j + \sum_ja_j^{(k)}\lf_j$, and the constant must vanish. This is the form used in \cref{sec:lea-sigma}. (iii) The three ingredients are separate: the pole structure (\cref{lem:lea-taucoef,lem:lea-tcoef}), the regularity theorem (\cref{rmk:lea-nopole}) and the coefficientwise limit at $z = \infty$ (\cref{lem:lea-tcoef}). That every pole $-E/j$, $j\leq k+1$, actually occurs is proved for the top pole when $\Psh_k\neq 0$ (\cref{cor:lea-top}). (iv) (F3) for $\Dil_1$: $\Dil_1(d) = (-1)^d\coef{z^{d-1}}\rr_{0,1} = (-1)^d\coef{z^{d-1}}(\rho^{-1}\Gk)$ lies in $\Kal\fps{t}$ (\cref{prop:lea-key}), with value $-E^{1-d}$ at $t = 0$, i.e.\ $\qt|_{t=0} = \qo = -E\psi/(E-\psi)$; in particular $\qt$ and $\rho\qt$ lie in the domain $\Dom = \qo + t\RRing\fps{\psi}$ of \cref{def:set-fpot}.

\subsection{The H-row in leaf form}
\label{sec:lea-sigma}

By \eqref{eq:lea-conj} and \eqref{eq:lea-Hrow},
\begin{equation}
\label{eq:lea-sigma1}
\begin{gathered}
\rho\,\rr_{1,1}(z) = \Bigl(\Lv_1+\frac{z}{t}\Bigr)\Gk + z\,\frac{d\Gk}{dt},\qquad
\brow(\psi) = \rho\,\vv_1(-\psi) = (\brow_0,\brow_1),\\
\brow_0 = \gl,\qquad \brow_1 = \Lv_1\gl - \frac{\psi}{t}\Dg = \Bigl(\Lv_1-\frac{\psi}{t}\Bigr)\gl - \psi\frac{d\gl}{dt}
\end{gathered}
\end{equation}
($z/t$ becomes $-\psi/t$ and $z\,d\Gk/dt$ becomes $-\psi\,d\gl/dt$ under $z = -\psi$; $\Dg = \frac{d}{dt}(t\gl) = \sum_k(k+1)\gl_kt^k$). Both components of $t\brow$ lie in $\Kal\fps{t}\fps{\psi}$, because $t\Lv_1\in\Kal\fps{t}$. Only (F1) enters \eqref{eq:lea-sigma1}: $t^2\rho^2\Delta_1 = -\epsilon_1\epsilon_2E^3$ holds by the definition of $\rho$, so the logarithmic derivative of $\Delta_1$ is that of $(t\rho)^{-2}$ and no square root of $\Delta_1$ is needed. By \cref{lem:set-eig},
\begin{equation}
\label{eq:lea-L1exp}
\Lv_1 = \frac{E}{t} + \lrem(t),\qquad \lrem = \frac{\wA}{\tau\sast} - \alpha_0 - \frac{E}{t}\in\Kal\fps{t},\qquad \lrem_0 = \frac{E}{2} - \alpha_0 - \frac{\alpha_1\alpha_2}{E}.
\end{equation}
The structural fact used below is that the coefficient of the pole of $\Lv_1$, $t\Lv_1(0) = E$, equals the leaf scale $E$ of $\lf_j = 1/(E-j\psi)$ and the first weight of $\bw$. It holds because $\Lv_1 = \wA/(\tau\sast) - \alpha_0$ with $\sast(0) = \wA/E$.

\begin{prop}[the leaf structure of $\brow$]\genusstatus{AI-proposed, unchecked}
\label{prop:lea-sigma}
\begin{enumerate}
\item $\brow^{(-1)} = \coef{t^{-1}}\brow = (0,E)$.
\item For $k\geq 0$, $\coef{t^k}\brow_0 = \gl_k$ and
\begin{equation}
\label{eq:lea-sigk}
\coef{t^k}\brow_1 = \bigl(E - (k+2)\psi\bigr)\gl_{k+1} + \sum_{m=0}^{k}\lrem_m\,\gl_{k-m}\in\Lsp_{k+1}.
\end{equation}
\end{enumerate}
Hence $\brow = t^{-1}(0,E) + \sum_{k\geq 0}t^k\brow^{(k)}$ with both components of $\brow^{(k)}$ in $\Lsp_{k+1}$; explicitly
\begin{equation}
\label{eq:lea-Sj}
\brow^{(k)} = \sum_{j=1}^{k+1}S_j^{(k)}\lf_j,\qquad S_j^{(k)} := \Bigl(g_j^{(k)},\ (k+2-j)\,a_j^{(k+1)} + \sum_{m=0}^{k}\lrem_m\,g_j^{(k-m)}\Bigr)
\end{equation}
($g_j^{(n)} := 0$ for $j > n+1$). In particular $\brow^{(0)} = (E\lf_1,\ \alpha_1\alpha_2\lf_1) = E\lf_1\,(1,\Lv_2(0))$, and $\gl_1 = \frac{\uq}{2E}(2\lf_2 - \lf_1)$.
\end{prop}

\begin{proof}
By \eqref{eq:lea-sigma1} and \eqref{eq:lea-L1exp}, $\coef{t^k}(\Lv_1\gl) = E\gl_{k+1} + \sum_{m=0}^{k}\lrem_m\gl_{k-m}$ and $\coef{t^k}((\psi/t)\Dg) = \psi\coef{t^{k+1}}\Dg = (k+2)\psi\gl_{k+1}$; this gives \eqref{eq:lea-sigk} for every $k\geq -1$ (empty sum for $k = -1$). For $k = -1$, $\coef{t^{-1}}\brow_1 = (E-\psi)\gl_0 = (E-\psi)E\lf_1 = E$ by (LAM3), and $\coef{t^{-1}}\brow_0 = 0$: this is (i).

Leaf structure. Write $\gl_{k+1} = \sum_{j=1}^{k+2}g_j^{(k+1)}\lf_j$. By \cref{thm:lea-L1} (as $k+1\geq 1$) and (LAM4), $\sum_jg_j^{(k+1)}/j = 0$. By (LAM3) with scale $E$, $(E-m\psi)\lf_j = m/j + E(1-m/j)\lf_j$, and $g_j^{(k+1)} = -(j/E)a_j^{(k+1)}$,
\[
\bigl(E-(k+2)\psi\bigr)\gl_{k+1} = (k+2)\sum_j\frac{g_j^{(k+1)}}{j} + E\sum_j\Bigl(1-\frac{k+2}{j}\Bigr)g_j^{(k+1)}\lf_j = \sum_{j=1}^{k+1}(k+2-j)\,a_j^{(k+1)}\lf_j :
\]
the factor $E-(k+2)\psi$ turns the top leaf $\lf_{k+2}$ into a constant, and all the constants add up to $(k+2)\sum_jg_j^{(k+1)}/j = 0$. The terms $\lrem_m\gl_{k-m}$ lie in $\Lsp_{k-m+1}\subseteq\Lsp_{k+1}$. Together with $\coef{t^k}\brow_0 = \gl_k$ this is \eqref{eq:lea-Sj}. (In the variable $z$: \eqref{eq:lea-sigk} at $\psi = -z$ is $(E+(k+2)z)\Gk_{k+1} + \sum_m\lrem_m\Gk_{k-m}$; the factor $E+(k+2)z$ kills the simple pole of $\Gk_{k+1}$ at $-E/(k+2)$; the result is regular at $0$, $O(1/z)$ at $\infty$, with simple poles in $\{-E/j : j\leq k+1\}$, so (PF) puts it in $\operatorname{span}\{1/(E+jz) : j\leq k+1\}$.)

The case $k = 0$, by hand. From \cref{rmk:row-first} and \eqref{eq:lea-key}, since $\tau/t = 1 - t/2 + O(t^2)$ and $\Hcoef_0 = 1$,
\[
\Gk_1 = \frac{E}{E+z}\Bigl(\Hcoef_1 - \frac12\Bigr) = \frac{2Ez + 2E(\alpha_2+\wA) - 4\alpha_2\wA - E^2 - 2Ez}{2(E+z)(E+2z)} = \frac{\uq}{2(E+z)(E+2z)},
\]
because $2E(\alpha_2+\wA) - 4\alpha_2\wA - E^2 = -(E-2\alpha_2)(E-2\wA) = -(E-2\alpha_2)(2\alpha_1-E) = \uq$. (Equivalently $\Gk_1 = \coef{t^1}\rr_{0,1} + \rho_1\Gk_0$ with $\rho_1 = \uq/(2E^2)$, since $\Gk = \rho\,\rr_{0,1}$.) So $\Rk_1 = z\Gk_1 = \frac{\uq}{2}\bigl[\frac{1}{E+z} - \frac{1}{E+2z}\bigr]$, $(a_1^{(1)},a_2^{(1)}) = \frac{\uq}{2}(1,-1)$, $q^{(1)} = \frac{\uq}{2}(\lf_1-\lf_2)$, and $\gl_1 = -E^{-1}(a_1^{(1)}\lf_1 + 2a_2^{(1)}\lf_2) = \frac{\uq}{2E}(2\lf_2-\lf_1)$. Then $(E-2\psi)\gl_1 = a_1^{(1)}\lf_1 = \frac{\uq}{2}\lf_1$ (by (LAM3): $(E-2\psi)\lf_2 = 1$ and $(E-2\psi)\lf_1 = 2 - E\lf_1$), $\lrem_0\gl_0 = (E^2/2 - \alpha_0E - \alpha_1\alpha_2)\lf_1$, and
\[
\coef{t^0}\brow_1 = \Bigl(\frac{\uq}{2} + \frac{E^2}{2} - \alpha_0E - \alpha_1\alpha_2\Bigr)\lf_1 = \alpha_1\alpha_2\lf_1,
\]
since $(\uq+E^2)/2 = \alpha_0E + 2\alpha_1\alpha_2$ (\cref{def:set-ND}). With $\Lv_2(0) = \alpha_1\alpha_2/E$ this is $\brow^{(0)} = E\lf_1(1,\Lv_2(0))$.
\end{proof}

Remarks. (i) $\brow^{(0)}_1 = \alpha_1\alpha_2\lf_1$ is the $z$-dependent part $\alpha_1\alpha_2/(E+z)$ of the H-row at order $t^0$: $\rho\,\rr_{1,1}(z) = E/t + \alpha_1\alpha_2/(E+z) + O(t)$, the second half of (F2). (ii) The two cancellations in the $k = 0$ computation, $2E(\alpha_2+\wA) - 4\alpha_2\wA - E^2 = \uq$ and $\uq/2 + \lrem_0E = \alpha_1\alpha_2$, use only $E = \alpha_1+\wA = \alpha_2+\wB$ and $\wA\wB = \alpha_1\alpha_2 + \alpha_0E$. (iii) $\brow_1$ is $O(1/\psi)$ but not $O(\psi^{-2})$; this does not matter for the leaf structure. (iv) Part (i) is the same statement as \cref{prop:lea-K1}(ii), proved there from Row 1 and \eqref{eq:lea-key}; here it comes out of the leaf expansion, using only (F2), \eqref{eq:lea-L1exp} and \eqref{eq:lea-sigma1}.

\begin{prop}[the projection of $\brow$ on the regular eigen-covector]\genusstatus{AI-proposed, unchecked}
\label{prop:lea-proj}
$\pim(\brow(\psi)) = -\Thf\,\frac{\psi}{t}\Dg(\psi) = \epsilon_1\epsilon_2\wB\tau\sast\Lv_2\,\frac{\psi}{t}\Dg(\psi)$ in $\Kal\fps{t}\fps{\psi}$, with $\Thf = \epsilon_1\epsilon_2(\alpha_1\alpha_2 - E\Lv_2)$.
\end{prop}

\begin{proof}
$\brow = \gl\,\chiv_1 + (\brow_1 - \Lv_1\gl)(0,1)$, since $\chiv_1 = (1,\Lv_1)$. Now $\pim(\chiv_1) = \langle\chiv_1,\chiv_2\rangle = 0$, $\pim((0,1)) = \epsilon_1\epsilon_2(\alpha_1\alpha_2 - E\Lv_2) = \Thf$ (\cref{prop:lea-A}(v)), and $\brow_1 - \Lv_1\gl = -(\psi/t)\Dg$ by \eqref{eq:lea-sigma1}.
\end{proof}

\subsection{The \texorpdfstring{$\Edge_{11}$}{E11} flow: (F3) for \texorpdfstring{$\Edge_{11}$}{E11}, (L3), and (L2) for \texorpdfstring{$j = 1$}{j = 1}}
\label{sec:lea-flow}

By \cref{lem:set-edgegen} and \cref{prop:lea-K1}(iv), in the leaf variables $(\psi,\psi') = (z,w)$,
\[
\rho^2\Edge_{11}(\psi,\psi') = \frac{\rho^2\Delta_1 - \langle\brow(\psi),\brow(\psi')\rangle}{\psi+\psi'},\qquad \rho^2\Delta_1 = -\frac{\epsilon_1\epsilon_2E^3}{t^2}.
\]
Put $\Yv(\psi) := t\brow(\psi)$, an element of $\Kal\fps{t}\fps{\psi}^2$ by \cref{prop:lea-sigma}, and
\begin{equation}
\label{eq:lea-Ebar}
\Ebar(\psi,\psi') := t^2\rho^2\Edge_{11}(\psi,\psi') = \frac{\Nbar(\psi,\psi')}{\psi+\psi'},\qquad \Nbar := -\epsilon_1\epsilon_2E^3 - \langle\Yv(\psi),\Yv(\psi')\rangle .
\end{equation}
$\Nbar$ lies in $\Kal\fps{t}\fps{\psi,\psi'}$, and $\Nbar(\psi,-\psi) = -\epsilon_1\epsilon_2E^3 - t^2\rho^2\langle\vv_1(-\psi),\vv_1(\psi)\rangle = -\epsilon_1\epsilon_2E^3 - t^2\rho^2\Delta_1 = 0$ by \cref{prop:lea-K1}(iii) and \eqref{eq:lea-rho}. So $\Nbar$ is divisible by $\psi+\psi'$ (in the variables $u = \psi$, $v = \psi+\psi'$ it vanishes at $v = 0$), $\Ebar$ lies in $\Kal\fps{t}\fps{\psi,\psi'}$ (quotients by the non-zero-divisor $\psi+\psi'$ are unique), and $\rho^2\Edge_{11}$ lies in $t^{-2}\Kal\fps{t}\fps{\psi,\psi'}$. Since $\langle\,,\rangle$ is symmetric, $\Ebar(\psi,\psi') = \Ebar(\psi',\psi)$, i.e.\ $\Edge_{11}(a,b) = \Edge_{11}(b,a)$. (Without $\Kone = 1$ the subtracted term would be $-\Kone(w)\Delta_1$ and the symmetry would not be manifest.)

\begin{lem}[the equation for $\Yv$]\genusstatus{AI-proposed, unchecked}
\label{lem:lea-Yqde}
$-\psi\,\frac{d\Yv(\psi)}{dt} = (\Aop-\Lv_1)\Yv(\psi)$ in $\Kal((t))\fps{\psi}^2$; equivalently $\frac{d\brow}{dt} = -\frac{\brow}{t} - \frac1\psi(\Aop-\Lv_1)\brow$.
\end{lem}

\begin{proof}
This is \cref{prop:lea-K1}(iv): the equation \eqref{eq:lea-qde} at $z = -\psi$, with $\leaVn(-\psi) = \Yv(\psi)$; $\Aop$ and $\Lv_1$ do not involve $z$. Sign: $\brow$ is built from $\vv_1(-\psi)$, so $z = -\psi$ and the $1/z$ of the equation becomes $-1/\psi$.
\end{proof}

\begin{thm}[the $\Edge_{11}$ flow]\genusstatus{AI-proposed, unchecked}
\label{thm:lea-flow}
In $\Kal\fps{t}\fps{\psi,\psi'}$,
\begin{align}
\label{eq:lea-flow}
\frac{d}{dt}\Ebar(\psi,\psi') &= \Thf(t)\,\Dg(\psi)\,\Dg(\psi'),\qquad \Dg = \frac{d}{dt}(t\gl),\\
\label{eq:lea-Theta}
\Thf &= \frac{\Delta_2}{\Lv_2-\Lv_1} = \epsilon_1\epsilon_2(\alpha_1\alpha_2 - E\Lv_2) = -\epsilon_1\epsilon_2\wB\tau\sast\Lv_2 = \epsilon_1\epsilon_2(E\dphi' - \alpha_2\wA).
\end{align}
\end{thm}

\begin{proof}
(i) Since $t^2\rho^2\Delta_1$ is constant, only the pairing term of $\Nbar$ depends on $t$: $d\Nbar/dt = -\langle d\Yv(\psi)/dt,\Yv(\psi')\rangle - \langle\Yv(\psi),d\Yv(\psi')/dt\rangle$. Multiply by $\psi\psi'$ and use \cref{lem:lea-Yqde} and self-adjointness (\cref{prop:lea-A}(iii)):
\[
\psi\psi'\frac{d\Nbar}{dt} = \psi'\langle(\Aop-\Lv_1)\Yv(\psi),\Yv(\psi')\rangle + \psi\langle\Yv(\psi),(\Aop-\Lv_1)\Yv(\psi')\rangle = (\psi+\psi')\langle(\Aop-\Lv_1)\Yv(\psi),\Yv(\psi')\rangle .
\]
(ii) $\psi+\psi'$ is $t$-independent, so $(\psi+\psi')\,d\Ebar/dt = d\Nbar/dt$, and after cancelling the non-zero-divisor $\psi+\psi'$, $\psi\psi'\,d\Ebar/dt = \langle(\Aop-\Lv_1)\Yv(\psi),\Yv(\psi')\rangle$. This is why the factor $\psi+\psi'$ disappears: the flows of the two arguments contribute $1/\psi$ and $1/\psi'$, self-adjointness makes the two terms equal, and $1/\psi + 1/\psi' = (\psi+\psi')/(\psi\psi')$.
(iii) By \eqref{eq:lea-proj} and \cref{prop:lea-proj}, $\langle(\Aop-\Lv_1)\Yv(\psi),\Yv(\psi')\rangle = \pim(\Yv(\psi))\pim(\Yv(\psi'))/\Thf$ with $\pim(\Yv(\psi)) = t\,\pim(\brow(\psi)) = -\Thf\,\psi\,\Dg(\psi)$; so $\psi\psi'\,d\Ebar/dt = \Thf\,\psi\psi'\,\Dg(\psi)\Dg(\psi')$; cancel $\psi\psi'$.
(iv) The forms of $\Thf$ are \cref{prop:lea-A}(v) and \cref{lem:set-Delta2}; and $\dphi' = \alpha_2 - \Lv_2$ (\cref{lem:set-N}) gives $\alpha_1\alpha_2 - E\Lv_2 = \alpha_1\alpha_2 - E\alpha_2 + E\dphi' = E\dphi' - \alpha_2\wA$, since $\alpha_2E - \alpha_1\alpha_2 = \alpha_2\wA$.
\end{proof}

Expansion of $\Thf$. $\Thf_0 = \epsilon_1\epsilon_2(E\dphi_1 - \alpha_2\wA) = 0$, since $\dphi_1 := \coef{t^1}\dphi = \alpha_2\wA/E$ (\cref{rmk:row-first}); $\Thf = \epsilon_1\epsilon_2E\dphi' - \epsilon_1\epsilon_2\alpha_2\wA$ gives $\Thf_k = (k+1)\epsilon_1\epsilon_2E\coef{t^{k+1}}\dphi$ for $k\geq 1$; and from the explicit form, $\Thf_1 = -\epsilon_1\epsilon_2\wB\sast(0)\Lv_2(0) = -\epsilon_1\epsilon_2\wB(\wA/E)(\alpha_1\alpha_2/E) = -\epsilon_1\epsilon_2\alpha_1\alpha_2\wA\wB/E^2$ (as $\tau = t + O(t^2)$); hence $\coef{t^2}\dphi = -\alpha_1\alpha_2\wA\wB/(2E^3)$. Two remarks. (a) $\Thf = \Delta_2/(\Lv_2-\Lv_1)$: the loop flow is governed by the regular idempotent; $\Delta_2$ is regular and non-zero while $\Lv_2-\Lv_1 = -\delta/\tau$ has a simple pole, so $\Thf = O(t)$. (b) Rank one: before expanding in $t$, the right side of \eqref{eq:lea-flow} is a product of a function of $(t,\psi)$ and a function of $(t,\psi')$, because $\Aop-\Lv_1$ has rank one (\cref{prop:lea-A}(v)). Its $t^k$-coefficients are not of rank one in general; they are sums of the rank-one terms $\Thf_{k_0}(k_1+1)(k_2+1)\gl_{k_1}(\psi)\gl_{k_2}(\psi')$ (\cref{thm:lea-L3}). What rank one gives is that every loop coefficient is built from products of the same leaf series $\gl_k$.

\paragraph{Integration from $t = 0$; (F3) for $\Edge_{11}$.}
By \cref{prop:lea-sigma}(i), $\coef{t^0}\Yv = (0,E)$, and $\langle(0,E),(0,E)\rangle = \epsilon_1\epsilon_2(-E)E^2 = -\epsilon_1\epsilon_2E^3$. Hence $\coef{t^0}\Nbar = 0$ and $\coef{t^0}\Ebar = 0$. This is where $\Kone = 1$ enters: for a $t$-independent $\Kone$ one would have $\Pij_{11}(-w,w) = -\Kone(w)\Delta_1$ and $\coef{t^0}\Ebar = -\epsilon_1\epsilon_2E^3(\Kone(\psi')-1)/(\psi+\psi')$. The kernel of $d/dt$ on $\Kal\fps{t}$ is $\Kal$ and $\coef{t^{n+1}}\Ebar = \coef{t^n}(d\Ebar/dt)/(n+1)$, so \eqref{eq:lea-flow} and $\coef{t^0}\Ebar = 0$ give
\begin{equation}
\label{eq:lea-int}
\Ebar = \int_0^t\Thf(u)\,\Dg(\psi;u)\,\Dg(\psi';u)\,du\ \ \text{(termwise)},\qquad\text{i.e.}\qquad \coef{t^{n+1}}\Ebar = \frac{1}{n+1}\coef{t^n}\bigl(\Thf\,\Dg(\psi)\Dg(\psi')\bigr).
\end{equation}
As $\Thf_0 = 0$, also $\coef{t^1}\Ebar = 0$. So $\Ebar$ lies in $t^2\Kal\fps{t}\fps{\psi,\psi'}$, and $\Edge_{11} = (t\rho)^{-2}\Ebar$ lies in $\Kal\fps{t}\fps{\psi,\psi'}$ since $\rho^{-2}\in\Kal\fps{t}$ (\cref{prop:lea-key}): $\Edge_{11}$ is regular at $t = 0$. This is (F3) for $\Edge_{11}$; in frame B the double pole of $\Delta_1$ is exactly compensated.

\begin{thm}[(L3), all $k$]\genusstatus{AI-proposed, unchecked}
\label{thm:lea-L3}
For every $k\geq 0$,
\begin{equation}
\label{eq:lea-L3}
\begin{aligned}
\coef{t^k}(\rho^2\Edge_{11})(\psi,\psi') &= \frac{1}{k+2}\sum_{\substack{k_0\geq 1,\ k_1,k_2\geq 0\\ k_0+k_1+k_2 = k+1}}\Thf_{k_0}(k_1+1)(k_2+1)\,\gl_{k_1}(\psi)\gl_{k_2}(\psi')\\
&= \sum_{\substack{i,l\geq 1\\ (i-1)+(l-1)\leq k}}e^{(k)}_{il}\,\lf_i(\psi)\lf_l(\psi'),\\
e^{(k)}_{il} &= \frac{1}{k+2}\sum_{\substack{k_0\geq 1\\ k_1+k_2 = k+1-k_0}}\Thf_{k_0}(k_1+1)(k_2+1)\,g_i^{(k_1)}g_l^{(k_2)},\qquad e^{(k)}_{il} = e^{(k)}_{li},
\end{aligned}
\end{equation}
with $g_i^{(m)} := 0$ for $i > m+1$. There is no $\lf_0$ in either variable. On the anti-diagonal $(i-1)+(l-1) = k$, i.e.\ $i = k_1+1$, $l = k_2+1$ with $k_1+k_2 = k$,
\begin{equation}
\label{eq:lea-anti}
e^{(k)}_{k_1+1,k_2+1} = (-1)^{k+1}\,\epsilon_1\epsilon_2\alpha_1\alpha_2\wA\wB\,\frac{(k_1+1)(k_2+1)\,\Psh_{k_1}\Psh_{k_2}}{(k+2)\,(k_1!)^2(k_2!)^2\,E^{2k}},
\end{equation}
which is non-zero if and only if $\Psh_{k_1}\Psh_{k_2}\neq 0$. At $t = 0$ the loop has rank one: $\coef{t^0}(\rho^2\Edge_{11}) = \frac12\Thf_1\gl_0(\psi)\gl_0(\psi') = -\frac{\epsilon_1\epsilon_2\alpha_1\alpha_2\wA\wB}{2}\lf_1(\psi)\lf_1(\psi')$, hence
\[
\Edge_{11}(z,w)\big|_{t=0} = -\frac{\epsilon_1\epsilon_2\alpha_1\alpha_2\wA\wB}{2}\,\lf_1(z)\lf_1(w)
\]
(with $\rho(0) = 1$).
\end{thm}

\begin{proof}
$\coef{t^k}(\rho^2\Edge_{11}) = \coef{t^{k+2}}\Ebar = \coef{t^{k+1}}(\Thf\,\Dg(\psi)\Dg(\psi'))/(k+2)$ by \eqref{eq:lea-int}; insert $\Thf = \sum_{k_0\geq 1}\Thf_{k_0}t^{k_0}$ and $\Dg = \sum_k(k+1)\gl_kt^k$. Support: every term has $k_1+k_2 = k+1-k_0\leq k$; by \cref{thm:lea-L1}, $\gl_{k_1}\in\Lsp_{k_1+1}$ and $\gl_{k_2}\in\Lsp_{k_2+1}$, so the term lies in $\operatorname{span}\{\lf_i(\psi)\lf_l(\psi') : i\leq k_1+1,\ l\leq k_2+1\}$, of total degree $(i-1)+(l-1)\leq k_1+k_2\leq k$. Symmetry: \eqref{eq:lea-L3} is invariant under $(\psi,k_1)\leftrightarrow(\psi',k_2)$. Anti-diagonal: a product $\lf_{k_1+1}(\psi)\lf_{k_2+1}(\psi')$ with $k_1+k_2 = k$ comes only from terms with $k_1'\geq k_1$, $k_2'\geq k_2$ and $k_1'+k_2'\leq k$, i.e.\ $k_1' = k_1$, $k_2' = k_2$ and $k_0 = 1$; its coefficient is $\frac{1}{k+2}\Thf_1(k_1+1)(k_2+1)g^{(k_1)}_{k_1+1}g^{(k_2)}_{k_2+1}$. Insert $\Thf_1 = -\epsilon_1\epsilon_2\alpha_1\alpha_2\wA\wB/E^2$ and $g^{(m)}_{m+1} = -\frac{m+1}{E}a^{(m)}_{m+1} = (-1)^m\Psh_m/((m!)^2E^{2m-1})$ (\cref{cor:lea-top}, proved in \cref{sec:lea-closed}, independently of this theorem); collecting, the coefficient is $(-1)^{k_1+k_2}(-1)\epsilon_1\epsilon_2\alpha_1\alpha_2\wA\wB(k_1+1)(k_2+1)\Psh_{k_1}\Psh_{k_2}/((k+2)(k_1!)^2(k_2!)^2E^{2+2k_1-1+2k_2-1})$, and $2k_1+2k_2 = 2k$. The value at $t = 0$ is the case $k = 0$ with $\gl_0 = E\lf_1$.
\end{proof}

Generating form: the $\lf_i$-coefficient of $\Dg$ is $-(i/E)\Ahat_i'$ (\cref{thm:lea-Aj}), so
\[
e_{il}(t) := \sum_ke^{(k)}_{il}t^k = \frac{il}{E^2t^2}\int_0^t\Thf(u)\Ahat_i'(u)\Ahat_l'(u)\,du
\]
(\cref{thm:lea-loops}). The case $k_1 = 0$ of \eqref{eq:lea-anti} is the corner $e^{(k)}_{1,k+1} = \frac{(k+1)^2}{k+2}\frac{\epsilon_1\epsilon_2\alpha_1\alpha_2\wA\wB}{E^2}a^{(k)}_{k+1}$.

\begin{cor}[powers of $\rho$; (L2) for $j = 1$]\genusstatus{AI-proposed, unchecked}
\label{cor:lea-rhoe}
For every integer $e$ and every $k\geq 0$, $\coef{t^k}(\rho^e\Edge_{11})$ lies in $\operatorname{span}\{\lf_i(\psi)\lf_l(\psi') : i,l\geq 1,\ (i-1)+(l-1)\leq k\}$. In particular ((L2) for $j = 1$), for every $b\geq 0$,
\begin{equation}
\label{eq:lea-L2one}
\coef{t^k}(\rho f_{1,b}) = \sum_{i=1}^{k+1}\Bigl(\sum_{m=0}^{k}\coef{t^m}(\rho^{-1})\sum_le^{(k-m)}_{il}\frac{l^b}{E^{b+1}}\Bigr)\lf_i(\psi)\in\Lsp_{k+1},
\end{equation}
and $\coef{t^k}f_{1,b}$ lies in $\Lsp_{k+1}$ as well. Explicitly $\coef{t^0}(\rho f_{1,b}) = -\frac{\epsilon_1\epsilon_2\alpha_1\alpha_2\wA\wB}{2}E^{-b-1}\lf_1$ for every $b$ (rank one at $t = 0$).
\end{cor}

\begin{proof}
$\rho^{e-2}$ lies in $\Kal\fps{t}$ (\cref{prop:lea-key}) with $\psi$-independent coefficients, so
\[
\coef{t^k}(\rho^e\Edge_{11}) = \sum_m\coef{t^m}(\rho^{e-2})\,\coef{t^{k-m}}(\rho^2\Edge_{11}),
\]
and the supports at the orders $k-m\leq k$ are contained in the support at order $k$. For the leg, $f_{1,b}(\psi) = \coef{\psi'^b}\Edge_{11}(\psi,\psi')$ and $\rho f_{1,b} = \coef{\psi'^b}(\rho\Edge_{11})$ (the case $e = 1$); by (LAM1), $\coef{\psi'^b}\lf_l(\psi') = l^b/E^{b+1}$, and the $\lf_i(\psi)$ that occur have $i-1\leq k$. The same with $e = 0$ gives the statement for $f_{1,b}$. The value at $t = 0$ is $e^{(0)}_{11}/E^{b+1}$ with $e^{(0)}_{11} = -\epsilon_1\epsilon_2\alpha_1\alpha_2\wA\wB/2$ (\cref{thm:lea-L3}).
\end{proof}

In the regularity theorem of \cref{sec:reg} every half-edge at a singular vertex of genus $h\geq 2$ carries one factor $\rho$. So a loop, or an edge between two such vertices, carries $\rho^2\Edge_{11}$; a leg, i.e.\ an $\Edge_{11}$-edge to a vertex that is not rescaled (singular of genus $0$ or $1$), with external power $b$ at the far end, carries $\rho f_{1,b}$; an $\Edge_{11}$-edge between two non-rescaled vertices carries $\Edge_{11}$ (\cref{rmk:lea-use}).

\subsection{The regular colour}
\label{sec:lea-regular}

\begin{lem}[the involution exchanges the idempotents; the same scalar at both]\genusstatus{AI-proposed, unchecked}
\label{lem:lea-conj}
Let $\iota$ be the involution $U\mapsto U_1+U_2-U$ of the quadratic algebra $\QA = \bbK(\cQ)[U]/(\Pq(U))$ (\cref{def:set-QA}), and let $\ev_i : \QA\to\bbK(\Ccov)$, $U\mapsto U_i$, be the two characters. Then $\iota$ is a ring automorphism of $\QA$ with $\ev_1\circ\iota = \ev_2$ and $\ev_2\circ\iota = \ev_1$; it commutes with the derivation $d/d\cQ$ of $\QA$ defined by implicit differentiation ($dU/d\cQ = -\Pq_{\cQ}/\Pq_U$), and with multiplication by scalar series in $\cQ$ and $z$. Consequently every rational expression in $U$, $\cQ$, $d/d\cQ$ and the $t$-independent scalars $\hnw$, $\hw^{-1}$ that defines a factor of idempotent $1$ defines the corresponding factor of idempotent $2$ when $U_1$ is replaced by $U_2$ throughout: $\rpt_{0,2}(z;t) = \sum_mz^m\wick_m(\cW_2(U_2),\cW_3(U_2),\dots)$, $\rr_{0,2} = h_{(\alpha_1,\wA,-E)}\rpt_{0,2}$, $\rr_{1,2} = (\Lv_2 + z\Dop_2)\rr_{0,2}$ with $d/dt$ the total derivative along $U_2(t)$, and $\Delta_2$ is the quadratic expression $\epsilon_1\epsilon_2(\alpha_0\alpha_1\alpha_2 + 2\alpha_1\alpha_2\Lv - E\Lv^2)$ of \cref{prop:set-idem} evaluated at $\Lv = \Lv_2$.
\end{lem}

\begin{proof}
$\ev_1,\ev_2$ are ring homomorphisms, jointly injective ($\Pq$ has distinct roots), and $\ev_1\circ\iota = \ev_2$, $\ev_2\circ\iota = \ev_1$ because $U_1+U_2$ is the sum of the roots; this gives the compatibility with the ring operations. For the derivation: $\ev_i(dU/d\cQ) = dU_i/d\cQ$ by implicit differentiation, so $\ev_i\circ d/d\cQ = (d/d\cQ)\circ\ev_i$; for $\{i,i'\} = \{1,2\}$, $\ev_i(\tfrac{d}{d\cQ}\iota C) = \tfrac{d}{d\cQ}\ev_{i'}(C) = \ev_{i'}(\tfrac{d}{d\cQ}C) = \ev_i(\iota\tfrac{d}{d\cQ}C)$, and joint injectivity gives $\tfrac{d}{d\cQ}\iota = \iota\tfrac{d}{d\cQ}$. The scalars $\hnw$ and $\hw^{-1}$ are the same series at both idempotents (\cref{def:set-rows,lem:set-hmult}), so multiplication by them commutes with $\iota$. The formulae for the factors of idempotent $2$ follow by applying $\ev_2 = \ev_1\circ\iota$ to the defining expressions of idempotent $1$, and $\ev_2$ commutes with the total derivative $d/d\cQ = \cQ^{-1}d/dt$.
\end{proof}

This is the only place where ``the same scalar at both colours'' is needed: with a colour-dependent constant the two columns would not be exchanged by $\iota$.

\begin{lem}[analyticity at the regular idempotent]\genusstatus{AI-proposed, unchecked}
\label{lem:lea-reg}
\begin{enumerate}
\item For $k\geq 1$, $\cW_k$ is a polynomial in $p_1 := (1-U)^{-1}$ and $p_{\cQ} := (1-\cQ U)^{-1}$ with coefficients in $\bbZ[\alpha_1,\alpha_2,\wA]$ ($\cW_1 = \alpha_1(1-p_{\cQ}) - \alpha_2 + \wA(1-p_1)$, and $\cW_k = -\alpha_1\Pol_{k-1}(p_{\cQ}) - \wA\Pol_{k-1}(p_1)$ for $k\geq 2$, \cref{lem:row-wr}). $U_2(t)$ is analytic at $t = 0$ with $U_2(0) = -\alpha_2/\wB$, and $1-U_2(0) = 1-\cQ(0)U_2(0) = E/\wB\neq 0$; so every $\cW_k(U_2(t))$ is analytic at $t = 0$, and $\cW_2(U_2)$ is a unit of $\Oan$ with $\cW_2(U_2(0)) = \alpha_2\wB/E$.
\item (R) Every $\rpt_{0,2}[m]$, $\rr_{0,2}[m]$, $\rr_{1,2}[m]$ is analytic at $t = 0$; $\rr_{0,2}[0] = 1$ and $\rr_{1,2}[0] = \Lv_2$.
\item $\Delta_2 = \epsilon_1\epsilon_2\wB\delta\sast\Lv_2 = \alpha_1\wA\epsilon_1\epsilon_2\cW_2(U_2)$ is analytic at $t = 0$ with $\Delta_2(0) = \epsilon_1\epsilon_2\alpha_1\alpha_2\wA\wB/E$ and $\Delta_2'(0) = 0$ ((F1) for $\Delta_2$).
\item $\Dil_2(n) = (-1)^n\rr_{0,2}[n-1]$ and every $\Edge_{22}(m,n)$ are analytic at $t = 0$ ((F3) for $\Dil_2$, $\Edge_{22}$).
\end{enumerate}
\end{lem}

\begin{proof}
(i) $\thU\log(1-aU) = 1-p_a$ and $\thU p_a^n = n(p_a^{n+1}-p_a^n)$ for $a\in\{1,\cQ\}$, and $\thU\log U = 1$, which gives the polynomial form. At $\cQ = 1$ the critical polynomial is $\Pq(U) = \wB U^2 - (\alpha_0+\alpha_1-\alpha_2)U - \alpha_2 = \wB U^2 - (\wB-\alpha_2)U - \alpha_2 = (\wB U+\alpha_2)(U-1)$, with the simple roots $1$ and $-\alpha_2/\wB$ (distinct, since $-\alpha_2/\wB = 1$ if and only if $E = 0$); so $U_2$ is a simple root at $\cQ = 1$ and analytic in $t$ near $0$. At $\cQ = 1$ the superpotential is $\cW_0(U) := E\log(1-U) - \alpha_2\log U$, $\thU^2\cW_0 = -EU/(1-U)^2$, and at $U = -\alpha_2/\wB$: $-E(-\alpha_2/\wB)/(E/\wB)^2 = \alpha_2\wB/E$.
(ii) By \cref{lem:lea-conj}, $\rpt_{0,2}[m]$ is a polynomial in $\cW_2(U_2)^{-1},\cW_3(U_2),\dots,\cW_{2m+2}(U_2)$, hence analytic by (i); $\hnw$, $\hw^{-1}$ are $t$-independent; $\Lv_2 = \wB\sast - \alpha_0$ lies in $\Oan$ (\cref{lem:set-eig}) and $\Delta_2$ is a unit of $\Oan$ by (iii), so $\rr_{1,2}[m] = \Lv_2\rr_{0,2}[m] + \Dop_2\rr_{0,2}[m-1]$ is analytic.
(iii) \Cref{lem:set-Delta2} and the fourth form of \cref{prop:set-idem}; the value $\alpha_1\wA\epsilon_1\epsilon_2\cdot\alpha_2\wB/E$ agrees with \cref{lem:set-Delta2}.
(iv) By \cref{lem:set-edgegen}, $\Edge_{22}(m,n)$ is a polynomial with constant coefficients in the entries of $\vv_2[p]$, $p\leq m+n+1$.
\end{proof}

Why a separate argument is needed for Row 2 at the regular idempotent. Row 1 at $i = 2$ is the definition of the H-row. Row 2 at $i = 2$ cannot be obtained as in \cref{prop:lea-row2}, because the closed form concerns the singular row only; nor does the route of \cref{prop:row-gaussode} apply, which multiplies the formal expansion by $\exp(\dphi/z)(f_2/f_2(0))^{-1/2}$ and expands at $\tau = 0$: at $\stil\to\infty$, $e^{\phi_\tau(\stil)/z}$ is $\tau^{-E/z}$ times a unit, $d\phi_\tau(\stil)/d\tau$ has a simple pole and $\phi_\tau''(\stil)$ a double zero. Instead we work coefficientwise in $z$ at fixed small $t\neq 0$, with no expansion at $\tau = 0$; analyticity at $\tau = 0$ enters only through units of $\Oan$ and values at $t = 0$.

\begin{lem}[data at the regular critical point]\genusstatus{AI-proposed, unchecked}
\label{lem:lea-regdata}
For $\tau\in\Ddisc''$ (small $t\neq 0$):
\begin{enumerate}
\item $\stil = (E+\alpha_0\tau+\delta)/(2\wB\tau) = \wA/(\wB\tau\sast)$ is a non-degenerate critical point of $\phi_\tau$, not in $\{0,1,1/\tau\}$; $\tau\stil$ lies in $\Oan$ with value $E/\wB$; $2\wB\tau\stil - (E+\alpha_0\tau) = \delta$; $\wB\stil = \Lv_1+\alpha_0$; $U_2 = 1-\tau\stil$ is a unit of $\Oan$ with $U_2(0) = -\alpha_2/\wB$.
\item $f_2 := \phi_\tau''(\stil) = \delta/(\stil(1-\stil)U_2) = \tau^2\Delta_2/(\alpha_1\wA\epsilon_1\epsilon_2U_2^2) = \tau^2f_{2}^{\mathrm{r}}$, where $f_2^{\mathrm{r}} := \Delta_2/(\alpha_1\wA\epsilon_1\epsilon_2U_2^2)$ is a unit of $\Oan$ with $f_2^{\mathrm{r}}(0) = \Delta_2(0)/(\alpha_1\wA\epsilon_1\epsilon_2U_2(0)^2) = \wB^3/(\alpha_2E)$.
\item $f_0(\tau) := \phi_\tau(\stil(\tau))$ (any branch) has $\partial_\tau f_0 = \alpha_2\stil/(1-\tau\stil)\in\tau^{-1}\Oan$ with leading term $-E/\tau$, and
\begin{equation}
\label{eq:lea-f0}
f_0' = (1-\tau)\frac{\alpha_2\stil}{1-\tau\stil} = \wA - \wB\stil = \alpha_2 - \Lv_1 = \Lv_2 - \muq .
\end{equation}
The canonical coordinate $u_2 := \cW(U_2)$ satisfies $du_2/dt = \Lv_2$ and $u_2 = \alpha_1t + E\log\tau + f_0 + \mathrm{const}$.
\end{enumerate}
\end{lem}

\begin{proof}
(i) \Cref{def:set-alg,lem:set-N,lem:set-eig}: $\stil$ is the second root of $\Ns$, hence a critical point of $\phi_\tau$ (\cref{lem:row-hess}), non-degenerate because $\Ns'(\stil) = \delta\neq 0$ on $\Ddisc$, not in $\{0,1,1/\tau\}$ on $\Ddisc''$; $\sast\stil = \wA/(\wB\tau)$ gives $\stil = \wA/(\wB\tau\sast)$ and $\wB\stil = \wA/(\tau\sast) = \Lv_1+\alpha_0$; $\tau\stil = (E+\alpha_0\tau+\delta)/(2\wB)$ has value $E/\wB$ at $0$; $U_2(0) = 1 - E/\wB = -\alpha_2/\wB$, and $U_2$ has no zero on $\Ddisc''$ since $\stil\neq 1/\tau$ there.
(ii) The Hessian formula $\phi_\tau''(s_c) = \Ns'(s_c)/(s_c(1-s_c)(1-\tau s_c))$ of \cref{lem:row-hess} at $s_c = \stil$ with $\Ns'(\stil) = \delta$, and the tensor rule $\cW_2(U_2) = U_2^2\phi_\tau''(\stil)/\tau^2$ with $\Delta_2 = \alpha_1\wA\epsilon_1\epsilon_2\cW_2(U_2)$ (\cref{prop:set-idem}). $f_2^{\mathrm{r}}$ is a unit since $\Delta_2$ and $U_2$ are, with $f_2^{\mathrm{r}}(0) = (\alpha_1\alpha_2\wA\wB/E)\cdot\wB^2/(\alpha_1\wA\alpha_2^2) = \wB^3/(\alpha_2E)$.
(iii) Envelope formula: $\frac{d}{d\tau}\phi_\tau(\stil(\tau)) = (\partial_\tau\phi_\tau)(\stil) + \phi_\tau'(\stil)\,\partial_\tau\stil = (\partial_\tau\phi_\tau)(\stil) = \alpha_2\stil/(1-\tau\stil)$. Since $\tau\stil\to E/\wB$ and $U_2\to-\alpha_2/\wB$, $\partial_\tau f_0 = \tau^{-1}(\tau\stil)\alpha_2/U_2$ has leading term $\alpha_2(E/\wB)/(\tau\cdot(-\alpha_2/\wB)) = -E/\tau$. For \eqref{eq:lea-f0}: the identity $(1-\tau)\alpha_2s = (\wA-\wB s)(1-\tau s)$ of \cref{lem:set-N}(ii) holds at both roots, being equivalent to $\Ns(s) = 0$; so $f_0' = (1-\tau)\partial_\tau f_0 = \wA-\wB\stil$; then $\wB\stil = \Lv_1+\alpha_0$ gives $\wA-\wB\stil = \wA-\alpha_0-\Lv_1 = \alpha_2-\Lv_1$, and $\Lv_2-\muq = \Lv_2 - E/\tau + \wA = \alpha_2-\Lv_1$ by $\Lv_1+\Lv_2 = E/\tau-\alpha_0$. The statement about $u_2$ is \cref{lem:row-canon} at the second root, with the Euler coordinate $\cW(1-\tau s) = \phi_\tau(s) + E\log\tau + \alpha_1\log\cQ + \alpha_1\log(-1)$ of \cref{lem:row-euler}.
\end{proof}

Fix $\tau_0\in\Ddisc''$ and an open disc $\leaDisc$ around $\tau_0$ whose closure lies in $\Ddisc''$. Fix on $\leaDisc$ branches of $\log s$, $\log(1-s)$, $\log(1-\tau s)$ near $s = \stil(\tau)$ (possible since $\stil$, $1-\stil$, $1-\tau\stil$ do not vanish on the closure of $\leaDisc$); then $\phi_\tau(s)$ is analytic near every point $(\tau,\stil(\tau))$, $\tau\in\leaDisc$. $\Oan(\leaDisc)$ denotes the analytic functions on $\leaDisc$ and $\Oan(\leaDisc)((z))$ the Laurent series in $z$, bounded below, with coefficients in $\Oan(\leaDisc)$ ($d/dt = (1-\tau)d/d\tau$ coefficientwise). Put, with $f_0$ and $f_2$ of \cref{lem:lea-regdata},
\[
\leaDt := \frac{d}{dt} + \frac{f_0'}{z} - \frac{f_2'}{2f_2},\qquad \leaDtau := \frac{\leaDt}{1-\tau} = \frac{d}{d\tau} + \frac{\partial_\tau f_0}{z} - \frac{\partial_\tau f_2}{2f_2}.
\]

\begin{thm}[Row 2 of the quantum differential equation at the regular idempotent]\genusstatus{AI-proposed, unchecked}
\label{thm:lea-row2reg}
In $\Man\fps{z}$ (in $\Oan\fps{z}$ after multiplication by $\tau$),
\begin{equation}
\label{eq:lea-row2reg}
(\Lv_2 + z\Dop_2)\rr_{1,2} = -\frac{\alpha_1\alpha_2}{\tau}\rr_{0,2} + \Bigl(\frac{E}{\tau}-\alpha_0\Bigr)\rr_{1,2},\qquad\text{equivalently}\qquad z\Dop_2\rr_{1,2} = \Lv_1(\rr_{1,2} - \Lv_2\rr_{0,2}),
\end{equation}
i.e.\ $(\Dop_2\rr_{1,2})[m-1] = \Lv_1(\rr_{1,2}[m] - \Lv_2\rr_{0,2}[m])$ for $m\geq 1$. With Row 1 this is $z\Dop_2\vv_2 = (\Aop-\Lv_2)\vv_2$, for the frame-B rows, the rows $\rZ$ and the point-frame rows alike.
\end{thm}

\begin{proof}
Step A (the regular row as a normalised Laplace expansion in the Euler coordinate). For every $\tau\in\leaDisc$,
\begin{equation}
\label{eq:lea-stepA}
\rpt_{0,2}(z;t) = U_2\,\leaLt(\tau),\qquad \leaLt(\tau) := \Lnorm[\phi_\tau,(1-\tau s)^{-1};\stil(\tau)]\in\Oan(\leaDisc)\fps{z}.
\end{equation}
This is the argument of \cref{prop:row-wickrow} with $(U_1,\sast)$ replaced by $(U_2,\stil)$. Let $y_2$ be a branch of $\log U_2$ and $\cW_y(y) := \cW(e^y)$; $\cW_y$ is analytic near $y_2$ with $\cW_y^{(k)}(y_2) = \cW_k(U_2)$, $\cW_y'(y_2) = 0$ and $\cW_y''(y_2) = \cW_2(U_2) = \delta U_2/(\tau^2\stil(1-\stil))\neq 0$ (\cref{lem:row-hess}); in this coordinate the integrand is $e^{\cW_y/z}dy$ with amplitude $1$, so by the Wick form (D1) of \cref{def:row-Lnorm} and \cref{lem:lea-conj}, $\rpt_{0,2} = \Lnorm[\cW_y,1;y_2] = \Lhat[\cW_y,1;y_2]$. The map $h(s) := \log(1-\tau s)$, with the branch for which $h(\stil) = y_2$, is analytic near $\stil$ with $h'(\stil) = -\tau/U_2\neq 0$. By \cref{lem:row-euler}, $\cW_y(h(s)) = \phi_\tau(s) + C(t)$ near $\stil$ with $C(t) = E\log\tau + \alpha_1t + \alpha_1\log(-1)$ independent of $s$, and the amplitude $1$ becomes $(1\circ h)h' = -\tau(1-\tau s)^{-1}$. By the coordinate invariance (P1)(iii) of \cref{prop:row-P1}, (D3) and linearity (the factor $-\tau$), $\Lhat[\cW_y,1;y_2] = \Lhat[\phi_\tau,(1-\tau s)^{-1};\stil] = U_2\leaLt$, since the amplitude takes the value $1/U_2$ at $\stil$. By (D2) the $z$-coefficients of $\leaLt$ are polynomials in $1/f_2$, in the $\phi_\tau^{(k)}(\stil)$ of \cref{lem:row-wr} and in $\partial_s^a(1-\tau s)^{-1}|_{\stil} = a!\,\tau^aU_2^{-a-1}$, so they lie in $\Oan(\leaDisc)$.

Step B (conjugation). As operators on $\Oan(\leaDisc)((z))$,
\begin{equation}
\label{eq:lea-stepB}
(\Lv_2 + z\Dop_2)\circ U_2 = U_2\circ(z\leaDt + \omq).
\end{equation}
Indeed $(\Lv_2+z\Dop_2)(U_2Y) = U_2\bigl[zY' + \Lv_2Y + z(U_2'/U_2 - \Delta_2'/(2\Delta_2))Y\bigr]$. By \cref{lem:lea-regdata}(ii), $\Delta_2 = \alpha_1\wA\epsilon_1\epsilon_2U_2^2f_2/\tau^2$, so $\Delta_2'/(2\Delta_2) = U_2'/U_2 + f_2'/(2f_2) - \tau'/\tau$ with $\tau'/\tau = (1-\tau)/\tau$; and $\Lv_2 = f_0' + \muq$ by \eqref{eq:lea-f0}. So the bracket equals $zY' + f_0'Y - z(f_2'/(2f_2))Y + (\muq + z(1-\tau)/\tau)Y = (z\leaDt+\omq)Y$, since $\muq + z(1-\tau)/\tau = \omq$. Only the logarithmic derivative of $\Delta_2$ enters, so no square root is involved. Step B is where the identity $\Delta_2 = \alpha_1\wA\epsilon_1\epsilon_2\cW_2(U_2)$ (the fourth form of \cref{prop:set-idem} at the second root) is used: if $\Delta_2$ were $\alpha_1\wA\epsilon_1\epsilon_2\cW_2(U_2)$ times a non-constant $p(t)$, the right side of \eqref{eq:lea-stepB} would contain the extra term $-(z/2)(p'/p)U_2Y$, and Step D would not give Row 2.

Step C (Gauss's equation for $\leaLt$). With $\ga = 1+\alpha_2/z$, $\gb = 1+\wA/z$, $\gc = 2+E/z$,
\begin{equation}
\label{eq:lea-stepC}
\Gauss_{\stil}[\leaLt] := \tau(1-\tau)\leaDtau^2\leaLt + (\gc - (\ga+\gb+1)\tau)\leaDtau\leaLt - \ga\gb\,\leaLt = 0\quad\text{in }\Oan(\leaDisc)((z)).
\end{equation}
The parameter derivative with moving expansion point, (P2) of \cref{prop:row-P2}, applies to $f(\tau,s) = \phi_\tau(s)$, to the analytic non-degenerate critical point $\stil(\tau)$ and to every amplitude $q(\tau,s)$ that is a Laurent polynomial in $z$ with coefficients analytic near $(\tau_1,\stil(\tau_1))$, $\tau_1\in\leaDisc$: $\leaDtau\Lnorm[\phi_\tau,q;\stil] = \Lnorm[\phi_\tau,\Dam q;\stil]$ with $\Dam q := \partial_\tau q + q\,(\partial_\tau\phi_\tau)/z$, and with the $f_0$ and $f_2$ of $\stil$. Apply it to $q = \amp_0 := (1-\tau s)^{-1}$ and to $q = \Dam\amp_0$, and use (D3) (the coefficients $\tau(1-\tau)$ etc.\ do not depend on $s$):
\[
\Gauss_{\stil}[\leaLt] = \Lnorm\bigl[\tau(1-\tau)\Dam^2\amp_0 + (\gc-(\ga+\gb+1)\tau)\Dam\amp_0 - \ga\gb\,\amp_0;\ \stil\bigr] = -\ga\,\Lnorm\bigl[k_E' + \phi_\tau'k_E/z;\ \stil\bigr] = 0,
\]
by the amplitude identity of \cref{lem:row-amp} ($k_E := s(1-s)(1-\tau s)^{-2}$), which is an identity of Laurent polynomials in $z$ with coefficients rational in $(s,\tau)$, valid identically and not only at a critical point, hence also near $\stil$; and by (P3) of \cref{prop:row-P3}, the normalised expansion of a total derivative vanishes. The amplitudes $\amp_0$, $\Dam\amp_0$, $\Dam^2\amp_0$ and $k_E$ are analytic near $\stil$ because $\stil\notin\{0,1,1/\tau\}$; that $\stil\to\infty$ as $\tau\to 0$ plays no role, since $\leaDisc$ does not contain $0$.

Step D (transfer by the operator identity). $z\leaDt = \Dz + a_0$ with the function $a_0 := f_0' - zf_2'/(2f_2)$ satisfies the Leibniz rule $z\leaDt(aY) = za'Y + a\,z\leaDt Y$, and $(z(1-\tau))^{-1}z\leaDt = \leaDtau$. So \cref{lem:lea-op} with $\Dd' = z\leaDt$ gives
\begin{equation}
\label{eq:lea-stepD}
(z\leaDt+\omq)^2 - \Bigl(\frac{E}{\tau}-\alpha_0\Bigr)(z\leaDt+\omq) + \frac{\alpha_1\alpha_2}{\tau} = \frac{(1-\tau)z^2}{\tau}\Gauss_{\stil}\quad\text{on }\Oan(\leaDisc)((z)).
\end{equation}
By \eqref{eq:lea-stepA}, \eqref{eq:lea-stepB} applied twice, \eqref{eq:lea-stepD} and \eqref{eq:lea-stepC},
\begin{align*}
\Bigl[(\Lv_2+z\Dop_2)^2 - \Bigl(\frac{E}{\tau}-\alpha_0\Bigr)(\Lv_2+z\Dop_2) + \frac{\alpha_1\alpha_2}{\tau}\Bigr]\rpt_{0,2}
&= U_2\Bigl[(z\leaDt+\omq)^2 - \Bigl(\frac{E}{\tau}-\alpha_0\Bigr)(z\leaDt+\omq) + \frac{\alpha_1\alpha_2}{\tau}\Bigr]\leaLt\\
&= U_2\,\frac{(1-\tau)z^2}{\tau}\,\Gauss_{\stil}[\leaLt] = 0 .
\end{align*}
The $t$-independent scalars $\hnw$ and $\hw^{-1}$ commute with everything, so the same identity holds for the rows $\rZ_{0,2}$ and $\rr_{0,2}$; with $\rr_{1,2} = (\Lv_2+z\Dop_2)\rr_{0,2}$ it is \eqref{eq:lea-row2reg} on $\leaDisc$.

Conclusion. Every $z$-coefficient of \eqref{eq:lea-row2reg} is an identity of analytic functions on $\leaDisc$; $\tau_0$ was arbitrary, so it holds on $\Ddisc''$. All its terms are meromorphic at $t = 0$ (\cref{lem:lea-reg}; $\alpha_1\alpha_2/\tau$ and $E/\tau$ have simple poles), so it holds in $\Man$. The second form uses $E/\tau-\alpha_0-\Lv_2 = \Lv_1$ and $\alpha_1\alpha_2/\tau = \Lv_1\Lv_2$. The components of $(\Aop-\Lv_2)\vv_2$ are $\rr_{1,2}-\Lv_2\rr_{0,2} = z\Dop_2\rr_{0,2}$ (Row 1) and $-(\alpha_1\alpha_2/\tau)\rr_{0,2} + (E/\tau-\alpha_0-\Lv_2)\rr_{1,2} = z\Dop_2\rr_{1,2}$ (Row 2).
\end{proof}

Why the normalised forms suffice: the point-frame row is itself normalised, so (P1)(iii) identifies it with $U_2\leaLt$ exactly, with no convention; the $t$-dependent factors that the unnormalised form of (P1) carries ($e^{C/z}$, the Jacobian, and $(-f_2)^{-1/2}$ in place of $\Delta_2^{-1/2}$) are accounted for in \eqref{eq:lea-stepB}, through $\muq$, $U_2$, $\tau'/\tau$ and \cref{lem:lea-regdata}(ii). An alternative proof of Row 2 at both idempotents, by the same calculation in the coordinate $U$ with Gauss's operator replaced by the Picard--Fuchs operator and the amplitude identity by Euler's identity, is given at indeterminate weights in \cref{prop:ide-rows}(ii); that statement is pointwise in $\cQ$ and in the weights and covers every weight point with (ND), $\alpha_0 = 0$ included.

\begin{cor}[orthogonality]\genusstatus{AI-proposed, unchecked}
\label{cor:lea-orth}
$\langle\vv_1(w),\vv_2(-w)\rangle = 0$ in $\Man\fps{w}$, to all orders in $w$.
\end{cor}

\begin{proof}
$Y_{12}(z) := \langle\vv_1(z),\vv_2(-z)\rangle$ lies in $\Man\fps{z}$. Use $z\Dop_1\vv_1 = (\Aop-\Lv_1)\vv_1$ (\cref{prop:lea-row2}), \eqref{eq:lea-row2reg} with $z$ replaced by $-z$ ($-z\Dop_2\vv_2(-z) = (\Aop-\Lv_2)\vv_2(-z)$; $\Aop$, $\Lv_2$, $\Dop_2$ do not involve $z$), $\vv_i' = \Dop_i\vv_i + (\Delta_i'/(2\Delta_i))\vv_i$, the $t$-independence of $\eta^{-1}$ and self-adjointness (\cref{prop:lea-A}(iii)):
\begin{align*}
zY_{12}' &= \langle(\Aop-\Lv_1)\vv_1(z),\vv_2(-z)\rangle - \langle\vv_1(z),(\Aop-\Lv_2)\vv_2(-z)\rangle + \frac{z}{2}\Bigl(\frac{\Delta_1'}{\Delta_1}+\frac{\Delta_2'}{\Delta_2}\Bigr)Y_{12}\\
&= \Bigl[(\Lv_2-\Lv_1) + \frac{z}{2}\Bigl(\frac{\Delta_1'}{\Delta_1}+\frac{\Delta_2'}{\Delta_2}\Bigr)\Bigr]Y_{12}.
\end{align*}
Write $Y_{12} = \sum_mY_mz^m$ with $Y_m\in\Man$. The $z^0$ coefficient gives $(\Lv_2-\Lv_1)Y_0 = 0$, and for $m\geq 1$ the $z^m$ coefficient gives $(\Lv_2-\Lv_1)Y_m = Y_{m-1}' - \frac12(\Delta_1'/\Delta_1 + \Delta_2'/\Delta_2)Y_{m-1}$. Since $\Lv_2-\Lv_1 = -\delta/\tau$ is a unit of $\Man$, $Y_0 = 0$ and inductively $Y_m = 0$.
\end{proof}

It is essential that $Y_{12}$ is a power series in $z$, bounded below: the formal symbol $e^{-(u_1-u_2)/z}$, an infinite series in $1/z$ (and not an element of $\Rz$, since $u_1-u_2$ contains $E\log\tau$), would formally give a non-zero flat section. $Y_0 = \langle\chiv_1,\chiv_2\rangle = 0$ is also direct. Row 2 at the regular idempotent is genuinely needed: Row 1 only determines $\rr_{1,2}$ from $\rr_{0,2}$, and the $z^m$ coefficient of $Y_{12}$ contains $\rr_{0,2}[m]$ only through $\langle\chiv_1,\chiv_2\rangle\rr_{0,2}[m] = 0$, so $Y_{12} = 0$ is a genuine constraint on the regular row, which the closed form of the singular row cannot supply. The symplectic relations $\langle\vv_i(w),\vv_j(-w)\rangle = \delta_{ij}\Delta_i$ are also proved at indeterminate weights in \cref{prop:ide-rows}(iv), the diagonal ones by evaluation at the semisimple point $\cQ = 0$, which needs $\alpha_0\neq 0$; the proofs given here (\cref{prop:lea-K1}, the present corollary and \cref{cor:lea-K2}) use the limit $t\to 0$ at $\cQ = 1$ instead and cover $\alpha_0 = 0$ (\cref{cor:lea-a0}).

\begin{prop}[the regular rows at $t = 0$]\genusstatus{AI-proposed, unchecked}
\label{prop:lea-regrows}
\begin{enumerate}
\item[(a)] (Analytic input.) $\rpt_{0,2}(z;0) = h_{(\alpha_2,-E,\wB)}(z)$; equivalently the frame-B value is
\[
\rr_{0,2}(z;0) = h_{(\alpha_1,\wA,-E)}(z)\,h_{(\alpha_2,-E,\wB)}(z) = h_{(\alpha_1,\alpha_2,\wA,\wB,-E,-E)}(z) =: \hreg(z),\qquad \hreg(z)\hreg(-z) = 1 .
\]
\item[(b)] (Uses \cref{thm:lea-row2reg}.) $\alpha_1\alpha_2\,\rr_{0,2}(z;0) = E\,\rr_{1,2}(z;0)$, i.e.
\[
\vv_2(z;0) = \rr_{0,2}(z;0)\,(1,\alpha_1\alpha_2/E) = \hreg(z)\,(1,\alpha_1\alpha_2/E);
\]
equivalently, $\coef{t^1}\rr_{0,2}(z;t) = 0$.
\item[(c)] (Uses (a), (b).)
\[
\Edge_{22}(z,w)\big|_{t=0} = \frac{\epsilon_1\epsilon_2\alpha_1\alpha_2\wA\wB}{E}\,\frac{1 - \hreg(z)^{-1}\hreg(w)^{-1}}{z+w},
\]
and, by (a) alone, $\Dil_2(n)|_{t=0} = (-1)^n\coef{z^{n-1}}\hreg(z)$.
\end{enumerate}
\end{prop}

\begin{proof}
(a) By \cref{lem:lea-reg}, each $\rpt_{0,2}[m]$ is analytic at $t = 0$ and its value there is $\wick_m$ evaluated on the $t = 0$ data $(\thU^k\cW_0)(-\alpha_2/\wB)$, $\cW_0(U) = E\log(1-U) - \alpha_2\log U$ the superpotential at $\cQ = 1$. By the Wick form (D1), $\rpt_{0,2}(z;0) = \Lnorm[\cW_{y,0},1;y_c]$ with $\cW_{y,0}(y) := \cW_0(e^y)$ and $y_c$ a branch of $\log(-\alpha_2/\wB)$: the normalised expansion, in the coordinate $y = \log U$, of the phase $a\log U + b\log(1-U)$ with $(a,b) := (-\alpha_2,E)$ at its critical point $U = a/(a+b) = -\alpha_2/\wB$. This is exactly the situation of \cref{lem:row-betaphase}, whose hypothesis $ab(a+b)\neq 0$ reads $\alpha_2E\wB\neq 0$ here; the lemma gives $\rpt_{0,2}(z;0) = h_{(-a,-b,a+b)} = h_{(\alpha_2,-E,\wB)}$. Frame B multiplies by $h_{(\alpha_1,\wA,-E)}$ (\cref{lem:set-hmult}), and the union of the multisets $(\alpha_1,\wA,-E)$ and $(\alpha_2,-E,\wB)$ is $(\alpha_1,\alpha_2,\wA,\wB,-E,-E)$; $\hreg(z)\hreg(-z) = 1$ because the exponent of $h_v$ is odd in $z$.

(b) First proof, from \cref{thm:lea-row2reg} alone. In $(\Dop_2\rr_{1,2})[m-1] = \Lv_1(\rr_{1,2}[m] - \Lv_2\rr_{0,2}[m])$ the left side is analytic at $t = 0$ (\cref{lem:lea-reg}), while $\Lv_1 = E/t + O(1)$ has a simple pole; so $\rr_{1,2}[m](0) = \Lv_2(0)\rr_{0,2}[m](0) = (\alpha_1\alpha_2/E)\rr_{0,2}[m](0)$ for $m\geq 1$, and for $m = 0$ by Row 1 ($\rr_{1,2}[0] = \Lv_2\rr_{0,2}[0]$). Second proof, from \cref{cor:lea-orth}: put $w = -z$ and multiply by $\rho$: $\langle\brow(z),\vv_2(z)\rangle = 0$; its $t^{-1}$ coefficient is $\langle(0,E),\vv_2(z;0)\rangle = \epsilon_1\epsilon_2E\bigl(\alpha_1\alpha_2\rr_{0,2}(z;0) - E\rr_{1,2}(z;0)\bigr)$, by $\langle(0,E),b\rangle = \epsilon_1\epsilon_2[\alpha_1\alpha_2Eb_0 - E^2b_1]$. Equivalent forms: by Row 1, $\rr_{1,2}[m](0) - (\alpha_1\alpha_2/E)\rr_{0,2}[m](0) = (\Dop_2\rr_{0,2}[m-1])(0) = (\rr_{0,2}[m-1])'(0)$, because $\Lv_2(0) = \alpha_1\alpha_2/E$ and $\Delta_2'(0) = 0$ (\cref{lem:set-Delta2}); so (b) is equivalent to $\coef{t^1}\rr_{0,2} = 0$.

(c) By \cref{lem:set-edgegen} and (a), (b),
\[
\Pij_{22}(z,w)\big|_{t=0} = -\langle\vv_2(-z;0),\vv_2(-w;0)\rangle = -\hreg(-z)\hreg(-w)\,\langle(1,\alpha_1\alpha_2/E),(1,\alpha_1\alpha_2/E)\rangle,
\]
and
\begin{align*}
\Bigl\langle\Bigl(1,\frac{\alpha_1\alpha_2}{E}\Bigr),\Bigl(1,\frac{\alpha_1\alpha_2}{E}\Bigr)\Bigr\rangle &= \epsilon_1\epsilon_2\Bigl[\alpha_0\alpha_1\alpha_2 + 2\alpha_1\alpha_2\frac{\alpha_1\alpha_2}{E} - E\frac{(\alpha_1\alpha_2)^2}{E^2}\Bigr]\\
&= \epsilon_1\epsilon_2\alpha_1\alpha_2\frac{\alpha_0E+\alpha_1\alpha_2}{E} = \frac{\epsilon_1\epsilon_2\alpha_1\alpha_2\wA\wB}{E} = \Delta_2(0);
\end{align*}
$\Pij_{22}(-w,w)|_{t=0} = -\hreg(w)\hreg(-w)\Delta_2(0) = -\Delta_2(0)$. Hence
\[
\Edge_{22}|_{t=0} = \Delta_2(0)\,\frac{1 - \hreg(z)^{-1}\hreg(w)^{-1}}{z+w},\qquad \Dil_2(n)|_{t=0} = (-1)^n\rr_{0,2}[n-1](0) = (-1)^n\coef{z^{n-1}}\hreg(z).
\]
\end{proof}

\begin{cor}[$\Ktwo = 1$]\genusstatus{AI-proposed, unchecked}
\label{cor:lea-K2}
$\Ktwo(z) := \langle\vv_2(z),\vv_2(-z)\rangle/\Delta_2$ lies in $\Oan\fps{z}$, is $t$-independent, and
\[
\Ktwo = \rr_{0,2}(z;0)\,\rr_{0,2}(-z;0) = \hreg(z)\hreg(-z) = 1 .
\] With \cref{prop:lea-K1}(iii) and \cref{cor:lea-orth}: $\langle\vv_i(w),\vv_j(-w)\rangle = \delta_{ij}\Delta_i$ for all $i,j$, to all orders in $w$. (Not used for \cref{thm:lea-main}.)
\end{cor}

\begin{proof}
Put $\nu_2 := 1/(\sqrt{\alpha_1\wA\epsilon_1\epsilon_2}\,U_2\,(f_2^{\mathrm{r}})^{1/2})$, a unit of $\Oan$ with $\nu_2^2 = \tau^2/(\alpha_1\wA\epsilon_1\epsilon_2U_2^2f_2) = 1/\Delta_2$ and $\nu_2'/\nu_2 = -\Delta_2'/(2\Delta_2)$ (any fixed square roots). Then $\leaVn_2 := \nu_2\vv_2$ lies in $\Oan\fps{z}^2$ and satisfies $z\leaVn_2' = (\Aop-\Lv_2)\leaVn_2$ by \cref{thm:lea-row2reg}. As in the proof of \cref{prop:lea-K1}(i), $\Ktwo = \langle\leaVn_2(z),\leaVn_2(-z)\rangle$ has zero $t$-derivative (self-adjointness), so $\Ktwo = \langle\vv_2(z;0),\vv_2(-z;0)\rangle/\Delta_2(0) = \rr_{0,2}(z;0)\rr_{0,2}(-z;0)\Delta_2(0)/\Delta_2(0)$ by \cref{prop:lea-regrows}(b) and the computation in its part (c), and $\hreg(z)\hreg(-z) = 1$.
\end{proof}

$\Ktwo = 1 + O(z^2)$ is automatic ($\Ktwo$ is even in $z$ and $\Ktwo(0) = \langle\chiv_2,\chiv_2\rangle/\Delta_2 = 1$); the vanishing of the higher coefficients is the content of the analytic input (a). $\Edge_{22}$ and $\Dil_2$ are regular whatever $\Ktwo$ is.

\begin{rmk}[$\Edge_{22}$ at $t = 0$ is not of rank one; the regular column is determined by the singular data]
\label{rmk:lea-E22}
(1) Write $\hreg(z)^{-1} = \exp\bigl(\sum_{k\geq 1}s_kz^{2k-1}\bigr)$ with $s_k := \bk_k\sum_{v\in(\alpha_1,\alpha_2,\wA,\wB,-E,-E)}v^{1-2k}$. Expanding \cref{prop:lea-regrows}(c),
\[
\frac{1-\hreg(z)^{-1}\hreg(w)^{-1}}{z+w} = -s_1 - \frac{s_1^2}{2}(z+w) - s_2(z^2-zw+w^2) - \frac{s_1^3}{6}(z+w)^2 + (\text{degree}\geq 3),
\]
so at $t = 0$: $\Edge_{22}(0,0) = -\Delta_2(0)s_1$, $\Edge_{22}(1,0) = \Edge_{22}(0,1) = -\Delta_2(0)s_1^2/2$, $\Edge_{22}(1,1) = \Delta_2(0)(s_2 - s_1^3/3)$, and the minor $\Edge_{22}(0,0)\Edge_{22}(1,1) - \Edge_{22}(0,1)\Edge_{22}(1,0) = \Delta_2(0)^2(s_1^4/12 - s_1s_2)$ is a rational function of the weights which is not identically zero (its restriction to the fibrewise torus $\alpha_0 = 0$, $\alpha_i = \epsilon_i$ is a non-zero rational function of $\alpha_1/\alpha_2$). So $\Edge_{22}|_{t=0}$ is not a product $\varphi(z)\varphi'(w)$, in contrast with $\Edge_{11}|_{t=0}$ and $\Edge_{12}|_{t=0}$ (\cref{thm:lea-L3,rmk:lea-legs}). Only the regularity of $\Edge_{22}$ is used in \cref{sec:reg}. (2) The regular column is determined recursively by the singular data. \Cref{cor:lea-orth} with $w = -z$, multiplied by $\rho$, reads $\langle\brow(z),\vv_2(z)\rangle = 0$; at order $t^k$ ($k\geq -1$), with $\vv_{2,n} := \coef{t^n}\vv_2(z)$, $\langle(0,E),b\rangle = \epsilon_1\epsilon_2E(\alpha_1\alpha_2b_0 - Eb_1)$, the H-row in the form $\alpha_1\alpha_2\rr_{0,2} - E\rr_{1,2} = (\alpha_1\alpha_2 - E\Lv_2)\rr_{0,2} - Ez\Dop_2\rr_{0,2}$, and $\alpha_1\alpha_2 - E\Lv_2 = -\wB\tau\sast\Lv_2 = O(t)$, this gives
\begin{multline*}
E(k+2)\,z\coef{t^{k+2}}\rr_{0,2} = \coef{t^{k+1}}\bigl((\alpha_1\alpha_2 - E\Lv_2)\rr_{0,2}\bigr) + Ez\coef{t^{k+1}}\Bigl(\rr_{0,2}\frac{\Delta_2'}{2\Delta_2}\Bigr)\\
+ \frac{1}{\epsilon_1\epsilon_2E}\sum_{k_1=0}^{k}\langle\brow^{(k_1)},\vv_{2,k-k_1}\rangle,
\end{multline*}
whose right side involves $\rr_{0,2}$ only up to order $t^{k+1}$. So, given $\rr_{0,2}(z;0) = \hreg$, the whole regular column, and with it every $\Edge_{12}$ leg and $\Edge_{22}$ to all orders, is determined recursively by the closed-form singular data and by $\Lv_2$, $\Delta_2$; e.g.\ $\coef{t^1}\rr_{0,2} = 0$. (3) In the language of \cref{rmk:lea-flat}, Steps A--D of \cref{thm:lea-row2reg} say that $\nu_2\rpt_{0,2}$ times the symbol $e^{u_2/z}$ is the second flat section of $z\,d/dt - \Aop$; its indicial exponent at $\tau = 0$ is $0$, while $(f_2)^{-1/2}\leaLt$ has $\tau$-order exactly $-1$ (the $z^{-2}$ coefficient of \eqref{eq:lea-stepC} is $\leaLt[0]$ times the eikonal polynomial of \cref{lem:row-eik} evaluated at $\partial_\tau f_0 = \stil/U_2$, the second root $-E/\tau + O(1)$ of the eikonal equation). Neither (2) nor (3) is used below.
\end{rmk}

\subsection{The legs: (L2) for \texorpdfstring{$j = 2$}{j = 2}, (F3) for \texorpdfstring{$\Edge_{12}$}{E12} and \texorpdfstring{$\Edge_{21}$}{E21}}
\label{sec:lea-legs}

\begin{lem}[extraction of a leg]\genusstatus{AI-proposed, unchecked}
\label{lem:lea-leg}
Let $\Edge\in\bbC\fps{\psi,w}$ and $\mathrm{Num} := (\psi+w)\Edge$. For every $b\geq 0$, in $\bbC((\psi))$,
\begin{equation}
\label{eq:lea-legext}
\coef{w^b}\Edge = \sum_{n=0}^{b}(-1)^n\psi^{-n-1}\coef{w^{b-n}}\mathrm{Num}.
\end{equation}
\end{lem}

\begin{proof}
Put $\Edge_m := \coef{w^m}\Edge\in\bbC\fps{\psi}$ and $\Edge_{-1} := 0$. Then $\coef{w^m}\mathrm{Num} = \psi\Edge_m + \Edge_{m-1}$, so
\[
\sum_{n=0}^{b}(-1)^n\psi^{-n-1}(\psi\Edge_{b-n} + \Edge_{b-n-1}) = \sum_{n=0}^{b}(-1)^n\psi^{-n}\Edge_{b-n} - \sum_{n=1}^{b+1}(-1)^n\psi^{-n}\Edge_{b-n} = \Edge_b .
\]
\end{proof}

Equivalently, \eqref{eq:lea-legext} is the expansion of $\mathrm{Num}/(\psi+w)$ in $\bbC((\psi))\fps{w}$, where $\psi+w$ is invertible with $1/(\psi+w) = \sum_n(-1)^nw^n\psi^{-n-1}$; $\bbC\fps{\psi,w}$ is contained in $\bbC((\psi))\fps{w}$. Being linear, \eqref{eq:lea-legext} holds coefficientwise in $t$.

\begin{prop}[(L2) for $j = 2$ and the regularity of $\Edge_{12}$]\genusstatus{AI-proposed, unchecked}
\label{prop:lea-L2}
Assume (R) of \cref{lem:lea-reg}(ii) (every $\vv_2[m]$ analytic at $t = 0$). Then:
\begin{enumerate}
\item $\rho\Edge_{12}$ and $\Edge_{12}$ are $O(t^0)$: every $\Edge_{12}(a,b)$ is regular at $t = 0$ ((F3) for $\Edge_{12}$);
\item for all $k\geq 0$ and $b\geq 0$, $\coef{t^k}(\rho f_{2,b})$ lies in $\Lsp_{k+1}$.
\end{enumerate}
\end{prop}

\begin{proof}
Put $\leaNtw(\psi,w) := \rho\Pij_{12}(\psi,w) = -\langle\brow(\psi),\vv_2(-w)\rangle$. Since $\rho$ depends on $t$ only, \cref{lem:set-edgegen} gives
\begin{equation}
\label{eq:lea-N12}
(\psi+w)\,\rho\Edge_{12}(\psi,w) = \leaNtw(\psi,w) - \leaNtw(-w,w).
\end{equation}
By \cref{prop:lea-sigma} and (R), $\leaNtw$ is a Laurent series in $t$ starting at $t^{-1}$, and each coefficient is a finite sum:
\begin{equation}
\label{eq:lea-N12coef}
\coef{t^k}\leaNtw(\psi,w) = -\sum_{k_1=-1}^{k}\langle\brow^{(k_1)}(\psi),\vv_{2,k-k_1}(-w)\rangle,\qquad \vv_{2,k_2}(-w) := \coef{t^{k_2}}\vv_2(-w),\ \text{a pair in }\bbC\fps{w}.
\end{equation}
(i) $\coef{t^{-1}}\leaNtw(\psi,w) = -\langle(0,E),\vv_{2,0}(-w)\rangle$ does not depend on $\psi$, so the $t^{-1}$ coefficient of the right side of \eqref{eq:lea-N12} is zero; $\bbC\fps{\psi,w}$ has no zero divisors, so $\coef{t^{-1}}(\rho\Edge_{12}) = 0$, and lower orders do not occur. So $\rho\Edge_{12} = O(t^0)$, and $\Edge_{12} = \rho^{-1}(\rho\Edge_{12}) = O(t^0)$.

(ii) Fix $k,b\geq 0$. Apply \eqref{eq:lea-legext} to $\Edge = \coef{t^k}(\rho\Edge_{12})$ and $\mathrm{Num} = \coef{t^k}(\leaNtw(\psi,w) - \leaNtw(-w,w))$:
\begin{equation}
\label{eq:lea-legrat}
\begin{gathered}
\coef{t^k}(\rho f_{2,b})(\psi) = \sum_{n=0}^{b}(-1)^n\psi^{-n-1}\,Y_{b-n}(\psi),\\
Y_m(\psi) := \coef{w^m}\coef{t^k}\bigl(\leaNtw(\psi,w) - \leaNtw(-w,w)\bigr) = -(-1)^m\sum_{k_1=-1}^{k}\langle\brow^{(k_1)}(\psi),\vv_{2,k-k_1}[m]\rangle - y_m,
\end{gathered}
\end{equation}
where $\vv_{2,k_2}[m] := \coef{t^{k_2}}\vv_2[m]\in\bbC^2$ (as $\coef{w^m}\vv_2(-w) = (-1)^m\vv_2[m]$) and $y_m := \coef{w^m}\coef{t^k}\leaNtw(-w,w)$ is a constant. By \cref{prop:lea-sigma}, both components of every $\brow^{(k_1)}$, $-1\leq k_1\leq k$, lie in $\Lsp_{k_1+1}^{+}$ (constants plus $\Lsp_{k_1+1}$), contained in $\Lsp_{k+1}^{+}$; hence $Y_m$ is the Taylor series of a rational function $\widehat{Y}_m\in\Lsp_{k+1}^{+}$. Consider $\mathrm{Rat}(\psi) := \sum_{n=0}^{b}(-1)^n\psi^{-n-1}\widehat{Y}_{b-n}(\psi)$.
\begin{itemize}
\item[(a)] Away from $\psi = 0$, $\mathrm{Rat}$ has at most simple poles, located at the points $E/j$ with $1\leq j\leq k+1$.
\item[(b)] $\mathrm{Rat}(\psi)\to 0$ as $\psi\to\infty$, because each $\widehat{Y}_m$ is bounded at $\infty$ and $\psi^{-n-1}\to 0$.
\item[(c)] The Laurent expansion of $\mathrm{Rat}$ at $\psi = 0$ is the right side of \eqref{eq:lea-legrat}, i.e.\ $\coef{t^k}(\rho f_{2,b})$, which lies in $\bbC\fps{\psi}$ because $\Edge_{12}(\psi,w)$ is a power series in $\psi$ (\cref{lem:set-edgegen}). Hence the principal part of $\mathrm{Rat}$ at $0$ vanishes: the apparent pole of order $\leq b+1$ cancels.
\end{itemize}
So $\mathrm{Rat} - \sum_{j=1}^{k+1}\mathrm{Rat}_j\lf_j$, with $\mathrm{Rat}_j := \lim_{\psi\to E/j}(E-j\psi)\mathrm{Rat}(\psi)$, has no poles in $\bbC$ and tends to $0$ at $\infty$, hence is $0$: $\coef{t^k}(\rho f_{2,b}) = \mathrm{Rat}$ lies in $\Lsp_{k+1}$.
\end{proof}

Explicitly, since $\psi^{-n-1}\lf_j = (j/E)^{n+1}\lf_j + $ (a polynomial in $1/\psi$ without constant term), $\mathrm{Rat}_j = \sum_{n=0}^{b}(-1)^n(j/E)^{n+1}[\lf_j]\widehat{Y}_{b-n}$. What is used about the regular rows: only (R). The rows $\vv_2$ enter \eqref{eq:lea-legrat} through the constants $\vv_{2,k_2}[m]$ and $y_m$; their $z$-dependence, their values, symplecticity and the orthogonality play no role. Analyticity is exactly what the degree bound needs: if some $\vv_2[m]$ had a pole of order $p$ at $t = 0$, \eqref{eq:lea-N12coef} would bring $\brow^{(k+p)}$ into the order $t^k$, i.e.\ $\lf_j$ with $j$ up to $k+p+1$. The $t = 0$ value of $\Edge_{12}$ is \cref{rmk:lea-legs}(b).

\begin{prop}[leaf form of the $\Edge_{12}$ edge]\genusstatus{AI-proposed, unchecked}
\label{prop:lea-E12}
Put $S_j(t) := \sum_{k\geq j-1}t^kS_j^{(k)}$ with the coefficient vectors \eqref{eq:lea-Sj}; $S_j = O(t^{j-1})$. Then
\begin{align}
\label{eq:lea-sigleaf}
\brow(\psi) &= \Bigl(0,\frac{E}{t}\Bigr) + \sum_{j\geq 1}S_j(t)\lf_j(\psi),\\
\label{eq:lea-E12leaf}
\rho\Edge_{12}(\psi,w) &= -\sum_{j\geq 1}j\,\langle S_j(t),\vv_2(-w;t)\rangle\,\lf_j(\psi)\lf_j(-w),\qquad \lf_j(-w) = \frac{1}{E+jw},
\end{align}
identities in $\bbC\fps{\psi,w}((t))$ in which only $j\leq k+1$ contribute at order $t^k$ (by (R)). Consequently, for all $k$ and $b$,
\begin{equation}
\label{eq:lea-E12coef}
\coef{t^k}(\rho f_{2,b})(\psi) = -\sum_{j=1}^{k+1}j\,\lf_j(\psi)\sum_{k_1=j-1}^{k}\Bigl\langle S_j^{(k_1)},\ \coef{w^b}\frac{\vv_{2,k-k_1}(-w)}{E+jw}\Bigr\rangle .
\end{equation}
\end{prop}

\begin{proof}
\eqref{eq:lea-sigleaf} is \cref{prop:lea-sigma}. Insert it in $\leaNtw$. The term $(0,E/t)$ contributes the same function of $w$ to $\leaNtw(\psi,w)$ and to $\leaNtw(-w,w)$, so by (LAM6) (with scale $E$: $\lf_j(\psi) - \lf_j(-w) = j(\psi+w)\lf_j(\psi)\lf_j(-w)$)
\begin{align*}
\leaNtw(\psi,w) - \leaNtw(-w,w) &= -\sum_j\langle S_j,\vv_2(-w)\rangle\bigl(\lf_j(\psi) - \lf_j(-w)\bigr)\\
&= -(\psi+w)\sum_jj\,\langle S_j,\vv_2(-w)\rangle\lf_j(\psi)\lf_j(-w).
\end{align*}
Divide by $\psi+w$ (\eqref{eq:lea-N12}); \eqref{eq:lea-E12coef} is the coefficient of $t^kw^b$.
\end{proof}

In \eqref{eq:lea-E12leaf} the cancellation at $\psi = 0$ of step (c) in the proof of \cref{prop:lea-L2} is explicit, and the subtracted term $\leaNtw(-w,w)$, which the orthogonality would set to zero, drops out without being evaluated. \Cref{prop:lea-L2}(ii) is the statement ``$j\leq k+1$'' in \eqref{eq:lea-E12coef}.

\begin{prop}[the transposition $\Edge_{21} = \Edge_{12}^{T}$; regularity of $\Edge_{21}$]\genusstatus{AI-proposed, unchecked}
\label{prop:lea-transp}
Assume \cref{cor:lea-orth}. Then $\Edge_{21}(z,w) = \Edge_{12}(w,z)$, i.e.\ $\Edge_{21}(b,a) = \Edge_{12}(a,b)$ for all $a,b$. Hence $\Edge_{21}$ is regular at $t = 0$, and an edge joining a singular vertex to a regular one gives the singular vertex the leg $f_{2,b}$ ($b$ the power at the regular end) for either orientation of the edge.
\end{prop}

\begin{proof}
$\varphi_{12}(u) := \Pij_{12}(-u,u) = -\langle\vv_1(u),\vv_2(-u)\rangle = 0$, so $\Edge_{12}(z,w) = \Pij_{12}(z,w)/(z+w)$. By the symmetry of $\langle\,,\rangle$, $\Pij_{21}(z,w) = -\langle\vv_2(-z),\vv_1(-w)\rangle = \Pij_{12}(w,z)$ and $\Pij_{21}(-w,w) = \Pij_{12}(w,-w) = \varphi_{12}(-w) = 0$. Hence $\Edge_{21}(z,w) = \Pij_{12}(w,z)/(z+w) = \Edge_{12}(w,z)$.
\end{proof}

What the orthogonality is needed for. Without it, $\Edge_{21}(z,w) - \Edge_{12}(w,z) = [\varphi_{12}(z) - \varphi_{12}(-w)]/(z+w)$, and by \eqref{eq:lea-legext} the singular-side leg of an $\Edge_{21}$ edge would be
\[
\coef{z^b}\Edge_{21}(z,\psi) = \rho^{-1}\Bigl[\sum_{n=0}^{b}(-1)^n\psi^{-n-1}\coef{z^{b-n}}\leaNtw(\psi,z) - (-1)^b\psi^{-b-1}\leaNtw(\psi,-\psi)\Bigr],
\]
which contains $\leaNtw(\psi,-\psi) = -\langle\brow(\psi),\vv_2(\psi)\rangle$, a pairing of $\brow$ with the full power series $\vv_2(\psi)$ (at $t = 0$, $\hreg(\psi)(1,\alpha_1\alpha_2/E)$, a divergent Stirling series); the rationality argument of \cref{prop:lea-L2} does not apply to it. With the orthogonality the leg is $f_{2,b}$ and is rational. Even the regularity of $\Edge_{21}$ needs the $t^{-1}$ part of the orthogonality: since $\rho\Pij_{21}(z,w) = -\langle\vv_2(-z),\brow(w)\rangle$ and $\brow^{(-1)} = (0,E)$,
\begin{gather*}
\coef{t^{-1}}(\rho\Edge_{21})(z,w) = \frac{f_e(z) - f_e(-w)}{z+w},\\
f_e(z) := -\langle\vv_2(-z;0),(0,E)\rangle = -\epsilon_1\epsilon_2E\bigl[\alpha_1\alpha_2\,\rr_{0,2}(-z;0) - E\,\rr_{1,2}(-z;0)\bigr],
\end{gather*}
and no lower powers of $t$ occur (\cref{lem:lea-reg}). This vanishes if and only if $f_e$ is constant; the constant is $f_e(0) = -\epsilon_1\epsilon_2E(\alpha_1\alpha_2 - E\Lv_2(0)) = 0$ ($\rr_{0,2}[0] = 1$, $\rr_{1,2}[0] = \Lv_2$, $\Lv_2(0) = \alpha_1\alpha_2/E$). So $\Edge_{21}$ is regular if and only if $\alpha_1\alpha_2\rr_{0,2}(z;0) = E\rr_{1,2}(z;0)$, which is \cref{prop:lea-regrows}(b) (from \cref{thm:lea-row2reg} alone).

\begin{cor}[all legs; unrescaled ends]\genusstatus{AI-proposed, unchecked}
\label{cor:lea-alllegs}
For all $k,b\geq 0$ and $j\in\{1,2\}$: $\coef{t^k}f_{j,b}$ and $\coef{t^k}(\rho f_{j,b})$ lie in $\Lsp_{k+1}$; $\coef{t^k}(\rho^e\Edge_{11}(\psi,\psi'))$ lies in $\operatorname{span}\{\lf_i(\psi)\lf_l(\psi') : i,l\geq 1,\ (i-1)+(l-1)\leq k\}$ for every integer $e$; all these series are $O(t^0)$. In particular $\Edge_{1j}(z,w)\to 0$ as $z\to\infty$ at fixed order in $w$, coefficientwise in $t$ (every element of $\Lsp_{k+1}$ vanishes at $\infty$).
\end{cor}

\begin{proof}
$j = 2$: \cref{prop:lea-L2} and multiplication by $\rho^{-1}\in\Kal\fps{t}$ (by (T): $\coef{t^k}(uX) = \sum_iu_i\coef{t^{k-i}}X$, and the spaces $\Lsp_{k+1}$ grow with $k$). $j = 1$ and the loops: \cref{cor:lea-rhoe}.
\end{proof}

\begin{rmk}[how the factors attach to a singular vertex]
\label{rmk:lea-use}
In the regularity theorem of \cref{sec:reg}, at a singular vertex $v$ of genus $h_v\geq 2$ the dilaton homogeneity rescales every half-edge at $v$ exactly once (one factor $\rho$ per half-edge at $v$). The factors attached to $v$ depend on the other end of the edge:
\begin{itemize}
\item the other end is a regular vertex: the edge factor $\Edge_{12}$ (or $\Edge_{21}$) gives $v$ the leg $\rho f_{2,b}$, $b$ the power at the regular end (\cref{prop:lea-L2}; either orientation by \cref{prop:lea-transp});
\item the other end is a singular vertex $v'$ of genus $0$ or $1$ (not rescaled): the edge factor is $\rho\Edge_{11} = \rho^{-1}(\rho^2\Edge_{11})$, and for each power $b'$ at $v'$ it gives $v$ the leg $\rho f_{1,b'}$ (\cref{cor:lea-rhoe}), of leaf degree $\leq k$ at order $t^k$; either orientation, since $\Edge_{11}$ is symmetric;
\item the other end is a singular vertex $v'$ of genus $\geq 2$ (also rescaled): the edge factor is $\rho^2\Edge_{11}$, which splits its leaf degree $\leq k$ at order $t^k$ between the two ends (\cref{thm:lea-L3});
\item a self-loop at $v$: the factor is $\rho^2\Edge_{11}$ (\cref{thm:lea-L3}).
\end{itemize}
At a vertex $v'$ of genus $0$ or $1$ itself only regularity is needed: $\Delta_1^{h-1}$ is $O(t^2)$ for $h = 0$ and equals $1$ for $h = 1$, and every factor attached to $v'$ is $O(t^0)$ ((F3), \cref{sec:lea-assembly}).
\end{rmk}

\subsection{Proof of Theorems A and \ref{thm:lea-qde}; (F1)--(F3) in one table}
\label{sec:lea-assembly}

\begin{proof}[Proof of \cref{thm:lea-main}]
(F1): $t^2\rho^2\Delta_1 = -\epsilon_1\epsilon_2E^3$, $\rho = \tau/(t\alg)$ with $\rho(0) = 1$, and the expansion: \cref{prop:lea-key}(ii), (iv). $\Delta_2 = \epsilon_1\epsilon_2\wB\delta\sast\Lv_2$ analytic with $\Delta_2(0) = \epsilon_1\epsilon_2\alpha_1\alpha_2\wA\wB/E$ and $\Delta_2'(0) = 0$: \cref{lem:set-Delta2,lem:lea-reg}(iii).
(F2): $\rr_{0,1}(z;0) = E/(E+z)$: \cref{cor:row-F2}. $\rho\,\rr_{1,1} = E/t + \alpha_1\alpha_2/(E+z) + O(t)$: \cref{prop:lea-sigma}(i), (ii) ($\brow^{(-1)} = (0,E)$, $\brow^{(0)}_1 = \alpha_1\alpha_2\lf_1$) at $\psi = -z$.
(F3): the table \eqref{eq:lea-table} below.
(L1): \cref{thm:lea-L1}.
(L2): $j = 1$ by \cref{cor:lea-rhoe}, $j = 2$ by \cref{prop:lea-L2}(ii); the unrescaled legs $\coef{t^k}f_{j,b}$ by \cref{cor:lea-alllegs}; $\Edge_{21}(b,a) = \Edge_{12}(a,b)$ by \cref{prop:lea-transp}.
(L3): \cref{thm:lea-L3}; the symmetry by \eqref{eq:lea-Ebar}; all powers $\rho^e$ by \cref{cor:lea-rhoe}.
Rationality of the coefficients: the passage (P) of \cref{lem:set-rings}. Sharpness: \cref{cor:lea-top,cor:lea-anti}. Large-$z$ form: \cref{lem:lea-tcoef} ($\lim z\,\rr_{0,1} = E/\rho$) and \cref{cor:lea-alllegs}.
\end{proof}

\begin{proof}[Proof of \cref{thm:lea-qde}]
(i) \cref{prop:lea-row2} and \cref{thm:lea-row2reg} (or \cref{prop:ide-rows}(ii)). (ii) \cref{prop:lea-K1}(iii), \cref{cor:lea-orth}, \cref{cor:lea-K2}. (iii) \cref{prop:lea-regrows}(a), (b).
\end{proof}

(F3), factor by factor (frame B; values at $t = 0$; $\kappa_2 := \epsilon_1\epsilon_2\alpha_1\alpha_2\wA\wB$, $\hreg = h_{(\alpha_1,\alpha_2,\wA,\wB,-E,-E)}$, $\chi_{12}(w) := -\kappa_2\hreg(w)^{-1}/(E+w)$):
\begingroup\small
\begin{equation}
\label{eq:lea-table}
\begin{array}{@{}lll@{}}
\text{factor} & \text{regularity at }t = 0 & \text{value at }t = 0\\[2pt]\hline\\[-8pt]
\Dil_1 & \text{\cref{thm:lea-L1}, remark (iv)} & \qo(\psi) = -E\psi/(E-\psi);\ \Dil_1(d;0) = -E^{1-d}\\
\Dil_2 & \text{\cref{lem:lea-reg}(iv)} & (-1)^n\coef{z^{n-1}}\hreg(z)\quad(\text{\cref{prop:lea-regrows}(a)})\\
\Edge_{11} & \text{\eqref{eq:lea-int}: }t^2\rho^2\Edge_{11} = O(t^2) & -\tfrac{\kappa_2}{2}\lf_1(\psi)\lf_1(\psi')\quad(\text{\cref{thm:lea-L3}})\\
\Edge_{12} & \text{\cref{prop:lea-L2}(i)} & \lf_1(\psi)\,\chi_{12}(w)\quad(\text{\cref{rmk:lea-legs}(b)})\\
\Edge_{21} & \text{\cref{prop:lea-regrows}(b) (}t^{-1}\text{ part); \cref{prop:lea-transp}} & \Edge_{21}(b,a) = \Edge_{12}(a,b)\\
\Edge_{22} & \text{\cref{lem:lea-reg}(iv)} & \tfrac{\kappa_2}{E}\,\dfrac{1-\hreg(z)^{-1}\hreg(w)^{-1}}{z+w}\quad(\text{\cref{prop:lea-regrows}(c)})
\end{array}
\end{equation}
\endgroup
By \eqref{eq:lea-sigleaf} the $t^{-1}$ part of $\brow$ is the $\psi$-independent covector $(0,E)$; hence the coefficients $(-1)^n\rr_{1,1}[n]$ of the H-row are regular for $n\geq 1$, while $\rr_{1,1}[0] = \Lv_1$ has a simple pole; they do not enter the graph sum without marked points.

\begin{rmk}[where each input is used; where (ND) is used]
\label{rmk:lea-inputs}
The analytic input of \cref{sec:row}, the Beta-phase lemma \cref{lem:row-betaphase} (through the constant $\Abeta$ of \cref{prop:row-beta} and the constant line of \cref{lem:set-hmult}, which also uses \cref{lem:ide-norm} for the definition of the rows), enters through \cref{thm:row-main} into the key formula \eqref{eq:lea-key} and from there into everything that uses $E/(E+z)$: (F2), (L1) at $t^0$ ($\qo$), $\brow^{(-1)} = (0,E)$, $\Kone = 1$, $\coef{t^0}\Ebar = 0$, (F3) for $\Edge_{11}$, the rank-one $t^0$ loop. The same lemma, with $(a,b) = (-\alpha_2,E)$, gives the value $\hreg$ of the regular row at $t = 0$ (\cref{prop:lea-regrows}(a)), used only for the $t = 0$ values at the regular idempotent and for $\Ktwo = 1$, not for \cref{thm:lea-main}. The identity $\Delta_i = \alpha_1\wA\epsilon_1\epsilon_2\cW_2(U_i)$ enters at $U_1$ in \cref{prop:lea-key} ($\rho = \tau/(t\alg)$) and at $U_2$ in \eqref{eq:lea-stepB}. The formal Laplace calculus ((P1)--(P3), homogeneity) enters \cref{prop:row-wickrow,prop:row-gaussode} and \cref{thm:lea-row2reg}. The condition (ND) is used as follows: $E\neq 0$ (the leaf scale; the recursion of \cref{lem:row-unique}; the distinct roots $1$ and $-\alpha_2/\wB$ at $\cQ = 1$); $\alpha_1\wA\neq 0$ (the non-degenerate Euler critical point $\phi_0''(\wA/E) = -E^3/(\alpha_1\wA)$; the Beta-phase lemma at $(\wA,\alpha_1)$); $\alpha_2\neq 0$ ($\sast\neq 1/\tau$; $U_2(0)\neq 0$; the Beta-phase lemma at $(-\alpha_2,E)$); $\wB\neq 0$ ($U_2(0)$ finite, $\Delta_2(0)\neq 0$); $\epsilon_1\epsilon_2\neq 0$ ($\eta$ invertible, $\Delta_i\neq 0$). Nothing uses $E = \epsilon_1+\epsilon_2$ (\cref{rmk:lea-cy}), and nothing uses the point $\cQ = 0$ except through the citations of \cref{lem:ide-norm} (the definition of the constant $\hnw$) and of \cref{prop:ide-rows}(iv) where it is offered as an alternative (\cref{cor:lea-a0}).
\end{rmk}

% s04b_closedforms. Section file (second part of s04, input from s04_leaves.tex); label prefix: lea.
% No new macros: those of s04_leaves.tex are in force (\TFsup, \leaVn, ...).

\subsection{Closed forms}
\label{sec:lea-closed}

Nothing in \cref{thm:lea-main} depends on this subsection except the sharpness statements and the explicit values quoted as examples. Recall $\Psh_k = \prod_{i=1}^{k}((k+1)\alpha_2 - iE)((k+1)\wA - iE)$, $\Cj_j = \Psh_{j-1}/(j!\,(j-1)!\,E^{2j-4})$, $\zq = \tau e^{\dphi/E}$, and the leaf coefficients $a_j^{(k)}$ of \eqref{eq:lea-pf}, with $\Aleaf_j(t) := \sum_ka_j^{(k)}t^k$ and $\Ahat_j := t\Aleaf_j$.

\begin{thm}[the dilaton leaf coefficients]\genusstatus{AI-proposed, unchecked}
\label{thm:lea-Aj}
For $j\geq 1$,
\[
\Cj_j = \frac{\Psh_{j-1}}{j!\,(j-1)!\,E^{2j-4}} = \frac{E^2\,(1-j\alpha_2/E)_{j-1}(1-j\wA/E)_{j-1}}{j!\,(j-1)!}\quad(\Cj_1 = E^2),\qquad \zq(t) = \tau e^{\dphi(\tau)/E} = t + O(t^2),
\]
\begin{equation}
\label{eq:lea-Aj}
\Aleaf_j(t) := (-1)^j\,\Cj_j\,t^{-1}\zq(t)^j\,{}_2F_1\Bigl(\frac{j\wB}{E},\frac{j\alpha_1}{E};\,j;\,\tau\Bigr),\qquad \Aleaf_j = O(t^{j-1}).
\end{equation}
Then $a_j^{(k)} = \coef{t^k}\Aleaf_j(t)$ for all $j\geq 1$, $k\geq 0$. Equivalently, as formal series in $t$ with finitely many $j$ at each order,
\[
z\Gk(z;t) = E + \sum_{j\geq 1}\frac{\Aleaf_j(t)}{E+jz},\qquad \rho\qt(\psi) = E + \sum_{j\geq 1}\Aleaf_j(t)\lf_j(\psi),\qquad \gl(\psi) = -\frac1E\sum_{j\geq 1}j\Aleaf_j(t)\lf_j(\psi).
\]
\end{thm}

\begin{proof}
Fix $j\geq 1$ and $k\geq 0$. By \cref{prop:lea-pf}, $\coef{t^k}(z\Gk)$ has at $-E/j$ at most a simple pole, contributing $a_j^{(k)}/(E+jz)$, so $a_j^{(k)} = j\Res_{z=-E/j}\coef{t^k}(z\Gk)$. Write $z\Gk = (\tau/t)\sum_n\tau^n(Ez/(E+z))\Ephi\Fhyp_n$; $\coef{t^k}(z\Gk)$ is a finite sum over $n\leq k$ of terms $(\text{scalar})\cdot(Ez/(E+z))\,z^{-p}\,\Fhyp_n(z)$ (each $\coef{\tau^m}\Ephi$ is a polynomial in $1/z$), and evaluation at $z = -E/j$ commutes with taking $t$-coefficients.

$j = 1$: $\Fhyp_n$ is regular at $z = -E$ (its poles are at $0$ and at $-E/l$, $l\geq 2$), so $\Res_{z=-E}\coef{t^k}(z\Gk) = \coef{t^k}\bigl(-E^2(\tau/t)e^{\dphi/E}\Fhyp(-E,\tau)\bigr)$, since $\Res_{z=-E}Ez/(E+z) = -E^2$ and $\Ephi(-E) = e^{\dphi/E}$; at $z = -E$ the parameters are $\ga = 1-\alpha_2/E = \wB/E$, $\gb = 1-\wA/E = \alpha_1/E$, $\gc = 2-1 = 1$, so $\Fhyp(-E,\tau) = {}_2F_1(\wB/E,\alpha_1/E;1;\tau)$. Hence $\Aleaf_1 = -E^2t^{-1}\zq\,{}_2F_1(\wB/E,\alpha_1/E;1;\tau) = (-1)^1\Cj_1t^{-1}\zq\,{}_2F_1$.

$j\geq 2$: $Ez/(E+z)$ and $z^{-p}$ are regular at $-E/j$; $\Fhyp_n$ is regular there for $n < j-1$ and has a simple pole for $n\geq j-1$, from the single factor $E+jz$ of its denominator. At $z = -E/j$: $\alpha_2+lz = (j\alpha_2-lE)/j$, $\wA+lz = (j\wA-lE)/j$, $z^n = (-E/j)^n$, $\prod_{l=2,\,l\neq j}^{n+1}(E+lz) = (E/j)^{n-1}\prod_{l=2,\,l\neq j}^{n+1}(j-l) = (E/j)^{n-1}(j-2)!\,(-1)^{n+1-j}(n+1-j)!$, and $\Res_{z=-E/j}1/(E+jz) = 1/j$; collecting the powers of $j$ ($j^{-2n}$ from the numerator, $j^{-(2n-1)}\cdot j^{-1}$ from the denominator) and the signs,
\[
\Res_{z=-E/j}\Fhyp_n = (-1)^{j-1}\,\frac{E\,(1-j\alpha_2/E)_n\,(1-j\wA/E)_n}{j^2\,n!\,(j-2)!\,(n+1-j)!},
\]
using $\prod_{l=1}^{n}(j\alpha_2-lE) = (-E)^n(1-j\alpha_2/E)_n$ and the same with $\wA$. Reindex $n = j-1+m$: $(1-j\alpha_2/E)_{j-1+m} = (1-j\alpha_2/E)_{j-1}(j\wB/E)_m$, since $1-j\alpha_2/E+j-1 = j(E-\alpha_2)/E = j\wB/E$; likewise $(1-j\wA/E)_{j-1+m} = (1-j\wA/E)_{j-1}(j\alpha_1/E)_m$; and $(j-1+m)! = (j-1)!\,(j)_m$. So
\[
\sum_{n\geq j-1}\tau^n\Res_{z=-E/j}\Fhyp_n = (-1)^{j-1}\,\frac{E\,(1-j\alpha_2/E)_{j-1}(1-j\wA/E)_{j-1}\,\tau^{j-1}}{j^2\,(j-2)!\,(j-1)!}\;{}_2F_1\Bigl(\frac{j\wB}{E},\frac{j\alpha_1}{E};j;\tau\Bigr).
\]
The remaining factors at $z = -E/j$ are $Ez/(E+z) = -E/(j-1)$ and $\Ephi = e^{j\dphi/E}$, together with $\tau/t$. Multiplying by $j$ and using $j(j-1)(j-2)!\,(j-1)! = j!\,(j-1)!$,
\[
a_j^{(k)} = \coef{t^k}\Bigl\{(-1)^j\,\frac{E^2(1-j\alpha_2/E)_{j-1}(1-j\wA/E)_{j-1}}{j!\,(j-1)!}\,t^{-1}\zq^j\,{}_2F_1\Bigl(\frac{j\wB}{E},\frac{j\alpha_1}{E};j;\tau\Bigr)\Bigr\},
\]
and $E^2(1-j\alpha_2/E)_{j-1}(1-j\wA/E)_{j-1} = \Psh_{j-1}/E^{2j-4}$, because $\prod_{i=1}^{j-1}(j\alpha_2-iE) = (-E)^{j-1}(1-j\alpha_2/E)_{j-1}$ and likewise for $\wA$. The three generating functions follow from \cref{prop:lea-pf}, \cref{thm:lea-L1} and $g_j^{(k)} = -(j/E)a_j^{(k)}$; $\Aleaf_j = O(t^{j-1})$ since $\zq^j = O(t^j)$.
\end{proof}

\begin{cor}[sum rule, $\Thf$, the H-component, dilaton consistency]\genusstatus{AI-proposed, unchecked}
\label{cor:lea-sumrule}
\begin{enumerate}
\item[(a)] Sum rule: $\sum_{j\geq 1}\Aleaf_j(t) = -E^2$, i.e.\ $\sum_j(-1)^j\Cj_j\zq^j\,{}_2F_1(j\wB/E,j\alpha_1/E;j;\tau) = -E^2t = E^2\log(1-\tau)$ (no pole of $\Gk$ at $z = 0$).
\item[(b)] $t\,z\Gk(z) = -E\log(1-\tau) + \sum_j\Ahat_j(\tau)/(E+jz)$, with $\Ahat_j = t\Aleaf_j = (-1)^j\Cj_j\zq^j\,{}_2F_1(j\wB/E,j\alpha_1/E;j;\tau)$ a function of $\tau$.
\item[(c)] $\Thf = \epsilon_1\epsilon_2(E\dphi' - \alpha_2\wA)$ \eqref{eq:lea-Theta}, $\Thf_1 = -\epsilon_1\epsilon_2\alpha_1\alpha_2\wA\wB/E^2$, $\Thf_k = (k+1)\epsilon_1\epsilon_2E\coef{t^{k+1}}\dphi$ for $k\geq 1$.
\item[(d)] The H-component: $\brow = (0,E/t) + \sum_jS_j(t)\lf_j$ with
\[
S_j(t) = \Bigl(-\frac{j}{E}\Aleaf_j,\ \Aleaf_j' - \frac{j-1}{t}\Aleaf_j - \frac{j}{E}\lrem\Aleaf_j\Bigr) = t^{j-1}\Bigl(-\frac{j}{E}\widetilde{\Aleaf}_j,\ \widetilde{\Aleaf}_j' - \frac{j}{E}\lrem\widetilde{\Aleaf}_j\Bigr),
\]
where $\widetilde{\Aleaf}_j := t^{1-j}\Aleaf_j = (-1)^j\Cj_j(\zq/t)^j\,{}_2F_1(j\wB/E,j\alpha_1/E;j;\tau)$ is analytic with $\widetilde{\Aleaf}_j(0) = (-1)^j\Cj_j$; the coefficient of $t^k$ is \eqref{eq:lea-Sj}, and $S_1(0) = (E,\alpha_1\alpha_2) = E(1,\alpha_1\alpha_2/E)$.
\item[(e)] Dilaton consistency: $\sum_{j\geq 1}j\Aleaf_j(t) = -E^2\rho(t)$, i.e.\ $\sum_{j\leq k+1}j\,a_j^{(k)} = -E^2\coef{t^k}\rho$ for every $k$ (not used).
\end{enumerate}
\end{cor}

\begin{proof}
(a) is the sum rule of \cref{prop:lea-pf} summed over $k$, with $a_1^{(0)} = -E^2$; multiply by $t = -\log(1-\tau)$. (b) is \cref{thm:lea-Aj} multiplied by $t$. (c) is the expansion after \cref{thm:lea-flow}. (d) By \cref{thm:lea-Aj}, $\gl = -E^{-1}\sum_jj\Aleaf_j\lf_j$, and by (LAM2), $\psi\lf_j = (E\lf_j-1)/j$, so $-\psi\gl = \sum_j\Aleaf_j\lf_j - E^{-1}\sum_j\Aleaf_j$ and $-\psi\gl' = \sum_j\Aleaf_j'\lf_j - E^{-1}\sum_j\Aleaf_j'$. With $\Lv_1 = E/t + \lrem$ \eqref{eq:lea-L1exp}, $\Lv_1\gl = -\sum_j(j/t)\Aleaf_j\lf_j - (\lrem/E)\sum_jj\Aleaf_j\lf_j$. Hence by \eqref{eq:lea-sigma1}
\[
\brow_1 = \Lv_1\gl - \frac{\psi}{t}\gl - \psi\gl' = \sum_j\Bigl[\Aleaf_j' - \frac{j-1}{t}\Aleaf_j - \frac{j}{E}\lrem\Aleaf_j\Bigr]\lf_j - \frac{1}{Et}\sum_j\bigl(\Aleaf_j + t\Aleaf_j'\bigr),
\]
and the constants add up to $-\frac{1}{Et}\frac{d}{dt}\sum_jt\Aleaf_j = -\frac{1}{Et}\frac{d}{dt}(-E^2t) = \frac{E}{t}$ by (a). With $\brow_0 = \gl$ this is the formula for $S_j$; the second form follows from $\Aleaf_j' - \frac{j-1}{t}\Aleaf_j = t^{j-1}\widetilde{\Aleaf}_j'$. $S_1(0) = (-\Aleaf_1(0)/E,\ \Aleaf_1'(0) - \lrem_0\Aleaf_1(0)/E) = (E,\ a_1^{(1)} + E\lrem_0) = (E,\ \uq/2 + E^2/2 - \alpha_0E - \alpha_1\alpha_2) = (E,\alpha_1\alpha_2)$. (e) $\Dil_1(1) = -\rr_{0,1}[0] = -1$ at every $t$ (\cref{def:set-factors}), so $\coef{\psi^1}(\rho\qt) = -\rho$; with $\coef{\psi^1}\lf_j = j/E^2$ (LAM1) this reads $\sum_jj\Aleaf_j/E^2 = -\rho$.
\end{proof}

Not given here: closed forms of $\rho$, $\alg$ and $e^{\dphi}$ free of branch ambiguity; generically $e^{\dphi} = (E\sast/\wA)^{\wA}\,(E(1-\sast)/\alpha_1)^{\alpha_1}\,U_1^{-\alpha_2}$ with $1-\sast = \alpha_1U_1/(E-\wB\tau\sast)$ (\cref{lem:set-N}(ii)), which is all that is needed to define $\zq$.

\begin{cor}[top and subleading coefficients; sharpness]\genusstatus{AI-proposed, unchecked}
\label{cor:lea-top}
For every $k\geq 0$, resp.\ $k\geq 1$,
\[
\begin{gathered}
a_{k+1}^{(k)} = (-1)^{k+1}\Cj_{k+1} = (-1)^{k+1}\frac{\Psh_k}{(k+1)!\,k!\,E^{2k-2}},\qquad
g_{k+1}^{(k)} = (-1)^k\frac{\Psh_k}{(k!)^2E^{2k-1}},\\
a_k^{(k)} = (-1)^{k+1}\,k\,\Cj_k\,\frac{\uq}{2E^2}.
\end{gathered}
\]
Consequently $\Cj_{k+1}\neq 0$ if and only if $\Psh_k\neq 0$ if and only if $E\notin\{(k+1)\alpha_2/i,\,(k+1)\wA/i : 1\leq i\leq k\}$: the top leaf $\lf_{k+1}$ occurs in $\coef{t^k}(\rho\qt)$ exactly then, and the bound $j\leq k+1$ of (L1) is sharp at generic weights.
\end{cor}

\begin{proof}
$\Aleaf_j = (-1)^j\Cj_jt^{j-1}(\zq/t)^j\,{}_2F_1(j\wB/E,j\alpha_1/E;j;\tau)$ with $(\zq/t)(0) = 1$ and ${}_2F_1 = 1 + O(\tau)$: the first claim is $j = k+1$ at order $t^k$; $g_{k+1}^{(k)} = -\frac{k+1}{E}a_{k+1}^{(k)}$ and $\frac{k+1}{(k+1)!} = \frac{1}{k!}$. For the third, with $j = k$ at order $t^k$ one needs $\coef{t^1}\bigl((\zq/t)^k\,{}_2F_1(k\wB/E,k\alpha_1/E;k;\tau)\bigr)$; now $\zq/t = (\tau/t)e^{\dphi/E} = 1 + (\dphi_1/E - \tfrac12)t + O(t^2)$ with $\dphi_1 = \alpha_2\wA/E$, and ${}_2F_1 = 1 + \frac{(k\wB/E)(k\alpha_1/E)}{k}\tau + O(\tau^2) = 1 + \frac{k\alpha_1\wB}{E^2}t + O(t^2)$, so
\[
\coef{t^1}\bigl((\zq/t)^k\,{}_2F_1\bigr) = k\Bigl(\frac{\alpha_2\wA}{E^2} - \frac12\Bigr) + \frac{k\alpha_1\wB}{E^2} = k\,\frac{2\alpha_2\wA + 2\alpha_1\wB - E^2}{2E^2} = -\frac{k\uq}{2E^2},
\]
since $2\alpha_2\wA + 2\alpha_1\wB - E^2 = 2E(\alpha_1+\alpha_2) - 4\alpha_1\alpha_2 - E^2 = -\uq$.
\end{proof}

Direct computation of the top coefficient without \cref{thm:lea-Aj} ($k\geq 1$): put $z_0 := -E/(k+1)$. Only the term $n = k$ of \eqref{eq:lea-Rk}, and in $\Hcoef_k$ only $\Fhyp_k$, has a pole at $z_0$ ($d_{k,k} = 1$), so $a_{k+1}^{(k)} = (k+1)\,\frac{Ez_0}{E+z_0}\Res_{z_0}\Fhyp_k$ with $Ez_0/(E+z_0) = -E/k$, $\alpha_2+iz_0 = ((k+1)\alpha_2-iE)/(k+1)$, $\wA+iz_0 = ((k+1)\wA-iE)/(k+1)$, $z_0^k = (-E)^k/(k+1)^k$, $\prod_{i=2}^{k}(E+iz_0) = E^{k-1}(k-1)!/(k+1)^{k-1}$ and $\Res_{z_0}1/(E+(k+1)z) = 1/(k+1)$; hence $\Res_{z_0}\Fhyp_k = (-1)^k\Psh_k/((k+1)^2k!\,(k-1)!\,E^{2k-1})$ and $a_{k+1}^{(k)} = (k+1)(-E/k)\Res_{z_0}\Fhyp_k = (-1)^{k+1}\Psh_k/((k+1)!\,k!\,E^{2k-2})$. For $k = 0$, $a_1^{(0)} = -E^2 = -\Psh_0E^2$.

\begin{thm}[the loops in closed form]\genusstatus{AI-proposed, unchecked}
\label{thm:lea-loops}
With $\Ahat_j = t\Aleaf_j$,
\begin{equation}
\label{eq:lea-loops}
\begin{gathered}
\rho^2\Edge_{11}(\psi,\psi') = \sum_{i,l\geq 1}e_{il}(t)\lf_i(\psi)\lf_l(\psi'),\\
e_{il}(t) = \frac{il}{E^2t^2}\int_0^t\Thf(u)\,\Ahat_i'(u)\,\Ahat_l'(u)\,du
= -\frac{il\,\epsilon_1\epsilon_2\wB}{E^2t^2}\int_0^{\tau}v(1-v)\,\sast(v)\Lv_2(v)\,\frac{d\Ahat_i}{dv}\frac{d\Ahat_l}{dv}\,dv,
\end{gathered}
\end{equation}
and $e_{il} = O(t^{(i-1)+(l-1)})$.
\end{thm}

\begin{proof}
By \cref{thm:lea-Aj}, $\Dg(\psi) = \frac{d}{dt}(t\gl(\psi)) = -\frac1E\sum_jj\Ahat_j'\lf_j(\psi)$. Insert this in \eqref{eq:lea-flow} and integrate from $0$ (\eqref{eq:lea-int}); the coefficient of $\lf_i(\psi)\lf_l(\psi')$ in $\Ebar = t^2\rho^2\Edge_{11}$ is $\frac{il}{E^2}\int_0^t\Thf\Ahat_i'\Ahat_l'\,du$. For the second form substitute $v = \tau(u)$: $du = dv/(1-v)$, $\Ahat' = (1-v)\,d\Ahat/dv$ and $\Thf = -\epsilon_1\epsilon_2\wB\,v\,\sast(v)\Lv_2(v)$ \eqref{eq:lea-Theta}, so $\Thf\Ahat_i'\Ahat_l'\,du = -\epsilon_1\epsilon_2\wB\,v(1-v)\sast\Lv_2\,\frac{d\Ahat_i}{dv}\frac{d\Ahat_l}{dv}\,dv$. The order follows from $\Thf = O(t)$ and $\Ahat_j' = O(t^{j-1})$.
\end{proof}

\begin{cor}[anti-diagonal and corner]\genusstatus{AI-proposed, unchecked}
\label{cor:lea-anti}
Since $\Thf = \Thf_1t + O(t^2)$ and $\Ahat_j' = j\,a_j^{(j-1)}t^{j-1} + O(t^j)$, for $i+l = k+2$:
\begin{equation}
\label{eq:lea-antidiag}
e_{il}^{(k)} = \frac{\Thf_1\,i^2l^2\,a_i^{(i-1)}a_l^{(l-1)}}{E^2(i+l)} = -\frac{\epsilon_1\epsilon_2\alpha_1\alpha_2\wA\wB\,i^2l^2\,a_i^{(i-1)}a_l^{(l-1)}}{E^4(i+l)} = -(-1)^{i+l}\frac{\epsilon_1\epsilon_2\alpha_1\alpha_2\wA\wB\,i^2l^2\,\Cj_i\Cj_l}{E^4(i+l)},
\end{equation}
which is \eqref{eq:lea-anti}; with $a_1^{(0)} = -E^2$ the corner is $e_{1,k+1}^{(k)} = \frac{(k+1)^2}{k+2}\frac{\epsilon_1\epsilon_2\alpha_1\alpha_2\wA\wB}{E^2}a_{k+1}^{(k)}$. Example $k = 1$: $e_{12}^{(1)} = e_{21}^{(1)} = -\frac{2\epsilon_1\epsilon_2\alpha_1\alpha_2\wA\wB\,\uq}{3E^2}$, $e_{22}^{(1)} = 0$. The anti-diagonal entries are non-zero if and only if $\Psh_{i-1}\Psh_{l-1}\neq 0$, so the support bound of (L3) is sharp entry by entry at generic weights.
\end{cor}

\begin{proof}
The leading term of the integrand in \eqref{eq:lea-loops} is $\Thf_1\,il\,a_i^{(i-1)}a_l^{(l-1)}u^{i+l-1}$; integrating gives $t^{i+l}/(i+l)$, and multiplying by $il/(E^2t^2)$ gives the coefficient of $t^{i+l-2} = t^k$. Insert $\Thf_1 = -\epsilon_1\epsilon_2\alpha_1\alpha_2\wA\wB/E^2$ and $a_j^{(j-1)} = (-1)^j\Cj_j$ (\cref{cor:lea-top}). For $k = 1$, $(i,l) = (1,2)$: $-\epsilon_1\epsilon_2\alpha_1\alpha_2\wA\wB\cdot 4\cdot(-E^2)(-\uq/2)/(3E^4)$, with $a_2^{(1)} = -\uq/2$ (\cref{prop:lea-sigma}); $e_{22}^{(1)} = 0$ by the support bound.
\end{proof}

\begin{rmk}[the legs in closed form]
\label{rmk:lea-legs}
\begin{enumerate}
\item[(a)] $\Edge_{11}$ legs: $\coef{t^k}(\rho f_{1,b}) = \sum_{k_1+k_2 = k}\coef{t^{k_1}}(\rho^{-1})\sum_{i,l}e_{il}^{(k_2)}\frac{l^b}{E^{b+1}}\lf_i(\psi)$ \eqref{eq:lea-L2one}: no $\lf_0$, and its top coefficient ($\lf_{k+1}$ at $t^k$, which requires $k_2 = k$, $i = k+1$, $l = 1$) is $e_{k+1,1}^{(k)}/E^{b+1} = \frac{(k+1)^2}{k+2}\frac{\epsilon_1\epsilon_2\alpha_1\alpha_2\wA\wB\,a_{k+1}^{(k)}}{E^{b+3}}$, for all $k$ and $b$. At $t^0$ the leg is $-\frac{\epsilon_1\epsilon_2\alpha_1\alpha_2\wA\wB}{2}E^{-b-1}\lf_1$.
\item[(b)] $\Edge_{12}$ at $t = 0$ (rank one):
\[
\Edge_{12}(\psi,w)\big|_{t=0} = \lf_1(\psi)\,\chi_{12}(w),\qquad \chi_{12}(w) := -\epsilon_1\epsilon_2\alpha_1\alpha_2\wA\wB\,\frac{\rr_{0,2}(-w;0)}{E+w} = -\epsilon_1\epsilon_2\alpha_1\alpha_2\wA\wB\,\frac{\hreg(w)^{-1}}{E+w}.
\]
Proof: at $t^0$, \eqref{eq:lea-E12leaf} has only $j = 1$ with $S_1^{(0)} = (E,\alpha_1\alpha_2) = E(1,\alpha_1\alpha_2/E)$, and for every covector $b$
\[
\Bigl\langle\Bigl(1,\frac{\alpha_1\alpha_2}{E}\Bigr),b\Bigr\rangle = \epsilon_1\epsilon_2\Bigl[\Bigl(\alpha_0\alpha_1\alpha_2 + \frac{(\alpha_1\alpha_2)^2}{E}\Bigr)b_0 + \Bigl(\alpha_1\alpha_2 - E\frac{\alpha_1\alpha_2}{E}\Bigr)b_1\Bigr] = \frac{\epsilon_1\epsilon_2\alpha_1\alpha_2\wA\wB}{E}\,b_0
\]
($\alpha_0E+\alpha_1\alpha_2 = \wA\wB$); so $\coef{t^0}(\rho\Edge_{12}) = -E\,\frac{\epsilon_1\epsilon_2\alpha_1\alpha_2\wA\wB}{E}\rr_{0,2}(-w;0)\lf_1(\psi)\lf_1(-w)$ and $\coef{t^0}\Edge_{12} = \coef{t^0}(\rho\Edge_{12})$ as $\rho(0) = 1$; then \cref{prop:lea-regrows}(a) with $\hreg(-w) = \hreg(w)^{-1}$. The orthogonality is not needed here.
\item[(c)] $\Edge_{12}$ top coefficients, given \cref{prop:lea-regrows}(a), (b): for all $k$, $b$,
\[
[\lf_{k+1}]\coef{t^k}(\rho f_{2,b}) = (k+1)^2\,\frac{\epsilon_1\epsilon_2\alpha_1\alpha_2\wA\wB}{E^2}\,a_{k+1}^{(k)}\,\coef{w^b}\frac{\hreg(w)^{-1}}{E+(k+1)w}.
\]
Proof: in \eqref{eq:lea-E12coef} only $j = k+1$ and $k_1 = k$ contribute to $\lf_{k+1}$, giving
\[
-(k+1)\Bigl\langle S_{k+1}^{(k)},\ \coef{w^b}\frac{\vv_{2,0}(-w)}{E+(k+1)w}\Bigr\rangle ;
\]
here $\vv_{2,0}(-w) = \hreg(w)^{-1}(1,\alpha_1\alpha_2/E)$, $\langle S,(1,\alpha_1\alpha_2/E)\rangle = \frac{\epsilon_1\epsilon_2\alpha_1\alpha_2\wA\wB}{E}S_0$ by (b), and $S_0 = g_{k+1}^{(k)} = -\frac{k+1}{E}a_{k+1}^{(k)}$ by \eqref{eq:lea-Sj}. The lower coefficients $[\lf_j]\coef{t^k}(\rho f_{2,b})$, $j\leq k$, involve $\vv_{2,k_2}$ with $k_2\geq 1$ and are determined recursively by \cref{rmk:lea-E22}(2); no closed form in $b$ is claimed for them.
\end{enumerate}
\end{rmk}

\begin{rmk}[low orders]
\label{rmk:lea-low}
All of the following has been derived above; $\kappa_2 := \epsilon_1\epsilon_2\alpha_1\alpha_2\wA\wB$.
\begin{itemize}
\item $\dphi = \frac{\alpha_2\wA}{E}t - \frac{\alpha_1\alpha_2\wA\wB}{2E^3}t^2 + O(t^3)$; $\alg = 1 - \frac{E^2+\uq}{2E^2}\tau + O(\tau^2)$; $\rho = 1 + \frac{\uq}{2E^2}t + O(t^2)$; $\Delta_1 = -\frac{\epsilon_1\epsilon_2E^3}{t^2}\bigl(1 - \frac{\uq}{E^2}t + O(t^2)\bigr)$; $\Delta_2(0) = \kappa_2/E$, $\Delta_2'(0) = 0$; $\lrem_0 = \frac{E}{2} - \alpha_0 - \frac{\alpha_1\alpha_2}{E}$; $\Thf_1 = -\kappa_2/E^2$.
\item $q^{(1)} = \frac{\uq}{2}(\lf_1-\lf_2)$; $\Gk_1 = \frac{\uq}{2(E+z)(E+2z)}$; $\gl_1 = \frac{\uq}{2E}(2\lf_2-\lf_1) = \frac{\uq}{2}\lf_1\lf_2$ (by (LAM5): $\lf_1\lf_2 = (2\lf_2-\lf_1)/E$); $\Gk_1 - \rho_1\Gk_0 = \coef{t^1}\rr_{0,1} = \frac{E}{E+z}\bigl(\Hcoef_1 - \frac{E^2+\uq}{2E^2}\bigr)$.
\item Top and subleading coefficients of $q^{(k)}$: $a_{k+1}^{(k)} = (-1)^{k+1}\Cj_{k+1}$, $a_k^{(k)} = (-1)^{k+1}k\Cj_k\uq/(2E^2)$ (\cref{cor:lea-top}); e.g.\ $a_3^{(2)} = -\Cj_3 = -\frac{\Psh_2}{12E^2} = -\frac{(3\alpha_2-E)(3\alpha_2-2E)(3\wA-E)(3\wA-2E)}{12E^2}$, $a_2^{(2)} = -\Cj_2\uq/E^2 = \frac{\uq^2}{2E^2}$ (since $\Cj_2 = \Psh_1/2 = \frac{(2\alpha_2-E)(2\wA-E)}{2} = \frac{(E-2\alpha_2)(2\alpha_1-E)}{2} = -\frac{\uq}{2}$), and $a_1^{(2)} = -a_2^{(2)} - a_3^{(2)}$ by the sum rule.
\item $\brow^{(0)} = (E\lf_1,\alpha_1\alpha_2\lf_1)$; $\coef{t^0}\Yv = (0,E)$.
\item Loops: $e_{11}^{(0)} = -\kappa_2/2$; at $k = 1$: $e_{12}^{(1)} = e_{21}^{(1)} = -\frac{2\kappa_2\uq}{3E^2}$, $e_{22}^{(1)} = 0$ ($e_{11}^{(1)}$ needs $\Thf_2 = 3\epsilon_1\epsilon_2E\coef{t^3}\dphi$, not computed here).
\item $\Edge_{11}$ legs at $t^0$: $-\frac{\kappa_2}{2}E^{-b-1}\lf_1$. $\Edge_{12}$ at $t^0$: $\lf_1(\psi)\chi_{12}(w)$.
\end{itemize}
Higher tables ($t^2$, $t^3$ of the loops and legs, $\brow^{(1)}$, $\Thf_2$, $\Thf_3$, $\lrem_1$, the explicit $\Edge_{12}$ leg polynomials in $b$) are not given.
\end{rmk}

\subsection{The case \texorpdfstring{$\alpha_0 = 0$}{alpha0 = 0}: rescaling to the fibrewise rows}
\label{sec:lea-a0}

The proof of \cref{thm:lea-main} given above does not use $\alpha_0\neq 0$ (see remark (iv) below). Independently of this, at $\alpha_0 = 0$ the objects of the theorem are related by explicit rescalings to those at a fibrewise point $\alpha_0 = 0$, $\alpha_i = \epsilon_i$, which makes the fibrewise case the general one for $\alpha_0 = 0$.

\begin{cor}[$\alpha_0 = 0$]\genusstatus{AI-proposed, unchecked}
\label{cor:lea-a0}
Let $\alpha_0 = 0$ and $\alpha_1\alpha_2E\epsilon_1\epsilon_2\neq 0$ (then $\wA = \alpha_2$, $\wB = \alpha_1$, $E = \alpha_1+\alpha_2$), and put $\xc := \alpha_1/\alpha_2$ (so $\xc\neq 0,-1$) and $\cc := 1+\xc = E/\alpha_2$. Let the superscript $\TFsup$ denote the frame-B objects of this note at the fibrewise weight point $(\alpha_0,\alpha_1,\alpha_2,\epsilon_1,\epsilon_2) = (0,\xc,1,\xc,1)$. Then, for $i = 1,2$ and all $t$:
\begin{gather*}
\rr_{0,i}(z;t) = \rr^{\TFsup}_{0,i}(z/\alpha_2;t),\qquad \rr_{1,i}(z;t) = \alpha_2\,\rr^{\TFsup}_{1,i}(z/\alpha_2;t),\qquad \Lv_i = \alpha_2\Lv_i^{\TFsup},\\
\Delta_i = \kap\,\Delta_i^{\TFsup},\quad \kap := \frac{\epsilon_1\epsilon_2\alpha_2^4}{\alpha_1},\qquad
\langle a,b\rangle = \kap\,\langle a^{\TFsup},b^{\TFsup}\rangle^{\TFsup}\ \text{for }a = (a_0,\alpha_2a_1^{\TFsup}),\ b\text{ likewise},\\
\Dil_i(d) = \alpha_2^{1-d}\Dil_i^{\TFsup}(d),\qquad \Edge_{ij}(m,n) = \kap\,\alpha_2^{-1-m-n}\,\Edge_{ij}^{\TFsup}(m,n),\qquad \rho = \rho^{\TFsup},\\
\lf_j(\psi) = \alpha_2^{-1}\lf_j^{\TFsup}(\psi/\alpha_2),\qquad \qt(\psi) = \alpha_2\,\qt^{\TFsup}(\psi/\alpha_2),\qquad f_{j,b}(\psi) = \kap\,\alpha_2^{-1-b}f^{\TFsup}_{j,b}(\psi/\alpha_2).
\end{gather*}
Consequently \cref{thm:lea-main} at $\alpha_0 = 0$ is \cref{thm:lea-main} at the fibrewise point $(0,\xc,1,\xc,1)$ read through these rescalings; no new row analysis is needed.
\end{cor}

\begin{proof}
(1) The superpotential. At $\alpha_0 = 0$, $\cW = \alpha_1\log(1-\cQ U) - \alpha_2\log U + \alpha_2\log(1-U) = \alpha_2\cW^{\TFsup}$ with $\cW^{\TFsup} := \xc\log(1-\cQ U) - \log U + \log(1-U)$, the superpotential at $(0,\xc,1,\xc,1)$; so $\cW_k = \alpha_2\cW_k^{\TFsup}$, the critical points coincide ($\Pq = \alpha_2\Pq^{\TFsup}$), and by the Wick form (D1) and the rescaling rule (D4) of \cref{def:row-Lnorm} ($\Lnorm[\lambda f,1](z) = \Lnorm[f,1](z/\lambda)$, applied in the coordinate $y = \log U$), $\rpt_{0,i}(z) = {\rpt}^{\TFsup}_{0,i}(z/\alpha_2)$.
(2) The twist. By \cref{lem:set-hmult}, $\rr_{0,i} = h_{(\alpha_1,\alpha_2,-E)}(z)\rpt_{0,i}$, and $h_{(\alpha_1,\alpha_2,-E)}(z) = h_{(\xc,1,-\cc)}(z/\alpha_2)$, because the exponent $-\sum_k\bk_kz^{2k-1}(\alpha_1^{1-2k}+\alpha_2^{1-2k}-E^{1-2k})$ equals $-\sum_k\bk_k(z/\alpha_2)^{2k-1}(\xc^{1-2k}+1-\cc^{1-2k})$ (rescaling the multiset by $\alpha_2$ rescales $z$); and $h_{(\xc,1,-\cc)}$ is the frame-B twist at the fibrewise point, where $\bw^{\TFsup} = (\cc,-\xc,-1)$ and $\hnw^{\TFsup} = h_{(\xc,1,-\xc,-1)} = 1$, so that $h_{(\xc,1,-\cc)} = (h_{\bw^{\TFsup}})^{-1}$ by $h_{-v} = h_v^{-1}$. So $\rr_{0,i}(z) = \rr^{\TFsup}_{0,i}(z/\alpha_2)$.
(3) $\Lv = -\alpha_1\cQ U/(1-\cQ U) = \alpha_2\Lv^{\TFsup}$. $\Dop_i$ is invariant under a constant rescaling of $\Delta_i$ (shown in (4)). So $\rr_{1,i}(z) = (\alpha_2\Lv_i^{\TFsup} + z\Dop_i)\rr^{\TFsup}_{0,i}(z/\alpha_2) = \alpha_2\bigl(\Lv_i^{\TFsup} + (z/\alpha_2)\Dop_i\bigr)\rr^{\TFsup}_{0,i}(z/\alpha_2) = \alpha_2\rr^{\TFsup}_{1,i}(z/\alpha_2)$.
(4) Pairing and $\Delta_i$. At $\alpha_0 = 0$, $\langle a,b\rangle = \epsilon_1\epsilon_2[\alpha_1\alpha_2(a_0b_1+a_1b_0) - Ea_1b_1]$ (\cref{prop:set-frob}); with $a_1 = \alpha_2a_1^{\TFsup}$, $b_1 = \alpha_2b_1^{\TFsup}$ and $\alpha_1\alpha_2 = \xc\alpha_2^2$, $E = \cc\alpha_2$,
\[
\langle a,b\rangle = \epsilon_1\epsilon_2\alpha_2^3\bigl[\xc(a_0b_1^{\TFsup} + a_1^{\TFsup}b_0) - \cc\,a_1^{\TFsup}b_1^{\TFsup}\bigr] = \frac{\epsilon_1\epsilon_2\alpha_2^3}{\xc}\langle a^{\TFsup},b^{\TFsup}\rangle^{\TFsup},
\]
since at the fibrewise point $\epsilon_1\epsilon_2 = \xc$ and $\langle a,b\rangle^{\TFsup} = \xc^2(a_0b_1+a_1b_0) - \cc\xc\,a_1b_1$; and $\epsilon_1\epsilon_2\alpha_2^3/\xc = \kap$. $\Delta_i = \epsilon_1\epsilon_2(\Lv_i-\Lv_j)(\alpha_1\alpha_2 - E\Lv_i) = \epsilon_1\epsilon_2\alpha_2^3(\Lv_i^{\TFsup}-\Lv_j^{\TFsup})(\xc - \cc\Lv_i^{\TFsup}) = (\epsilon_1\epsilon_2\alpha_2^3/\xc)\Delta_i^{\TFsup}$, since $\Delta_i^{\TFsup} = \xc(\Lv_i^{\TFsup}-\Lv_j^{\TFsup})(\xc-\cc\Lv_i^{\TFsup})$ by the second form of \cref{prop:set-idem}. Hence $\Delta_i'/\Delta_i = (\Delta_i^{\TFsup})'/\Delta_i^{\TFsup}$ and $\Dop_i = \Dop_i^{\TFsup}$.
(5) $\Dil_i(d) = (-1)^d\rr_{0,i}[d-1] = \alpha_2^{1-d}\Dil_i^{\TFsup}(d)$ by (2). By (3), $\vv_i(z) = (\rr^{\TFsup}_{0,i}(z/\alpha_2),\,\alpha_2\rr^{\TFsup}_{1,i}(z/\alpha_2))$ has the form of (4), so $\Pij_{ij}(z,w) = -\langle\vv_i(-z),\vv_j(-w)\rangle = \kap\,\Pij_{ij}^{\TFsup}(z/\alpha_2,w/\alpha_2)$, and by \cref{lem:set-edgegen}
\[
\Edge_{ij}(z,w) = \frac{\Pij_{ij}(z,w) - \Pij_{ij}(-w,w)}{z+w} = \frac{\kap}{\alpha_2}\,\frac{\Pij_{ij}^{\TFsup}(z/\alpha_2,w/\alpha_2) - \Pij_{ij}^{\TFsup}(-w/\alpha_2,w/\alpha_2)}{z/\alpha_2 + w/\alpha_2} = \frac{\kap}{\alpha_2}\Edge_{ij}^{\TFsup}\Bigl(\frac{z}{\alpha_2},\frac{w}{\alpha_2}\Bigr),
\]
i.e.\ $\Edge_{ij}(m,n) = \kap\,\alpha_2^{-1-m-n}\Edge_{ij}^{\TFsup}(m,n)$.
(6) $\rho$ and the leaves. $\Delta_1t^2/(-\epsilon_1\epsilon_2E^3) = \kap\Delta_1^{\TFsup}t^2/(-\epsilon_1\epsilon_2\alpha_2^3\cc^3) = \Delta_1^{\TFsup}t^2/(-\xc\cc^3)$, which is the defining expression of $\rho^{\TFsup}$ (at the fibrewise point $-\epsilon_1\epsilon_2E^3 = -\xc\cc^3$); so $\rho = \rho^{\TFsup}$. $\lf_j(\psi) = 1/(\alpha_2\cc - j\psi) = \alpha_2^{-1}/(\cc - j\psi/\alpha_2) = \alpha_2^{-1}\lf_j^{\TFsup}(\psi/\alpha_2)$. $\qt(\psi) = -\psi\rr_{0,1}(-\psi) = \alpha_2\cdot(-\psi/\alpha_2)\rr^{\TFsup}_{0,1}(-\psi/\alpha_2) = \alpha_2\qt^{\TFsup}(\psi/\alpha_2)$. Finally $f_{j,b}(\psi) = \sum_a\Edge_{1j}(a,b)\psi^a = \sum_a\kap\alpha_2^{-1-a-b}\Edge_{1j}^{\TFsup}(a,b)\psi^a = \kap\alpha_2^{-1-b}\sum_a\Edge^{\TFsup}_{1j}(a,b)(\psi/\alpha_2)^a = \kap\alpha_2^{-1-b}f^{\TFsup}_{j,b}(\psi/\alpha_2)$.
(7) Each statement of \cref{thm:lea-main} at $\alpha_0 = 0$ is now the corresponding statement at the fibrewise point: membership in $\Lsp_{k+1}$ is preserved by $\psi\mapsto\psi/\alpha_2$ and by the overall constants ($\lf_j^{\TFsup}(\psi/\alpha_2) = \alpha_2\lf_j(\psi)$); (F1) is preserved since $\rho = \rho^{\TFsup}$ and $\Delta_2 = \kap\Delta_2^{\TFsup}$; (F2) since $E/(E+z) = \cc/(\cc+z/\alpha_2)$ and $E/t + \alpha_1\alpha_2/(E+z) = \alpha_2[\cc/t + \xc/(\cc+z/\alpha_2)]$; (F3) since the rescalings are by constants; the symmetries $\Edge_{11}(a,b) = \Edge_{11}(b,a)$ and $\Edge_{21}(b,a) = \Edge_{12}(a,b)$ are preserved since the factor $\alpha_2^{-1-m-n}$ is symmetric.
\end{proof}

Remarks. (i) The fibrewise case needs $\xc\neq 0,-1$, i.e.\ $\alpha_1\neq 0$ and $E\neq 0$; $\alpha_2\neq 0$ for the rescaling; $\epsilon_1\epsilon_2\neq 0$ for $\kap\neq 0$: exactly (ND) at $\alpha_0 = 0$ (where $\wA = \alpha_2$, $\wB = \alpha_1$). (ii) The $\cQ = 0$ constant at $\alpha_0 = 0$ is $\hnw = h_{(\alpha_1,\alpha_2,-\epsilon_1,-\epsilon_2)}$, not $1$ unless $\alpha_i = \epsilon_i$; it is the frame change $\hw^{-1}\hnw = h_{(\alpha_1,\alpha_2,-E)}$ that is independent of $\epsilon_1,\epsilon_2$. At the fibrewise point both are trivial. (iii) The vertex class $\Hg_{\bw}$, $\bw = (E,-\epsilon_1,-\epsilon_2) = \alpha_2(\cc,-\epsilon_1/\alpha_2,-\epsilon_2/\alpha_2)$, differs from the fibrewise class $(\cc,-\xc,-1)$ unless $\epsilon_1 = \alpha_1$; this concerns the regularity theorem of \cref{sec:reg}, not \cref{thm:lea-main}. (iv) The direct proof of \cref{sec:row} and of \cref{sec:lea-qde,sec:lea-L1,sec:lea-sigma,sec:lea-flow,sec:lea-regular,sec:lea-legs,sec:lea-assembly} covers $\alpha_0 = 0$ as well: the point $\cQ = 0$ is never used there, and $\hnw$ enters only through the multiset identity of \cref{lem:set-hmult}. The only statements of this section that need $\alpha_0\neq 0$ are the citations of \cref{lem:ide-norm} for the definition of the constant $\hnw$ and of \cref{prop:ide-rows}(iv) as an alternative source for $\Ktwo = 1$; both extend to $\alpha_0 = 0$ by rationality in the weights (the passage (P) of \cref{lem:set-rings}) once the factors are defined there. This is what the three-torus corollary of \cref{sec:rig} uses (\cref{cor:rig-3T}).

\begin{rmk}[the fibrewise torus as a special case]
\label{rmk:lea-fibrewise}
At the fibrewise point $(\alpha_0,\alpha_1,\alpha_2,\epsilon_1,\epsilon_2) = (0,\xc,1,\xc,1)$ with $\cc = 1+\xc$, \cref{thm:lea-main} is the leaf-structure theorem of the fibrewise torus $\TF$, and the objects of this section take the following form: $\wA = 1$, $\wB = \xc$, $E = \cc$, $\uq = -(1-\xc)^2$, $\epsilon_1\epsilon_2 = \xc$, $-\epsilon_1\epsilon_2E^3 = -\xc\cc^3$; the Euler phase is $\phi_\tau(s) = \log s + \xc\log(1-s) - \log(1-\tau s)$ with $\phi_0''(1/\cc) = -\cc^3/\xc$; $\delta^2 = \cc^2 - 4\xc\tau$, $\sast = 2/(\cc+\delta)$, $\stil = (\cc+\delta)/(2\xc\tau)$, $U_2(0) = -1/\xc$; the class is $\bw = (\cc,-\xc,-1)$, $\hnw = 1$, and the frame-B rows are $\hw^{-1}\rpt_{0,i}$; $\Lv_1 = 1/(\tau\sast) = \cc/\tau - \xc\sast$, $\Lv_2 = \xc\sast$, $\muq = \cc/\tau - 1$, $\lrem_0 = \cc/2 - \xc/\cc$; the pairing is $\langle a,b\rangle = \xc^2(a_0b_1+a_1b_0) - \cc\xc\,a_1b_1$; $\Delta_1 = -\xc\delta/(\tau^2{\sast}^2)$, $\Delta_2 = \xc^3\delta{\sast}^2$, $\Delta_2(0) = \xc^3/\cc$; $\Thf = -\xc^3\tau{\sast}^2$, $\Thf_1 = -\xc^3/\cc^2$; $\kappa_2 = \xc^3$; $\hreg = h_{(\xc,1,1,\xc,-\cc,-\cc)} = \hw^{-2}$, so $\vv_2(z;0) = \hw(z)^{-2}(1,\xc/\cc)$; $\Psh_k = \prod_{i=1}^{k}(k+1-i\cc)^2$, $\Cj_j = \cc^2[(1-j/\cc)_{j-1}]^2/(j!\,(j-1)!)$, and $\Aleaf_j = (-1)^j\Cj_jt^{-1}\zq^j\,{}_2F_1(j\xc/\cc,j\xc/\cc;j;\tau)$ (as $\wB/E = \alpha_1/E = \xc/\cc$). The identity $(1-\tau)\alpha_2s = (\wA-\wB s)(1-\tau s)$ of \cref{lem:set-N}(ii) becomes $1-s = \xc s(1-\tau s)$ at the critical points, and the numerator of $\Fhyp_n$ becomes $\prod_{j=1}^{n}(1+jz)^2$, with double zeros at $-1/j$; neither special form is used in the proofs above. By \cref{cor:lea-a0}, every weight point with $\alpha_0 = 0$ and (ND) is a rescaling of such a fibrewise point.
\end{rmk}


% s05_identification. Section file; new macros at the top; label prefix: ide.
% New macros (section-local; none of them clashes with main.tex):
\newcommand{\ideKt}{\widetilde{\bbK}_{\mathrm{CY}}}      % \KCY with e(wt_+)^{1/2}, e(wt_-)^{1/2} adjoined
\newcommand{\ideFz}{\bbF_{\circ}}                        % \ideKt((Q)), the field of the flatness proof
\newcommand{\ideone}{\mathbf{1}}                         % the unit of the state space; also the trivial rank-one CohFT
\newcommand{\ideV}{\mathcal{V}}                          % state space \KCY 1 + \KCY H (with scalars extended)
\newcommand{\idePF}{\mathsf{PF}}                         % the second-order operator of Euler's identity
\newcommand{\idekm}{\mathsf{k}_{\mathrm{m}}}             % its amplitude -z\alpha_1 Q(1-U)/(1-QU)
\newcommand{\ideCh}{\mathcal{X}}                         % \ideCh_h = \sum_k b_k (2k)!/B_{2k} ch_{2k-1}(E)
\newcommand{\ideCurve}{\mathfrak{C}}                     % universal curve over the open moduli space
\newcommand{\idesec}{\varsigma}                          % its sections \idesec_l
\newcommand{\ideLam}{\boldsymbol{\Lambda}}               % diag(L_+, L_-)
\newcommand{\idecan}{\boldsymbol{u}}                     % diag(u_+, u_-), canonical coordinates
\newcommand{\ideS}{\mathscr{S}}                          % fundamental solution R \Psi e^{u/z}
\newcommand{\idelg}[1]{\mathrm{a}^{#1}}                  % leg functions \idelg{i}(\psi;x) = f_i^*(B(\psi)x)
\newcommand{\ideP}{\mathfrak{P}}                         % \ideP_i = \Delta_i^{(2g-2+n)/2}
\newcommand{\ideN}{\mathfrak{N}}                         % kappa series of the translation legs
\newcommand{\ideG}{\mathfrak{G}}                         % kappa series of the pushforward
\newcommand{\ideW}{\mathsf{W}_{\Psi}}                    % \Psi^{-1}\thQ\Psi
\newcommand{\ideRf}{\mathsf{R}^{\mathrm{f}}}             % \Psi^{-1}\Rm\Psi
\newcommand{\ideY}{\mathsf{Y}}                           % the pairing <v_+(z), v_-(-z)> in the symplecticity step
\newcommand{\ideZ}{\mathsf{Z}}                           % \Psi^{-1}\Rm_1^{-1}\Rm_2\Psi in the ambiguity lemma
\newcommand{\idesp}{\mathfrak{s}}                        % the scalar series \idesp_p = \exp\sum_k a_k^{p} z^{2k-1}
\newcommand{\ideQU}{\mathrm{V}}                          % \ideQU = Q U_- , the regular coordinate at Q = 0
\newcommand{\ideT}{\mathsf{T}_{\!\Rm}}                   % translation z(1 - R^{-1}(z) 1)

\section{Identification of the graph sum with the Gromov--Witten potential}
\label{sec:ide}

This section proves Theorem~B: for $g\geq 2$ the frame-B Hodge-vertex graph sum $\Gsum_g$ of \cref{def:set-graphsum} is the Gromov--Witten potential $\sF_g=\sum_{d\geq 0}\GW_{g,d}\cQ^d$ of $Z$ on the Calabi--Yau torus. The argument runs as follows. The equivariant classes of $Z$ form a cohomological field theory which is semisimple at $\cQ=0$ on $\TCY$, because $\alpha_0\neq 0$ there; Teleman's classification writes it as $\Rm.\TFT$ for a unique symplectic $\Rm$, which we show to be a flat solution of the quantum differential equation. The rows built from the normalised stationary-phase expansions of the Euler integrand, multiplied by one scalar constant $\hnw$, form another such solution; the two agree up to constants, which are fixed at $\cQ=0$ by the Hodge integrals of constant maps. The change to frame~B is a scalar in the group of \cite{PPZ15}. Throughout, the weights are the generic point of $\TCY$: the field is $\KCY=\bbQ(\alpha_0,\alpha_1,\alpha_2,\epsilon_1)$ with $\epsilon_2\coloneqq E-\epsilon_1$, and every hypothesis on the weights used below is a non-vanishing or non-constancy statement which holds in $\KCY$. The inputs of this section are \cite{Tel12}, \cite{PPZ15}, \cite{GP99}, \cite{CG07}, Mumford's relation and formula \cite{Mum83,FP00}, \cite[Theorem~4, p.~3]{FP00} and, through \cref{lem:row-betaphase}, Stirling's series. Theorem~A is not used; \cref{thm:lea-qde} is cited in a remark only. The nine steps of the flatness proof are named (Fl1)--(Fl9) here, to keep them apart from the properties (F1)--(F3) of Theorem~A.

\subsection{Statement}

\begin{prop}[Theorem~B: the frame-B formula on $\TCY$]\genusstatus{AI-proposed, unchecked}
\label{prop:ide-main}
Let $g\geq 2$. Then $\sF_g$ lies in $\KCY(\cQ)$, is regular at $\cQ=0$ with power series expansion $\sum_{d\geq 0}\GW_{g,d}\cQ^d$ there, and
\begin{equation}
\label{eq:ide-main}
\sF_g=\Gsum_g=\sum_{\Gamma}\frac{1}{|\Aut\Gamma|}\sum_{\col\colon\Ver(\Gamma)\to\{1,2\}}\GCont(\Gamma,\col)
\end{equation}
in $\KCY(\cQ)$, the sum over the connected stable graphs of genus $g$ without legs, with
\[
\GCont(\Gamma,\col)=\sum_{\mathrm{dec}}\prod_v\Delta_{\col_v}^{g_v-1}\int\Hg_{\bw}\,\psi^{a}\prod\Dil_{\col_v}(d)\ \prod_e\Edge_{\col\col'}(a_e,b_e)
\]
as in \cref{eq:set-graphsum}: the vertex class is $\Hg_{\bw}=\prod_{a=1}^3\Lamv_{g_v}(w_a)/w_a^{g_v}$ with $\bw=(E,-\epsilon_1,-\epsilon_2)$, and the frame-B factors $\Delta_i$, $\Dil_i(d)$, $\Edge_{ij}(a,b)$ are those of \cref{def:set-factors}, computed from the rows $\vv_i=(\rr_{0,i},\rr_{1,i})$ with
\[
\rr_{0,i}=\hw^{-1}\hnw\,\rpt_{0,i}=h_{(\alpha_1,\wA,-E)}\,\rpt_{0,i},\qquad \rr_{1,i}=(\Lv_i+z\Dop_i)\,\rr_{0,i},\qquad \Dop_i=\frac{\dd}{\dd t}-\frac{\Delta_i'}{2\Delta_i},
\]
where $\rpt_{0,i}$ is the normalised stationary-phase expansion of $\re^{\cW/z}\dd U/U$ at the root $U_i$ (\cref{def:set-rows}) and $\nw=(\alpha_1,\wA,-\epsilon_1,-\epsilon_2)$. Moreover each coloured term $\GCont(\Gamma,\col)$ is an element of the function field $\KCY(\Ccov)$ of the double cover $\ydc^2=\Disc(\cQ)$, regular over every $\cQ\notin\{0,1,\infty\}$ with $\Disc(\cQ)\neq 0$.
\end{prop}

Two remarks on the statement. The field $\KCY$ is the function field of $\mathrm{Lie}(\TCY)$: on $\TCY$ the elements $\alpha_0,\alpha_1,\alpha_2,\wA,\wB,\epsilon_1,\epsilon_2,E,\alpha_1-\alpha_2$ are non-zero (\cref{def:set-ND}), and this is all that is used. The proof treats the four weights as indeterminates; it never passes through a non-semisimple point, because $\cQ=0$ is semisimple on $\TCY$ (\cref{lem:ide-ss}); a specialisation step would be needed only on the fibrewise torus, where $\alpha_0=0$, and that case is treated in \cref{sec:rig} by a ring map (\cref{lem:rig-3Tid}, \cref{rmk:ide-scope}).

\subsection{The cohomological field theory and its genus-zero data}

Let $\Omega=(\Omega_{g,n})$ denote the $\TCY$-equivariant Gromov--Witten classes of $Z$ with the Novikov variable $\cQ$ of degree $d[\bbP^1]$, defined by localisation on the fixed loci of the proper moduli space of maps to $\bbP^1$ (\cref{def:set-inv}), so that $\sF_g=\int_{\Mbar_g}\Omega_{g,0}$. Let $H$ be the equivariant lift of the hyperplane class of $\bbP^1$ with $H|_{p_+}=0$, $H|_{p_-}=-\alpha_0$, so that $H(H+\alpha_0)=0$ in $H^*_{\TCY}(\bbP^1)$, and let $\ideV\coloneqq\KCY\ideone\oplus\KCY H$ be the state space, with scalars extended whenever needed.

\begin{lem}[Fixed loci, the twisted theory, the axioms]\genusstatus{AI-proposed, unchecked}
\label{lem:ide-fixed}
\begin{enumerate}
\item A $\TCY$-fixed stable map to $Z$ has image in the zero section; the $\TCY$-fixed loci of $\Mbar_{g,n}(Z,d)$ are those of the proper space $\Mbar_{g,n}(\bbP^1,d)$, and $R\pi_*f^*\Nb$, $\Nb=\cO(-1)_{\alpha_1}\oplus\cO(-1)_{\alpha_2}\oplus\cO_{-\epsilon_1}\oplus\cO_{-\epsilon_2}$, is moving on them.
\item $\Omega$ consists of the classes of the $e_{\TCY}(R\pi_*f^*\Nb)^{-1}$-twisted theory of $\bbP^1$ in the sense of \cite[Section~1]{CG07}:
\[
\Omega_{g,n}(\gamma_1,\dots,\gamma_n)=\sum_{d\geq 0}\cQ^d\,\mathrm{st}_*\Big([\Mbar_{g,n}(\bbP^1,d)]^{\vir}_{\TCY}\cap\prod_k\ev_k^*\gamma_k\cap e_{\TCY}(R\pi_*f^*\Nb)^{-1}\Big).
\]
\item $\Omega$ is a cohomological field theory with flat unit $\ideone$ over $\KCY\fps{\cQ}$ for the pairing
\[
\eta(a,b)=\int_{\bbP^1}\frac{ab}{e_{\TCY}(\Nb)}=\int_Z ab=\frac{(ab)|_{p_+}}{\eul(\wts_+)}+\frac{(ab)|_{p_-}}{\eul(\wts_-)},
\]
with $\eul(\wts_+)=\alpha_0\alpha_1\alpha_2\epsilon_1\epsilon_2$ and $\eul(\wts_-)=-\alpha_0\wA\wB\epsilon_1\epsilon_2$; in the basis $(\ideone,H)$ this is the pairing of \cref{prop:set-frob}.
\end{enumerate}
\end{lem}
\begin{proof}
(i) A compact irreducible $\TCY$-invariant curve in $Z$ maps to the fixed point of $\bbC^2$, that is, the origin, since $\epsilon_1,\epsilon_2\neq 0$, and a compact curve in the resolved conifold lies in the zero section, the conifold being affine away from it. Hence the image of a $\TCY$-fixed map lies in the zero section $\bbP^1$, on which $\TCY$ acts by the rotation of weight $\alpha_0\neq 0$, and the $\TCY$-fixed loci of $\Mbar_{g,n}(Z,d)$ are the $\alpha_0$-fixed loci of the proper space $\Mbar_{g,n}(\bbP^1,d)$. On such a locus the restriction of $R\pi_*f^*\Nb$ is moving: its weights are the fibre characters $\alpha_1,\alpha_2,-\epsilon_1,-\epsilon_2$ at $p_+$ and $\wB,\wA,-\epsilon_1,-\epsilon_2$ at $p_-$ (vertex and flag terms), and, on an edge of degree $d_e$ covering $\bbP^1$, the weights $\alpha_1+k\alpha_0/d_e$, $\alpha_2+k\alpha_0/d_e$, $1\leq k\leq d_e-1$, of $H^1(\bbP^1,\cO(-d_e))$; every one of these is a non-zero linear form on $\mathrm{Lie}(\TCY)$. \draftnote{The sources state ``$R\pi_*f^*\Nb$ is moving there'' without listing the weights; the list is the normalisation-sequence computation of the companion paper for $\cO(-3)$, read for $\cO(-1)_{\alpha_1}\oplus\cO(-1)_{\alpha_2}$; to be checked.}

(ii) By (i) and the localisation formula of \cite{GP99} applied on $\Mbar_{g,n}(\bbP^1,d)$ (to a generic one-parameter subgroup through which $\TCY$ acts with the same fixed loci; both sides are rational in the weights and such subgroups are Zariski dense), the fixed part of the obstruction theory of $Z$ on a fixed locus is that of $\bbP^1$, the virtual classes of the fixed loci agree, and the virtual normal bundle of $Z$ is that of $\bbP^1$ plus $R\pi_*f^*\Nb$; so the localisation definition of the classes of $Z$ is the twisted integral displayed. \draftnote{BS label: defn: refined GW invariant; the definition by localisation on the proper fixed locus.}

(iii) The $S_n$-symmetry is clear. Unit: $\Omega_{g,n+1}(\gamma,\ideone)=p^*\Omega_{g,n}(\gamma)$ by the fundamental-class axiom of the virtual class, because the twisting class is pulled back along the forgetful map, $R\pi_*f^*\Nb$ on $\Mbar_{g,n+1}(\bbP^1,d)$ being the pull-back of the one on $\Mbar_{g,n}(\bbP^1,d)$ \cite[App.~1~(iv)]{CG07}; and $\Omega_{0,3}(a,b,\ideone)=\eta(a,b)$, since the terms with $d>0$ are pulled back along the forgetful map $\Mbar_{0,3}(\bbP^1,d)\to\Mbar_{0,2}(\bbP^1,d)$ and vanish. Gluing: the splitting axiom of the virtual class, together with the statement that the restriction of $R\pi_*f^*\Nb$ to a boundary divisor is the sum of the two sides (or the one side, for a non-separating node) minus $\ev^*\Nb$ at the node \cite[App.~1~(v), (vi)]{CG07} (equivalently, the normalisation sequence); the node factor $e(\Nb)$ turns the diagonal class of $\bbP^1$ into the bivector $\eta^{-1}$. This replaces the Poincar\'e pairing of $\bbP^1$ by $\eta$, and the last formula is localisation on $Z$, whose tangent weights at $p_\pm$ are $\wts_\pm$. \draftnote{The fundamental-class and splitting axioms of the virtual class are used in standard form; a reference (Behrend--Manin) is to be added to the bibliography.}
\end{proof}

\begin{lem}[The genus-zero potential]\genusstatus{AI-proposed, unchecked}
\label{lem:ide-g0}
Let $t_0,t_1$ be the flat coordinates along $\ideone,H$. The genus-zero potential is
\[
\cF_0=\cF_0^{\mathrm{cl}}+\frac{\Li_3(\re^{t_1})}{\epsilon_1\epsilon_2},\qquad
\cF_0^{\mathrm{cl}}=\frac16\Big[\frac{t_0^3}{\alpha_0\alpha_1\alpha_2\epsilon_1\epsilon_2}-\frac{(t_0-\alpha_0t_1)^3}{\alpha_0\wA\wB\epsilon_1\epsilon_2}\Big],
\]
with $\cQ=\re^{t_1}$: that is, $\GW_{0,d}=1/(d^3\epsilon_1\epsilon_2)$ for $d\geq 1$ and $\langle H,H,H\rangle_{0,d}=d^3\GW_{0,d}=1/(\epsilon_1\epsilon_2)$. Consequently
\begin{gather*}
H\qprod H=(\Lv_++\Lv_-)\,H-\Lv_+\Lv_-\,\ideone,\\
\Lv_++\Lv_-=-\frac{\alpha_0+(\alpha_1+\alpha_2)\cQ}{1-\cQ}=\frac{E}{\tau}-\alpha_0,\qquad
\Lv_+\Lv_-=-\frac{\alpha_1\alpha_2\cQ}{1-\cQ}=\frac{\alpha_1\alpha_2}{\tau},
\end{gather*}
where $\Lv_\pm$ denote the two eigenvalues of $H\qprod$.
\end{lem}
\begin{proof}
The classical part is the localisation formula $\frac16\int_Z(t_0+t_1H)^3=\frac16\sum_{p=\pm}(t_0+t_1H|_p)^3/\eul(\wts_p)$ with $H|_{p_+}=0$, $H|_{p_-}=-\alpha_0$ and the values of $\eul(\wts_\pm)$ of \cref{lem:ide-fixed}(iii). For $d>0$: a section of $f^*\cO(-1)$ on a connected curve $C$ vanishes on the non-contracted components, where the bundle has negative degree, and is constant on the contracted ones, which meet them, so $H^0(C,f^*\cO(-1))=0$; hence $\Mbar_{0,0}(Z,d)=\Mbar_{0,0}(\bbP^1,d)\times\bbC^2$ with obstruction bundle $R^1\pi_*f^*(\cO(-1)\oplus\cO(-1))$, whose rank $2d-2$ is the virtual dimension of $\Mbar_{0,0}(\bbP^1,d)$. Localising on the factor $\bbC^2$ gives $\GW_{0,d}=C(0,d)/(\epsilon_1\epsilon_2)$ with $C(0,d)=\int_{[\Mbar_{0,0}(\bbP^1,d)]^{\vir}}e(R^1\pi_*f^*(\cO(-1)\oplus\cO(-1)))=1/d^3$, the Aspinwall--Morrison formula \cite[p.~3]{FP00}. The divisor equation of the twisted theory of $\bbP^1$ (the statement (Fl4) in the proof of \cref{prop:ide-flat}, which uses only the fundamental-class axiom and \cite[App.~1~(iv)]{CG07} and nothing else of that proof), applied in genus $0$, gives $\langle H,H,H\rangle_{0,d}=d^3\GW_{0,d}$, so the degree-$d$ part of the potential is $\GW_{0,d}\re^{dt_1}$ and the sum over $d\geq 1$ is $\Li_3(\re^{t_1})/(\epsilon_1\epsilon_2)$. The quantum product is read off from the third derivatives of $\cF_0$ through $\eta$: $\partial_{t_1}^3\Li_3(\re^{t_1})=\re^{t_1}/(1-\re^{t_1})=\cQ/(1-\cQ)=-1/\tau$, and the computation of the structure constants, carried out in \cref{prop:set-frob}, gives $H\qprod H$ as displayed; $\Lv_\pm$ are the roots of $\Lv^2-(\Lv_++\Lv_-)\Lv+\Lv_+\Lv_-$.
\end{proof}

\begin{lem}[Semisimplicity at $\cQ=0$]\genusstatus{AI-proposed, unchecked}
\label{lem:ide-ss}
The eigenvalues $\Lv_\pm$ of $H\qprod$ lie in $\KCY\fps{\cQ}$, with $\Lv_+(0)=0$, $\Lv_-(0)=-\alpha_0$, and $\Lv_+-\Lv_-$ is a unit of $\KCY\fps{\cQ}$. The idempotents $\eid_\pm=(H-\Lv_\mp)/(\Lv_\pm-\Lv_\mp)$ lie in $\ideV\otimes\KCY\fps{\cQ}$, and $\eid_\pm|_{\cQ=0}$ are the point classes $\eid_p|_q=\delta_{pq}$ of $p_\pm$. With $\Delta_i\coloneqq 1/\eta(\eid_i,\eid_i)$,
\begin{equation}
\label{eq:ide-Delta}
\Delta_i=\epsilon_1\epsilon_2\big(\alpha_0\alpha_1\alpha_2+2\alpha_1\alpha_2\Lv_i-E\Lv_i^2\big)=\epsilon_1\epsilon_2(\Lv_i-\Lv_j)(\alpha_1\alpha_2-E\Lv_i)
=\alpha_1\wA\epsilon_1\epsilon_2\,\cW_2(U_i),
\end{equation}
$j\neq i$, and $\Delta_+(0)=\eul(\wts_+)=\alpha_0\alpha_1\alpha_2\epsilon_1\epsilon_2$, $\Delta_-(0)=\eul(\wts_-)=-\alpha_0\wA\wB\epsilon_1\epsilon_2$, both non-zero in $\KCY$. So the cohomological field theory $\Omega$ is semisimple over $\KCY\fps{\cQ}$, and the ratio $\Delta_+(0)/\Delta_-(0)=-\alpha_1\alpha_2/(\wA\wB)$ is a non-constant element of $\KCY$.
\end{lem}
\begin{proof}
By \cref{lem:ide-g0} the eigenvalues are the roots of $\Lv^2-(\Lv_++\Lv_-)\Lv+\Lv_+\Lv_-$, whose coefficients lie in $\KCY\fps{\cQ}$ and which reduces modulo $\cQ$ to $\Lv^2+\alpha_0\Lv=\Lv(\Lv+\alpha_0)$, with the distinct roots $0$ and $-\alpha_0$ since $\alpha_0\neq 0$; by Hensel's lemma both roots lie in $\KCY\fps{\cQ}$, and their difference has the value $\alpha_0$ at $\cQ=0$, hence is a unit. Modulo $\cQ$ the quantum product is the classical product of $H^*_{\TCY}(Z)\otimes\KCY=\KCY\eid_+\oplus\KCY\eid_-$, in which $H=-\alpha_0\eid_-$; so $\eid_\pm(0)$ are the point classes. The formulas \cref{eq:ide-Delta} are \cref{prop:set-idem} and \cref{lem:set-crit}; we recall their derivation in \cref{prop:ide-rows}(i). The values at $\cQ=0$ follow from the first form with $\Lv_+(0)=0$ and $\Lv_-(0)=-\alpha_0$: $\Delta_-(0)=\epsilon_1\epsilon_2(\alpha_0\alpha_1\alpha_2-2\alpha_0\alpha_1\alpha_2-E\alpha_0^2)=-\alpha_0\epsilon_1\epsilon_2(\alpha_1\alpha_2+\alpha_0E)=-\alpha_0\wA\wB\epsilon_1\epsilon_2$, as $\wA\wB=\alpha_0E+\alpha_1\alpha_2$. They agree with $\eta(\eid_p(0),\eid_p(0))=1/\eul(\wts_p)$ from \cref{lem:ide-fixed}(iii). The ratio $-\alpha_1\alpha_2/((\alpha_0+\alpha_1)(\alpha_0+\alpha_2))$ is a non-constant rational function of $(\alpha_0,\alpha_1,\alpha_2)$.
\end{proof}

\subsection{Hodge classes as an R-matrix action}

The two lemmas of this subsection are statements about $\Mbar_{g,n}$ alone. They are placed here because the freeness lemma below uses them. For a scalar series $h=\sum_{k\geq 1}b_kz^{2k-1}$ put
\[
\ideCh_h\coloneqq\sum_{k\geq 1}b_k\frac{(2k)!}{B_{2k}}\,\mathrm{ch}_{2k-1}(\Hdg)\in H^*(\Mbar_{g,n};\bbQ),
\]
and write $\exp(h)$ for the rank-one R-matrix with logarithm $h$, acting on the trivial rank-one theory $\ideone$ by the unit-preserving action of \cite[Definitions~2.2, 2.13]{PPZ15}; thus $h_{\mathsf{v}}.\ideone$ denotes the action of the series $h_{\mathsf{v}}$ of \cref{def:set-rows}, whose logarithm is $-\sum_k\bk_kz^{2k-1}\sum_av_a^{1-2k}$. Recall $\Hg_{\mathsf{v}}=\prod_{a}\Lamv(v_a)/v_a^{g}$ for a multiset $\mathsf{v}$ of non-zero weights.

\begin{lem}[Hodge classes as an R-action]\genusstatus{AI-proposed, unchecked}
\label{lem:ide-mumford}
For all stable $(g,n)$, in the conventions $\kappa_{2l-1}=p_*(\psi_{n+1}^{2l})$ of Arbarello--Cornalba,
\begin{equation}
\label{eq:ide-mumford}
\mathrm{ch}_{2l-1}(\Hdg)=\frac{B_{2l}}{(2l)!}\Big[\kappa_{2l-1}-\sum_{i=1}^n\psi_i^{2l-1}+\frac12\xi_*\Big(\sum_{a+a'=2l-2}(-1)^{a'}\psi^a\psi'^{a'}\Big)\Big],
\end{equation}
including $g=0$, where both sides vanish ($\xi$ is the gluing map of the boundary with the two branches $\psi,\psi'$ at the node). Further $\exp(h).\ideone=\exp(\ideCh_h)$, and $\Hg_{\mathsf{v}}=h_{\mathsf{v}}.\ideone$.
\end{lem}
\begin{proof}
\emph{Mumford's formula with markings.} Mumford's formula as stated in \cite[Section~1.1]{FP00} applies to the family $\pi\colon\Mbar_{g,n+1}\to\Mbar_{g,n}$ of curves with the marked points as the divisor $D=\sum_iD_i$. With $K_\pi\coloneqq c_1(\omega_\pi)$, $K_\pi\cdot D_i=\psi_iD_i$, $D_i^2=-\psi_iD_i$ and $D_iD_j=0$ for $i\neq j$,
\[
(K_\pi+D_i)^{2l}=K_\pi^{2l}+\psi_i^{2l-1}D_i\sum_{r=1}^{2l}(-1)^{r-1}\binom{2l}{r}=K_\pi^{2l}+\psi_i^{2l-1}D_i,
\]
so $\pi_*(c_1(\omega_\pi(D))^{2l})=\pi_*(c_1(\omega_\pi)^{2l})+\sum_i\psi_i^{2l-1}$, which turns Mumford's formula for $\pi_*(c_1(\omega_\pi)^{2l})$ into \cref{eq:ide-mumford}.

\emph{The Hodge class.} The relation $c(\Hdg)c(\Hdg^\vee)=1$ \cite[(5.4), (5.5)]{Mum83} gives $\mathrm{ch}_{2k}(\Hdg)=0$ for $k\geq 1$, hence
\[
\frac{\Lamv(v)}{v^g}=c_{-1/v}(\Hdg)=\exp\Big(-\sum_{k\geq 1}(2k-2)!\,\mathrm{ch}_{2k-1}(\Hdg)\,v^{1-2k}\Big).
\]

\emph{The action.} By \cite[Prop.~2.14]{PPZ15}, $\frac{\dd}{\dd s}(\re^{sh}.\ideone)=\mathcal{D}_h(\re^{sh}.\ideone)$, where, from \cite[Defs.~2.2, 2.13]{PPZ15},
\[
\mathcal{D}_h\Omega_{g,n}=\sum_kb_k\Big[-\sum_i\psi_i^{2k-1}\Omega_{g,n}+p_*\big(\psi_{n+1}^{2k}\Omega_{g,n+1}\big)+\frac12\xi_*\Big(\sum_{a+a'=2k-2}(-1)^{a'}\psi^a\psi'^{a'}\Omega'\Big)\Big]:
\]
the legs give $-\psi^{2k-1}$, from $\Rm^{-1}(\psi)$; the translation $\ideT(z)=b_kz^{2k}$ gives the $\kappa$ term; and the edge gives $(z^{2k-1}+w^{2k-1})/(z+w)$. On $\exp(s\ideCh_h)$, which is pulled back under $p$ and additive under $\xi^*$ for $k\geq 1$ (loops included), the projection formula gives
\[
\mathcal{D}_h\exp(s\ideCh_h)=\exp(s\ideCh_h)\sum_kb_k\Big[\kappa_{2k-1}-\sum_i\psi_i^{2k-1}+\frac12\xi_*(\cdots)\Big]=\ideCh_h\exp(s\ideCh_h)
\]
by \cref{eq:ide-mumford}. So $\exp(s\ideCh_h)$ solves the same graded differential equation as $\re^{sh}.\ideone$ with the same initial value $1$ at $s=0$, which proves $\exp(h).\ideone=\exp(\ideCh_h)$.

\emph{Comparison.} The series $\log h_{\mathsf{v}}$ has $b_k=-\bk_k\sum_av_a^{1-2k}$ with $\bk_k=B_{2k}/(2k(2k-1))$, so $b_k(2k)!/B_{2k}=-(2k-2)!\sum_av_a^{1-2k}$ and $\ideCh_{\log h_{\mathsf{v}}}=-\sum_k(2k-2)!\,\mathrm{ch}_{2k-1}(\Hdg)\sum_av_a^{1-2k}=\log\Hg_{\mathsf{v}}$. This agrees with \cite[Prop.~8.7(i)]{Tel12} and \cite[Section~2.3, Lemma]{Giv01}.
\end{proof}

\begin{lem}[Non-vanishing of $\mathrm{ch}_{2k-1}$]\genusstatus{AI-proposed, unchecked}
\label{lem:ide-ch}
For $k\geq 2$, $\mathrm{ch}_{2k-1}(\Hdg)\neq 0$ on $\Mbar_k$, and $\mathrm{ch}_{2k-1}(\Hdg)\neq 0$ on $\Mbar_{k,1}$ for $k\geq 1$.
\end{lem}
\begin{proof}
For $c_t=\sum_i\lambda_it^i$ one has $c_t'c_{-t}=\sum_{n\geq 1}(-1)^{n-1}p_nt^{n-1}$ with $p_n=n!\,\mathrm{ch}_n$; the coefficient of $t^{2k-2}$, using $c(\Hdg)c(\Hdg^\vee)=1$ and $\lambda_i=0$ for $i>k$ on $\Mbar_k$ (the Hodge bundle has rank $k$), gives $(2k-1)!\,\mathrm{ch}_{2k-1}(\Hdg)=(-1)^{k-1}\lambda_k\lambda_{k-1}$. Further, by the coefficient of $t^{2k-2}$ in $c_tc_{-t}=1$ on $\Mbar_k$, $\lambda_{k-1}^2=2\lambda_k\lambda_{k-2}$, so $\int_{\Mbar_k}\lambda_k\lambda_{k-1}\lambda_{k-2}=\frac12\int_{\Mbar_k}\lambda_{k-1}^3\neq 0$ for $k\geq 2$ by \cite[Theorem~4]{FP00}; hence $\mathrm{ch}_{2k-1}(\Hdg)\neq 0$ on $\Mbar_k$. For the second statement, $\mathrm{ch}_{2k-1}(\Hdg)$ on $\Mbar_{k,1}$ is the pull-back $\pi^*$ of the class on $\Mbar_k$ for $k\geq 2$, and $\pi^*\colon H^*(\Mbar_k;\bbQ)\to H^*(\Mbar_{k,1};\bbQ)$ is injective because $\pi_*(\psi_1\pi^*y)=(2k-2)y$. For $k=1$, $\int_{\Mbar_{1,1}}\lambda_1=1/24$.
\end{proof}

\subsection{Teleman's classification over the field}

Let $\TFT$ be the semisimple topological field theory over $\KCY\fps{\cQ}$ with $\TFT_{g,n}(\eid_i^{\otimes n})=\Delta_i^{g-1}$ and $\TFT_{g,n}(x_1,\dots,x_n)=0$ for mixed insertions; its genus-zero part is the Frobenius algebra $(\ideV\otimes\KCY\fps{\cQ},\qprod,\eta)$. An R-matrix is an element $\Rm(z)=1+\Rm_1z+\Rm_2z^2+\cdots$ of $\End(\ideV)\fps{z}$ (scalars extended to a field containing $\KCY\lser{\cQ}$) which is symplectic, $\eta(\Rm(z)a,\Rm(-z)b)=\eta(a,b)$, equivalently $\Rm^{-1}(z)^\dagger=\Rm(-z)$ with $\dagger$ the $\eta$-adjoint \cite[(5)]{PPZ15}; $\Rm.\TFT$ denotes the unit-preserving action of \cite[Definitions~2.2, 2.13]{PPZ15}: a sum over stable graphs with the legs $\Rm^{-1}(\psi_l)$, the edges $(\eta^{-1}-(\Rm^{-1}(\psi)\otimes\Rm^{-1}(\psi'))\eta^{-1})/(\psi+\psi')$ and the translation $\ideT(z)=z(\ideone-\Rm^{-1}(z)\ideone)$. We write $\fid_i\coloneqq\Delta_i^{1/2}\eid_i$ for the normalised idempotents, which are $\eta$-orthonormal, and $\fid_i^*=\eta(\fid_i,\cdot)=\Delta_i^{-1/2}\eid_i^*$, where $\eid_i^*=\Delta_i\eta(\eid_i,\cdot)$ is the dual basis. The square roots $\Delta_i^{1/2}$ lie in $\ideKt\fps{\cQ}$, $\ideKt\coloneqq\KCY\big(\eul(\wts_+)^{1/2},\eul(\wts_-)^{1/2}\big)$, because $\Delta_i\in\KCY\fps{\cQ}$ has the non-zero value $\eul(\wts_i)$ at $\cQ=0$ (\cref{lem:ide-ss}); they are used only in the flatness proof and in the ambiguity lemma. Put $\ideFz\coloneqq\ideKt\lser{\cQ}$, a finite extension of $\KCY\lser{\cQ}$.

\begin{lem}[Freeness of the action]\genusstatus{AI-proposed, unchecked}
\label{lem:ide-free}
Let $\TFT$ be a semisimple topological field theory over a field of characteristic $0$ with $\TFT(\eid_i^{\otimes n})=\Delta_i^{g-1}$ (of any rank). If $\Rm$ is symplectic and $\Rm.\TFT=\TFT$, then $\Rm=1$.
\end{lem}
\begin{proof}
Suppose $\Rm\neq 1$ and write $\Rm=1+\Rm_mz^m+\cdots$ with $\Rm_m\neq 0$, in the basis $\fid_i$ (over the field with the $\Delta_i^{1/2}$ adjoined).

Evaluate the degree-$m$ part of $(\Rm.\TFT)_{0,m+3}(\fid_i,\fid_j,\dots,\fid_j)$ with $i\neq j$. A graph with an edge would need a vertex carrying $\fid_i$ and the edge alone, which is unstable, so only the one-vertex graphs occur; translations give mixed insertions, which vanish; so only the leg term at the $\fid_i$-leg survives, which is $-(\Rm_m)_{ji}\Delta_j^{(m+1)/2}\psi_1^m$, and $\psi_1^m\neq 0$ on $\Mbar_{0,m+3}$. Hence $\Rm_m$ is diagonal, $\Rm_m=\mathrm{diag}(r_i)$.

Symplecticity makes $m$ odd: $\Rm(z)\Rm(-z)^\dagger=1$ gives $\Rm_m+(-1)^m\Rm_m^\dagger=0$ at order $z^m$, and a diagonal $\Rm_m$ in the orthonormal basis is self-adjoint, so $m$ even would force $\Rm_m=0$. Write $m=2k-1$, $k\geq 1$.

\emph{Colour isolation.} Work on $\Mbar_{k,1}$ with the insertion $\eid_i$ (on $\Mbar_k$ there is no marking and $\TFT_{k,0}$ is a sum over colours, so one colour cannot be separated there). The degree-$m$ part of $(\Rm.\TFT)_{k,1}(\eid_i)$ is linear in $\Rm_m$, monochromatic (a half-edge at a vertex carrying $\eid_i$, or joined to it, must have colour $i$, because $\Rm_m$ is diagonal and mixed insertions vanish), and equals $\Delta_i^{k-1}$ times the degree-$m$ part of $(\exp(r_iz^m).\ideone)_{k,1}$, by $\sum_v(g_v-1)+|\Edg(\Gamma)|=g-1$. By \cref{lem:ide-mumford} this is $\Delta_i^{k-1}r_i\frac{(m+1)!}{B_{m+1}}\mathrm{ch}_m(\Hdg)$, and by \cref{lem:ide-ch} $\mathrm{ch}_{2k-1}(\Hdg)\neq 0$ on $\Mbar_{k,1}$. Hence $r_i=0$ for every $i$, a contradiction.
\end{proof}

\begin{prop}[Existence and field of definition of $\Rm$]\genusstatus{AI-proposed, unchecked}
\label{prop:ide-teleman}
There is a unique symplectic $\Rm\in\End(\ideV)\otimes\KCY\lser{\cQ}\fps{z}$ with $\Omega=\Rm.\TFT$.
\end{prop}
\begin{proof}
\emph{Field of definition.} The classification of \cite{Tel12} is stated over $\bbC$ (it allows the generic point of the Novikov variables, \cite[(1.6)]{Tel12}). The field $\KCY\lser{\cQ}$ has characteristic $0$ and cardinality $2^{\aleph_0}$, so it embeds in $\bbC$ as a field: send a transcendence basis over $\bbQ$ injectively into one of $\bbC$ and extend, $\bbC$ being algebraically closed. Fix such an embedding $\iota$. Since $H^*(\Mbar_{g,n};\bbQ)$ is finite dimensional, $\iota(\Omega)$ is a cohomological field theory with flat unit in $H^*(\Mbar_{g,n};\bbQ)\otimes\bbC$, semisimple because $\Omega$ is semisimple over $\KCY\lser{\cQ}$ (\cref{lem:ide-ss}), with the topological field theory $\iota(\TFT)$. By \cite[(1.15), Theorem~(2), Prop.~8.7(ii)]{Tel12}, $\iota(\Omega)=\Rm.\iota(\TFT)$ for a symplectic $\Rm$ with coefficients in $\bbC$; here \cite[Defs.~2.2, 2.5, 2.13 and Prop.~2.14]{PPZ15} are identities of graph sums whose proofs hold over any field of characteristic $0$. By \cref{lem:ide-free} and the group property \cite[Prop.~2.14]{PPZ15}, $\Rm$ is unique: if $\Rm_1.\TFT=\Rm_2.\TFT$ then $(\Rm_1^{-1}\Rm_2).\TFT=\TFT$, so $\Rm_1=\Rm_2$. The embedding forgets the derivation $\thQ$, which is recovered in \cref{prop:ide-flat} after the descent.

\emph{Descent.} Read $\Rm$ in $\End(\ideV\otimes\bbC)$, where $\ideV$ has the basis $(\ideone,H)$ over $\KCY\lser{\cQ}$. The idempotents $\eid_\pm=(H-\Lv_\mp)/(\Lv_\pm-\Lv_\mp)$ lie in $\ideV\otimes\KCY\fps{\cQ}$, because $\Lv_+-\Lv_-$ is a unit of $\KCY\fps{\cQ}$ (\cref{lem:ide-ss}), so $\TFT$ and $\Omega$ are defined over $\KCY\lser{\cQ}$, and no square root of $\Delta_i$ is needed for this step. Let $\phi\in\Aut(\bbC/\KCY\lser{\cQ})$. The action $\Rm.\TFT$ is given by polynomial formulas over $\bbQ$ in the entries of $\Rm$ and the values of $\TFT$, so $\phi(\Rm).\TFT=\phi(\Rm.\TFT)=\Omega=\Rm.\TFT$; hence $\Rm^{-1}\phi(\Rm)$ fixes $\TFT$ \cite[Prop.~2.14]{PPZ15}, and $\phi(\Rm)=\Rm$ by \cref{lem:ide-free}. The fixed field of $\Aut(\bbC/F)$ is $F$ for every subfield $F$ of $\bbC$: an element algebraic over $F$ and not in $F$ has a conjugate different from itself, an element $y$ transcendental over $F$ can be sent to $y+1$, and these automorphisms of $F(y)$, respectively of $F(\text{conjugates})$, extend to $\bbC$. So $\Rm$ has coefficients in $\KCY\lser{\cQ}$, and, by uniqueness over $\KCY\lser{\cQ}$, it does not depend on the embedding.
\end{proof}

Let $\Psi\colon\ideFz^2\to\ideV\otimes\ideFz$ be the operator $\mathbf{e}_i\mapsto\fid_i$, so that $\Psi^T\eta\Psi=1$, that is, $\eta(\Psi\mathsf{u},\Psi\mathsf{u}')=\mathsf{u}\cdot\mathsf{u}'$, and $\Psi^{-1}y=(\fid_+^*(y),\fid_-^*(y))$; let $\ideLam\coloneqq\mathrm{diag}(\Lv_+,\Lv_-)$, so that $H\qprod\Psi=\Psi\ideLam$. The derivation $\thQ=\cQ\,\dd/\dd\cQ=\dd/\dd t$ extends coefficientwise to $\ideFz$, with constants $\ideKt$; it kills $\eta$ and the elements of $\ideV$ over $\KCY$. Let $\idecan=\mathrm{diag}(u_+,u_-)$ be canonical coordinates, $\thQ\idecan=\ideLam$ (any choice; only $\thQ\idecan$ enters).

\begin{lem}[Ambiguity of flat symplectic solutions]\genusstatus{AI-proposed, unchecked}
\label{lem:ide-ambig}
Let $\Rm_1,\Rm_2\in\End(\ideV)\otimes\ideFz\fps{z}$ be symplectic, with $\Rm_1(0)=\Rm_2(0)=1$, both solving the flatness equation
\begin{equation}
\label{eq:ide-flat}
z\,\thQ\big(\Rm(z)\Psi\big)=\big(H\qprod\Rm(z)-\Rm(z)\,H\qprod\big)\Psi .
\end{equation}
Then $\Rm_2=\Rm_1\Psi\exp\big(\mathrm{diag}\sum_{k\geq 1}a_k^{(i)}z^{2k-1}\big)\Psi^{-1}$ with constants $a_k^{(i)}\in\ideKt$; that is, in the frame $\fid_i$ the two matrices differ by right multiplication by a diagonal matrix of exponentials of odd series with constant coefficients.
\end{lem}
\begin{proof}
Equation \cref{eq:ide-flat} says $z\thQ\Rm=H\qprod\Rm-\Rm H\qprod-z\Rm\,\thQ\Psi\,\Psi^{-1}$, since $\Psi\ideLam\Psi^{-1}=H\qprod$. Hence for $\mathfrak{Y}\coloneqq\Rm_1^{-1}\Rm_2$, using $z\thQ(\Rm_1^{-1})=-\Rm_1^{-1}(z\thQ\Rm_1)\Rm_1^{-1}$,
\[
z\thQ\mathfrak{Y}=[H\qprod,\mathfrak{Y}]+z\,[\thQ\Psi\,\Psi^{-1},\mathfrak{Y}],
\]
and for $\ideZ\coloneqq\Psi^{-1}\mathfrak{Y}\Psi$ the terms in $\thQ\Psi$ cancel: $z\thQ\ideZ=[\ideLam,\ideZ]$. Write $\ideZ=1+\sum_{n\geq 1}\ideZ_nz^n$. At order $z^n$, $\thQ\ideZ_{n-1}=[\ideLam,\ideZ_n]$, that is, $(\Lv_i-\Lv_j)(\ideZ_n)_{ij}=(\thQ\ideZ_{n-1})_{ij}$ for $i\neq j$ and $(\thQ\ideZ_{n-1})_{ii}=0$. By induction on $n$: $\ideZ_0=1$ is diagonal and constant; if $\ideZ_{n-1}$ is diagonal with $\thQ$-constant entries, then the off-diagonal entries of $\ideZ_n$ vanish because $\Lv_+-\Lv_-$ is a unit (\cref{lem:ide-ss}), and the next equation says that the diagonal entries of $\ideZ_n$ are $\thQ$-constant. The constants of $\thQ$ on $\ideFz=\ideKt\lser{\cQ}$ are $\ideKt$. So $\ideZ=\mathrm{diag}(\zeta_+(z),\zeta_-(z))$ with $\zeta_i\in\ideKt\fps{z}$, $\zeta_i(0)=1$. Since $\Rm_1^{-1}\Rm_2$ is symplectic and $\Psi$ is an isometry, $\ideZ(z)^T\ideZ(-z)=1$, that is, $\zeta_i(z)\zeta_i(-z)=1$, so $\log\zeta_i$ is odd: $\zeta_i=\exp\sum_ka_k^{(i)}z^{2k-1}$.
\end{proof}

\subsection{Flatness of the R-matrix}

\begin{prop}[Flatness of $\Rm$]\genusstatus{AI-proposed, unchecked}
\label{prop:ide-flat}
The R-matrix $\Rm$ of \cref{prop:ide-teleman} satisfies \cref{eq:ide-flat}. Equivalently, $\ideS\coloneqq\Rm\Psi\re^{\idecan/z}$ satisfies $z\thQ\ideS=H\qprod\ideS$: it is a fundamental solution of the quantum differential equation, unitary in the sense that $\Rm$ is symplectic. In terms of the covectors $\vv_i\coloneqq\eid_i^*\circ\Rm^{-1}(-z)$, $i=\pm$, the equation reads
\begin{equation}
\label{eq:ide-roweq}
z\Dop_i\vv_i=(\Aop-\Lv_i)\vv_i,\qquad \Dop_i=\frac{\dd}{\dd t}-\frac{\Delta_i'}{2\Delta_i},
\end{equation}
where $\Aop$ is the transpose of the matrix of $H\qprod$, acting on covectors by $(\Aop a)(x)=a(H\qprod x)$; in the basis $(\ideone,H)$, $\Aop=\begin{pmatrix}0&1\\-\Lv_+\Lv_-&\Lv_++\Lv_-\end{pmatrix}$ on $a=(a_0,a_1)$ (\cref{prop:set-frob}). The equation is for the $\Rm$ whose legs are $\Rm^{-1}(\psi)$, with the same $z$; it is not an equation for $\Rm^{-1}(-z)$.
\end{prop}

\begin{proof}
\emph{Set-up.} Work over $\ideFz$ with the derivation $\thQ$, using $\Omega=\Rm.\TFT$ (\cref{prop:ide-teleman}; the coefficients of $\Rm$ lie in $\KCY\lser{\cQ}\subset\ideFz$). On $\ideV_0\coloneqq\KCY\ideone\oplus\KCY H$ the derivation $\thQ$ vanishes and $\eta$ takes values in $\KCY$. In this proof state-space elements are denoted $x,x_l\in\ideV_0$, so that $\thQ x=0$. Put $\Rinv(\psi)\coloneqq\Rm^{-1}(\psi)$ and
\[
\idelg{i}(\psi;x)\coloneqq\fid_i^*\big(\Rinv(\psi)x\big)\in\ideFz\fps{\psi}.
\]
Let $M_{g,n}$ denote the moduli space of smooth curves ($\Mbar_{g,n}$ is its compactification). Each of the following nine steps is a statement about curves or about a semisimple Frobenius algebra of any rank, except (Fl4) and (Fl8), where the target enters.

\emph{(Fl1) Edgeless form.} On $M_{g,n}$ every graph with an edge vanishes, being pushed forward from the boundary. The translation $\ideT(z)=z(\ideone-\Rinv(z)\ideone)$ has no $z^0$ and $z^1$ terms, since $\Rinv(z)=1-\Rm_1z+\cdots$; so the translation legs, pushed forward along the forgetful maps, give a series in the classes $\kappa_j$, $j\geq 1$, finite in each degree, with no boundary correction, because the corrections from the divisors where a leg $l$ bubbles off with translation points vanish by the dimension count of (Fl3) (a translation leg on a genus-$0$ bubble with $m$ translation points has degree $\geq 2m>m-1$). Only the one-vertex graph remains, and only its monochromatic terms, because $\TFT$ vanishes on mixed insertions and $\TFT_{g,N}(\fid_i^{\otimes N})=\Delta_i^{(2g-2+N)/2}$. Distributing one factor $\Delta_i^{1/2}$ to each translation leg,
\begin{equation}
\label{eq:ide-Fl1}
\begin{gathered}
\Omega_{g,n}(x_1,\dots,x_n)=\sum_{i=\pm}\ideP_i\,\ideN_i(\kappa)\prod_{l=1}^n\idelg{i}(\psi_l;x_l),\\
\ideP_i\coloneqq\Delta_i^{(2g-2+n)/2},\qquad
\ideN_i(\kappa)\coloneqq\sum_{m\geq 0}\frac{1}{m!}\,p_{m*}\prod_{k=1}^m\Big(\Delta_i^{1/2}\fid_i^*\ideT(\psi_{n+k})\Big),
\end{gathered}
\end{equation}
where $p_m\colon\Mbar_{g,n+m}\to\Mbar_{g,n}$ forgets the last $m$ points; $\ideN_i$ is a series in the $\kappa_j$, $j\geq 1$, with coefficients in $\ideFz$, finite in each degree, with constant term $\ideN_i(0)=1$. Only $\TFT(\fid_i^{\otimes N})=\Delta_i^{(2g-2+N)/2}$ and the vanishing of mixed insertions are used.

\emph{(Fl2) Pushforward.} Let $p\colon\ideCurve\to M_{g,n}$ be the restriction of $\Mbar_{g,n+1}\to\Mbar_{g,n}$ over $M_{g,n}$, with sections $\idesec_l$, $l=1,\dots,n$. On $\ideCurve$: $\psi_l=p^*\psi_l+\idesec_l$, $\psi_l\idesec_l=\psi_{n+1}\idesec_l=0$, hence $\psi_l^k=(p^*\psi_l)^k+\idesec_l(p^*\psi_l)^{k-1}$; $\idesec_l\idesec_{l'}=0$ for $l\neq l'$; $\kappa_j=p^*\kappa_j+\psi_{n+1}^j$; $p_*\psi_{n+1}^k=\kappa_{k-1}$, where $\kappa_0=2g-2+n$ is a number; and $p_*(\idesec_l\,p^*c)=c$. These are statements about curves, with no target.

\emph{(Fl3) Bubble.} The fibres of $p$ are smooth, so a point of $\ideCurve$ has at most one node: graphs with two or more edges miss $\ideCurve$, and the one-edge graphs meeting $\ideCurve$ are the bubbles, a vertex $(0,3)$ carrying the legs $l$ and $n+1$ joined by one edge to the genus-$g$ vertex, supported on $\idesec_l$. A translation leg on the bubble vanishes, since $p_{m*}$ of a class of degree $\geq 2m$ on $\Mbar_{0,3+m}$, of dimension $m$, vanishes for $m\geq 1$; the legs of the bubble carry $\Rinv(0)=1$, since $\psi=0$ on $\Mbar_{0,3}$. The bubble vertex is $\TFT_{0,3}(x_l,H,\cdot)=\eta(H\qprod x_l,\cdot)$, which, contracted with the bivector of the edge, puts $H\qprod x_l$ on the genus-$g$ side; the edge factor at $\psi_h=0$ on the bubble side is $(\eta^{-1}-(1\otimes\Rinv(\psi_{h'}))\eta^{-1})/\psi_{h'}$, so the bubble contributes
\[
\ideP_i\ideN_i\,\fid_i^*\Big(\frac{1-\Rinv(\psi_l)}{\psi_l}\,H\qprod x_l\Big)\prod_{l'\neq l}\idelg{i}(\psi_{l'};x_{l'})
\]
on leg $l$ of the genus-$g$ vertex, after $p_*$. The edgeless graph \cref{eq:ide-Fl1} for $(g,n+1)$ with insertions $(x,H)$, restricted to $\idesec_l$, where $\kappa_j\mapsto p^*\kappa_j$, $\idelg{i}(\psi_{n+1};H)\mapsto\fid_i^*(H)$ and $\Rinv(\psi_l)=\Rinv(p^*\psi_l)+\idesec_l(\Rinv(p^*\psi_l)-1)/p^*\psi_l$, contributes after $p_*$
\[
\Delta_i^{(2g-1+n)/2}\ideN_i\,\fid_i^*(H)\,\fid_i^*\Big(\frac{\Rinv(\psi_l)-1}{\psi_l}x_l\Big)\prod_{l'\neq l}\idelg{i}(\psi_{l'};x_{l'}).
\]
The identities $H=\sum_i\Lv_i\eid_i$, that is, $\fid_i^*(H)=\Lv_i\Delta_i^{-1/2}$, and $\fid_i^*(a\qprod b)=\Delta_i^{1/2}\fid_i^*(a)\fid_i^*(b)$ hold in any semisimple Frobenius algebra; so $\Delta_i^{1/2}\fid_i^*(H)\fid_i^*(y)=\Lv_i\fid_i^*(y)=\fid_i^*(H\qprod y)$, and the two terms add up to $\ideP_i\ideN_i\fid_i^*\big(\frac{1}{\psi_l}[H\qprod,\Rinv(\psi_l)]x_l\big)$ on leg $l$. The part of the edgeless graph not supported on the sections gives, after $p_*$, a series in the $\kappa_j$ times $\prod_l\idelg{i}(\psi_l;x_l)$. Hence
\begin{equation}
\label{eq:ide-Fl3}
p_*\Omega_{g,n+1}(x,H)=\sum_i\ideP_i\ideN_i\sum_{l=1}^n\fid_i^*\Big(\frac{1}{\psi_l}[H\qprod,\Rinv(\psi_l)]\,x_l\Big)\prod_{l'\neq l}\idelg{i}(\psi_{l'};x_{l'})+\sum_i\ideG_i(\kappa)\prod_l\idelg{i}(\psi_l;x_l),
\end{equation}
with $\ideG_i$ a series in the $\kappa_j$, $j\geq 1$ (and the number $\kappa_0$), finite in each degree, containing no $\psi_l$.

\emph{(Fl4) Divisor equation} (the first of the two places where the target enters). For $x\in\ideV_0^{\otimes n}$ and every $d\geq 0$,
\[
p_*\Omega_{g,n+1,d}(x,H)=d\,\Omega_{g,n,d}(x),\qquad\text{hence}\qquad p_*\Omega_{g,n+1}(x,H)=\thQ\Omega_{g,n}(x).
\]
Indeed $\mathrm{ft}_*(\ev_{n+1}^*H)=d$, the fibre degree of $f^*\cO(1)$, an equivariant class of degree $0$, for the equivariant lift of $H$ used here as for any other; the virtual class is pulled back (fundamental-class axiom) and so is the twisting class \cite[App.~1~(iv)]{CG07}. At $d=0$ both sides vanish. So the left side of \cref{eq:ide-Fl3} is $\thQ\Omega_{g,n}(x)$, with $\thQ$ acting coefficientwise on \cref{eq:ide-Fl1}.

\emph{(Fl5) Defect.} Since $x\mapsto(\idelg{+}(\psi;x),\idelg{-}(\psi;x))=\Psi^{-1}\Rinv(\psi)x$ is invertible over $\ideFz\fps{\psi}$, there is a unique $\Mdef(\psi)\in\mathrm{Mat}_2(\ideFz\fps{\psi})$ with
\[
\thQ\idelg{i}(\psi;x)-\fid_i^*\Big(\frac{1}{\psi}[H\qprod,\Rinv(\psi)]\,x\Big)=\sum_k\Mdef(\psi)^i{}_k\,\idelg{k}(\psi;x)\qquad\text{for all }x\in\ideV_0 .
\]
$\Mdef$ is built from $\Rm$, $\Psi$ and $H\qprod$ alone, so it does not depend on $g$. Apply $\thQ$ to \cref{eq:ide-Fl1} at $n=2$, where $\ideP_i=\Delta_i^g$, and compare with \cref{eq:ide-Fl3} through (Fl4): with $\mathrm{a}_l\coloneqq(\idelg{\pm}(\psi_l;x_l))$,
\begin{equation}
\label{eq:ide-Fl5}
\sum_i\ideP_i\ideN_i\Big[(\Mdef(\psi_1)\mathrm{a}_1)^i\,\mathrm{a}_2^i+\mathrm{a}_1^i\,(\Mdef(\psi_2)\mathrm{a}_2)^i\Big]=\sum_i\widetilde{\ideG}_i(\kappa)\,\mathrm{a}_1^i\mathrm{a}_2^i,\qquad
\widetilde{\ideG}_i\coloneqq\ideG_i-\thQ(\ideP_i\ideN_i),
\end{equation}
an identity in $H^*(M_{g,2})\otimes\ideFz$, bilinear in $(\mathrm{a}_1,\mathrm{a}_2)$ over $\ideFz\fps{\psi_1,\psi_2}$; since $x\mapsto\mathrm{a}(\psi;x)$ is invertible over $\ideFz$, it holds for arbitrary vectors $\mathrm{a}_1\in\ideFz\fps{\psi_1}^2$, $\mathrm{a}_2\in\ideFz\fps{\psi_2}^2$.

\emph{(Fl6) Stability.} By \cite[2.17--2.19]{Tel12}, resting on \cite{Loo96}, \cite{MW07}, \cite{Har85} and \cite{Iva90} as reported there, in cohomological degree $<(g-1)/2$ (with $\deg\psi_l=2$, $\deg\kappa_j=2j$) the monomials in $\psi_1,\psi_2$ and in the classes $\kappa_j^{M}$, $j\geq 1$, pulled back from $M_g$, are linearly independent in $H^*(M_{g,2};\bbQ)$. The $\kappa_j$ used here are the Arbarello--Cornalba classes, and on $M_{g,2}$ they satisfy $\kappa_j=\kappa_j^M+\psi_1^j+\psi_2^j$ \cite[(2.15)]{Tel12}, a triangular degree-preserving change of generators; so the monomials in $\kappa_j$ ($j\geq 1$), $\psi_1,\psi_2$ are independent in the same range, and over $\ideFz$ by base change, $H^*(M_{g,2};\ideFz)=H^*(M_{g,2};\bbQ)\otimes_{\bbQ}\ideFz$ with $\thQ$ acting on the second factor. Only independence in some range growing with $g$ is used. Take in \cref{eq:ide-Fl5} the coefficient of $\kappa^0\psi_2^0$, where $\ideN_i=1$, $\mathrm{a}_2$ is a constant vector and $\mathrm{a}_1\in\ideFz\fps{\psi_1}^2$ is arbitrary: with $\ideP\coloneqq\mathrm{diag}(\Delta_+^g,\Delta_-^g)$, as bilinear forms in $(\mathrm{a}_1,\mathrm{a}_2)$,
\[
\ideP\,\Mdef(\psi)+\Mdef_0^T\ideP=\widetilde{\ideG}(0)\quad\text{diagonal, modulo degree}\ \geq(g-1)/2\ \text{in}\ \psi,
\]
$\Mdef_0\coloneqq\Mdef(0)$. Since $\Mdef$ does not depend on $g$, the coefficient of $\psi^n$ for $n\geq 1$ gives $\ideP\Mdef_n=0$ once $g>4n+1$ (then $\psi_1^n$ has cohomological degree $2n<(g-1)/2$), and $\ideP$ is invertible; so $\Mdef_n=0$ for all $n\geq 1$. At $\psi^0$ the off-diagonal entry reads
\[
\Delta_+^g(\Mdef_0)_{+-}+\Delta_-^g(\Mdef_0)_{-+}=0\qquad\text{for every }g\geq 2 .
\]

\emph{(Fl7) Skewness of $\Mdef_0$.} With $\ideW\coloneqq\Psi^{-1}\thQ\Psi$ and $\ideRf\coloneqq\Psi^{-1}\Rm\Psi$, one has $\Mdef_0=-\ideW+[\ideLam,\ideRf_1]$. Indeed $\mathrm{a}(\psi;x)=\Psi^{-1}\Rinv(\psi)x$, so $\thQ\mathrm{a}=-\ideW\mathrm{a}+\Psi^{-1}(\thQ\Rinv)x$, and
\[
\Psi^{-1}\frac1\psi[H\qprod,\Rinv]x=\frac1\psi\big(\ideLam\mathrm{a}-\Psi^{-1}\Rinv\Psi\,\ideLam\,(\Psi^{-1}\Rinv\Psi)^{-1}\mathrm{a}\big);
\]
at $\psi^0$, using $\Psi^{-1}\Rinv(\psi)\Psi=1-\ideRf_1\psi+\cdots$ and $\thQ\Rinv=O(\psi)$, this gives $\Mdef_0=-\ideW-\big(-[\ideLam,\ideRf_1]\big)$. $\ideW$ is skew, because $\Psi^T\eta\Psi=1$ and $\thQ\eta=0$ give $(\Psi^{-1}\thQ\Psi)^T+\Psi^{-1}\thQ\Psi=0$. $\ideRf_1$ is symmetric: symplecticity gives $\Rm_1=\Rm_1^\dagger$ at order $z$, and in the $\eta$-orthonormal frame $\fid_i$ the adjoint is the transpose; its commutator with the diagonal $\ideLam$ is skew. No rank enters. So $(\Mdef_0)_{-+}=-(\Mdef_0)_{+-}$, and the off-diagonal relation of (Fl6) becomes $(\Delta_+^g-\Delta_-^g)(\Mdef_0)_{+-}=0$.

\emph{(Fl8) Non-degeneracy} (the second place where the target enters). By \cref{lem:ide-ss}, $\Delta_+(0)/\Delta_-(0)=-\alpha_1\alpha_2/(\wA\wB)$ is a non-constant element of $\KCY$. Since $\bbQ$ is algebraically closed in the purely transcendental extension $\KCY=\bbQ(\alpha_0,\alpha_1,\alpha_2,\epsilon_1)$, this element is transcendental over $\bbQ$, hence not a root of unity. So $\Delta_+^g-\Delta_-^g$ has the non-zero value $\Delta_-(0)^g\big((\Delta_+(0)/\Delta_-(0))^g-1\big)$ at $\cQ=0$, for every $g\geq 1$, and $(\Mdef_0)_{+-}=0$; with skewness and the vanishing of the diagonal part of a skew matrix, $\Mdef_0=0$. Hence $\Mdef=0$:
\[
\thQ\big(\Psi^{-1}\Rinv(\psi)\big)=\Psi^{-1}\frac1\psi\big[H\qprod,\Rinv(\psi)\big]
\]
as operators, since both sides agree on $\ideV_0$ and are $\ideFz$-linear.

\emph{(Fl9) The equation.} Take $\eta$-adjoints, with $\Rinv(\psi)^\dagger=\Rm(-\psi)$ (the symplectic condition), $\Psi^\dagger=\Psi^{-1}$ ($\Psi$ is an isometry), $(H\qprod)^\dagger=H\qprod$ (\cref{prop:set-frob}) and $\thQ\eta=0$, so that $\thQ$ commutes with $\dagger$:
\[
\thQ\big(\Rm(-\psi)\Psi\big)=\frac1\psi\big(\Rm(-\psi)\,H\qprod-H\qprod\Rm(-\psi)\big)\Psi .
\]
Putting $z=-\psi$ gives \cref{eq:ide-flat}. With $H\qprod\Psi=\Psi\ideLam$ and $\thQ\idecan=\ideLam$, \cref{eq:ide-flat} is equivalent to $z\thQ\ideS=H\qprod\ideS$ for $\ideS=\Rm\Psi\re^{\idecan/z}$, the form \cite[Section~1.3, (3)--(5)]{Giv01} of the system. The $\Rm$ here is the one of the legs $\Rm^{-1}(\psi_l)$ of \cite[Def.~2.2]{PPZ15}, so no comparison of conventions is needed. For the rows: by symplecticity and $\eid_i^*=\Delta_i\eta(\eid_i,\cdot)$,
\[
\vv_i(x)=\eid_i^*\big(\Rm^{-1}(-z)x\big)=\Delta_i\,\eta\big(\Rm^{-1}(-z)^\dagger\eid_i,x\big)=\Delta_i^{1/2}\,\eta\big(\Rm(z)\fid_i,x\big),
\]
so $\Delta_i^{-1/2}\vv_i$ is the $i$-th column of $\Rm(z)\Psi$ lowered by the constant pairing $\eta$. The $i$-th column of \cref{eq:ide-flat} is $z\thQ(\Rm\Psi\mathbf{e}_i)=H\qprod(\Rm\Psi\mathbf{e}_i)-\Lv_i\Rm\Psi\mathbf{e}_i$; lowering by $\eta$ turns $H\qprod$ into $\Aop$, and $z\thQ(\Delta_i^{-1/2}\vv_i)=\Delta_i^{-1/2}\big(z\vv_i'-z\frac{\Delta_i'}{2\Delta_i}\vv_i\big)=\Delta_i^{-1/2}z\Dop_i\vv_i$. This is \cref{eq:ide-roweq}, with the same $z$.
\end{proof}

Beyond the semisimple Frobenius data ($\Lv_+-\Lv_-$ a unit, for the $\eid_i$, $\Psi$ and the ambiguity lemma; $\Delta_i(0)\neq 0$, for $\fid_i$ and the invertibility of $\ideP$), the target enters only through (Fl4) and (Fl8).

\begin{rmk}[The route through Teleman's Proposition~7.7; not used]
\label{rmk:ide-tel77}
Let $\mathsf{sh}\coloneqq\exp(s\thQ)\colon\KCY\lser{\cQ}\to\KCY\lser{\cQ}\fps{s}$ be the ring map $\cQ\mapsto\cQ\re^{s}$. By (Fl4), $\sum_m\frac{s^m}{m!}p^m_*\Omega_{g,n+m}(x,H^m)=\mathsf{sh}(\Omega_{g,n}(x))$. The definition \cite[Def.~7.1]{Tel12} carries the factor $(-1)^m$, so, applying the embedding of \cref{prop:ide-teleman} coefficientwise in $s$, $\mathsf{sh}(\Omega)$ is Teleman's family at $u=-sH$, not at $u=sH$ (against the divisor equation of (Fl4), the one printed in \cite[Example~7.17]{Tel12} has the opposite sign; that printed sign is not relied on). Applying $\mathsf{sh}$ to coefficients gives $\mathsf{sh}(\Rm).\mathsf{sh}(\TFT)=\mathsf{sh}(\Omega)$, so $\mathsf{sh}(\Rm)$ is the datum of this family, unique by \cref{lem:ide-free} over $\bbC\lser{s}$. Equation \cite[(7.7.a)]{Tel12} with $v=-H$ gives at $s=0$ the equation \cref{eq:ide-flat}, and \cite[Prop.~8.2(ii)]{Tel12} along $v=-H$ is $z\thQ\ideS=H\qprod\ideS$. The legs of \cite[(7.6)]{Tel12} carry $\mathbf{E}^{-1}(\psi_i)$ and those of \cite[Def.~2.2]{PPZ15} carry $\Rm^{-1}(\psi_i)$, so in this comparison Teleman's $\mathbf{E}$ is the $\Rm$ of \cite{PPZ15}, with the same $z$. No step of the proof above uses \cite[Prop.~7.7, Prop.~8.2(ii) or Example~7.17]{Tel12}.
\end{rmk}

\subsection{The rows}

We now construct a second flat symplectic solution from the normalised stationary-phase expansions of \cref{sec:row}. Recall from \cref{def:set-rows} that
\[
\rpt_{0,i}(z;t)=\sum_mz^m\wick_m\big(\cW_2(U_i),\cW_3(U_i),\dots\big)
\]
is the normalised expansion of $\re^{\cW/z}\dd U/U$ at the root $U_i$ of the critical equation $\Pq(U)=0$ of \eqref{eq:set-P} (\cref{lem:set-crit}), where
\[
\cW(U)=\alpha_1\log(1-\cQ U)-\alpha_2\log U+\wA\log(1-U),\quad\cW_k=\thU^k\cW,\quad
\Lv(U)=-\frac{\alpha_1\cQ U}{1-\cQ U}=\frac{\de\cW}{\de t}\Big|_U ,
\]
and the eigenvalues are $\Lv_i=\Lv(U_i)$. At $\cQ=0$ the two roots are labelled by $U_+\to-\alpha_2/\alpha_0$ and $\ideQU\coloneqq\cQ U_-\to\alpha_0/\wB$ (\cref{lem:set-limits}): indeed $\Pq(U)|_{\cQ=0}=-\alpha_0U-\alpha_2$ has the simple root $-\alpha_2/\alpha_0$, and in terms of $\ideQU=\cQ U$ the equation $\cQ\Pq(U)=\wB\ideQU^2-(\alpha_0+\cQ(\alpha_1-\alpha_2))\ideQU-\alpha_2\cQ=0$ reduces at $\cQ=0$ to $\ideQU(\wB\ideQU-\alpha_0)=0$ with the simple root $\alpha_0/\wB$. By Hensel's lemma $U_+$ and $\ideQU$ lie in $\KCY\fps{\cQ}$ (both roots are simple because $\alpha_0\neq 0$). This labelling of the two roots is the one in force in this subsection; the labelling $U_1\to 1$, $U_2\to-\alpha_2/\wB$ at $\cQ\to 1$ of \cref{lem:set-crit} is a different one, and no matching between the two is needed, because the colour sum is symmetric under the exchange of the two roots (\cref{lem:ide-terms}).

\begin{lem}[Normalisation constants at $\cQ=0$]\genusstatus{AI-proposed, unchecked}
\label{lem:ide-norm}
The $z$-coefficients of $\rpt_{0,\pm}$ lie in $\KCY\fps{\cQ}$, and
\[
\rpt_{0,+}(z;\cQ=0)=h_{(\alpha_2,\,-\wA,\,\alpha_0)}(z),\qquad
\rpt_{0,-}(z;\cQ=0)=h_{(-\alpha_0,\,-\alpha_1,\,\wB)}(z).
\]
Hence
\[
\frac{h_{\wts_+}}{\rpt_{0,+}(z;0)}=\frac{h_{\wts_-}}{\rpt_{0,-}(z;0)}=\hnw,\qquad\nw=(\alpha_1,\wA,-\epsilon_1,-\epsilon_2),
\]
the same scalar at both colours.
\end{lem}
\begin{proof}
At $U_+$: for $k\geq 1$, $\thU^k\log(1-\cQ U)$ is a polynomial without constant term in $\cQ U/(1-\cQ U)$, which lies in $\cQ\KCY\fps{\cQ}$ at $U=U_+$; $\thU\log U=1$ and $\thU^k\log U=0$ for $k\geq 2$; and $\thU^k\log(1-U)$ is a polynomial in $(1-U)^{-1}$, with $(1-U_+)^{-1}\in\KCY\fps{\cQ}$ of value $\alpha_0/\wA$. So every $\cW_k(U_+)$, $k\geq 2$, lies in $\KCY\fps{\cQ}$ and its value at $\cQ=0$ is the $k$-th derivative in $y=\log U$ of $f=a\log U+b\log(1-U)$, $(a,b)=(-\alpha_2,\wA)$, at $U_+(0)=-\alpha_2/\alpha_0=a/(a+b)$, which is the critical point of $f$ (here $a\neq 0$, $b\neq 0$, $a+b=\alpha_0\neq 0$). At $U_-$: write $\cW(U)=\alpha_1\log(1-\ideQU)-\alpha_2\log\ideQU+\alpha_2\log\cQ+\wA\log(1-\ideQU/\cQ)$ with $\ideQU=\cQ U$, and $\thU=\theta_{\ideQU}$; at $\ideQU=\ideQU(\cQ)\in\KCY\fps{\cQ}$ with $\ideQU(0)=\alpha_0/\wB\neq 0$, one has $\theta_{\ideQU}\log(1-\ideQU/\cQ)=\ideQU/(\ideQU-\cQ)\in\KCY\fps{\cQ}$ with value $1$, and $\theta_{\ideQU}^{k}\log(1-\ideQU/\cQ)=\theta_{\ideQU}^{k-1}\big(\ideQU/(\ideQU-\cQ)\big)\in\cQ\KCY\fps{\cQ}$ for $k\geq 2$; further $(1-\ideQU)^{-1}\in\KCY\fps{\cQ}$ with value $\wB/\alpha_1$. So $\cW_k(U_-)$ lies in $\KCY\fps{\cQ}$ for $k\geq 2$, and its value at $\cQ=0$ is the $k$-th derivative in $\log\ideQU$ of $\alpha_1\log(1-\ideQU)$, which for $k\geq 2$ is also that of $f=a\log\ideQU+b\log(1-\ideQU)$, $(a,b)=(\alpha_0,\alpha_1)$, at $\ideQU(0)=\alpha_0/\wB=a/(a+b)$, the critical point of $f$ (the critical equation $\cW_1(U_-)=0$ reads, at $\cQ=0$, $\alpha_0+\alpha_1\theta_{\ideQU}\log(1-\ideQU)=0$; here $a,b,a+b=\wB$ are non-zero). In both cases $\cW_2\to-a(a+b)/b\neq 0$ (for $f(y)=ay+b\log(1-\re^y)$ one has $f''=-b\re^y/(1-\re^y)^2$, and at $\re^y=a/(a+b)$ this is $-a(a+b)/b$). By the Wick form (D1) of \cref{def:row-Lnorm} and the characterisation of $\wick_m$ as a polynomial in $\cW_2^{-1},\cW_3,\dots$ (\cref{def:set-rows}), each $z$-coefficient of $\rpt_{0,\pm}$ lies in $\KCY\fps{\cQ}$ and its value at $\cQ=0$ is the corresponding coefficient of $\Lnorm[f,1;y_0]$ in the variable $y=\log U$ at the critical point $y_0=\log(a/(a+b))$. By \cref{lem:row-betaphase}, applied with $(a,b)=(-\alpha_2,\wA)$ and with $(a,b)=(\alpha_0,\alpha_1)$ (in both cases $ab(a+b)\neq 0$ in $\KCY$), this limit is $h_{(-a,-b,a+b)}$, which gives the two displayed values. Finally, as multisets, $\wts_+=(\alpha_0,\alpha_1,\alpha_2,-\epsilon_1,-\epsilon_2)=\nw+(\alpha_2,-\wA,\alpha_0)$ and $\wts_-=(-\alpha_0,\wB,\wA,-\epsilon_1,-\epsilon_2)=\nw+(-\alpha_0,-\alpha_1,\wB)$, so $h_{\wts_+}=\hnw h_{(\alpha_2,-\wA,\alpha_0)}$ and $h_{\wts_-}=\hnw h_{(-\alpha_0,-\alpha_1,\wB)}$ by $h_{\mathsf{u}+\mathsf{v}}=h_{\mathsf{u}}h_{\mathsf{v}}$ (\cref{def:set-rows}).
\end{proof}

Define the rows with the constant,
\[
\rZ_{0,i}\coloneqq\hnw\,\rpt_{0,i},\qquad \rZ_{1,i}\coloneqq(\Lv_i+z\Dop_i)\,\rZ_{0,i},\qquad i=\pm,
\]
and let $\RZ$ be the endomorphism with $\eid_i^*\circ\RZ^{-1}(-z)=(\rZ_{0,i},\rZ_{1,i})$, $i=\pm$: it exists and is unique since $\eid_\pm^*$ is a basis of the dual, and $\RZ(0)=1$ because $\rZ_{0,i}[0]=1$, $\rZ_{1,i}[0]=\Lv_i$, that is, $(\rZ_{0,i},\rZ_{1,i})|_{z=0}=(1,\Lv_i)=\chiv_i=\eid_i^*$.

\begin{prop}[The rows form a flat unitary solution]\genusstatus{AI-proposed, unchecked}
\label{prop:ide-rows}
\begin{enumerate}
\item (Frobenius data.) $\Lv_i=\Lv(U_i)=\thQ u_i$ for $u_i\coloneqq\cW(U_i)$; $U_i/(1-U_i)=(\Lv_i-\alpha_2)/\wA$; and \cref{eq:ide-Delta} holds.
\item (Row~2 of the quantum differential equation.) With $\dd/\dd t=\thQ$ at fixed $U$ and
\begin{equation}
\label{eq:ide-PF}
\idePF\coloneqq z^2(\cQ-1)\Big(\frac{\dd}{\dd t}\Big)^2-z\big(\alpha_0+(\alpha_1+\alpha_2)\cQ\big)\frac{\dd}{\dd t}+\alpha_1\alpha_2\cQ
=(\cQ-1)\Big[\Big(z\frac{\dd}{\dd t}\Big)^2-(\Lv_++\Lv_-)\,z\frac{\dd}{\dd t}+\Lv_+\Lv_-\Big],
\end{equation}
Euler's integration by parts is the identity in $U$ and $\cQ$
\begin{equation}
\label{eq:ide-euler}
\idePF\,\re^{\cW/z}=\thU\big(\idekm\,\re^{\cW/z}\big),\qquad\idekm\coloneqq-z\alpha_1\cQ\,\frac{1-U}{1-\cQ U},
\end{equation}
and consequently $z\Dop_i(\rZ_{0,i},\rZ_{1,i})=(\Aop-\Lv_i)(\rZ_{0,i},\rZ_{1,i})$, as an identity in $\KCY(\Ccov)\fps{z}$ (\cref{lem:ide-terms}) and in $\KCY\fps{\cQ}\fps{z}$.
\item ($\cQ=0$.) The rows lie in $\KCY\fps{\cQ}\fps{z}^2$, and $(\rZ_{0,i},\rZ_{1,i})|_{\cQ=0}=h_{\wts_i}(z)\,(1,\Lv_i(0))$: that is, $\RZ|_{\cQ=0}=\mathrm{diag}(h_{\wts_+},h_{\wts_-})$ on the idempotents $\eid_\pm(0)$.
\item (Symplecticity.) $\langle\vv^Z_+(z),\vv^Z_-(-z)\rangle=0$ and $\langle\vv^Z_i(z),\vv^Z_i(-z)\rangle=\Delta_i$, where $\vv^Z_i=(\rZ_{0,i},\rZ_{1,i})$ and $\langle\ ,\ \rangle$ is the pairing of covectors by $\eta^{-1}$. Hence $\RZ$ is symplectic, $\RZ\Psi\re^{\idecan/z}$ is a unitary fundamental solution of the quantum differential equation, and
\[
\Rm=\RZ\,\Psi\exp\Big(\mathrm{diag}\sum_{k\geq 1}a_k^{(i)}z^{2k-1}\Big)\Psi^{-1}\qquad\text{with constants }a_k^{(i)}\in\ideKt .
\]
\end{enumerate}
\end{prop}
\begin{proof}
(i) $\de_t\cW|_U=\alpha_1\thQ\log(1-\cQ U)=-\alpha_1\cQ U/(1-\cQ U)=\Lv(U)$, and at a root $U_i$ of $\Pq$, which is the critical equation $\thU\cW=0$ (\cref{lem:set-crit}), $\thQ u_i=\de_t\cW|_{U_i}+\cW_1(U_i)\thQ\log U_i=\Lv(U_i)$. The critical equation is $\cW_1=\Lv(U)-\alpha_2-\wA U/(1-U)=0$, which gives $U_i/(1-U_i)=(\Lv_i-\alpha_2)/\wA$. That $\Lv(U_\pm)$ are the eigenvalues $\Lv_\pm$ of $H\qprod$ is \cref{lem:set-eig}(ii); directly, $\Lv(U)=(\wB U+\alpha_2)/\tau$ modulo $\Pq$ (\cref{lem:set-crit}(iii)) and $\tau^2\cQ\big[\Lv(U)^2-(\Lv_++\Lv_-)\Lv(U)+\Lv_+\Lv_-\big]=\wB\,\Pq(U)$ identically in $U$, with the sum and product of \cref{lem:ide-g0}. For \cref{eq:ide-Delta}: $\eta(\eid_i,\eid_i)=\eta(\eid_i,\ideone)=\big(\eta(H,\ideone)-\Lv_j\eta(\ideone,\ideone)\big)/(\Lv_i-\Lv_j)=(\alpha_1\alpha_2-E\Lv_j)/\big(\alpha_1\alpha_2\wA\wB\epsilon_1\epsilon_2(\Lv_i-\Lv_j)\big)$ by the values $\eta(\ideone,\ideone)=E/(\alpha_1\alpha_2\wA\wB\epsilon_1\epsilon_2)$, $\eta(\ideone,H)=1/(\wA\wB\epsilon_1\epsilon_2)$ of \cref{lem:ide-fixed}(iii); the identity $(\alpha_1\alpha_2-E\Lv_+)(\alpha_1\alpha_2-E\Lv_-)=\alpha_1^2\alpha_2^2-E\alpha_1\alpha_2(\Lv_++\Lv_-)+E^2\Lv_+\Lv_-=\alpha_1\alpha_2(\alpha_1\alpha_2+\alpha_0E)=\alpha_1\alpha_2\wA\wB$ turns the inverse into $\epsilon_1\epsilon_2(\Lv_i-\Lv_j)(\alpha_1\alpha_2-E\Lv_i)$, and $\Lv_j=E/\tau-\alpha_0-\Lv_i$, $\Lv_i\Lv_j=\alpha_1\alpha_2/\tau$ give the first form. For the third form, $\cW_2=\thU\cW_1=-\alpha_1\cQ U/(1-\cQ U)^2-\wA U/(1-U)^2$; with $-\alpha_1\cQ U/(1-\cQ U)^2=\Lv(1+\cQ U/(1-\cQ U))=\Lv-\Lv^2/\alpha_1$ and $-\wA U/(1-U)^2=-\wA\,\frac{U}{1-U}\big(1+\frac{U}{1-U}\big)=-(\Lv-\alpha_2)(\Lv+\alpha_0)/\wA$ at $U_i$, one gets $\alpha_1\wA\cW_2(U_i)=\alpha_1\wA\Lv_i-\wA\Lv_i^2-\alpha_1(\Lv_i-\alpha_2)(\Lv_i+\alpha_0)=-E\Lv_i^2+2\alpha_1\alpha_2\Lv_i+\alpha_0\alpha_1\alpha_2$, as $\wA+\alpha_1=E$ and $\wA-\alpha_0+\alpha_2=2\alpha_2$. \draftnote{Compare BS, Appendix A, for the Frobenius data of the conifold at generic weights.}

(ii) \emph{The identity \cref{eq:ide-euler}.} Divide both sides by $\re^{\cW/z}$. On the left, with $\de_t\cW=\Lv(U)$ and $\de_t^2\cW=\thQ\Lv(U)=-\alpha_1\cQ U/(1-\cQ U)^2=\Lv(U)/(1-\cQ U)$,
\[
\re^{-\cW/z}\idePF\re^{\cW/z}=(\cQ-1)\Big[\Lv^2+\frac{z\Lv}{1-\cQ U}\Big]-\big(\alpha_0+(\alpha_1+\alpha_2)\cQ\big)\Lv+\alpha_1\alpha_2\cQ .
\]
On the right, $\re^{-\cW/z}\thU(\idekm\re^{\cW/z})=\thU\idekm+\idekm\cW_1/z$, with
\[
\thU\idekm=-z\alpha_1\cQ U\,\frac{-(1-\cQ U)+\cQ(1-U)}{(1-\cQ U)^2}=-z\alpha_1\cQ U\,\frac{\cQ-1}{(1-\cQ U)^2}=(\cQ-1)\frac{z\Lv}{1-\cQ U},
\]
and, writing $\alpha_1\cQ/(1-\cQ U)=-\Lv/U$,
\[
\frac{\idekm\cW_1}{z}=-\alpha_1\cQ\frac{1-U}{1-\cQ U}\Big[\Lv-\alpha_2-\wA\frac{U}{1-U}\Big]=\Lv^2\frac{1-U}{U}-\alpha_2\Lv\frac{1-U}{U}-\wA\Lv .
\]
The $z$-terms agree, and the difference of the $z^0$-terms, multiplied by $U$, is
\begin{multline*}
\Lv^2\big((\cQ-1)U-(1-U)\big)+\Lv\big(-(\alpha_0+(\alpha_1+\alpha_2)\cQ)U+\alpha_2(1-U)+\wA U\big)+\alpha_1\alpha_2\cQ U\\
=-(1-\cQ U)\Lv^2+\alpha_2(1-\cQ U)\Lv-\alpha_1\cQ U\Lv+\alpha_1\alpha_2\cQ U,
\end{multline*}
using $-\alpha_0U-\alpha_2U+\wA U=0$. This equals $(\alpha_2-\Lv)\big((1-\cQ U)\Lv+\alpha_1\cQ U\big)=0$, by the definition of $\Lv(U)$. So \cref{eq:ide-euler} holds identically in $U$ and $\cQ$, not only at the critical points.

\emph{Row~2.} Fix complex weights in (ND) and a complex $\cQ_0\notin\{0,1,\infty\}$ with $\Disc(\cQ_0)\neq 0$. Then the two roots are distinct, hence simple, and lie outside $\{0,1,1/\cQ\}$, by $U_+U_-=-\alpha_2/(\wB\cQ)$ and $\Pq(1)=\wA(\cQ-1)$, $\Pq(1/\cQ)=\alpha_1(1-\cQ)/\cQ$ (\cref{lem:set-crit}); each $U_i(\cQ)$ is analytic near $\cQ_0$, $y_i\coloneqq\log U_i$ is defined, and $\cW^y(y;t)\coloneqq\cW(\re^y)$ is analytic near $(y_i,t_0)$ with $\de_y^k\cW^y(y_i)=\cW_k(U_i)$, $\de_y\cW^y(y_i)=0$ and $\de_y^2\cW^y(y_i)=\cW_2(U_i)\neq 0$ (by \cref{eq:ide-Delta}, $\cW_2(U_i)=0$ would force $\Lv_i=\Lv_j$ or $\alpha_1\alpha_2-E\Lv_i=0$, and the latter is excluded by $(\alpha_1\alpha_2-E\Lv_+)(\alpha_1\alpha_2-E\Lv_-)=\alpha_1\alpha_2\wA\wB\neq 0$). In the coordinate $y$ the integrand is $\re^{\cW^y/z}\dd y$ with amplitude $1$, so $\rpt_{0,i}=\Lnorm[\cW^y,1;y_i]$ by (D1) of \cref{def:row-Lnorm} and the Wick form of $\rpt_{0,i}$ (this is step~(1) of the proof of \cref{prop:row-wickrow}). Now \cref{prop:row-P2} (the parameter derivative with a moving expansion point, with the parameter $t$ in place of $\tau$, the coordinate $y$ in place of $s$ and the critical point $y_i(t)$) gives, for every amplitude $q(t,y)$ which is a Laurent polynomial in $z$ with coefficients analytic near $(t_0,y_i(t_0))$,
\[
\Big(\frac{\dd}{\dd t}+\frac{\thQ u_i}{z}-\frac{\thQ\cW_2(U_i)}{2\cW_2(U_i)}\Big)\Lnorm[\cW^y,q;y_i]=\Lnorm\Big[\cW^y,\ \de_tq+\frac{q\,\de_t\cW^y}{z};y_i\Big],
\]
where $u_i=\cW^y(y_i)$ is the critical value and $\de_t\cW^y=\Lv(\re^y)$. By (i) and \cref{eq:ide-Delta}, $\thQ u_i=\Lv_i$ and $\thQ\cW_2(U_i)/\cW_2(U_i)=\Delta_i'/\Delta_i$, so the operator on the left is $z^{-1}(\Lv_i+z\Dop_i)$. Applying it twice to $\rpt_{0,i}$ (first with $q=1$, which gives the amplitude $\Lv(\re^y)/z$, then with $q=\Lv(\re^y)/z$, which gives $\thQ\Lv(\re^y)/z+\Lv(\re^y)^2/z^2$), and using the linearity (D3) of $\Lnorm$ over functions independent of $y$,
\begin{align*}
\Big[(\Lv_i+z\Dop_i)^2-(\Lv_++\Lv_-)(\Lv_i+z\Dop_i)+\Lv_+\Lv_-\Big]\rpt_{0,i}
&=\Lnorm\big[\cW^y,\ z\thQ\Lv+\Lv^2-(\Lv_++\Lv_-)\Lv+\Lv_+\Lv_-;\,y_i\big]\\
&=\frac{1}{\cQ-1}\Lnorm\big[\cW^y,\ \re^{-\cW/z}\idePF\re^{\cW/z};\,y_i\big],
\end{align*}
by \cref{eq:ide-PF}. By \cref{eq:ide-euler} the amplitude is $\thU\idekm+\idekm\cW_1/z=\de_y\idekm(\re^y)+\idekm(\re^y)\,\de_y\cW^y/z$, the amplitude of the total derivative $\de_y(\re^{\cW^y/z}\idekm(\re^y))$, with $\idekm(\re^y)$ analytic near $y_i$ because $1-\cQ U_i\neq 0$. By \cref{prop:row-P3} its normalised expansion vanishes. Since $\hnw$ does not depend on $t$, the same identity holds for $\rZ_{0,i}=\hnw\rpt_{0,i}$, and with $\rZ_{1,i}=(\Lv_i+z\Dop_i)\rZ_{0,i}$ it reads
\[
z\Dop_i\rZ_{1,i}=\Lv_j\big(\rZ_{1,i}-\Lv_i\rZ_{0,i}\big),\qquad j\neq i,
\]
which, together with $z\Dop_i\rZ_{0,i}=\rZ_{1,i}-\Lv_i\rZ_{0,i}$ (the definition of $\rZ_{1,i}$), is $z\Dop_i\vv^Z_i=(\Aop-\Lv_i)\vv^Z_i$ componentwise: $(\Aop a)_0=a_1$ and $(\Aop a)_1=-\Lv_+\Lv_-a_0+(\Lv_++\Lv_-)a_1$, so $((\Aop-\Lv_i)\vv_i^Z)_1=-\Lv_i\Lv_j\rZ_{0,i}+\Lv_j\rZ_{1,i}$. The $z$-coefficients of both sides are, for indeterminate weights, elements of $\KCY(\Ccov)$ (\cref{lem:ide-terms}, whose proof uses only the definitions), and they agree as analytic functions of $\cQ$ near every $\cQ_0$ as above at every complex weight point of (ND); by the principle (P) of \cref{lem:set-rings}, in the form: an element of $\KCY(\Ccov)$ vanishes if its evaluations at all complex weights with (ND) and all complex $\cQ$ in an open set vanish, the identity holds in $\KCY(\Ccov)\fps{z}$, hence, expanding at either point of $\Ccov$ over $\cQ=0$, in $\KCY\fps{\cQ}\fps{z}$. \draftnote{The sources state the conclusion ``at every $\cQ$ with distinct roots not in $\{0,1,1/\cQ\}$''; the passage to the function field is the principle (P), applied here with the parameter $t$ in place of $\tau$. To be checked that \cref{prop:row-P2} is stated in \cref{sec:row} with a general parameter, as in its source.} This is also the argument of \cref{thm:lea-row2reg}, Steps C--D, with Gauss's operator replaced by $\idePF$.

(iii) By the proof of \cref{lem:ide-norm}, every $\cW_k(U_\pm)$, $k\geq 2$, lies in $\KCY\fps{\cQ}$ and $\cW_2(U_\pm)$ is a unit there; so the $z$-coefficients of $\rpt_{0,\pm}$, polynomials in $\cW_2^{-1},\cW_3,\dots$, lie in $\KCY\fps{\cQ}$. Further $\Lv_+=-\alpha_1\cQ U_+/(1-\cQ U_+)\in\cQ\KCY\fps{\cQ}$, $\Lv_-=-\alpha_1\ideQU/(1-\ideQU)\in\KCY\fps{\cQ}$ with value $-\alpha_1(\alpha_0/\wB)/(\alpha_1/\wB)=-\alpha_0$, and $\Delta_\pm\in\KCY\fps{\cQ}$ with the non-zero values $\eul(\wts_\pm)$ (\cref{lem:ide-ss}), so $\Delta_i'/\Delta_i\in\KCY\fps{\cQ}$ and $\rZ_{1,i}=(\Lv_i+z\Dop_i)\rZ_{0,i}$ has coefficients in $\KCY\fps{\cQ}$. At $\cQ=0$ the derivation $\thQ=\cQ\,\dd/\dd\cQ$ vanishes, so $\rZ_{1,i}(z;0)=\Lv_i(0)\rZ_{0,i}(z;0)$, and $\rZ_{0,i}(z;0)=\hnw\rpt_{0,i}(z;0)=h_{\wts_i}$ by \cref{lem:ide-norm}. Since $\eid_i^*(0)=(1,\Lv_i(0))$, this says $\eid_i^*\circ\RZ^{-1}(-z)|_{\cQ=0}=h_{\wts_i}(z)\eid_i^*$, that is, $\RZ^{-1}(-z)|_{\cQ=0}=\mathrm{diag}(h_{\wts_i}(z))$ on the $\eid_i(0)$, and $h_{\wts_i}(z)^{-1}=h_{\wts_i}(-z)$ gives $\RZ|_{\cQ=0}=\mathrm{diag}(h_{\wts_i})$.

(iv) Write $\vv_i=\vv^Z_i$ and $\ideY(z)\coloneqq\langle\vv_+(z),\vv_-(-z)\rangle\in\KCY\fps{\cQ}\fps{z}$. The pairing has constant coefficients, so $\dd\langle a,b\rangle/\dd t=\langle a',b\rangle+\langle a,b'\rangle$, and $\Aop$ is self-adjoint for it because $H\qprod$ is $\eta$-self-adjoint (\cref{prop:set-frob}). From (ii), $z\vv_i'=(\Aop-\Lv_i+z\frac{\Delta_i'}{2\Delta_i})\vv_i$, and replacing $z$ by $-z$, $z\,\vv_-(-z)'=-(\Aop-\Lv_-)\vv_-(-z)+z\frac{\Delta_-'}{2\Delta_-}\vv_-(-z)$. Hence
\begin{align*}
z\ideY'&=\langle(\Aop-\Lv_+)\vv_+,\vv_-(-z)\rangle-\langle\vv_+,(\Aop-\Lv_-)\vv_-(-z)\rangle+\frac z2\Big(\frac{\Delta_+'}{\Delta_+}+\frac{\Delta_-'}{\Delta_-}\Big)\ideY\\
&=\Big[(\Lv_--\Lv_+)+\frac z2\Big(\frac{\Delta_+'}{\Delta_+}+\frac{\Delta_-'}{\Delta_-}\Big)\Big]\ideY .
\end{align*}
Write $\ideY=\sum_n\ideY_nz^n$ with $\ideY_n\in\KCY\fps{\cQ}$. At $z^0$: $(\Lv_--\Lv_+)\ideY_0=0$, so $\ideY_0=0$ since $\Lv_+-\Lv_-$ is a unit; at $z^n$: $\ideY_{n-1}'=(\Lv_--\Lv_+)\ideY_n+\frac12(\ldots)\ideY_{n-1}$, so $\ideY_{n-1}=0$ implies $\ideY_n=0$. Thus $\ideY=0$. Likewise, writing $\langle\vv_i(z),\vv_i(-z)\rangle=\Delta_i\,\mathsf{K}_i(z)$ (so that $\mathsf{K}_1,\mathsf{K}_2$ are the symplectic constants $\Kone,\Ktwo$ of \cref{def:set-factors}), the terms in $\Aop-\Lv_i$ cancel and $z(\Delta_i\mathsf{K}_i)'=z\frac{\Delta_i'}{\Delta_i}\Delta_i\mathsf{K}_i$, so $\mathsf{K}_i'=0$ and $\mathsf{K}_i$ is independent of $\cQ$; its value at $\cQ=0$ is, by (iii), $h_{\wts_i}(z)h_{\wts_i}(-z)\langle\chiv_i,\chiv_i\rangle/\Delta_i$ at $\cQ=0$, and $\langle\chiv_i,\chiv_i\rangle=\epsilon_1\epsilon_2(\alpha_0\alpha_1\alpha_2+2\alpha_1\alpha_2\Lv_i-E\Lv_i^2)=\Delta_i$ by the pairing of \cref{prop:set-frob} and \cref{eq:ide-Delta}, while $h_{\wts_i}(z)h_{\wts_i}(-z)=1$. So $\mathsf{K}_i=1$. In terms of $\RZ$: $\langle\vv_i(z),\vv_j(-z)\rangle=\Delta_i\Delta_j\,\eta\big(\RZ^{-1}(z)\RZ^{-1}(-z)^\dagger\eid_i,\eid_j\big)$, so the two identities say $\RZ^{-1}(z)\RZ^{-1}(-z)^\dagger=1$, which is the symplectic condition. Then, exactly as in (Fl9), $\vv_i=\Delta_i^{1/2}\eta(\RZ(z)\fid_i,\cdot)$ and the row equations of (ii) are the columns of \cref{eq:ide-flat} for $\RZ$, lowered by $\eta$; so $\RZ$ is a flat symplectic solution with $\RZ(0)=1$, and $\RZ\Psi\re^{\idecan/z}$ is a unitary fundamental solution. By \cref{prop:ide-flat} the same holds for $\Rm$, and \cref{lem:ide-ambig} gives the last statement.
\end{proof}

\subsection{The constants}

\begin{prop}[$\Rm=\RZ$]\genusstatus{AI-proposed, unchecked}
\label{prop:ide-const}
The constants of \cref{prop:ide-rows}(iv) vanish: $\Rm=\RZ$. Consequently, for $g\geq 2$, $\sF_g=\int_{\Mbar_g}\Omega_{g,0}$ is the point-frame graph sum for the rows $\vv^Z_i=(\rZ_{0,i},\rZ_{1,i})$, that is, the sum \cref{eq:set-graphsum} with the vertex class $\Hg_{\bw}$ replaced by $1$ and the factors $\Delta_i$, $\Dil_i(d)$, $\Edge_{ij}(a,b)$ of \cref{def:set-factors} computed from the rows $\vv^Z_i$.
\end{prop}
\begin{proof}
\emph{Reduction at $\cQ=0$.} Write $\idesp_p\coloneqq\exp\sum_ka_k^{(p)}z^{2k-1}$, $p=\pm$. The objects $\Rm$, $\Rm^{-1}$, $\eid_i$, $\Delta_i^{\pm1}$ and $\eta^{-1}$ are regular at $\cQ=0$ (\cref{prop:ide-rows}(iii), (iv), \cref{lem:ide-ss}), and the graph sum $\Rm.\TFT$ is polynomial in them, so $\Omega|_{\cQ=0}=\Rm|_{\cQ=0}.\TFT|_{\cQ=0}$ with $\Rm|_{\cQ=0}=\mathrm{diag}(h_{\wts_p}\idesp_p)$ on the $\eid_p(0)$ by \cref{prop:ide-rows}(iii), (iv). A diagonal R-matrix has diagonal edges, so every graph is monochromatic, and by $\sum_v(g_v-1)+|\Edg(\Gamma)|=g-1$ the $(\eid_p,\dots,\eid_p)$ component of $\Rm|_{\cQ=0}.\TFT|_{\cQ=0}$ is $\Delta_p(0)^{g-1}\big((h_{\wts_p}\idesp_p).\ideone\big)_{g,n}$, while mixed insertions vanish.

\emph{Constant maps.} At $\cQ=0$ only constant maps contribute: the $\TCY$-fixed loci of $\Mbar_{g,n}(Z,0)=\Mbar_{g,n}\times Z$ are $\Mbar_{g,n}\times\{p_\pm\}$, with virtual class $[\Mbar_{g,n}]$ and virtual normal bundle $T_pZ-\Hdg^\vee\otimes T_pZ$; by \cite{GP99}, with $e(\Hdg^\vee\otimes\bbC_w)=\Lamv(w)$ and $e(T_pZ)=\eul(\wts_p)=\prod_{w\in\wts_p}w$,
\[
\Omega_{g,n}(\eid_p,\dots,\eid_p)|_{\cQ=0}=\frac{\prod_{w\in\wts_p}\Lamv(w)}{\eul(\wts_p)}=\eul(\wts_p)^{g-1}\,\Hg_{\wts_p},
\]
mixed insertions vanish because $\eid_p|_{q}=\delta_{pq}$, and $\Delta_p(0)=\eul(\wts_p)$ (\cref{lem:ide-ss}). Comparing the two descriptions, and using $\Hg_{\wts_p}=h_{\wts_p}.\ideone$ (\cref{lem:ide-mumford}) and the group property \cite[Prop.~2.14]{PPZ15}, $(h_{\wts_p}\idesp_p).\ideone=h_{\wts_p}.(\idesp_p.\ideone)$, we get $\idesp_p.\ideone=\ideone$ for all stable $(g,n)$ with $n\geq 1$.

\emph{Induction.} Suppose $\idesp_p\neq 1$ and let $k$ be minimal with $a_k^{(p)}\neq 0$. Work on $\Mbar_{k,1}$ with the insertion $\eid_p$: on $\Mbar_k$, which has no marking, $\Omega_{k,0}$ is a sum over the colours and one colour cannot be separated. By the flat unit, $(\idesp_p.\ideone)_{k,1}=\pi^*(\idesp_p.\ideone)_{k,0}$, and the degree-$(2k-1)$ part of $(\idesp_p.\ideone)_{k,1}$ is, by \cref{lem:ide-mumford}, $a_k^{(p)}\frac{(2k)!}{B_{2k}}\mathrm{ch}_{2k-1}(\Hdg)$ on $\Mbar_{k,1}$, which is non-zero by \cref{lem:ide-ch} (for $k\geq 2$ because $\pi^*\colon H^*(\Mbar_k;\bbQ)\to H^*(\Mbar_{k,1};\bbQ)$ is injective, $\pi_*(\psi_1\pi^*y)=(2k-2)y$; for $k=1$ directly on $\Mbar_{1,1}$). This contradicts $\idesp_p.\ideone=\ideone$. So $a^{(p)}_k=0$ for all $k$ and $\Rm=\RZ$.

\emph{The graph sum.} Integrating $\Omega_{g,0}=\RZ.\TFT$ over $\Mbar_g$ gives a sum over the connected stable graphs of genus $g$ without legs, in which the vertex at $v$ with colour $i$ contributes $\TFT(\eid_i^{\otimes N})=\Delta_i^{g_v-1}$ integrated against the psi classes, the translation leg at a colour-$i$ vertex contributes $\eid_i^*(\ideT(\psi))=\psi\big(1-\eid_i^*(\RZ^{-1}(\psi)\ideone)\big)=\psi\big(1-\rZ_{0,i}(-\psi)\big)=\sum_{d\geq 2}\Dil_i(d)\psi^d$ (the terms $\psi^0,\psi^1$ cancel because $\rZ_{0,i}[0]=1$, and $\Dil_i(d)=(-1)^d\rZ_{0,i}[d-1]$), and the edge between colours $i,j$ contributes
\[
\frac{\Delta_i\delta_{ij}-\langle\vv_i(-z),\vv_j(-w)\rangle}{z+w}=\Edge_{ij}(z,w),
\]
because $(\eid_i^*\otimes\eid_j^*)(\eta^{-1})=\Delta_i\delta_{ij}$ and $(\eid_i^*\otimes\eid_j^*)\big((\RZ^{-1}(z)\otimes\RZ^{-1}(w))\eta^{-1}\big)=\sum_k\Delta_k\,\vv_i(-z)(\eid_k)\vv_j(-w)(\eid_k)=\langle\vv_i(-z),\vv_j(-w)\rangle$, and by the symplecticity (iv) this is the edge factor of \cref{lem:set-edgegen}, $\Edge_{ij}(z,w)=[\Pij_{ij}(z,w)-\Pij_{ij}(-w,w)]/(z+w)$ with $\Pij_{ij}(-w,w)=-\langle\vv_i(w),\vv_j(-w)\rangle=-\Delta_i\delta_{ij}$. With the decorations of \cref{def:set-graphsum} (exponents on half-edges, ordered tuples of dilaton degrees $d\geq 2$ with weight $1/m_v!$) this is the point-frame sum. (This is Givental's formula \cite[(1), (6), (7)]{Giv01}, with the vertex weights $\Delta_i^{g_v-1}$, the dilaton leaves $T^i_d$, $d\geq 2$, and the edges $\Delta_i^{1/2}\Delta_j^{1/2}V^{ij}$, written in terms of the rows $\vv_i=\Delta_i^{1/2}\re^{-u_i/z}\,\eta(\ideS\mathbf{e}_i,\cdot)$.)
\end{proof}

\subsection{Frame B}

\begin{prop}[The frame-B rows]\genusstatus{AI-proposed, unchecked}
\label{prop:ide-frameB}
Let $\RB\coloneqq\RZ\hw^{-1}$ with $\bw=(E,-\epsilon_1,-\epsilon_2)$, whose rows are $\eid_i^*\circ\RB^{-1}(-z)=\hw(z)^{-1}\vv^Z_i=(\rr_{0,i},\rr_{1,i})$, the frame-B rows of \cref{def:set-rows}. Then, for $g\geq 2$, $\sF_g$ is the Hodge-vertex graph sum \cref{eq:ide-main} with vertex class $\Hg_{\bw}$ and the factors computed from the rows $(\rr_{0,i},\rr_{1,i})$. The exponent $-1$ is derived, not chosen.
\end{prop}
\begin{proof}
By \cref{lem:ide-mumford}, $\Hg_{\bw}=\hw.\ideone$ as rank-one theories, where $\hw$ is a scalar series (the same at both idempotents) with $\hw(z)\hw(-z)=1$, hence a symplectic R-matrix. Since $\hw$ is scalar, the legs, edges and translations of the graph sum defining $\hw.\TFT$ are those of $\hw.\ideone$ multiplied by the identity, respectively by the bivector $\eta^{-1}$ and the unit; by the gluing and unit axioms of $\TFT$ the vertex values reassemble, on each stable graph, to the pull-back of $\TFT_{g,n}$ along the gluing map, and the projection formula gives $(\hw.\TFT)_{g,n}=(\hw.\ideone)_{g,n}\,\TFT_{g,n}=(\Hg_{\bw}\otimes\TFT)_{g,n}$. \draftnote{The sources state ``$\Hg_{\bw}\otimes\TFT=\hw.\TFT$ ($\hw$ is scalar)'' in one phrase; the projection-formula sentence is the expansion of that phrase, to be checked.} By the group property \cite[Prop.~2.14]{PPZ15},
\[
\RB.(\Hg_{\bw}\otimes\TFT)=\RB.(\hw.\TFT)=(\RB\hw).\TFT=\RZ.\TFT=\Omega .
\]
Integrated over $\Mbar_g$, the right side is $\sF_g$ by \cref{prop:ide-const}. The left side is the graph sum of \cref{prop:ide-const}, last paragraph, for the rows of $\RB$, with the vertex values $\Delta_i^{g_v-1}\Hg_{\bw}$ in place of $\Delta_i^{g_v-1}$: this is \cref{eq:ide-main}, the sum \cref{eq:set-graphsum} with the vertex class $\Hg_{\bw}$ and the frame-B factors of \cref{def:set-factors}, which are computed from the rows $\eid_i^*\circ\RB^{-1}(-z)$. Finally, right multiplication of an R-matrix by the scalar $\hw^s$ multiplies the rows by $\hw(z)^s$: $\eid_i^*\circ(\RZ\hw^s)^{-1}(-z)=\hw(-z)^{-s}\,\eid_i^*\circ\RZ^{-1}(-z)=\hw(z)^s\vv^Z_i$, by $\hw(z)\hw(-z)=1$. So the rows of $\RB=\RZ\hw^{-1}$ are $\hw^{-1}\vv_i^Z=\hw^{-1}\hnw(\rpt_{0,i},(\Lv_i+z\Dop_i)\rpt_{0,i})$, which are the frame-B rows of \cref{def:set-rows}, since $\Lv_i+z\Dop_i$ is linear over $\bbC\fps{z}$ with $t$-independent coefficients and so commutes with the multiplication by $\hw^{-1}\hnw$; and $\hw^{-1}\hnw=h_{(\alpha_1,\wA,-E)}$ by \cref{lem:set-hmult}.
\end{proof}

\subsection{The coloured terms as elements of the function field}

\begin{lem}[Description of the coloured terms]\genusstatus{AI-proposed, unchecked}
\label{lem:ide-terms}
Let $\Ccov$ be the double cover $\ydc^2=\Disc(\cQ)$ of $\bbP^1_{\cQ}$, with function field $\KCY(\Ccov)=\KCY(\cQ)[U]/(\Pq(U))$, so that $2\wB\cQ U-(\alpha_0+\cQ(\alpha_1-\alpha_2))=\pm\ydc$ (\cref{lem:set-disc}(v)); the two colours are the two roots of $\Pq$, exchanged by the deck involution of $\Ccov$. Then:
\begin{enumerate}
\item every $z$-coefficient of $\rr_{0,i}$, $\rr_{1,i}$, every $\Dil_i(d)$, every $\Edge_{ij}(a,b)$ and $\Delta_i^{\pm1}$ is an element of $\KCY(\Ccov)$, whose denominators are products of powers of $\cW_2(U_i)$, $1-U_i$, $1-\cQ U_i$, $\ydc$, $\wB\cQ$ and the constants $\alpha_1,\wA,\epsilon_1,\epsilon_2,E$;
\item hence, for $g\geq 2$, each coloured term $\GCont(\Gamma,\col)$ of \cref{eq:ide-main} is an element of $\KCY(\Ccov)$, regular over every $\cQ\notin\{0,1,\infty\}$ with $\Disc(\cQ)\neq 0$;
\item the factors of colour $2$ are the same expressions as those of colour $1$ evaluated at the other root, so the colour sum $\sum_{\col}\GCont(\Gamma,\col)$ is invariant under the deck involution and lies in $\KCY(\cQ)$;
\item the two points of $\Ccov$ over $\cQ=0$ are unramified, and the expansion of every coloured term at either of them lies in $\KCY\fps{\cQ}$; so the total \cref{eq:ide-main} is regular at $\cQ=0$.
\end{enumerate}
\end{lem}
\begin{proof}
(i) Each $z$-coefficient of $\rpt_{0,i}$ is a polynomial with rational coefficients in $\cW_3(U_i),\cW_4(U_i),\dots$ and $1/\cW_2(U_i)$ (\cref{def:set-rows}); $\cW_k(U)$ is, for $k\geq 2$, a polynomial over $\KCY$ in $U/(1-U)$ and $\cQ U/(1-\cQ U)$, and $\cW_1=\Lv(U)-\alpha_2-\wA U/(1-U)$ (\cref{lem:set-crit}). The operator $\Dop_i$ involves $\dd U_i/\dd t=\cQ\,\dd U_i/\dd\cQ=-\cQ(\de_{\cQ}\Pq)/(\de_U\Pq)(U_i)$, whose denominator $\de_U\Pq(U_i)=2\wB\cQ U_i-(\alpha_0+\cQ(\alpha_1-\alpha_2))=\pm\ydc$, and $\Delta_i'/\Delta_i$ with $\Delta_i=\alpha_1\wA\epsilon_1\epsilon_2\cW_2(U_i)$ (\cref{eq:ide-Delta}); $\hnw$, $\hw^{-1}$ and $\eta^{-1}$ are constants of $\KCY\fps{z}$ with the denominators $\alpha_1,\wA,\epsilon_1,\epsilon_2,E$. The dilaton and edge factors are polynomials in the row coefficients and $\eta^{-1}$ (\cref{def:set-factors}, \cref{lem:set-edgegen}). Since $U_i$ is finite exactly where $\wB\cQ\neq 0$, the list of denominators is as stated.

(ii) In genus $g$ only finitely many $z$-coefficients enter: the dilaton degrees and the edge exponents of a decoration are bounded by $3g_v-3+n_v+m_v$. The denominators of (i) vanish only as follows: $\cW_2(U_i)=0$ iff $\Delta_i=0$ iff, by \cref{eq:ide-Delta}, $\Lv_i=\Lv_j$ or $\Lv_i=\infty$, that is, $\Disc(\cQ)=0$ or $\cQ=1$ (\cref{lem:set-disc}(i), $\Disc=\cQ^2\tau^2(\Lv_+-\Lv_-)^2$, and $1-\cQ U_i=0$ only over $\cQ=1$); $1-U_i$ and $1-\cQ U_i$ vanish only over $\cQ=1$, by $\Pq(1)=\wA(\cQ-1)$, $\Pq(1/\cQ)=\alpha_1(1-\cQ)/\cQ$ (\cref{lem:set-crit}); $\ydc=0$ iff $\Disc=0$; $\wB\cQ=0$ iff $\cQ=0$; and at $\cQ=\infty$ see \cref{lem:set-limits}. So each coloured term is regular over every $\cQ\notin\{0,1,\infty\}$ with $\Disc(\cQ)\neq 0$.

(iii) Every factor of colour $i$ is a universal expression in $U_i$, $\cQ$ and the weights; the deck involution $U\mapsto U_++U_--U$ of the quadratic algebra $\KCY(\cQ)[U]/(\Pq)$ (\cref{def:set-QA}) exchanges the two roots and hence the two colours, and the scalar constants $\hnw$, $\hw^{-1}$ are the same at both colours (\cref{lem:ide-norm}). The sum over all colourings of a graph is therefore invariant under the deck involution, and an element of $\KCY(\Ccov)$ fixed by the involution lies in $\KCY(\cQ)$.

(iv) $\Disc(0)=\alpha_0^2\neq 0$, so $\cQ=0$ is unramified (\cref{lem:set-disc}(v)), and the two points over it are $U\to-\alpha_2/\alpha_0$ and $\cQ U\to\alpha_0/\wB$, exchanged by the deck involution together with the colours. By \cref{prop:ide-rows}(iii) the rows at both colours lie in $\KCY\fps{\cQ}\fps{z}$ when expanded at the corresponding point, and so do $\Delta_i^{\pm1}$ (\cref{lem:ide-ss}); hence every coloured term expands in $\KCY\fps{\cQ}$ at either point. A rational function of $\cQ$ over $\KCY$ which, as an element of $\KCY(\Ccov)$, is regular over a point of $\bbP^1$ that is unramified is regular there; so the total is regular at $\cQ=0$.
\end{proof}

\subsection{Proof of Theorem~B}

\begin{proof}[Proof of \cref{prop:ide-main}]
By \cref{lem:ide-fixed} the classes $\Omega$ form a cohomological field theory with flat unit over $\KCY\fps{\cQ}$, semisimple at $\cQ=0$ by \cref{lem:ide-ss}, because $\alpha_0\neq 0$ in $\KCY$. By \cref{prop:ide-teleman}, $\Omega=\Rm.\TFT$ for a unique symplectic $\Rm$ over $\KCY\lser{\cQ}$, and by \cref{prop:ide-flat} $\Rm$ is a flat solution of the quantum differential equation. The rows with the constant $\hnw$ (\cref{lem:ide-norm}) form a flat unitary solution $\RZ$ (\cref{prop:ide-rows}), so $\Rm=\RZ$ up to constants (\cref{lem:ide-ambig}), and the constants vanish by the Hodge integrals of constant maps (\cref{prop:ide-const}): $\sF_g=\int_{\Mbar_g}\Omega_{g,0}$ is the point-frame sum for the rows $\vv^Z_i$, an identity in $\KCY\fps{\cQ}$ whose left side is $\sum_{d\geq 0}\GW_{g,d}\cQ^d$ by definition. By \cref{prop:ide-frameB} this is the frame-B Hodge-vertex sum \cref{eq:ide-main}, read in $\KCY\fps{\cQ}$ at either point of $\Ccov$ over $\cQ=0$. By \cref{lem:ide-terms}, every coloured term lies in $\KCY(\Ccov)$, regular over $\cQ\notin\{0,1,\infty\}$ with $\Disc\neq 0$, the colour sum lies in $\KCY(\cQ)$ and is regular at $\cQ=0$, and its expansion there is the series just identified. So $\sF_g\in\KCY(\cQ)$, regular at $\cQ=0$ with expansion $\sum_d\GW_{g,d}\cQ^d$, and \cref{eq:ide-main} holds in $\KCY(\cQ)$.

Every hypothesis on the weights used above holds in $\KCY$: in \cref{lem:ide-fixed}, that the fixed loci are those of $\bbP^1$ and the twisting bundle is moving (all weights are non-zero linear forms, $\alpha_0\neq 0$); in \cref{lem:ide-ss}, $\Lv_+(0)-\Lv_-(0)=\alpha_0\neq 0$ and $\Delta_\pm(0)=\eul(\wts_\pm)\neq 0$; in \cref{prop:ide-teleman}, that $\KCY\lser{\cQ}$ is a field of characteristic $0$ and cardinality $2^{\aleph_0}$; in (Fl8), that $\Delta_+(0)/\Delta_-(0)=-\alpha_1\alpha_2/(\wA\wB)$ is non-constant; in \cref{lem:ide-norm}, that the limiting pairs $(a,b)=(-\alpha_2,\wA)$ and $(\alpha_0,\alpha_1)$ have $a,b,a+b$ non-zero, that is, $\alpha_2,\wA,\alpha_0$ and $\alpha_0,\alpha_1,\wB$ non-zero; in \cref{prop:ide-rows}(ii), (iv), the identities of \cref{lem:ide-g0} and the invertibility of $\Lv_+-\Lv_-$ in $\KCY\fps{\cQ}$; in \cref{prop:ide-const}, $\mathrm{ch}_{2k-1}(\Hdg)\neq 0$ on $\Mbar_{k,1}$ and the non-vanishing of the five weights of $\wts_\pm$, so that $\Hg_{\wts_p}=h_{\wts_p}.\ideone$ makes sense; in \cref{prop:ide-frameB}, that $\bw=(E,-\epsilon_1,-\epsilon_2)$ has non-zero entries.
\end{proof}

\begin{rmk}[The alternative route by specialisation]
\label{rmk:ide-alt}
The same proof applies at the five indeterminate weights, over $\Kal=\bbQ(\alpha_0,\alpha_1,\alpha_2,\epsilon_1,\epsilon_2)$, with the frame $\bw=(\alpha_0+\alpha_1+\alpha_2,-\epsilon_1,-\epsilon_2)$, and gives \cref{eq:ide-main} for the five-torus invariants $\sF_g^{\Tmax}$ in $\Kal(\cQ)$. The identity on $\TCY$ then also follows by specialising along $\epsilon_2\mapsto E-\epsilon_1$: the colour-summed total is a quotient of polynomials in the weights and $\cQ$ whose denominator does not vanish on $\TCY$, because every per-term denominator is one of the factors listed in \cref{lem:ide-terms}(i), none of which involves $E-\epsilon_1-\epsilon_2$; and on the Gromov--Witten side $\GW^{\TCY}_{g,d}=\GW^{\Tmax}_{g,d}|_{\TCY}$ for $d\geq 1$ by the polynomiality and restriction statement of \cref{lem:rig-poly}, and for $d=0$ by the constant-map formula of \cref{rmk:rig-d0}, whose denominators $\alpha_0\alpha_1\alpha_2\epsilon_1\epsilon_2$ and $\alpha_0\wA\wB\epsilon_1\epsilon_2$ do not vanish on $\TCY$. Either route gives \cref{eq:ide-main}; the one followed above works over $\KCY$ directly and uses no specialisation.
\end{rmk}

\begin{rmk}[What is and is not an analytic continuation; the fibrewise torus]
\label{rmk:ide-scope}
What is proved about the rows is this: they are formal series in $z$ whose coefficients are elements of $\KCY(\Ccov)$; the genus-$g$ graph sum uses finitely many of them; the identity \cref{eq:ide-main} is an identity in $\KCY(\cQ)$, and the equality of its expansion at $\cQ=0$ with $\sum_d\GW_{g,d}\cQ^d$ is \cref{prop:ide-main}. No analytic continuation in $\cQ$ is used anywhere; everything is algebraic in $\KCY(\Ccov)$. In particular the Laurent expansion of \cref{eq:ide-main} at the conifold point $t=0$, which is the object of \cref{sec:reg,sec:rig}, is the expansion of an element of $\KCY(\cQ)$ at $\cQ=1$.

On the fibrewise torus $\TF$ ($\alpha_0=0$, $\alpha_i=\epsilon_i$), and more generally on the three-torus $\Tthree$ ($\alpha_0=0$), the point $\cQ=0$ is not semisimple: both roots of $\Pq$ tend to infinity and both $\Lv_i\to 0$, and the coloured terms have poles at $\cQ=0$ which cancel only in the total. The argument of this section therefore does not apply there directly. What holds instead is the following. The $\TCY$-invariants are the restrictions of the five-torus invariants, and the $\Tthree$- and $\TF$-invariants are the restrictions of the $\TCY$-invariants, by the contravariance of the invariants under restriction of the torus, together with the polynomiality of $\epsilon_1\epsilon_2\GW_{g,d}$ for $d\geq 1$ (\cref{lem:rig-poly}); \draftnote{BS label: lem: GW contravariance.} this is not used in the proof of \cref{eq:ide-main} but identifies $\sF_g$ with the series of \cite{BS}. On the graph side, \cref{eq:ide-main} is restricted to $\alpha_0=0$ by a ring map applied to the identity at indeterminate weights, with the rows rescaled to the fibrewise rows (\cref{lem:rig-3Tid}, \cref{cor:lea-a0}); the constant $\hnw$ becomes $h_{(\alpha_1,\alpha_2,-\epsilon_1,-\epsilon_2)}$, which on $\TF$ equals $1$ because $\nw$ is then of the form $(\mathsf{v},-\mathsf{v})$.
\end{rmk}

\begin{rmk}[Unitarity and $t$-independence of the frame-B rows; the dilaton shift]
\label{rmk:ide-frame}
\begin{enumerate}
\item $\hw$ is scalar (the same at both idempotents), independent of $t$, and $\hw(z)\hw(-z)=1$. So $\RB=\RZ\hw^{-1}$ is again a flat symplectic solution of the quantum differential equation, the ambiguity being exactly that of \cref{lem:ide-ambig}; the frame-B rows satisfy \cref{eq:ide-roweq} and the symplecticity $\langle\vv_i(z),\vv_j(-z)\rangle=\delta_{ij}\Delta_i$ (that is, $\Kone=\Ktwo=1$ and the orthogonality), which \cref{thm:lea-qde} proves directly from the Euler integral. All dependence on $t$ in \cref{eq:ide-main} sits in $\Delta_i$, $\Dil_i(d)$, $\Edge_{ij}(a,b)$; the vertex class $\Hg_{\bw}$ is independent of $t$. The identity \cref{eq:ide-main} holds for every multiset class $\mathsf{v}$ in place of $\bw$, with $\mathsf{v}$-dependent frame factors; the class $\bw=(E,-\epsilon_1,-\epsilon_2)$ is the one for which the singular frame-B row has the value $E/(E+z)$ at $t=0$ (\cref{cor:row-F2}, \cref{lem:set-hmult}).
\item The translation of \cite[Def.~2.13]{PPZ15} for $\RB$, $z(\ideone-\RB^{-1}(z)\ideone)$, has no $z^0$ and $z^1$ terms; its colour-$i$ part is $\psi(1-\rr_{0,i}(-\psi))=\sum_{d\geq 2}\Dil_i(d)\psi^d$, and the omitted term $\Dil_i(1)\psi=-\psi$ is the dilaton shift of \cref{def:set-graphsum}.
\end{enumerate}
\end{rmk}

% s06_regularity. Section file; new macros at the top; label prefix: reg.
% Local macros (used in this section only):
\providecommand{\regbra}[2]{\langle #1\mid #2\rangle}      % \regbra{f}{\msU}: one extra insertion f(\psi_0)
\providecommand{\regD}{\dg}                                  % total leaf degree at a vertex, \regD_v
\section{Regularity at the conifold point and the polar part}
\label{sec:reg}

This section proves Theorem~C of the introduction (\cref{thm:int-reg}) in the form of \cref{cor:reg-main}. It combines three inputs: the leaf structure of \cref{thm:lea-main} (Theorem~A), which gives the regularity properties (F1), (F3) and the leaf memberships (L1)--(L3) of every frame-B factor to all orders in $t$; the vanishing theorem for the Hodge integrals $\Xiz_h(\msM)$ of \cref{def:set-fpot}, cited from the companion paper as \cref{thm:reg-L}; and the value of the single integral $\ZZ[g,0]$, cited from the companion paper as \cref{thm:reg-VI}. Everything up to and including \cref{thm:reg-main} and parts (a), (b) of \cref{cor:reg-main} is a statement about the decoration sum $\GCont(\Gamma,\col)$ of \cref{def:set-graphsum}, as defined; the identification of the graph sum $\Gsum_g$ with the potential $\sF_g$ (\cref{prop:ide-main}) is used only in the second half of \cref{cor:reg-main}(c), and the proof of regularity and of the value of the pole does not depend on it. As a statement about the decoration sum, Theorem~C therefore rests on Theorem~A and on the two cited theorems of \cite{LP2}, and on nothing else. \draftnote{LP2 numbering provisional.}

The scope is the Calabi--Yau torus $\TCY$, at the generic point (field $\KCY$) or at any weight point with (ND) of \cref{def:set-ND} (field $\bbC$); every statement about a graph is per colouring, and every identity is one of formal series in $z$, $w$, $\psi$ to all orders in $t$. Throughout, $g\geq 2$ is the genus of the graph sum and $h\geq 2$ the genus of a vertex. We write $\RRing\coloneqq\KCY\fps{t}$ at the generic point and $\RRing\coloneqq\bbC\fps{t}$ at a weight point.

\emph{Regularity at $t=0$.} By \cref{def:set-graphsum} and \cref{lem:set-rings}, the decoration sum $\GCont(\Gamma,\col)$ of a coloured stable graph is a finite sum (the decorations being bounded by the dimension of the moduli spaces) of products of coefficients of the frame-B factors $\Delta_{\col_v}^{g_v-1}$, $\Dil_{\col_v}(d)$, $\Edge_{\col_v\col_{v'}}(a,b)$, each of which is a Laurent series in $t$ with coefficients in $\KCY$ (in $\bbC$ at a weight point), by the passage (T) between $\tau$-orders and $t$-orders of \cref{lem:set-rings}. We say that $\GCont(\Gamma,\col)$ is \emph{regular at $t=0$ after summing over decorations} if this Laurent series in $t$ has no negative powers. Its \emph{polar part} is the sum of its terms of negative degree in $t$.

\subsection{The Hodge side: domain, dilaton homogeneity, range of genera}
\label{sec:reg-hodge}

We use the objects of \cref{def:set-fpot}: for $h\geq 2$ the class $\Hg_{\bw}=\prod_{a=1}^3\Lamv(\bw_a)/\bw_a^{h}$ with $\bw=(E,-\epsilon_1,-\epsilon_2)$, the base-point derivatives $\dbase^s$, the neighbourhood $\Dom=\qo+t\RRing\fps{\psi}$ of $\qo$, the derivatives $\dd^n\Fpot_h(q)[f_1,\dots,f_n]$ for $q\in\Dom$, the dilaton-shifted form, and the numbers $\Xiz_h(\msU)$, $\dg(\msU)=\sum_i(j_i-1)$, $J_{\msU}=\sum_ij_i$, $\ZZ[h,n]=\Xiz_h(1^n)$, so that $\ZZ[h,0]=\Fpot_h(\qo)$. Recall $\lf_j=1/(E-j\psi)$, $\qo=E-E^2\lf_1$ and $\Tdil=-\psi^2/(E-\psi)=\qo+\psi$. The class $\Hg_{\bw}$ is pulled back along the forgetful maps, so the string and dilaton equations hold for it; this is the only property of $\Hg_{\bw}$ used in this subsection and in the next, and it holds for every triple $\bw$. The dilaton series of the singular idempotent, $\qt(\psi)=\sum_{d\geq 1}\Dil_1(d)\psi^d=-\psi\,\rr_{0,1}(-\psi;t)$ (\cref{def:set-factors}), lies in $\Dom$: its coefficients lie in $\RRing$ and $\qt|_{t=0}=\qo$, by (F3) and by remark~(iv) of \cref{thm:lea-L1}. Since $\rho\in 1+t\RRing$ by (F1) of \cref{thm:lea-main}, also $\rho\,\qt\in\Dom$, although its $\psi$-coefficient is $-\rho$ rather than $-1$; this is why the domain $\Dom$, and not only the dilaton-shifted form, has to be used. At the regular idempotent, $\qt^{(2)}|_{t=0}=-\psi\,\rr_{0,2}(-\psi;0)$ is not $\qo$ (\cref{prop:lea-regrows}), so $\qt^{(2)}\notin\Dom$; there only the dilaton-shifted form is used, and nothing is rescaled.

\begin{lem}[Dilaton homogeneity on $\Dom$]\genusstatus{AI-proposed, unchecked}
\label{lem:reg-H0}
Let $h\geq 2$, $q\in\Dom$, and let $\nu\in 1+t\RRing$ be independent of $\psi$. Then $\nu q\in\Dom$, and for $f_1,\dots,f_n\in\RRing\fps{\psi}$,
\[
\dd^n\Fpot_h(\nu q)[f_1,\dots,f_n]=\nu^{2-2h-n}\,\dd^n\Fpot_h(q)[f_1,\dots,f_n],
\]
equivalently
\begin{equation}
\label{eq:reg-H0}
\dd^n\Fpot_h(q)[f_1,\dots,f_n]=\nu^{2h-2}\,\dd^n\Fpot_h(\nu q)[\nu f_1,\dots,\nu f_n].
\end{equation}
\end{lem}

Only $\nu=\rho$ is used below.

\begin{proof}
That $\nu q\in\Dom$ is immediate: $\nu q=\qo+\upsilon+(\nu-1)q$ with $(\nu-1)q\in t\RRing\fps{\psi}$.

\emph{Step 1 (base point).} Let $\mathsf{b}_1,\dots,\mathsf{b}_s\in\RRing\fps{\psi}$, and let
\[
\mathsf{a}_m\coloneqq\frac{1}{m!}\int_{\Mbar_{h,s+m}}\Hg_{\bw}\prod_{i=1}^s\mathsf{b}_i(\psi_i)\prod_{k=1}^{m}\Tdil(\psi_{s+k})
\]
be the summand of $\dbase^s(\mathsf{b}_1,\dots,\mathsf{b}_s)$ with $m$ insertions of $\Tdil$. The dilaton equation gives
\[
\int_{\Mbar_{h,s+m+1}}\Hg_{\bw}\,\psi_{s+m+1}\prod_{i}\mathsf{b}_i(\psi_i)\prod_{k\leq m}\Tdil(\psi_{s+k})=(2h-2+s+m)\int_{\Mbar_{h,s+m}}\Hg_{\bw}\prod_i\mathsf{b}_i(\psi_i)\prod_{k\leq m}\Tdil(\psi_{s+k}).
\]
Since $\qo=\Tdil-\psi$, the insertion of $\qo$ in the last slot splits into an insertion of $\Tdil$, which turns the summand with $m$ insertions of $\Tdil$ into $(m+1)\mathsf{a}_{m+1}$, and an insertion of $-\psi$, which gives $-(2h-2+s+m)\mathsf{a}_m$. Reindexing the first sum,
\[
\dbase^{s+1}(\mathsf{b}_1,\dots,\mathsf{b}_s,\qo)=\sum_{m\geq 0}\big[m-(2h-2+s+m)\big]\mathsf{a}_m=(2-2h-s)\,\dbase^s(\mathsf{b}_1,\dots,\mathsf{b}_s).
\]
The dependence on $m$ cancels. (At $s=0$ this reads $E^2\Xiz_h(0)-E^2\Xiz_h(1)=(2-2h)\ZZ[h,0]$, which is the sum of the identities of \cref{rmk:reg-L0} and \cref{lem:reg-C} below.)

\emph{Step 2 (Euler identity on $\Dom$).} Let $q=\qo+\upsilon$ with $\upsilon\in t\RRing\fps{\psi}$, and abbreviate $f=(f_1,\dots,f_n)$. By the definition of $\dd^{n+1}\Fpot_h$ and Step~1 with $s=n+r$,
\begin{align*}
\dd^{n+1}\Fpot_h(q)[f,q]
&=\sum_{r\geq 0}\frac{1}{r!}\Big[\dbase^{n+r+1}(f,\upsilon^{r},\qo)+\dbase^{n+r+1}(f,\upsilon^{r+1})\Big]\\
&=\sum_{r\geq 0}\frac{1}{r!}(2-2h-n-r)\,\dbase^{n+r}(f,\upsilon^{r})+\sum_{r\geq 0}\frac{1}{r!}\,\dbase^{n+r+1}(f,\upsilon^{r+1})\\
&=(2-2h-n)\,\dd^n\Fpot_h(q)[f],
\end{align*}
because, after the reindexing $r\to r-1$ in the last sum, the terms carrying the factor $r$ in the first sum cancel against it. Here $\upsilon^r$ stands for $r$ slots filled with $\upsilon$.

\emph{Step 3 (integration).} Put $\mu\coloneqq\nu-1\in t\RRing$. Let $\mathsf{r}$ be a formal variable, $q_{\mathsf{r}}\coloneqq(1+\mathsf{r}\mu)q$, and $\mathsf{G}(\mathsf{r})\coloneqq\dd^n\Fpot_h(q_{\mathsf{r}})[f]$. Each $t$-coefficient of $\mathsf{G}$ is a polynomial in $\mathsf{r}$, because at fixed $t$-order only finitely many $r$ contribute to the Taylor sum. The Taylor form of the definition gives $\dd\mathsf{G}/\dd\mathsf{r}=\mu\,\dd^{n+1}\Fpot_h(q_{\mathsf{r}})[f,q]$, and Step~2 at the point $q_{\mathsf{r}}=(1+\mathsf{r}\mu)q$ of $\Dom$ gives
\[
(1+\mathsf{r}\mu)\,\frac{\dd\mathsf{G}}{\dd\mathsf{r}}=\mu\,(2-2h-n)\,\mathsf{G}.
\]
The solution of this equation in $\RRing\fps{\mathsf{r}}$ with value $\mathsf{G}(0)$ at $\mathsf{r}=0$ is unique (compare the coefficients of $\mathsf{r}^k$, which is possible in characteristic $0$), so $\mathsf{G}(\mathsf{r})=(1+\mathsf{r}\mu)^{2-2h-n}\mathsf{G}(0)$. Put $\mathsf{r}=1$, which is legitimate since each $t$-coefficient is a polynomial in $\mathsf{r}$. This is the first form of the statement; the second follows by replacing $f_i$ by $\nu f_i$ and using multilinearity.
\end{proof}

\begin{lem}[Range of genera]\genusstatus{AI-proposed, unchecked}
\label{lem:reg-range}
Let $\Gamma$ be a connected stable graph of genus $g$, without legs, with at least one edge. Then $g_v\leq g-1$ for every vertex $v$. Consequently \cref{thm:reg-main}(a) uses $\Mirr(h)$ only for $2\leq h\leq g-1$, and at $g=2$ a graph with an edge has no vertex of genus $\geq 2$ at all.
\end{lem}

\begin{proof}
If $h^1(\Gamma)\geq 1$ this follows from $g=h^1(\Gamma)+\sum_vg_v$. If $\Gamma$ is a tree with at least two vertices, it has at least two vertices of valence $1$, each of genus $\geq 1$ by stability, so no vertex carries all of $g$.
\end{proof}

The stability count concerns half-edges only: $\Gsum_g$ has no external markings, and the legs of step~(ii) of the proof of \cref{thm:reg-main} appear only after the graph has been fixed. The single-vertex graph does need $h=g$.

\subsection{String and dilaton identities}
\label{sec:reg-string}

For $h\geq 2$, a multiset $\msU=\{j_1,\dots,j_r\}$ of integers and a series $f\in\RRing\fps{\psi}$, write
\[
\regbra{f}{\msU}\coloneqq\sum_{m\geq 0}\frac{1}{m!}\int_{\Mbar_{h,1+r+m}}\Hg_{\bw}\;f(\psi_0)\prod_{i=1}^{r}\lf_{j_i}(\psi_i)\prod_{k=1}^{m}\Tdil(\psi_{r+k}),
\]
one extra insertion in the integrals defining $\Xiz_h(\msU)=\dbase^r(\lf_{j_1},\dots,\lf_{j_r})$. The string equation, valid since $h\geq 2$ and $\Hg_{\bw}$ is pulled back along the forgetful map, removes an insertion of $1$ and acts on a series $f$ at another marking by $f\mapsto(f-f(0))/\psi$. On a leaf this gives
\begin{equation}
\label{eq:reg-S}
\frac{\lf_j-\lf_j(0)}{\psi}=\frac{1}{\psi}\Big(\frac{1}{E-j\psi}-\frac1E\Big)=\frac{j}{E}\,\lf_j,
\qquad
\frac{\Tdil-\Tdil(0)}{\psi}=\frac{\Tdil}{\psi}=-\frac{\psi}{E-\psi}=-\psi\lf_1 .
\end{equation}

\begin{lem}[String identity]\genusstatus{AI-proposed, unchecked}
\label{lem:reg-string}
For $h\geq 2$ and every multiset $\msU$ of integers, $\Xiz_h(1,\msU)=(J_{\msU}/E^2)\,\Xiz_h(\msU)$. Iterating,
\[
\Xiz_h(1^p,\msU)=\prod_{i=0}^{p-1}(J_{\msU}+i)\cdot E^{-2p}\,\Xiz_h(\msU).
\]
\end{lem}

\begin{proof}
By (LAM2) of \cref{lem:set-leafid} at $j=1$, $\psi\lf_1=E\lf_1-1$, that is, $\lf_1=E^{-1}(1+\psi\lf_1)$. Hence $\Xiz_h(1,\msU)=E^{-1}\big[\regbra{1}{\msU}+\regbra{\psi\lf_1}{\msU}\big]$. Apply the string equation to the insertion $1$ in $\regbra{1}{\msU}$: by \eqref{eq:reg-S} each probe $\lf_{j_i}$ contributes $(j_i/E)$ times the same integral without the extra marking, and each of the $m$ dilaton leaves contributes an integral in which one $\Tdil$ is replaced by $-\psi\lf_1$; the factor $m/m!=1/(m-1)!$ makes the sum of the latter over $m$ equal to $-\regbra{\psi\lf_1}{\msU}$. So
\[
\regbra{1}{\msU}=\frac{J_{\msU}}{E}\,\Xiz_h(\msU)-\regbra{\psi\lf_1}{\msU},
\]
and the second term cancels in $\Xiz_h(1,\msU)$. The iteration uses $J_{(1^i,\msU)}=J_{\msU}+i$.
\end{proof}

\begin{lem}[Vanishing on $1^p$]\genusstatus{AI-proposed, unchecked}
\label{lem:reg-C}
For $h\geq 2$ and $p\geq 1$, $\Xiz_h(1^p)=0$; that is, $\ZZ[h,n]=0$ for $n\geq 1$.
\end{lem}

\begin{proof}
Apply \cref{lem:reg-string} with $\msU$ empty: the factor $J_{\emptyset}=0$ kills every $p\geq 1$. The proof uses only the string equation and the shapes $\lf_j=1/(E-j\psi)$, $\Tdil=-\psi^2/(E-\psi)$ of the leaves; no weight enters.
\end{proof}

\begin{rmk}[The $\lf_0$ identity; not used]
\label{rmk:reg-L0}
Put $\lf_0\coloneqq 1/E$. For $h\geq 2$ and every multiset $\msU$ of integers,
\[
\Xiz_h(0,\msU)=\frac{\dg(\msU)-2h+2}{E^2}\,\Xiz_h(\msU).
\]
Indeed, from $\Tdil=\qo+\psi=E-E^2\lf_1+\psi$ and (LAM2) at $j=1$ one gets $\psi\lf_1=E\lf_1-1=(\psi-\Tdil)/E$. The dilaton equation gives $\regbra{\psi}{\msU}-\regbra{\Tdil}{\msU}=(2h-2+r)\,\Xiz_h(\msU)$, because after reindexing the sum over the number of dilaton leaves the factor $1/m!$ cancels the dependence on $m$ exactly as in Step~1 of the proof of \cref{lem:reg-H0}. Together with the display in the proof of \cref{lem:reg-string},
\[
\regbra{1}{\msU}=\frac{J_{\msU}}{E}\Xiz_h(\msU)-\regbra{\psi\lf_1}{\msU}=\frac{J_{\msU}-(2h-2+r)}{E}\,\Xiz_h(\msU),
\]
and multiplication by $\lf_0=1/E$ with $J_{\msU}-r=\dg(\msU)$ gives the identity. In particular $\Xiz_h(0)=((2-2h)/E^2)\,\ZZ[h,0]$: a part $0$ lowers $\dg$ by one and reaches the empty multiset, so it is not covered by \cref{thm:reg-L}. The identity is recorded only to explain the next remark.
\end{rmk}

\begin{rmk}[Why a constant leaf would be fatal, and why none occurs]
\label{rmk:reg-nol0}
In the proof of \cref{thm:reg-main} the coefficients of the factors attached to a singular vertex are expanded in the leaves $\lf_j$, $j\geq 1$, and the counting of step~(iv) rests on the bound $\dg\leq k$ at order $t^k$. A constant component of such a coefficient, that is, a multiple of $\lf_0=1/E$, could not be stripped by \cref{lem:reg-string}; by \cref{rmk:reg-L0} its contribution to $\dg$ is $-1$, so it would lower $\dg$ instead of being absorbed by the counting, and when the remaining multiset is empty it would leave $((2-2h)/E^2)\ZZ[h,0]$, a multiple of the value of \cref{thm:reg-VI}, at the order $t^{2-2h+k}$. So a constant leaf would be a genuine obstruction. None occurs: by \cref{thm:lea-main}, $[t^k](\rho\qt)\in\Lsp_{k+1}$ for $k\geq 1$ (L1), $[t^k](\rho f_{j,b})\in\Lsp_{k+1}$ for every $k\geq 0$ (L2, \cref{prop:lea-L2}, \cref{cor:lea-rhoe}), and $[t^k](\rho^2\Edge_{11}(\psi,\psi'))$ lies in the span of the products $\lf_i(\psi)\lf_l(\psi')$ with $i,l\geq 1$ (L3, \cref{thm:lea-L3}); the spaces $\Lsp_m$ contain no constants. The constant $E$ in $\qo=E-E^2\lf_1=[t^0](\rho\qt)$ is the base point of the Taylor expansion on $\Dom$, not an insertion.
\end{rmk}

\subsection{The two cited theorems}
\label{sec:reg-cited}

The companion paper \cite{LP2} works at scale $1$: for an indeterminate $\tau$ it puts $\st\coloneqq\tau(1+\tau)$ and, on $\Mbar_{h,n}$,
\[
\Hg_{(1,\tau,-1-\tau)}=\frac{\Lamv(1)\Lamv(\tau)\Lamv(-1-\tau)}{(-\st)^h},\qquad
\lf^{\tau}_j(\psi)=\frac{1}{1-j\psi},\qquad
\Tdil^{\tau}(\psi)=-\frac{\psi^2}{1-\psi},
\]
and for a finite multiset $\msU$ of integers with $r=\#\msU$ and $2h-2+r>0$
\[
\Xit_h(\msU)\coloneqq\sum_{m\geq 0}\frac{1}{m!}\int_{\Mbar_{h,r+m}}\Hg_{(1,\tau,-1-\tau)}\prod_{j\in\msU}\lf^{\tau}_j(\psi)\prod_{k=1}^m\Tdil^{\tau}(\psi_{r+k}),
\]
a finite sum because $\Tdil^{\tau}$ has $\psi$-degree at least $2$. The symbol $\st$ is used only in this subsection, and $\tau$ means here the indeterminate of \cite{LP2}, not $1-1/\cQ$; the point at which it is evaluated below is written $\tauw$.

\begin{lem}[Scale law]\genusstatus{AI-proposed, unchecked}
\label{lem:reg-scale}
Let $E,\mathsf{v}_2,\mathsf{v}_3$ lie in a field of characteristic $0$, with $E+\mathsf{v}_2+\mathsf{v}_3=0$ and $E\mathsf{v}_2\mathsf{v}_3\neq 0$, and let $\Xiz^{\mathsf{v}}_h(\msU)$ be defined as $\Xiz_h(\msU)$ in \cref{def:set-fpot} with the class $\Hg_{\mathsf{v}}=\prod_a\Lamv(\mathsf{v}_a)/\mathsf{v}_a^h$, $\mathsf{v}=(E,\mathsf{v}_2,\mathsf{v}_3)$, the leaves $\lf_j=1/(E-j\psi)$ and $\Tdil=-\psi^2/(E-\psi)$. Then, with $\tau\coloneqq\mathsf{v}_2/E$, one has $\mathsf{v}=E(1,\tau,-1-\tau)$, $\st=-\mathsf{v}_2\mathsf{v}_3/E^2\neq 0$, and for every finite multiset $\msU$ with $r=\#\msU$ and $2h-2+r>0$
\begin{equation}
\label{eq:reg-scale}
\Xiz^{\mathsf{v}}_h(\msU)=E^{-(3h-3+2r)}\,\Xit_h(\msU)\big|_{\tau=\mathsf{v}_2/E}.
\end{equation}
For $\mathsf{v}=\bw=(E,-\epsilon_1,-\epsilon_2)$ on $\TCY$ one has $\tauw=-\epsilon_1/E$, $\st=-\epsilon_1\epsilon_2/E^2$, and $\Xiz^{\bw}_h=\Xiz_h$ is the integral of \cref{def:set-fpot}.
\end{lem}

\begin{proof}
This is a weight count \cite[Corollary~8.4]{LP2}. Replace $\psi$ by $\psi/E$. Write $\Hg_{\mathsf{v}}=\sum_i(\Hg_{\mathsf{v}})_i$ with $(\Hg_{\mathsf{v}})_i$ of cohomological degree $i$; because of the normalisation by $\mathsf{v}_a^h$, $(\Hg_{\mathsf{v}})_i$ is homogeneous of degree $-i$ in $\mathsf{v}$, so it equals $E^{-i}$ times the corresponding part of $\Hg_{(1,\tau,-1-\tau)}$. The degree-$b$ part of a leaf $1/(E-j\psi)=E^{-1}\sum_b(j\psi/E)^b$ scales as $E^{-1-b}$, and the degree-$b$ part of a dilaton leaf $-\psi^2/(E-\psi)$ scales as $E^{1-b}$. On $\Mbar_{h,r+m}$ the total cohomological degree is $3h-3+r+m$, so the total exponent of $E$ is $-(3h-3+r+m)-r+m=-(3h-3+2r)$, independent of $m$. For $\mathsf{v}=\bw$ the sum-zero condition is $E=\epsilon_1+\epsilon_2$, the defining relation of $\TCY$.
\end{proof}

Without the normalisation by $\mathsf{v}_a^h$ the exponent would be $3-2r$; the vanishing below is not affected by the normalisation.

\begin{thm}[Vanishing of the Hodge integrals; cited]\genusstatus{AI-proposed, unchecked}
\label{thm:reg-L}
\begin{enumerate}
\item[(i)] \textup{\cite[Theorem~8.1]{LP2}} For every $h\geq 2$ and every non-empty multiset $\msM$ of integers $\geq 2$ with $\dg(\msM)\leq 2h-3$, $\Xit_h(\msM)=0$ in $\bbQ(\tau)$.
\item[(ii)] \textup{\cite[Corollary~8.2]{LP2}} With $\Lambda(u)=\sum_iu^i\lambda_i$ and the Chern roots $\mathsf{r}_i$ of the Hodge bundle $\Hdg$, one has $(1-\mathsf{r}/\tau)(1+\mathsf{r}/(1+\tau))=1-(\mathsf{r}+\mathsf{r}^2)/\st$, hence $\Hg_{(1,\tau,-1-\tau)}=\Lambda(-1)\sum_{k=0}^h(-1/\st)^k\,e_k(\mathsf{r}+\mathsf{r}^2)$, and for every multiset $\msU$ (the empty one included)
\begin{equation}
\label{eq:reg-laurent}
\Xit_h(\msU;\tau)=\sum_{k=0}^h\Big(-\frac{1}{\st}\Big)^k\,\Xi^{(k)}_h(\msU),
\end{equation}
with $\Xi^{(k)}_h(\msU)\in\bbQ$ independent of $\tau$. For $\msM$ as in \textup{(i)} every $\Xi^{(k)}_h(\msM)$ vanishes, so $\Xit_h(\msM;\tau_0)=0$ at every $\tau_0$ in a field of characteristic $0$ with $\st(\tau_0)\neq 0$, the diagonal $\tau_0=-1/2$ included.
\item[(iii)] \textup{\cite[Corollary~8.3]{LP2}} The same vanishing, in $\bbQ(\tau)$ and at every such $\tau_0$, holds for every non-empty multiset $\msU$ of positive integers with $\dg(\msU)\leq 2h-3$. No part $0$ is allowed.
\item[(iv)] \textup{(\cref{lem:reg-scale} and \cite[Corollary~8.4]{LP2})} Consequently $\Mirr(h)$ of \cref{def:set-mirr} holds for every $h\geq 2$, for the class $\Hg_{\bw}$, $\bw=(E,-\epsilon_1,-\epsilon_2)$, at the generic point of $\TCY$ (in $\KCY$, indeed in $\bbQ(\epsilon_1,\epsilon_2)$) and at every weight point with $\epsilon_1\epsilon_2E\neq 0$, the diagonal $\epsilon_1=\epsilon_2$ included: for every non-empty $\msM$ with parts $\geq 2$ and $\dg(\msM)\leq 2h-3$,
\[
\Xiz_h(\msM)=E^{-(3h-3+2r)}\,\Xit_h(\msM)\big|_{\tau=\tauw=-\epsilon_1/E}=0,\qquad \st=-\frac{\epsilon_1\epsilon_2}{E^2}\neq 0 .
\]
Together with \cref{lem:reg-string,lem:reg-C}, the same vanishing holds for every non-empty multiset with parts $\geq 1$ in the same range.
\end{enumerate}
\end{thm}

The theorem rests on the Mari\~no--Vafa formula \cite[eq.~(1)--(3)]{LLZ03} and textbook facts; its proof is not reproduced here. Part (iv) is where the Calabi--Yau condition is used, and it is the only place: the hypothesis of (iv) is that $\bw$ is a sum-zero triple whose \emph{first} weight is the leaf scale $E$ of the leaves $\lf_j=1/(E-j\psi)$ produced by Theorem~A. On $\TCY$ the triple $\bw=(E,-\epsilon_1,-\epsilon_2)$ sums to zero precisely because $E=\epsilon_1+\epsilon_2$, so that $\bw=E(1,\tauw,-1-\tauw)$ with $\tauw=-\epsilon_1/E$; and the leaf scale of Theorem~A is $E=\alpha_0+\alpha_1+\alpha_2$, the first weight of $\bw$, by (F2) of \cref{thm:lea-main} and \cref{lem:set-hmult}. Theorem~A itself does not use the Calabi--Yau condition (\cref{rmk:lea-cy}). The specialisation in (ii) and (iv) requires no sign, rationality or integrality condition on the weights: only $\st\neq 0$ and $E\neq 0$ are needed, which on $\TCY$ with (ND) hold.

\begin{thm}[The value of the vertex; cited]\genusstatus{AI-proposed, unchecked}
\label{thm:reg-VI}
\begin{enumerate}
\item[(i)] \textup{\cite[Theorem~9.3]{LP2}} Put $\ZZ_h(\tau)\coloneqq\Xit_h(\emptyset)$ and $\mathsf{a}_1\coloneqq-\tau$, $\mathsf{a}_2\coloneqq 1+\tau$, so that $\mathsf{a}_1\mathsf{a}_2=-\st$. For every $h\geq 2$, identically in $\bbQ(\tau)$,
\[
\st^{\,h-1}\,\ZZ_h(\tau)=\beta_{2h-2}(\mathsf{a}_1,\mathsf{a}_2),\qquad
\beta_n(a,b)\coloneqq(n-1)!\,\coef{\mathsf{t}^n}\,\frac{1}{4\sinh(a\mathsf{t}/2)\sinh(b\mathsf{t}/2)} .
\]
\item[(ii)] \textup{\cite[Theorem~9.1]{LP2}} For $\bw=(E,-\epsilon_1,-\epsilon_2)$ with $E=\epsilon_1+\epsilon_2$ and every $g\geq 2$,
\begin{equation}
\label{eq:reg-VI}
(-\epsilon_1\epsilon_2E^3)^{g-1}\,\ZZ[g,0]=\beta_{2g-2}(\epsilon_1,\epsilon_2),
\end{equation}
identically in $(\epsilon_1,\epsilon_2)$, and at every weight point with $\epsilon_1\epsilon_2E\neq 0$.
\end{enumerate}
\end{thm}

Part (ii) is the form in which the theorem is used; it is literally the situation of this note, with the same class $\Hg_{\bw}$, the same leaves and the same scale $E$, so no conversion is needed. For the record we transcribe how (ii) follows from (i), since this is also what gives the statement at every weight point. By \cref{lem:reg-scale} at $r=0$ with $\mathsf{v}=\bw$, $\ZZ[g,0]=E^{-(3g-3)}\ZZ_g(\tauw)$ at $\tauw=-\epsilon_1/E$, where $\mathsf{a}_i=\epsilon_i/E$. Further $-\epsilon_1\epsilon_2E^3=\st E^5$ and $\beta_n(c\,a,c\,b)=c^{n}\beta_n(a,b)$, since $(n-1)!\coef{\mathsf{t}^n}$ of a function of $c\mathsf{t}$ picks up $c^n$. So
\[
(-\epsilon_1\epsilon_2E^3)^{g-1}\ZZ[g,0]=\st^{\,g-1}E^{5g-5-3g+3}\ZZ_g=E^{2g-2}\beta_{2g-2}(\mathsf{a}_1,\mathsf{a}_2)=\beta_{2g-2}(\epsilon_1,\epsilon_2).
\]
The identity (i) specialises to every $\tau_0$ with $\st(\tau_0)\neq 0$: by \eqref{eq:reg-laurent} at the empty multiset, $\ZZ_g(\tau)$ is a polynomial in $1/\st$ with rational coefficients, and $\beta_{2g-2}(\mathsf{a}_1,\mathsf{a}_2)$ is $(2g-3)!/(\mathsf{a}_1\mathsf{a}_2)$ times a symmetric polynomial in $\mathsf{a}_1^2,\mathsf{a}_2^2$ (see \cref{lem:reg-beta}), hence a polynomial in $\st$ and $1/\st$, because $\mathsf{a}_1^2+\mathsf{a}_2^2=1+2\st$ and $\mathsf{a}_1^2\mathsf{a}_2^2=\st^2$; an identity in $\bbQ(\tau)$ between two Laurent polynomials in $\st$ holds wherever $\st\neq 0$. This gives \eqref{eq:reg-VI} at every weight point with $\epsilon_1\epsilon_2E\neq 0$. The proof of (i) in \cite{LP2} uses the Mari\~no--Vafa formula only at the profile $1^a$ and nothing from \cref{thm:reg-L}; it is not reproduced here.

\begin{lem}[The polynomials $\beta_{2g-2}$ and $\Pg$]\genusstatus{AI-proposed, unchecked}
\label{lem:reg-beta}
Let $\Ssh(y)\coloneqq y/(2\sinh(y/2))$, an even power series in $y$ with constant term $1$, and recall $\Pg=\coef{t^{2g}}\,\epsilon_1\epsilon_2t^2/(4\sinh(\epsilon_1t/2)\sinh(\epsilon_2t/2))$ from \cref{def:set-inv}. Then $\Pg=\coef{t^{2g}}\,\Ssh(\epsilon_1t)\Ssh(\epsilon_2t)$ is a symmetric homogeneous polynomial of degree $2g$ in $\epsilon_1,\epsilon_2$, a polynomial in $\epsilon_1^2,\epsilon_2^2$, and for $g\geq 2$
\begin{equation}
\label{eq:reg-beta}
\beta_{2g-2}(\epsilon_1,\epsilon_2)=\frac{(2g-3)!\,\Pg}{\epsilon_1\epsilon_2}.
\end{equation}
\end{lem}

\begin{proof}
From the definitions,
\[
\Ssh(\epsilon_1t)\,\Ssh(\epsilon_2t)=\frac{\epsilon_1\epsilon_2t^2}{4\sinh(\epsilon_1t/2)\sinh(\epsilon_2t/2)},
\]
so $\Pg=\coef{t^{2g}}\Ssh(\epsilon_1t)\Ssh(\epsilon_2t)$. Since $\Ssh$ is even, write $\Ssh(y)=\sum_{i\geq 0}s_iy^{2i}$ with $s_0=1$; then
\[
\Pg=\sum_{i+j=g}s_is_j\,\epsilon_1^{2i}\epsilon_2^{2j},
\]
which is the stated polynomial. Next,
\begin{multline*}
\beta_{2g-2}(\epsilon_1,\epsilon_2)=(2g-3)!\,\coef{t^{2g-2}}\frac{1}{4\sinh(\epsilon_1t/2)\sinh(\epsilon_2t/2)}\\
=\frac{(2g-3)!}{\epsilon_1\epsilon_2}\,\coef{t^{2g}}\frac{\epsilon_1\epsilon_2t^2}{4\sinh(\epsilon_1t/2)\sinh(\epsilon_2t/2)}
=\frac{(2g-3)!}{\epsilon_1\epsilon_2}\,\Pg,
\end{multline*}
both sides being $(2g-3)!/(\epsilon_1\epsilon_2)$ times $\coef{t^{2g}}\Ssh(\epsilon_1t)\Ssh(\epsilon_2t)$.
\end{proof}

\begin{rmk}[The first values]
\label{rmk:reg-betavalues}
From $2\sinh(y/2)=y\,(1+y^2/24+y^4/1920+y^6/322560+\cdots)$ one finds $\Ssh(y)=1-y^2/24+7y^4/5760-31y^6/967680+\cdots$, so that
\begin{align*}
\mathsf{P}_2&=\frac{7\epsilon_1^4+10\epsilon_1^2\epsilon_2^2+7\epsilon_2^4}{5760},&
\beta_2(\epsilon_1,\epsilon_2)&=\frac{7\epsilon_1^4+10\epsilon_1^2\epsilon_2^2+7\epsilon_2^4}{5760\,\epsilon_1\epsilon_2},\\
\mathsf{P}_3&=-\frac{31\epsilon_1^6+49\epsilon_1^4\epsilon_2^2+49\epsilon_1^2\epsilon_2^4+31\epsilon_2^6}{967680},&
\beta_4(\epsilon_1,\epsilon_2)&=-\frac{31\epsilon_1^6+49\epsilon_1^4\epsilon_2^2+49\epsilon_1^2\epsilon_2^4+31\epsilon_2^6}{161280\,\epsilon_1\epsilon_2},
\end{align*}
using $s_0s_2=7/5760$, $s_1^2=1/576=10/5760$, $s_0s_3=-31/967680$, $s_1s_2=-7/138240=-49/967680$ and $(2g-3)!=1,6$ at $g=2,3$. These are the first two of the polynomials $\beta_n$ of \cite{BS}. \draftnote{BS label: eq:betapol.} At real $\epsilon_1,\epsilon_2\neq 0$ all terms of $\beta_{2g-2}$ have the sign $(-1)^g$, since the sign of $s_i$ is $(-1)^i$, so $\beta_{2g-2}(\epsilon_1,\epsilon_2)\neq 0$ there; it may vanish at special complex weights. No proof step uses the non-vanishing.
\end{rmk}

\subsection{Theorem R: regularity and the polar part of the decoration sums}
\label{sec:reg-thmR}

The hypotheses $\Mirr(h)$ in the following statement are discharged by \cref{thm:reg-L}(iv) for every $h\geq 2$; they are kept in the statement so that the use of each instance is visible. Recall (\cref{def:set-rho}, \cref{prop:lea-key}) that $\rho=(\Delta_1t^2/(-\epsilon_1\epsilon_2E^3))^{-1/2}$, the branch with $\rho(0)=1$, is an element of $1+t\RRing$ free of $z$, and that (F1) of \cref{thm:lea-main} is the exact identity
\begin{equation}
\label{eq:reg-F1}
\Delta_1\,\rho^2=-\frac{\epsilon_1\epsilon_2E^3}{t^2}.
\end{equation}

\begin{thm}[Theorem R on the Calabi--Yau torus]\genusstatus{AI-proposed, unchecked}
\label{thm:reg-main}
Fix $g\geq 2$. Work in frame B with the class $\bw=(E,-\epsilon_1,-\epsilon_2)$, on $\TCY$ at the generic point or at any weight point with (ND).
\begin{enumerate}
\item[(a)] If $\Mirr(h)$ holds for every $2\leq h\leq g-1$, then every coloured stable graph $(\Gamma,\col)$ of genus $g$ with at least one edge is regular at $t=0$ after summing over decorations.
\item[(b)] If $\Mirr(g)$ holds, then the single-vertex graph is regular at colour $2$, and at colour $1$ its decoration sum has the polar part
\begin{equation}
\label{eq:reg-pole}
(-\epsilon_1\epsilon_2E^3)^{g-1}\,\Fpot_g(\qo)\,t^{2-2g}=(-\epsilon_1\epsilon_2E^3)^{g-1}\,\ZZ[g,0]\,t^{2-2g},
\end{equation}
and no other polar term: the coefficients of $t^{2-2g+k}$ vanish for $1\leq k\leq 2g-3$.
\end{enumerate}
\end{thm}

\begin{proof}
\emph{Step 0 (the vertex factor at a singular vertex).} By \cref{def:set-graphsum} and the dilaton-shifted form of \cref{def:set-fpot}, the sum over the dilaton decorations (degrees $d\geq 2$, weight $1/m!$) at a vertex $v$ of genus $h\geq 2$ and colour $i$, whose $n$ half-edges carry the series $f_1,\dots,f_n\in\RRing\fps{\psi}$ coming from the edge factors, equals
\begin{equation}
\label{eq:reg-vertex}
\Delta_i^{\,h-1}\,\dd^n\Fpot_h\big(\qt^{(i)}\big)[f_1,\dots,f_n],
\end{equation}
multilinear in the slots $f_1,\dots,f_n$. The vertex class $\Hg_{\bw}$ is independent of $t$ in frame B; all dependence on $t$ sits in $\Delta_i$, in the dilaton series $\qt^{(i)}$ and in the edge factors. For a vertex of genus $0$ or $1$ the decoration sum is the finite sum of \cref{def:set-graphsum} itself, and \eqref{eq:reg-vertex} is not used there. At colour $1$, $\qt\in\Dom$ and $\rho\qt\in\Dom$ (\cref{sec:reg-hodge}), so the Taylor form of \cref{def:set-fpot} applies.

Edges do not factorise in general: a self-loop at $v$, or an edge between two singular vertices of genus $\geq 2$, is a series in two variables $\psi,\psi'$. Since $\dbase^s$ is $\RRing$-multilinear and depends only on a finite jet of each slot (the number of dilaton leaves is at most $3h-3+s$, and the degree in each slot is cut off by the dimension), it is $t$-adically continuous in each slot and may be applied term by term to a $t$-series of leaves and coefficient by coefficient in the variable $\psi'$ of the far end. At a far end which is a vertex $u$ of genus $0$ or $1$, only the powers $\psi'^{\,b}$ with $b\leq 3g_u-3+n_u$ occur, so only finitely many coefficients enter; at fixed $t$-order the number of Taylor insertions is bounded by that order, so only finitely many multisets occur. Everything below is to be read in this sense.

\emph{Step (i) (homogeneity).} Call a colour-$1$ vertex of genus $h\geq 2$ \emph{rescaled}. At a rescaled vertex $v$ with half-edge series $f_1,\dots,f_n$, \eqref{eq:reg-vertex} with \eqref{eq:reg-H0} at $\nu=\rho$ and with \eqref{eq:reg-F1} gives
\begin{equation}
\label{eq:reg-resc}
\Delta_1^{\,h-1}\,\dd^n\Fpot_h(\qt)[f_1,\dots,f_n]
=\Big(\frac{-\epsilon_1\epsilon_2E^3}{t^2}\Big)^{h-1}\dd^n\Fpot_h(\rho\qt)[\rho f_1,\dots,\rho f_n].
\end{equation}
Three properties of $\rho$ are used: it lies in $\RRing$ and is free of $z$ and $\psi$; $\rho(0)=1$, so that $\nu=\rho$ is admissible in \cref{lem:reg-H0} and $[t^0](\rho\qt)=\qo$ is the base point of $\Dom$; and the exact identity \eqref{eq:reg-F1}. So every half-edge at a rescaled vertex is rescaled by exactly one factor $\rho$, and the prefactor $\Delta_1^{h-1}$ becomes $(-\epsilon_1\epsilon_2E^3)^{h-1}t^{2-2h}$ exactly.

\emph{Step (ii) (leaf expansion).} By (L1) of \cref{thm:lea-main} (\cref{thm:lea-L1}),
\[
\rho\qt=\qo+\sum_{k\geq 1}t^kq^{(k)},\qquad q^{(k)}\in\Lsp_{k+1}=\operatorname{span}(\lf_1,\dots,\lf_{k+1}).
\]
The factor attached to a half-edge of the rescaled vertex $v$ depends on the other end of its edge, as described in \cref{rmk:lea-use}:
\begin{itemize}
\item if the other end is a regular vertex, the edge factor is $\Edge_{12}$ or $\Edge_{21}$, and for each power $w^b$ at the regular end it gives $v$ the leg $\rho f_{2,b}(\psi)=\rho\,[w^b]\Edge_{12}(\psi,w)$; by (L2) (\cref{prop:lea-L2}, \cref{prop:lea-transp} for the orientation $\Edge_{21}(b,a)=\Edge_{12}(a,b)$), $[t^k](\rho f_{2,b})\in\Lsp_{k+1}$ for every $k\geq 0$ and every $b$;
\item if the other end is a singular vertex $v'$ of genus $0$ or $1$, which is not rescaled, the edge factor is $\rho\,\Edge_{11}=\rho^{-1}(\rho^2\Edge_{11})$, and for each power $\psi'^{\,b'}$ at $v'$ it gives $v$ the leg $\rho f_{1,b'}$; by (L2) at $j=1$ (\cref{cor:lea-rhoe}), $[t^k](\rho f_{1,b'})\in\Lsp_{k+1}$; either orientation, since $\Edge_{11}$ is symmetric (L3);
\item if the other end is a rescaled vertex $v'$, the edge factor is $\rho^2\Edge_{11}(\psi,\psi')$, and by (L3) (\cref{thm:lea-L3}) its coefficient at order $t^k$ lies in $\operatorname{span}\{\lf_i(\psi)\lf_l(\psi'):i,l\geq 1,\ (i-1)+(l-1)\leq k\}$; the bound on the leaf degree is split between the two ends;
\item for a self-loop at $v$ the factor is $\rho^2\Edge_{11}(\psi,\psi')$ with both ends at $v$, and the whole bound is assigned to $v$.
\end{itemize}
In every case, the order-$t^k$ coefficient of each factor attached to $v$ is a finite combination of leaves $\lf_j$, $j\geq 1$, of total degree $\dg\leq k$, where $\dg(\lf_{j_1}\cdots\lf_{j_s})=\sum(j_i-1)$; the factor $\rho^{-1}=1+O(t)$ at an unrescaled far end raises the $t$-order and does not raise $\dg$; and no constant leaf occurs (\cref{rmk:reg-nol0}). The unrescaled ends themselves are covered by \cref{cor:lea-alllegs}.

By the definition of $\Fpot_h$ on $\Dom$, $\dd^n\Fpot_h(\qo+\upsilon)[\rho f_1,\dots,\rho f_n]$ with $\upsilon=\sum_{k\geq 1}t^kq^{(k)}$ is the Taylor series $\sum_r(1/r!)\,\dbase^{n+r}(\rho f_1,\dots,\rho f_n,\upsilon^r)$ at $\qo$, finite at each order in $t$ because $\upsilon=O(t)$ and each $\dbase$ is a finite sum. Expanding every slot in leaves by multilinearity, a term of total $t$-order $k$ is a product of coefficients whose leaf degrees add up to at most $k$. Hence the coefficient of $t^{2-2h+k}$ of the vertex factor \eqref{eq:reg-resc} is a finite $\KCY$-combination (a $\bbC$-combination at a weight point) of numbers
\[
\Xiz_h(\msU)=\dbase^{\,r}(\lf_{j_1},\dots,\lf_{j_r}),\qquad \msU=\{j_1,\dots,j_r\},\quad\text{all parts }j_i\geq 1,\quad \dg(\msU)\leq k,
\]
and each half-edge at $v$ contributes at least one leaf, so $r\geq n\geq 1$ whenever $v$ has a half-edge. All leaves occurring at order $t^k$ have integer scales $j\leq k+1$. The $\Edge_{11}$ block between two rescaled vertices is not factorised: it is expanded in the $t$-independent products $\lf_i(\psi)\lf_l(\psi')$ with coefficients in $\RRing$, and the degree bound is split between the ends.

At an unrescaled vertex only regularity is needed, and it is supplied by (F1) and (F3) of \cref{thm:lea-main} (the table \cref{eq:lea-table}): at a singular vertex of genus $0$, $\Delta_1^{-1}=t^2\rho^2/(-\epsilon_1\epsilon_2E^3)=O(t^2)$; at a singular vertex of genus $1$, $\Delta_1^{0}=1$; at a regular vertex of any genus $h\geq 0$, $\Delta_2^{h-1}$ is analytic at $t=0$ because $\Delta_2$ is analytic with $\Delta_2(0)=\epsilon_1\epsilon_2\alpha_1\alpha_2\wA\wB/E\neq 0$ (F1); and every dilaton factor $\Dil_i(d)$ and every edge coefficient $\Edge_{ij}(a,b)$ attached to such a vertex is analytic at $t=0$ (F3).

\emph{Step (iii) (vanishing).} Let $v$ be a rescaled vertex of genus $h$, and consider a number $\Xiz_h(\msU)$ occurring at order $t^{2-2h+k}$ in Step~(ii), with all parts $\geq 1$ and $\dg(\msU)\leq k$. Strip the parts equal to $1$ by \cref{lem:reg-string}: writing $\msU=(1^p,\msM)$ with $\msM$ the parts $\geq 2$,
\[
\Xiz_h(1^p,\msM)=\prod_{i=0}^{p-1}(J_{\msM}+i)\,E^{-2p}\,\Xiz_h(\msM),\qquad\dg(\msM)=\dg(\msU).
\]
If $\msM$ is empty, then $p=r\geq 1$ and the factor $J_{\emptyset}=0$ kills the term (\cref{lem:reg-C}). If $\msM$ is non-empty, all its parts are $\geq 2$ and $\dg(\msM)=\dg(\msU)\leq k$, so $\Xiz_h(\msM)=0$ by $\Mirr(h)$ whenever $k\leq 2h-3$, whatever the $t$-order at which the term occurs; in a graph with an edge, $h\leq g-1$ by \cref{lem:reg-range}, so the hypothesis of (a) covers every rescaled vertex. Consequently a non-zero term in the expansion needs $\regD_v\geq 2h_v-2$ at every rescaled vertex $v$ with at least one leaf, where $\regD_v$ denotes the total leaf degree at $v$ in that term.

\emph{Step (iv) (counting).} This proves (a). In a connected graph with an edge every vertex has at least one half-edge, so $r_v\geq 1$ at every rescaled vertex. Consider a non-zero term in the expansion of $\GCont(\Gamma,\col)$. Each factor contributes a $t$-order, and by Step~(ii) its leaf degree at any rescaled end is bounded by that order, with the bound of an edge between two rescaled vertices split between its ends and that of a loop assigned wholly to its vertex. Hence the $t$-order of the term is at least
\[
\sum_{v\ \mathrm{rescaled}}\big(2-2h_v+\regD_v\big)\;+\;2\cdot\#\{\text{singular vertices of genus }0\}\;+\;0,
\]
where the $2$ comes from $\Delta_1^{-1}=O(t^2)$ at a singular vertex of genus $0$, and the final $0$ accounts for the singular vertices of genus $1$ ($\Delta_1^0=1$), for all regular vertices ($\Delta_2^{h-1}$, $\Dil_2$, $\Edge_{2j}$ analytic), and for every order of an edge or dilaton factor not allocated to a $\regD_v$, all of which contribute at least $0$ by (F1), (F3). By Step~(iii) each summand $2-2h_v+\regD_v$ is at least $0$. Hence the order is at least $0$: the Laurent series of $\GCont(\Gamma,\col)$ has no negative powers.

\emph{Step (v) (the single-vertex graph).} This proves (b). The single-vertex graph has no edges and trivial automorphism group. Its colour-$2$ term is $\Delta_2^{\,g-1}$ times a finite sum of products of the dilaton factors $\Dil_2(d)$ (the dilaton-shifted form of \cref{def:set-fpot}), which is analytic at $t=0$ by (F1) and (F3). At colour $1$ the term is, by Steps~(i) and (ii) with no half-edges,
\[
\Big(\frac{-\epsilon_1\epsilon_2E^3}{t^2}\Big)^{g-1}\Fpot_g(\rho\qt)
=(-\epsilon_1\epsilon_2E^3)^{g-1}\,t^{2-2g}\sum_{k\geq 0}t^k\,\coef{t^k}\Fpot_g\Big(\qo+\sum_{k'\geq 1}t^{k'}q^{(k')}\Big),
\]
and the coefficient of $t^k$ of the Taylor series is a combination of numbers $\Xiz_g(\msU)$ with $\msU$ empty only at $k=0$, where the value is $\Xiz_g(\emptyset)=\Fpot_g(\qo)=\ZZ[g,0]$, and otherwise with $\msU$ non-empty, all parts $\geq 1$ and $\dg(\msU)\leq k$. For $1\leq k\leq 2g-3$ these vanish by Step~(iii) at $h=g$, the case in which only parts $1$ occur being killed by \cref{lem:reg-C} and the other cases by $\Mirr(g)$. So the polar part is exactly \eqref{eq:reg-pole}, and the coefficients of $t^{3-2g},\dots,t^{-1}$ are zero.
\end{proof}

Three remarks on the logic. The leaves occurring at order $t^k$ have integer scales $j\leq k+1$, and only the vanishing of the $\Xiz_h(\msM)$ is used, never a value other than $\ZZ[g,0]$. The ranges match exactly: the polar orders $k\leq 2h-3$ at a vertex of genus $h$ need exactly $\dg\leq 2h-3$, which is the range of $\Mirr(h)$, and the bound $\dg\leq k$ of Theorem~A is sharp for generic weights (\cref{cor:lea-top}, \cref{cor:lea-anti}), so there is no slack in either direction. Finally, no colouring sum is used: every step is per colouring, and the single-vertex graph is excluded from (a) only through $r\geq 1$ in Step~(iii) and through $h\leq g-1$ in \cref{lem:reg-range}.

\subsection{Theorem C}
\label{sec:reg-thmC}

\begin{cor}[Theorem C: regularity and the polar part]\genusstatus{AI-proposed, unchecked}
\label{cor:reg-main}
Fix $g\geq 2$. In frame B with $\bw=(E,-\epsilon_1,-\epsilon_2)$, on $\TCY$ at the generic point or at any weight point with (ND):
\begin{enumerate}
\item[(a)] every coloured stable graph of genus $g$ with at least one edge is regular at $t=0$ after summing over decorations;
\item[(b)] the single-vertex graph is regular at colour $2$, and at colour $1$ its polar part is exactly
\begin{equation}
\label{eq:reg-gap}
\gapg\,t^{2-2g},\qquad
\gapg\coloneqq(-\epsilon_1\epsilon_2E^3)^{g-1}\,\ZZ[g,0]=\beta_{2g-2}(\epsilon_1,\epsilon_2)=\frac{(2g-3)!\,\Pg}{\epsilon_1\epsilon_2},
\end{equation}
with no polar term of order $t^{3-2g},\dots,t^{-1}$;
\item[(c)] hence the expansion of the graph sum $\Gsum_g$ of \cref{eq:set-graphsum} in $\KCY\lser{t}$ (in $\bbC\lser{t}$ at a weight point) has polar part exactly $\gapg\,t^{2-2g}$; and, by \cref{prop:ide-main}, the same holds for $\sF_g$: the image of $\sF_g\in\KCY(\cQ)$ under $\cQ=\re^{t}$ has polar part exactly $\gapg\,t^{2-2g}$.
\end{enumerate}
\end{cor}

\begin{proof}
By \cref{thm:reg-L}(iv), $\Mirr(h)$ holds for every $h\geq 2$ at the generic point of $\TCY$ and at every weight point with $\epsilon_1\epsilon_2E\neq 0$, in particular at every weight point with (ND). So the hypotheses of \cref{thm:reg-main}(a) and (b) are satisfied, and (a) is \cref{thm:reg-main}(a). For (b), \cref{thm:reg-main}(b) gives the polar part $(-\epsilon_1\epsilon_2E^3)^{g-1}\ZZ[g,0]\,t^{2-2g}$ with no other polar term, \cref{thm:reg-VI}(ii) gives $(-\epsilon_1\epsilon_2E^3)^{g-1}\ZZ[g,0]=\beta_{2g-2}(\epsilon_1,\epsilon_2)$, identically and at every weight point with $\epsilon_1\epsilon_2E\neq 0$, and \cref{lem:reg-beta} gives $\beta_{2g-2}=(2g-3)!\Pg/(\epsilon_1\epsilon_2)$.

(c), first half. The graph sum $\Gsum_g=\sum_\Gamma|\Aut\Gamma|^{-1}\sum_{\col}\GCont(\Gamma,\col)$ is a finite sum of Laurent series in $t$ (\cref{def:set-graphsum} and the opening of this section). By (a) every summand belonging to a graph with an edge has no negative powers; by (b) the single-vertex graph, whose automorphism group is trivial, contributes no negative powers at colour $2$ and exactly $\gapg t^{2-2g}$ at colour $1$. So the polar part of $\Gsum_g$ is $\gapg t^{2-2g}$.

(c), second half. By \cref{lem:ide-terms}, each coloured term $\GCont(\Gamma,\col)$ is an element of the function field $\KCY(\Ccov)$ of the double cover $\ydc^2=\Disc(\cQ)$, regular over every $\cQ\notin\{0,1,\infty\}$ with $\Disc(\cQ)\neq 0$; since $\Disc(1)=E^2\neq 0$ (\cref{lem:set-disc}), the point $\cQ=1$ is not a branch point, and the expansion of each term at the colour-$1$ or colour-$2$ point of $\Ccov$ over $\cQ=1$, in the local parameter $t=\log\cQ$, is an element of $\KCY\lser{t}$. By \cref{rmk:set-scope} and \cref{lem:set-rings} these expansions are the formal series in $t$ to which Theorem~A and \cref{thm:reg-main} refer, so the first half of (c) applies to them. By \cref{prop:ide-main}, $\sF_g\in\KCY(\cQ)$ equals $\Gsum_g$ in $\KCY(\Ccov)$. The substitution $\cQ=\re^{t}$ maps $\KCY(\cQ)$, localised away from the poles at $\cQ=1$, into $\KCY\lser{t}$: $\cQ-1=t(1+t/2+\cdots)$ is $t$ times a unit of $\KCY\fps{t}$, and $\cQ$ and $\Disc(\cQ)$ are units at $\cQ=1$. The image of $\sF_g$ is the sum of the expansions of the terms, hence has polar part exactly $\gapg t^{2-2g}$.
\end{proof}

As a statement about the decoration sum, parts (a), (b) and the first half of (c) use Theorem~A (\cref{thm:lea-main}), the two cited theorems of \cite{LP2} (\cref{thm:reg-L}, \cref{thm:reg-VI}), and the elementary identities of \cref{sec:reg-hodge,sec:reg-string}; they use no result of \cref{sec:ide}, no input from \cite{BS}, no topological recursion and no finite-order computation. The identification with $\sF_g$ in the second half of (c) is \cref{prop:ide-main}, with its own inputs. The value $\gapg$ is the coefficient of $t^{2-2g}$ in $\Pg\Li_{3-2g}(\re^{t})/(\epsilon_1\epsilon_2)$, since the only singular term of $\Li_{3-2g}(\re^{t})$ at $t=0$ is $(2g-3)!\,t^{2-2g}$ (\cref{lem:rig-li}); this is how Theorem~C enters the proof of Theorem~D in \cref{sec:rig}.

\begin{rmk}[Specialisation to weight points]
\label{rmk:reg-special}
At a point of $\mathrm{Lie}(\TCY)$ with $\alpha_1\alpha_2\wA\wB\epsilon_1\epsilon_2E\neq 0$, with complex or rational weights, the same proof applies word for word: Theorem~A holds there (\cref{thm:lea-main} is stated at every weight point with (ND)), $\Mirr(h)$ holds there by \cref{thm:reg-L}(ii), (iv), and the value \eqref{eq:reg-VI} holds there because $\ZZ_g(\tau)$ is a polynomial in $1/\st$ (the paragraph after \cref{thm:reg-VI}). The value $\gapg$ is then a number. It is non-zero at real weights (\cref{rmk:reg-betavalues}) but may vanish at special complex weights; no proof step uses its non-vanishing. The regularity statements (a) and (c) do not depend on the value of $\gapg$ at all.
\end{rmk}

\begin{rmk}[Which instances of the vanishing theorem are used, by genus]
\label{rmk:reg-genus}
By \cref{lem:reg-range} and Step~(v) of the proof of \cref{thm:reg-main}:
\begin{itemize}
\item at $g=2$, the graphs with an edge use no instance of $\Mirr$ at all, since every vertex has genus $\leq 1$; the single-vertex graph uses $\Mirr(2)$;
\item at $g=3$, the graphs with an edge use only $\Mirr(2)$; the single-vertex graph uses $\Mirr(3)$;
\item at $g\geq 4$, the graphs with an edge use $\Mirr(h)$ for $2\leq h\leq g-1$, and the single-vertex graph also $\Mirr(g)$.
\end{itemize}
In every genus, the value of the pole uses \cref{thm:reg-VI} and nothing from \cref{thm:reg-L}. \draftnote{The sources record that $\Mirr(2)$ admits a proof independent of the Mari\~no--Vafa formula; that proof is not among the sources of this note and is not reproduced. The authors may wish to cite it, so that genus $2$ rests on \cref{thm:reg-VI} alone.}
\end{rmk}

% s07_rigidity. Section file; new macros at the top; label prefix: rig.
\section{Rigidity}
\label{sec:rig}

This section proves Theorem D (\cref{thm:int-rig}) and draws its consequences: the rigidity conjecture of \cite{BS} for the resolved conifold (\cref{cor:rig-conj}), a flop identity for the constant-map terms (\cref{prop:rig-flop}), and the three-torus $\Tthree$ (\cref{cor:rig-3T}). Throughout, the weights are those of $\TCY$, the field is $\KCY=\bbQ(\alpha_0,\alpha_1,\alpha_2,\epsilon_1)$ with $\epsilon_2=E-\epsilon_1$, and $E=\alpha_0+\alpha_1+\alpha_2=\epsilon_1+\epsilon_2$. In \cref{lem:rig-poly} and \cref{lem:rig-g1} identities on the five-torus $\Tmax$ occur; there the sums $\alpha_0+\alpha_1+\alpha_2$ and $\epsilon_1+\epsilon_2$ are written out, and $E$ is used only after the restriction to $\TCY$. The letter $N$ denotes the numerator of a rational function of $\cQ$ (the critical polynomial of \cref{sec:set} does not occur in this section), and $\sQ\coloneqq 1/\cQ$ is the local coordinate at $\cQ=\infty$.

\subsection{Statement and reduction}
\label{sec:rig-plan}

\begin{thm}[Theorem D: rigidity for $g\geq 2$]\genusstatus{AI-proposed, unchecked}
\label{thm:rig-main}
Let $g\geq 2$. In $\KCY(\cQ)$,
\begin{equation}
\label{eq:rig-main}
\sF_g(\cQ)-\sF_g(0)=\frac{\Pg\,\Li_{3-2g}(\cQ)}{\epsilon_1\epsilon_2},
\qquad\text{equivalently}\qquad
\epsilon_1\epsilon_2\,\GW_{g,d}=\Pg\,d^{\,2g-3}\quad\text{for every } d\geq 1 .
\end{equation}
Both sides of the second identity are polynomials in $(\alpha_0,\alpha_1,\alpha_2,\epsilon_1)$ on $\mathrm{Lie}(\TCY)$ (\cref{lem:rig-poly}), so it holds at every point of $\mathrm{Lie}(\TCY)$, including $\alpha_1=\alpha_2$, $E=0$, $\epsilon_1=0$ or $\epsilon_2=0$, and $\alpha_0=0$. In particular $\epsilon_1\epsilon_2\GW_{g,d}$ depends on $(\epsilon_1,\epsilon_2)$ only, symmetrically and through $\epsilon_1^2,\epsilon_2^2$.
\end{thm}

Here $\sF_g(\cQ)=\sum_{d\geq 0}\GW_{g,d}\cQ^d$ is the series of \cref{def:set-inv}, with the constant-map term $\sF_g(0)=\GW_{g,0}$ included (\cref{rmk:set-d0}), $\Pg=\coef{t^{2g}}\,\epsilon_1\epsilon_2t^2/(4\sinh(\epsilon_1t/2)\sinh(\epsilon_2t/2))$, and $\Li_{-n}(\cQ)=\sum_{d\geq 1}d^n\cQ^d$ (\cref{lem:rig-li}). The equivalence of the two forms of \eqref{eq:rig-main} is the comparison of the coefficients of $\cQ^d$, $d\geq 1$, in $\KCY\fps{\cQ}$: $\sF_g$ is regular at $\cQ=0$ with expansion $\sum_d\GW_{g,d}\cQ^d$ (\cref{lem:rig-poles}), and subtracting the constant $\sF_g(0)\in\KCY$ changes no pole.

\emph{Inputs.} Theorem D uses Theorem B (\cref{prop:ide-main}, through \cref{lem:ide-terms}) and Theorem C (\cref{cor:reg-main}), hence all the inputs on which these rest (\cref{sec:int-inputs}); the only further input is the localisation formula of Graber--Pandharipande \cite{GP99}, used in \cref{lem:rig-poly}. The constant term $\sF_g(0)$ is not determined (\cref{rmk:rig-F0}).

\emph{Plan.} Put
\[
\sG_g\coloneqq\sF_g-\sF_g(0)=\sum_{d\geq 1}\GW_{g,d}\cQ^d\in\KCY(\cQ).
\]
The proof has four steps: (i) the poles of $\sF_g$ lie over $\cQ\in\{1,\infty\}$ or over the two roots of $\Disc(\cQ)$ (\cref{lem:rig-poles}); (ii) there is no pole on $\Disc=0$ (\cref{lem:rig-disc}, using the polynomiality \cref{lem:rig-poly} and the small root \cref{lem:rig-root}); (iii) there is no pole at $\cQ=\infty$ (\cref{lem:rig-inf}); (iv) the comparison with $\Li_{3-2g}$ at $\cQ=1$ (\cref{lem:rig-li}) and the conclusion. Steps (i)--(iii) use Theorem B, the genus-zero data of \cref{sec:set} and the polynomiality lemma \cref{lem:rig-poly} (localisation on a proper space) only; step (iv) is the only place where Theorem C, and through it Theorem A, enters.

\subsection{Step (i): the possible poles}

\begin{lem}[the possible poles]\genusstatus{AI-proposed, unchecked}
\label{lem:rig-poles}
$\sF_g$ lies in $\KCY(\cQ)$, is regular at $\cQ=0$ with expansion $\sum_{d\geq 0}\GW_{g,d}\cQ^d$ there, and its poles lie over $\cQ=1$, over $\cQ=\infty$, or over the two roots of $\Disc(\cQ)$.
\end{lem}

\begin{proof}
This is the description of the coloured terms in \cref{lem:ide-terms}: each coloured term of the graph sum \eqref{eq:set-graphsum} is an element of the function field $\KCY(\Ccov)$ of the double cover $\ydc^2=\Disc(\cQ)$ (\cref{lem:set-disc}(v)), regular over every $\cQ\notin\{0,1,\infty\}$ with $\Disc(\cQ)\neq 0$; the colour sum is invariant under the deck involution of $\Ccov$, hence lies in $\KCY(\cQ)$; and the expansion of every coloured term at either of the two points of $\Ccov$ over $\cQ=0$ lies in $\KCY\fps{\cQ}$ (the rows at both colours lie in $\KCY\fps{\cQ}\fps{z}$ by \cref{prop:ide-rows}(iii), $\cQ=0$ being semisimple because $\alpha_0\neq 0$, \cref{lem:ide-ss}; and $\cQ=0$ is unramified for $\Ccov$ by \cref{lem:set-disc}(v)). By \cref{prop:ide-main} this colour sum is $\sF_g$, and its expansion at $\cQ=0$ is $\sum_d\GW_{g,d}\cQ^d$. Finally, a rational function of $\cQ$ over $\KCY$ which, as an element of $\KCY(\Ccov)$, is regular over a point of $\bbP^1_{\cQ}$ that is unramified for $\Ccov$, is regular at that point.
\end{proof}

\begin{rmk}[what is proved about the rows]
\label{rmk:rig-rows}
The rows of \cref{def:set-rows} are formal series in $z$ whose coefficients are elements of $\KCY(\Ccov)$; the genus-$g$ graph sum uses finitely many of them; the identity \eqref{eq:ide-main} is an identity in $\KCY(\cQ)$, and the equality of its expansion at $\cQ=0$ with $\sum_d\GW_{g,d}\cQ^d$ is Theorem B. No analytic continuation in $\cQ$ is used anywhere in this section: every statement is algebraic in $\KCY(\Ccov)$ (compare \cref{rmk:ide-scope}).
\end{rmk}

\subsection{Step (ii): no pole on the discriminant}

The argument has three parts: the polynomiality of $\epsilon_1\epsilon_2\GW_{g,d}$ in the weights for $d\geq 1$, the existence of a root of $\Disc$ of positive $\alpha_0$-adic valuation, and an $\alpha_0$-adic evaluation. Recall $\Kred=\bbQ(\alpha_1,\alpha_2,\epsilon_1)$, so that $\KCY=\Kred(\alpha_0)$.

\begin{lem}[polynomiality]\genusstatus{AI-proposed, unchecked}
\label{lem:rig-poly}
For $d\geq 1$ and all $g\geq 0$, the five-torus invariant $\epsilon_1\epsilon_2\GW^{\Tmax}_{g,d}$ is a polynomial: it lies in the ring $\bbQ[\alpha_0,\alpha_1,\alpha_2,\epsilon_1,\epsilon_2]$, is homogeneous of degree $2g$, and is symmetric under $\epsilon_1\leftrightarrow\epsilon_2$, under $\alpha_1\leftrightarrow\alpha_2$ and under $(\alpha_0,\alpha_1,\alpha_2)\mapsto(-\alpha_0,\alpha_0+\alpha_1,\alpha_0+\alpha_2)$. The $\TCY$-invariant $\GW_{g,d}$ is its restriction $\epsilon_2\mapsto\alpha_0+\alpha_1+\alpha_2-\epsilon_1$. Hence, on $\TCY$,
\[
\GW_{g,d}\in\frac{1}{\epsilon_1\epsilon_2}\,\bbQ[\alpha_0,\alpha_1,\alpha_2,\epsilon_1]\subset\Kred\fps{\alpha_0}\qquad(d\geq 1),
\]
the inclusion because $1/\epsilon_2=1/\big((\alpha_1+\alpha_2-\epsilon_1)+\alpha_0\big)$ expands in $\Kred\fps{\alpha_0}$, $\alpha_1+\alpha_2-\epsilon_1$ being a non-zero element of $\Kred$.
\end{lem}

\begin{proof}
Let $M\coloneqq\Mbar_g(X,d)$, $X$ the resolved conifold and $d\geq 1$. Every compact curve in $X$ lies in the zero section $\bbP^1$, so $M=\Mbar_g(\bbP^1,d)$ as a space, and $M$ is proper: the class $d[\bbP^1]$ is rigid and $X$ is equirigid in the sense of \cite{BS}. \draftnote{BS label: defn: rigid; the sentence that the resolved conifold is equirigid.} The $\Tmax$-fixed loci of $\Mbar_g(Z,d)$ are those of $M$ (\cref{lem:ide-fixed}). For $d\geq 1$ and $C$ connected, $H^0(C,f^*\cO(-1))=0$: a section vanishes on the non-contracted components, where $f^*\cO(-1)$ has negative degree, and is constant on the contracted components, each of which meets a non-contracted one. So $R^1_{a}\coloneqq R^1\pi_*f^*\cO(-1)_{a}$ is a vector bundle of rank $d+g-1$ on $M$ for $a=\alpha_1,\alpha_2$, and
\[
R\pi_*f^*\Nb=-R^1_{\alpha_1}-R^1_{\alpha_2}+(\cO-\Hdg^\vee)_{-\epsilon_1}+(\cO-\Hdg^\vee)_{-\epsilon_2}.
\]
Thus
\[
e_{\Tmax}(R\pi_*f^*\Nb)^{-1}=\frac{e_{\Tmax}(R^1_{\alpha_1})\,e_{\Tmax}(R^1_{\alpha_2})\,e_{\Tmax}(\Hdg^\vee_{-\epsilon_1})\,e_{\Tmax}(\Hdg^\vee_{-\epsilon_2})}{\epsilon_1\epsilon_2}
\]
is $(\epsilon_1\epsilon_2)^{-1}$ times an honest $\Tmax$-equivariant class on $M$ ($e_{\Tmax}(R^1_a)$ depends on $\alpha_0$ through the rotation of the base). The localisation definition of the invariant (\cref{def:set-inv}) equals the $\Tmax$-equivariant integral of this class over the proper space $M$ by the localisation formula of \cite{GP99}: that formula is stated for a $\bbC^*$-action, and one applies it to a generic one-parameter subgroup of $\Tmax$ with the same fixed loci; both sides are rational functions of the weights and such subgroups are Zariski dense, so the formula holds for $\Tmax$. The equivariant pushforward of a class on the proper $M$ lands in $H^*_{\Tmax}(\mathrm{pt})=\bbQ[\alpha_0,\alpha_1,\alpha_2,\epsilon_1,\epsilon_2]$, which gives the polynomiality. The degree is $2(d+g-1)+2g-2(d+g-1)=2g$, the virtual dimension of $M$ being $2d+2g-2$. The symmetries are those of \cref{lem:set-sym}: they are induced by automorphisms of $Z$, and the expression is symmetric in $\epsilon_1,\epsilon_2$. Restriction to $\TCY$: the fixed loci are the same, the Euler classes of the fixed and moving parts restrict, and no denominator vanishes on $\TCY$ (the weights of the virtual normal bundles are of the form (a fibre weight at $p_\pm$)$+j\alpha_0/k$ and $\alpha_0/k$ with integers $j$ and $k\geq 1$, none of which is a multiple of $\alpha_0+\alpha_1+\alpha_2-\epsilon_1-\epsilon_2$); so $\GW^{\TCY}_{g,d}=\GW^{\Tmax}_{g,d}|_{\TCY}$. More generally, let $\mgp{T}'\subset\Tmax$ be a subtorus on which none of the weights of $R\pi_*f^*\Nb$ on the fixed loci vanishes: these weights are of the form (a fibre weight of $\Nb$ at $p_\pm$)$+j\alpha_0/k$ with integers $0\leq j\leq k$, $k\geq 1$, and when $\alpha_0=0$ on $\mgp{T}'$ they are the four fibre weights $\alpha_1,\alpha_2,-\epsilon_1,-\epsilon_2$ of $\Nb$ at $p_+$ (equal to those at $p_-$), which are then the weights of $R\pi_*f^*\Nb$ on all of $M$. Then the $\mgp{T}'$-fixed loci of $\Mbar_g(Z,d)$ are those of $M$ (every curve of class $d[\bbP^1]$ lies in the zero section, and the two directions of $\bbC^2$ carry the non-zero weights $-\epsilon_1,-\epsilon_2$), $R\pi_*f^*\Nb$ is moving on them, and the $\mgp{T}'$-invariant is defined by localisation as in \cref{def:set-inv}. Since $M$ is proper, the $\Tmax$-equivariant integral over $M$ restricts to the $\mgp{T}'$-equivariant integral of the restricted class, which equals the localisation definition by \cite{GP99} applied to $\mgp{T}'$; so $\epsilon_1\epsilon_2\GW^{\mgp{T}'}_{g,d}$ is the restriction of the polynomial $\epsilon_1\epsilon_2\GW^{\Tmax}_{g,d}$. This covers $\Tthree$ and $\TF$ (\cref{lem:rig-g1}, \cref{lem:rig-3Tid}) and the subtori of $\TCY$ in \cref{cor:rig-conj}. \draftnote{The sources state the restriction for $\TCY$ and, in the three-torus argument, for $\Tthree$ (where the normal weights of the zero section are the four fibre weights); the general subtorus statement is written out here. The hypothesis is stated as the non-vanishing of the weights of $R\pi_*f^*\Nb$ on the fixed loci, not only of the four fibre weights at $p_\pm$: for a subtorus with $\alpha_0\neq 0$ a weight $\alpha_1+j\alpha_0/k$ with $0<j<k$ can vanish although $\alpha_1$ and $\wB$ do not.} This is the contravariance of the invariants of \cite{BS}, and the polynomiality is the polynomiality of the invariants for a rigid class there. \draftnote{BS label: lem: shape of refGW series; lem: GW contravariance.}
\end{proof}

\begin{rmk}[the constant-map term]
\label{rmk:rig-d0}
For $d=0$ the fixed locus is $\Mbar_g\times\{p_+,p_-\}$ and localisation gives
\[
\GW_{g,0}=\sum_{p=\pm}\int_{\Mbar_g}\prod_{a=1}^{5}\frac{\Lamv_g(w^{(p)}_a)}{w^{(p)}_a},
\]
a rational function of the weights with denominators $\alpha_0\alpha_1\alpha_2\epsilon_1\epsilon_2$ and $\alpha_0\wA\wB\epsilon_1\epsilon_2$ (this is the value at $\cQ=0$ of the graph sum, \cref{prop:ide-const}). It is a constant in $\cQ$ and plays no role in step (ii); only $\GW_{g,d}\in\Kred\fps{\alpha_0}$ for $d\geq 1$ is used there. A weaker form of \cref{lem:rig-poly}, $\GW^{\Tmax}_{g,d}\in\bbQ[\alpha_1,\alpha_2,\epsilon_1,\epsilon_2][1/(\alpha_1\alpha_2\epsilon_1\epsilon_2)]\fps{\alpha_0}$, obtained by expanding $e_{\Tmax}^{-1}$ in $\alpha_0$, would also suffice for step (ii); the polynomial form is what makes the extension to the excluded loci in the proof of \cref{thm:rig-main} a one-line statement.
\end{rmk}

\begin{lem}[the small root]\genusstatus{AI-proposed, unchecked}
\label{lem:rig-root}
$\Disc$ has a unique root $\cQ_1$ in $\Kred\fps{\alpha_0}$ with $\cQ_1=O(\alpha_0)$; it has $\alpha_0$-adic valuation $2$:
\[
\cQ_1=-\frac{\alpha_0^2}{4\alpha_1\alpha_2}+\frac{(\alpha_1+\alpha_2)\,\alpha_0^3}{8\alpha_1^2\alpha_2^2}-\frac{(5\alpha_1^2+6\alpha_1\alpha_2+5\alpha_2^2)\,\alpha_0^4}{64\,\alpha_1^3\alpha_2^3}+O(\alpha_0^5).
\]
\end{lem}

\begin{proof}
Modulo $\alpha_0$, $\Disc=4\alpha_1\alpha_2\cQ+(\alpha_1-\alpha_2)^2\cQ^2$ has the simple root $\cQ=0$ (its derivative there is $4\alpha_1\alpha_2\neq 0$ in $\Kred$), which lifts uniquely to a root in the complete local ring $\Kred\fps{\alpha_0}$ by successive approximation (simple-root lifting). Write $\cQ_1=c_1\alpha_0+c_2\alpha_0^2+c_3\alpha_0^3+c_4\alpha_0^4+\dots$ and expand $\Disc(\cQ_1)=\alpha_0^2+(2\alpha_0(\alpha_1+\alpha_2)+4\alpha_1\alpha_2)\cQ_1+(\alpha_1-\alpha_2)^2\cQ_1^2$. The coefficient of $\alpha_0^1$ is $4\alpha_1\alpha_2c_1$, so $c_1=0$. Then: at order $\alpha_0^2$, $1+4\alpha_1\alpha_2c_2=0$; at order $\alpha_0^3$, $4\alpha_1\alpha_2c_3+2(\alpha_1+\alpha_2)c_2=0$; at order $\alpha_0^4$, $4\alpha_1\alpha_2c_4+2(\alpha_1+\alpha_2)c_3+(\alpha_1-\alpha_2)^2c_2^2=0$, giving
\[
c_2=-\frac{1}{4\alpha_1\alpha_2},\qquad c_3=\frac{\alpha_1+\alpha_2}{8\alpha_1^2\alpha_2^2},\qquad
c_4=-\frac{4(\alpha_1+\alpha_2)^2+(\alpha_1-\alpha_2)^2}{64\,\alpha_1^3\alpha_2^3}=-\frac{5\alpha_1^2+6\alpha_1\alpha_2+5\alpha_2^2}{64\,\alpha_1^3\alpha_2^3}.
\]
The other root of $\Disc$ is $-4\alpha_1\alpha_2/(\alpha_1-\alpha_2)^2+O(\alpha_0)$, a unit of $\Kred\fps{\alpha_0}$.
\end{proof}

\begin{lem}[no pole on $\Disc=0$]\genusstatus{AI-proposed, unchecked}
\label{lem:rig-disc}
$\sG_g$, and hence $\sF_g$, has no pole at the roots of $\Disc$.
\end{lem}

\begin{proof}
By \cref{lem:rig-poles} and the irreducibility of $\Disc$ over $\KCY$ (\cref{lem:set-disc}(iv)), in the unique factorisation domain $\KCY[\cQ]$ we can write
\begin{equation}
\label{eq:rig-NG}
\sG_g=\frac{N}{(1-\cQ)^n\,\Disc^m},\qquad N\in\KCY[\cQ],\quad n,m\geq 0,\quad \Disc\nmid N
\end{equation}
($\sG_g$ is regular at $\cQ=0$ and has no other finite poles; $m$ is then the order of the pole of $\sG_g$ along $\Disc$, and we must show $m=0$). Suppose $m\geq 1$.

\emph{1. The coefficients of $N$.} The identity $N=(1-\cQ)^n\Disc^m\,\sG_g$ holds in $\KCY\fps{\cQ}$ (expand the rational functions at $\cQ=0$; $\sG_g$ is the series $\sum_{d\geq 1}\GW_{g,d}\cQ^d$ by \cref{lem:rig-poles}). Embed $\KCY=\Kred(\alpha_0)$ into the field $\Kred\lser{\alpha_0}$ of formal Laurent series in $\alpha_0$ with finitely many negative powers; the identity then holds in $\Kred\lser{\alpha_0}\fps{\cQ}$. On the right, $(1-\cQ)^n\Disc^m$ has coefficients in $\Kred[\alpha_0]$, and $\sG_g$ has coefficients in $\Kred\fps{\alpha_0}$ by \cref{lem:rig-poly}, so the right side lies in $\Kred\fps{\alpha_0}\fps{\cQ}$. Hence every coefficient of $N$ lies in $\Kred(\alpha_0)\cap\Kred\fps{\alpha_0}$, and $N\in\Kred\fps{\alpha_0}[\cQ]$.

\emph{2. Evaluation at $\cQ_1$.} Let $\ev\colon\Kred\fps{\alpha_0}\fps{\cQ}\to\Kred\fps{\alpha_0}$ be the substitution $\cQ\mapsto\cQ_1$. It is well defined and a ring homomorphism because $\cQ_1$ has valuation $\val(\cQ_1)=2>0$: for a series $\sum_dc_d\cQ^d$ with $c_d\in\Kred\fps{\alpha_0}$, the terms $c_d\cQ_1^d$ have valuation at least $2d$, so the sum converges $\alpha_0$-adically, and products are respected ($\Kred\fps{\alpha_0}$ is complete, and the substitution of a topologically nilpotent element is multiplicative). On the polynomial $N$, $\ev$ is the ordinary evaluation $N(\cQ_1)$. Applying $\ev$ to the identity of step 1,
\[
N(\cQ_1)=(1-\cQ_1)^n\,\Disc(\cQ_1)^m\,\ev(\sG_g)=0,
\]
since $\Disc(\cQ_1)=0$ and $m\geq 1$. The point is that the series $\sG_g\in\Kred\fps{\alpha_0}\fps{\cQ}$ can be evaluated at $\cQ_1$ only because $\val(\cQ_1)>0$, and that the value of the rational function $\sG_g$ at $\cQ_1$ is not used: what is used is that the identity $N=(1-\cQ)^n\Disc^m\sG_g$ of $\Kred\fps{\alpha_0}\fps{\cQ}$ is preserved by the ring homomorphism $\ev$, and that $N(\cQ_1)$ is the evaluation of a polynomial.

\emph{3. Divisibility.} $N$ lies in $\Kred(\alpha_0)[\cQ]$ and vanishes at $\cQ_1$ in the extension field $\Kred\lser{\alpha_0}$ of $\Kred(\alpha_0)$. The minimal polynomial of $\cQ_1$ over $\Kred(\alpha_0)$ is $\Disc/(\alpha_1-\alpha_2)^2$, which is irreducible over $\Kred(\alpha_0)=\KCY$ by \cref{lem:set-disc}(iv) and has $\cQ_1$ as a root. Hence $\Disc$ divides $N$ in $\Kred(\alpha_0)[\cQ]=\KCY[\cQ]$, contradicting \eqref{eq:rig-NG}. Therefore $m=0$.
\end{proof}

\subsection{Step (iii): regularity at \texorpdfstring{$\cQ=\infty$}{Q = infinity}}

\begin{lem}[regularity at infinity]\genusstatus{AI-proposed, unchecked}
\label{lem:rig-inf}
$\sF_g$ is regular at $\cQ=\infty$ (here $\alpha_1\neq\alpha_2$ in $\KCY$ is used).
\end{lem}

\begin{proof}
The point $\cQ=\infty$ is unramified for $\Ccov$ (\cref{lem:set-disc}(v): the leading coefficient $(\alpha_1-\alpha_2)^2$ of $\Disc$ is non-zero), so the two points of $\Ccov$ over $\sQ=1/\cQ=0$ are the two colours, and regularity has to be shown at both of them. In the coordinate $\sQ$, the critical equation \cref{lem:set-crit} divided by $\cQ$ reads
\[
\wB\,U^2-(\alpha_0\sQ+\alpha_1-\alpha_2)\,U-\alpha_2\sQ=0,
\]
with the simple roots $U\to(\alpha_1-\alpha_2)/\wB$ and $U\to 0$ at $\sQ=0$ (simple because $\alpha_1\neq\alpha_2$); for the second root put $V\coloneqq\cQ U$, so that $\wB\sQ\,V^2-(\alpha_0\sQ+\alpha_1-\alpha_2)\,V-\alpha_2=0$, with the simple root $V\to-\alpha_2/(\alpha_1-\alpha_2)$. Hence (\cref{lem:set-limits}) the critical points $U_1(\sQ)$ and $V_2(\sQ)\coloneqq\cQ U_2$ are power series in $\sQ$ over $\KCY$ with
\[
U_1(0)=\frac{\alpha_1-\alpha_2}{\wB},\qquad V_2(0)=-\frac{\alpha_2}{\alpha_1-\alpha_2}.
\]
The ingredients of the factors at colour $1$ are the $\cW_k(U_1)$, $k\geq 2$, polynomials in $U_1/(1-U_1)$ and $\cQ U_1/(1-\cQ U_1)=1/(\sQ/U_1-1)$ (\cref{lem:set-crit}); both are power series in $\sQ$, since $1-U_1(0)=\wA/\wB\neq 0$ and $\sQ/U_1-1\to-1$. At colour $2$, $U_2/(1-U_2)=\sQ V_2/(1-\sQ V_2)$ and $\cQ U_2/(1-\cQ U_2)=V_2/(1-V_2)$ are power series in $\sQ$, since $1-V_2(0)=\alpha_1/(\alpha_1-\alpha_2)\neq 0$. Further $\cW_2(U_i)(\sQ=0)\neq 0$ by the values below, so $1/\cW_2(U_i)$ is a power series in $\sQ$; $\Lv_i=\Lv(U_i)=-\alpha_1\cQ U_i/(1-\cQ U_i)$ is a power series in $\sQ$ with $\Lv_1(0)=\alpha_1$, $\Lv_2(0)=\alpha_2$; the derivation $\dd/\dd t=\cQ\,\dd/\dd\cQ=-\sQ\,\dd/\dd\sQ$ maps power series in $\sQ$ to power series in $\sQ$; $\Delta_i=\alpha_1\wA\epsilon_1\epsilon_2\cW_2(U_i)$ (\cref{prop:set-idem}) is a power series in $\sQ$ with non-zero value, so $\Delta_i^{-1}$ and $\Delta_i'/\Delta_i$ are power series in $\sQ$; and $\hnw$, $\hw^{-1}$, $\eta^{-1}$ are constants. Hence every $z$-coefficient of $\rr_{0,i}$ and $\rr_{1,i}$, every $\Dil_i(d)$, every $\Edge_{ij}(a,b)$ and every $\Delta_i^{\pm 1}$ is a power series in $\sQ$ (the normalised expansion of \cref{def:set-rows} has leading coefficient $1$ and involves no square root, so no branch question arises at $\sQ=0$). Therefore every coloured term of \eqref{eq:ide-main}, and $\sF_g$ itself, is a power series in $\sQ$: $\sF_g$ has no pole at $\cQ=\infty$.

\emph{The values at $\sQ=0$} (\cref{lem:set-limits}; they are used again in \cref{prop:rig-flop}). From $\Lv_1+\Lv_2=E/\tau-\alpha_0\to E-\alpha_0=\alpha_1+\alpha_2$ and $\Lv_1\Lv_2=\alpha_1\alpha_2/\tau\to\alpha_1\alpha_2$ at $\tau=1$, the eigenvalues tend to $\{\alpha_1,\alpha_2\}$, and by $\Delta_i=\epsilon_1\epsilon_2(\Lv_i-\Lv_j)(\alpha_1\alpha_2-E\Lv_i)$ (\cref{prop:set-idem}), with $\alpha_2-E=-\wB$ and $\alpha_1-E=-\wA$,
\begin{align}
\Delta_1(\infty)&=\epsilon_1\epsilon_2(\alpha_1-\alpha_2)(\alpha_1\alpha_2-E\alpha_1)=-\alpha_1\wB(\alpha_1-\alpha_2)\epsilon_1\epsilon_2,\notag\\
\Delta_2(\infty)&=\epsilon_1\epsilon_2(\alpha_2-\alpha_1)(\alpha_1\alpha_2-E\alpha_2)=\alpha_2\wA(\alpha_1-\alpha_2)\epsilon_1\epsilon_2,
\label{eq:rig-Dinf}
\end{align}
both non-zero in $\KCY$. The same values follow from $\Delta_i=\alpha_1\wA\epsilon_1\epsilon_2\cW_2(U_i)$: at colour $1$, $\cW_2(U_1)\to-\wA U/(1-U)^2$ at $U=(\alpha_1-\alpha_2)/\wB$, $1-U=\wA/\wB$, which gives $-(\alpha_1-\alpha_2)\wB/\wA$; at colour $2$, $\cW_2(U_2)\to-\alpha_1V/(1-V)^2$ at $V=-\alpha_2/(\alpha_1-\alpha_2)$, $1-V=\alpha_1/(\alpha_1-\alpha_2)$, which gives $\alpha_2(\alpha_1-\alpha_2)/\alpha_1$. In particular $\cW_2(U_i)(0)\neq 0$, as used above.
\end{proof}

\subsection{Step (iv): comparison with the polylogarithm}

\begin{lem}[the polylogarithms $\Li_{-n}$]\genusstatus{AI-proposed, unchecked}
\label{lem:rig-li}
For $n\geq 1$,
\[
\Li_{-n}(\cQ)=\sum_{d\geq 1}d^n\cQ^d=\frac{\Aeul_n(\cQ)}{(1-\cQ)^{n+1}},\qquad \Aeul_n\in\bbZ[\cQ],\quad\deg\Aeul_n=n,\quad\Aeul_n(0)=0,
\]
with $\Aeul_0=\cQ$ and $\Aeul_{m+1}=\cQ\big[(1-\cQ)\Aeul_m'+(m+1)\Aeul_m\big]$. Hence $\Li_{-n}$ is a rational function of $\cQ$ whose only pole is at $\cQ=1$, of order $n+1$, vanishing at $\cQ=0$ and at $\cQ=\infty$, and
\[
\Li_{-n}(1/\cQ)=(-1)^{n+1}\Li_{-n}(\cQ).
\]
Its expansion at $t=0$ under $\cQ=\re^{t}$ is
\[
\Li_{-n}(\re^{t})=(-1)^{n+1}\,n!\,t^{-n-1}+(\text{a power series in }t),
\]
with no other negative power of $t$. For $n=2g-3$ (odd) the polar part is $(2g-3)!\,t^{2-2g}$.
\end{lem}

\begin{proof}
The recursion follows from $\Li_{-m-1}=\thQ\Li_{-m}$, $\thQ=\cQ\,\dd/\dd\cQ$: if $\Li_{-m}=\Aeul_m/(1-\cQ)^{m+1}$ then $\thQ\Li_{-m}=\cQ\big[(1-\cQ)\Aeul_m'+(m+1)\Aeul_m\big]/(1-\cQ)^{m+2}$. By induction $\Aeul_n$ has integer coefficients, vanishes at $0$ and has degree exactly $n$ (the coefficient of $\cQ^{m+1}$ in $\Aeul_{m+1}$ is $-m+(m+1)=1$ times the leading coefficient of $\Aeul_m$). So the only pole of $\Li_{-n}$ is at $\cQ=1$, of order $n+1$, and $\Li_{-n}$ vanishes at $\cQ=\infty$ since $\deg\Aeul_n=n<n+1$. The functional equation: for $n=1$, $\Li_{-1}(\cQ)=\cQ/(1-\cQ)^2$ and $\Li_{-1}(1/\cQ)=\cQ/(\cQ-1)^2=\Li_{-1}(\cQ)$; and $\thQ$ anticommutes with $\cQ\mapsto 1/\cQ$ (if $\varphi(\cQ)=f(1/\cQ)$ then $\thQ\varphi=-(\thQ f)(1/\cQ)$), so $\Li_{-n-1}(1/\cQ)=-\thQ\big[\Li_{-n}(1/\cQ)\big]=-(-1)^{n+1}\Li_{-n-1}(\cQ)=(-1)^{n+2}\Li_{-n-1}(\cQ)$. For $n=0$ instead $\Li_0(1/\cQ)=1/(\cQ-1)=-1-\Li_0(\cQ)$, so $\Li_0$ does not vanish at $\infty$; this is why $g\geq 2$ is needed below. The expansion at $t=0$: $\Li_0(\re^t)=\re^t/(1-\re^t)=-1/t-1/2-t/12+\dots$, and since $\thQ=\dd/\dd t$, $\Li_{-n}(\re^t)=(\dd/\dd t)^n\Li_0(\re^t)=(-1)^{n+1}n!\,t^{-n-1}+(\text{power series})$, with no other negative power. For $n=2g-3$, $(-1)^{2g-2}=1$.
\end{proof}

\subsubsection*{Proof of \cref{thm:rig-main}}
\label{sec:rig-proof}

Let $g\geq 2$ and put
\[
D(\cQ)\coloneqq\sF_g(\cQ)-\sF_g(0)-\frac{\Pg\,\Li_{3-2g}(\cQ)}{\epsilon_1\epsilon_2}\in\KCY(\cQ).
\]
\begin{itemize}
\item \emph{Poles of $D$ at finite $\cQ\neq 1$.} $\sF_g$ has none: it is regular at $\cQ=0$ and has no pole on $\Disc=0$ by \cref{lem:rig-poles,lem:rig-disc}. $\Li_{3-2g}$ has none by \cref{lem:rig-li}.
\item \emph{Pole at $\cQ=\infty$.} None, by \cref{lem:rig-inf} for $\sF_g$ and \cref{lem:rig-li} for $\Li_{3-2g}$.
\item \emph{Pole at $\cQ=1$.} Under $\cQ=\re^t$, $D$ maps to an element of $\KCY\lser{t}$ whose polar part is, by Theorem C (\cref{cor:reg-main}(c)), $\gapg\,t^{2-2g}$ minus, by \cref{lem:rig-li} and \cref{lem:reg-beta},
\[
\frac{\Pg}{\epsilon_1\epsilon_2}\,(2g-3)!\,t^{2-2g}=\beta_{2g-2}(\epsilon_1,\epsilon_2)\,t^{2-2g}=\gapg\,t^{2-2g};
\]
the polar part of $D$ is zero. Here it matters that Theorem C gives the polar part of $\sF_g$ exactly, with no terms $t^{3-2g},\dots,t^{-1}$ (these are what step (v) of \cref{thm:reg-main} kills by $\Mirr(g)$), and that $\Li_{3-2g}(\re^t)$ has none either. A rational function of $\cQ$ over $\KCY$ whose image in $\KCY\lser{t}$ has no negative powers is regular at $\cQ=1$: a pole of order $r\geq 1$ at $\cQ=1$ would give the leading term $(\cQ-1)^{-r}\mapsto t^{-r}$ times a unit of $\KCY\fps{t}$.
\end{itemize}
So $D$ is a rational function of $\cQ$ over $\KCY$ without poles on $\bbP^1$, hence a constant $c\in\KCY$, and $D(0)=\sF_g(0)-\sF_g(0)-0=0$ gives $c=0$. This is \eqref{eq:rig-main} in $\KCY(\cQ)$; comparing the coefficients of $\cQ^d$ gives $\epsilon_1\epsilon_2\GW_{g,d}=\Pg\,d^{2g-3}$ in $\KCY$ for every $d\geq 1$.

\emph{Extension to the excluded loci.} Both sides of $\epsilon_1\epsilon_2\GW_{g,d}=\Pg d^{2g-3}$ are polynomials in $(\alpha_0,\alpha_1,\alpha_2,\epsilon_1)$ on $\mathrm{Lie}(\TCY)$: the left side by \cref{lem:rig-poly}, the right side because $\Pg\in\bbQ[\epsilon_1^2,\epsilon_2^2]$ with $\epsilon_2=\alpha_0+\alpha_1+\alpha_2-\epsilon_1$. Two polynomials that are equal in $\KCY$ are equal, so the identity holds at every point of $\mathrm{Lie}(\TCY)$: at $\alpha_1=\alpha_2$ (where the argument of \cref{lem:rig-inf} would fail), at $E=0$ (outside the frame of \cref{sec:lea}), at $\epsilon_1=0$ or $\epsilon_2=0$ (where $\GW_{g,d}$ itself is undefined but $\epsilon_1\epsilon_2\GW_{g,d}$ is not), and at $\alpha_0=0$ (where $\cQ=0$ is not semisimple). Nothing else is needed: the proof is an identity in the function field, and the excluded loci are reached by polynomial continuation, not by taking limits of the graph sum. \qed

\subsection{Genus zero and genus one}

\begin{prop}[genus zero]\genusstatus{AI-proposed, unchecked}
\label{prop:rig-g0}
For every $d\geq1$, $\GW_{0,d}=1/(d^3\epsilon_1\epsilon_2)$ at every weight point with $\epsilon_1\epsilon_2\neq0$, i.e.\ $\epsilon_1\epsilon_2\GW_{0,d}=\Ph{0}\,d^{-3}$ with $\Ph{0}=1$ at every point of $\mathrm{Lie}(\TCY)$.
\end{prop}

\begin{proof}
This is the Aspinwall--Morrison formula, \cite[p.~3]{FP00}, as used in \cref{lem:ide-g0}: for $d\geq 1$, $H^0(C,f^*\cO(-1))=0$ gives $\Mbar_{0,0}(Z,d)=\Mbar_{0,0}(\bbP^1,d)\times\bbC^2$ with obstruction bundle $R^1\pi_*f^*(\cO(-1)\oplus\cO(-1))$ of rank equal to the virtual dimension, so that, localising on $\bbC^2$, $\GW_{0,d}=1/(d^3\epsilon_1\epsilon_2)$ at every weight. And $\Ph{0}=1$, the constant term of $\epsilon_1\epsilon_2t^2/(4\sinh(\epsilon_1t/2)\sinh(\epsilon_2t/2))$ (\cref{def:set-inv}).
\end{proof}

\begin{lem}[genus one]\genusstatus{AI-proposed, unchecked}
\label{lem:rig-g1}
For every $d\geq 1$, on $\TCY$,
\[
\epsilon_1\epsilon_2\,\GW_{1,d}=\frac{\Ph{1}}{d}=-\frac{\epsilon_1^2+\epsilon_2^2}{24\,d}.
\]
For $d\geq 2$ this uses, as an input, the formula of \cite{BS} for the resolved conifold on the subtori on which $\alpha_i=\epsilon_j$; for $d=1$ it is proved by direct localisation, without that input. \draftnote{BS label: prop:ResConifold, eq:resconGWref.}
\end{lem}

\begin{proof}
In genus one, $\Lamv_1(-\epsilon)=-\epsilon-\lambda_1$ and $\lambda_1^2=0$ (Mumford's relation $c(\Hdg)c(\Hdg^\vee)=1$, \cite{Mum83}), so, in the integral of \cref{lem:rig-poly} on the five-torus,
\[
\epsilon_1\epsilon_2\,\GW^{\Tmax}_{1,d}=\epsilon_1\epsilon_2\,N_d+(\epsilon_1+\epsilon_2)\,L_d,
\]
\[
N_d\coloneqq\int_{[\Mbar_1(\bbP^1,d)]^{\vir}}e(R^1_{\alpha_1})\,e(R^1_{\alpha_2}),\qquad
L_d\coloneqq\int_{[\Mbar_1(\bbP^1,d)]^{\vir}}e(R^1_{\alpha_1})\,e(R^1_{\alpha_2})\,\lambda_1 .
\]
By the degree count of \cref{lem:rig-poly} (the two bundles $R^1_a$ have rank $d$ in genus one, against the virtual dimension $2d$), $N_d\in\bbQ$ and $L_d\in\bbQ[\alpha_0,\alpha_1,\alpha_2]_1$ is a linear form; neither involves $\epsilon_1,\epsilon_2$. By the symmetries of \cref{lem:set-sym}, $L_d=c\,(\alpha_0+\alpha_1+\alpha_2)$ for some $c\in\bbQ$: invariance under $\alpha_1\leftrightarrow\alpha_2$ gives $L_d=c_0\alpha_0+c_1(\alpha_1+\alpha_2)$, and invariance under $(\alpha_0,\alpha_1,\alpha_2)\mapsto(-\alpha_0,\alpha_0+\alpha_1,\alpha_0+\alpha_2)$ gives $-c_0+2c_1=c_0$, i.e.\ $c_0=c_1$.

\emph{$d\geq 2$.} On the fibrewise torus $\TF$ ($\alpha_0=0$, $\alpha_i=\epsilon_i$), a subtorus of $\TCY$ on which $\alpha_1=\epsilon_1$ and on which the invariants are the restrictions of the five-torus ones (\cref{lem:rig-poly}, last paragraph of the proof), the formula of \cite{BS} gives, for all $d\geq 1$, $\sum_g\GW_{g,d}=\frac{1}{d}\big(2\sinh\frac{d\epsilon_1}{2}\big)^{-1}\big(2\sinh\frac{d\epsilon_2}{2}\big)^{-1}$, whose genus-one part is $\epsilon_1\epsilon_2\GW_{1,d}=\coef{t^2}\Ssh(\epsilon_1t)\Ssh(\epsilon_2t)/d=-(\epsilon_1^2+\epsilon_2^2)/(24d)=\epsilon_1\epsilon_2/(12d)-(\epsilon_1+\epsilon_2)^2/(24d)$. On $\TF$ one has $(\epsilon_1+\epsilon_2)(\alpha_0+\alpha_1+\alpha_2)=(\epsilon_1+\epsilon_2)^2$, so the restriction of the five-torus identity reads $\epsilon_1\epsilon_2N_d+c\,(\epsilon_1+\epsilon_2)^2$; the quadratic forms $\epsilon_1\epsilon_2$ and $(\epsilon_1+\epsilon_2)^2$ being linearly independent over $\bbQ$, comparison gives $N_d=1/(12d)$ and $c=-1/(24d)$. Hence, on the five-torus,
\[
\epsilon_1\epsilon_2\,\GW^{\Tmax}_{1,d}=\frac{\epsilon_1\epsilon_2}{12d}-\frac{(\epsilon_1+\epsilon_2)(\alpha_0+\alpha_1+\alpha_2)}{24d},
\]
which on $\TCY$ is $\epsilon_1\epsilon_2/(12d)-E^2/(24d)=\Ph{1}/d$, using $E=\epsilon_1+\epsilon_2=\alpha_0+\alpha_1+\alpha_2$.

\emph{$d=1$ by direct localisation.} For $d=1$ the $\Tmax$-fixed loci of $\Mbar_g(\bbP^1,1)$ are indexed by $g_++g_-=g$: the zero section covered once, with a contracted component of genus $g_p$ attached at $p$ (absent if $g_p=0$), with fixed locus $\Mbar_{g_+,1}\times\Mbar_{g_-,1}$. The factors of the localisation formula of \cite{GP99} on the proper space $\Mbar_g(\bbP^1,1)$ are: at a vertex $p$ with a contracted component, $H^1$ and $H^0$ of $T_Z$ on that component, i.e.\ $\prod_a\Lamv(w^{(p)}_a)/\prod_aw^{(p)}_a$ over the five tangent weights $w^{(p)}_a$ at $p$; at the node, $\prod_aw^{(p)}_a$ from the normalisation sequence and $1/(w^{(p)}_1-\psi)$ from the smoothing of the node, $w^{(p)}_1=\pm\alpha_0$ being the tangent weight of $\bbP^1$ at $p$; on the edge, $H^0=H^1=0$ for each $\cO(-1)$, $H^0(\bbP^1,\cO_{-\epsilon_i})$ giving $1/(\epsilon_1\epsilon_2)$, and, of $H^0(\bbP^1,T\bbP^1)$ with weights $(\alpha_0,0,-\alpha_0)$, the part not cancelled by the automorphisms of the edge fixing its nodes, i.e.\ $1/w^{(p)}_1$ for each node. Collecting,
\begin{equation}
\label{eq:rig-loc}
\epsilon_1\epsilon_2\,\GW^{\Tmax}_{g,1}=\sum_{g_++g_-=g}V_+(g_+)\,V_-(g_-),\qquad V_\pm(0)\coloneqq 1,\qquad
V_\pm(h)\coloneqq\int_{\Mbar_{h,1}}\frac{\prod_{a=1}^{5}\Lamv_h(w^{(\pm)}_a)}{w^{(\pm)}_1\,(w^{(\pm)}_1-\psi)}\quad(h\geq 1),
\end{equation}
with $\Lamv_h(w)=\sum_i(-1)^i\lambda_iw^{h-i}$ and
\[
(w^{(+)}_a)_{a=1}^5=\wts_+=(\alpha_0,\alpha_1,\alpha_2,-\epsilon_1,-\epsilon_2),\qquad (w^{(-)}_a)_{a=1}^5=\wts_-=(-\alpha_0,\wB,\wA,-\epsilon_1,-\epsilon_2).
\]
At $g=0$, \eqref{eq:rig-loc} gives $\GW_{0,1}=1/(\epsilon_1\epsilon_2)$, the value of \cref{prop:rig-g0}. \draftnote{Formula \eqref{eq:rig-loc} is the standard localisation formula with the ingredients just listed; its normalisation agrees with the genus-zero value. The authors may wish to confirm the sign conventions of the node and edge factors against their reference for localisation on $\Mbar_g(\bbP^1,1)$.}

At $h=1$, $\Lamv_1(w)=w-\lambda_1$ and $\lambda_1^2=0$, so with $\Pi_\pm\coloneqq\prod_{a=1}^5w^{(\pm)}_a$ and $1/(w_1(w_1-\psi))=w_1^{-2}(1+\psi/w_1+\dots)$, the degree-one part of the integrand of $V_\pm(1)$ is $(\Pi_\pm/w_1^2)\big[\psi/w_1-\lambda_1\sum_{a=1}^5 1/w_a\big]$, and $\int_{\Mbar_{1,1}}\psi=\int_{\Mbar_{1,1}}\lambda_1=1/24$ gives
\[
V_\pm(1)=\frac{\Pi_\pm}{24\,w_1^2}\Big[\frac{1}{w_1}-\sum_{a=1}^{5}\frac{1}{w_a}\Big]=-\frac{\Pi_\pm}{24\,w_1^2}\sum_{a=2}^{5}\frac{1}{w_a}.
\]
At $p_+$, $\Pi_+=\alpha_0\alpha_1\alpha_2\epsilon_1\epsilon_2$ and $w_1=\alpha_0$, so
\[
V_+(1)=-\frac{1}{24\alpha_0}\Big[(\alpha_1+\alpha_2)\epsilon_1\epsilon_2-\alpha_1\alpha_2(\epsilon_1+\epsilon_2)\Big];
\]
at $p_-$, $\Pi_-=-\alpha_0\wA\wB\epsilon_1\epsilon_2$ and $w_1^2=\alpha_0^2$, so
\[
V_-(1)=\frac{1}{24\alpha_0}\Big[(\wA+\wB)\epsilon_1\epsilon_2-\wA\wB(\epsilon_1+\epsilon_2)\Big].
\]
Since $\wA+\wB-\alpha_1-\alpha_2=2\alpha_0$ and $\wA\wB-\alpha_1\alpha_2=\alpha_0(\alpha_0+\alpha_1+\alpha_2)$,
\[
\epsilon_1\epsilon_2\,\GW^{\Tmax}_{1,1}=V_+(1)+V_-(1)=\frac{\epsilon_1\epsilon_2}{12}-\frac{(\alpha_0+\alpha_1+\alpha_2)(\epsilon_1+\epsilon_2)}{24}
\]
on the five-torus, i.e.\ $N_1=1/12$ and $L_1=-(\alpha_0+\alpha_1+\alpha_2)/24$, without the input of \cite{BS}; on $\TCY$ this is $\epsilon_1\epsilon_2\GW_{1,1}=\Ph{1}$.
\end{proof}

\begin{rmk}
\label{rmk:rig-g1input}
The input of \cite{BS} in \cref{lem:rig-g1} determines the two numbers $N_d$ and $c$ for $d\geq 2$. A direct evaluation of the linear form $L_d=-(\alpha_0+\alpha_1+\alpha_2)/(24d)$ for all $d$ (as done for $d=1$ above), or a genus-one analogue of Theorem B, would remove it; neither is given here.
\end{rmk}

\subsection{The rigidity conjecture of \texorpdfstring{\cite{BS}}{[BS]} for the resolved conifold}

The conjecture of \cite{BS}, stated by content: let $X$ be a Calabi--Yau threefold, $Z=X\times\bbC^2$, and $\mgp{G}$ a connected reductive group acting on $Z$ with maximal torus $\mgp{T}$ such that the $\mgp{T}$-fixed loci of the moduli spaces of stable maps are proper, so that the $\mgp{G}$-equivariant invariants $\GW_g(Z,\mgp{G};\beta)$ are defined by localisation. If $\beta$ is a rigid curve class in $X$ (i.e.\ $\Mbar_g(X,\beta)$ is proper for every $g$) and the $\mgp{G}$-action on $Z$ is Calabi--Yau, then the fixed-class, all-genus series $\GW(Z,\mgp{G};\beta)\coloneqq\sum_{g\geq 0}\GW_g(Z,\mgp{G};\beta)$ depends only on $\epsilon_1$ and $\epsilon_2$, the weights of the two trivial directions:
\[
\epsilon_1\epsilon_2\,\GW(Z,\mgp{G};\beta)\in\bbQ\fps{\epsilon_-^2,\epsilon_+^2},\qquad \epsilon_\pm\coloneqq\tfrac12(\epsilon_1\pm\epsilon_2),
\]
i.e.\ it depends on $(\epsilon_1,\epsilon_2)$ symmetrically and evenly. The sum is over genera at fixed $\beta$, and $\beta=0$ is excluded: for non-compact $X$ the class $0$ is not rigid, $\Mbar_g(X,0)=\Mbar_g\times X$ being non-proper, and the fixed-genus series of \cite{BS} sums over non-zero classes only). \draftnote{BS label: conj: rigidity; defn: rigid; defn: refined GW series.}

\begin{cor}[the rigidity conjecture for the resolved conifold]\genusstatus{AI-proposed, unchecked}
\label{cor:rig-conj}
Let $X$ be the resolved conifold and $Z=X\times\bbC^2$ with its five-torus $\Tmax$ and Calabi--Yau four-torus $\TCY$ (\cref{def:set-ND}). Then for every $d\geq 1$ and every $g\geq 0$,
\[
\epsilon_1\epsilon_2\,\GW_g(Z,\TCY;d[\bbP^1])=\Pg\,d^{\,2g-3},
\]
so that
\[
\sum_{g\geq 0}\GW_g(Z,\TCY;d[\bbP^1])=\frac{1}{d}\Big(2\sinh\frac{d\epsilon_1}{2}\Big)^{-1}\Big(2\sinh\frac{d\epsilon_2}{2}\Big)^{-1}.
\]
In particular the conjecture holds for $X$ the resolved conifold, every non-zero class $d[\bbP^1]$, and $\mgp{G}=\TCY$; by contravariance it holds for every subtorus of $\TCY$ on which $R\pi_*f^*\Nb$ is moving on the fixed loci, so that the invariants are defined by localisation as in \cref{def:set-inv} (the weights to check are listed in the last paragraph of the proof of \cref{lem:rig-poly}), and for every connected $\mgp{G}$ whose maximal torus acts on $Z$ through such a subtorus (in particular for every Calabi--Yau subtorus of $\Tmax$ with this property).
\end{cor}

\begin{proof}
The identity $\epsilon_1\epsilon_2\GW_{g,d}=\Pg d^{2g-3}$ is \cref{thm:rig-main} for $g\geq 2$, \cref{lem:rig-g1} for $g=1$ and \cref{prop:rig-g0} for $g=0$; $\TCY$ is the maximal Calabi--Yau torus of $Z$ inside $\Tmax$, and its fixed loci are proper (\cref{lem:ide-fixed}), so the invariants are those of \cref{def:set-inv}. Summing over $g$,
\[
\sum_{g\geq 0}\frac{\Pg\,d^{2g-3}}{\epsilon_1\epsilon_2}=\frac{1}{d^3\epsilon_1\epsilon_2}\sum_{g\geq 0}\coef{t^{2g}}\Big(\frac{\epsilon_1\epsilon_2t^2}{4\sinh(\epsilon_1t/2)\sinh(\epsilon_2t/2)}\Big)\,d^{2g}
=\frac{1}{d}\Big(2\sinh\frac{d\epsilon_1}{2}\Big)^{-1}\Big(2\sinh\frac{d\epsilon_2}{2}\Big)^{-1},
\]
which is the formula of \cite{BS} for the subtori $\alpha_i=\epsilon_j$, now on all of $\TCY$. \draftnote{BS label: eq:resconGWref.} Each $\Pg$ is a symmetric polynomial in $\epsilon_1^2,\epsilon_2^2$, hence a polynomial in $\epsilon_+^2,\epsilon_-^2$. For a subtorus $\mgp{T}'\to\TCY$ the invariants are the images of those of $\TCY$ under the induced map of coefficient rings (contravariance, \cref{lem:rig-poly}), so the polynomial identity persists; for a connected $\mgp{G}$ with maximal torus $\mgp{T}'$ and Weyl group $W$, the $\mgp{G}$-invariant is the element of $R_{\mgp{G}}=R_{\mgp{T}'}^{W}$ whose image in $R_{\mgp{T}'}$ is the $\mgp{T}'$-invariant. \draftnote{BS label: lem: GW contravariance; the identification $R_{\mgp{G}}=R_{\mgp{T}}^{W}$ used in the definition of the invariants.}
\end{proof}

\emph{What is not implied.} The conjecture for group actions on $Z$ that are not conjugate into the five-torus $\Tmax$ is not addressed. The constant-map term $d=0$ is not part of the conjecture and is not determined here (\cref{rmk:rig-F0}). The inputs by genus: for $g\geq 2$ the result rests on Theorems A, B, C and \cite{GP99}; for $g=1$ it uses in addition the formula of \cite{BS} on the subtori $\alpha_i=\epsilon_j$ when $d\geq 2$; for $g=0$ it is \cite[p.~3]{FP00}.

\begin{rmk}[relation to the results of \cite{BS}]
\label{rmk:rig-BS}
(i) The formula of \cite{BS} for the resolved conifold on the subtori on which $\alpha_i=\epsilon_j$, for all $d>0$ and all genera, is proved there by degeneration of the base curve and the relative invariants of local curves. \Cref{cor:rig-conj} contains it for $g\geq 2$ and proves it without degeneration; it uses it at $g=1$ for $d\geq 2$ (\cref{lem:rig-g1}); it does not touch $g=0$. \draftnote{BS label: prop:ResConifold.}
(ii) The case in which the action on $X$ is also Calabi--Yau, i.e.\ $\epsilon_+=0$, $\epsilon_2=-\epsilon_1$, is proved in \cite{BS} for every rigid class $\beta$ in any $X$: there $\GW_g=(\epsilon_1\epsilon_2)^{g-1}N_{g,d}$ with $N_{g,d}$ the non-equivariant invariants, by Mumford's relation and the properness of $\Mbar_g(X,d)$. On $\TCY$ this is the hyperplane $E=0$ (so also $\alpha_0+\alpha_1+\alpha_2=0$), outside the frame of \cref{sec:lea}; \cref{thm:rig-main} reaches it by the polynomial continuation of its last paragraph, and there \eqref{eq:rig-main} reads, with $\epsilon\coloneqq\epsilon_1=-\epsilon_2$,
\[
\GW_{g,d}=-\frac{\Pg(\epsilon,-\epsilon)\,d^{2g-3}}{\epsilon^2}=\text{the }\epsilon^{2g-2}\text{ term of }-\frac{1}{d}\Big(2\sinh\frac{d\epsilon}{2}\Big)^{-2},
\]
since $\Pg(\epsilon,-\epsilon)=\coef{t^{2g}}\Ssh(\epsilon t)^2=\epsilon^{2g}\coef{t^{2g}}\Ssh(t)^2$ and $(2\sinh(d\epsilon/2))^{-2}=(d\epsilon)^{-2}\Ssh(d\epsilon)^2$: the genus-$g$ part of the unrefined formula for the resolved conifold. \draftnote{BS label: prop: rigidity unref limit; eq: res conifold unref.}
(iii) \cite{BS} verified the formula on $\TCY$ by computer for $d\leq 3$ up to contributions of genus five; \cref{cor:rig-conj} proves it for all $d$ and $g$ on $\TCY$.
(iv) \cite{BS} remark that the conjecture implies that the restriction to any admissible two-dimensional subtorus recovers the fully equivariant series. For the resolved conifold on $\TCY$ this is now a theorem.
\end{rmk}

\subsection{The flop}

\begin{prop}[the flop identity $\sF_g(\infty)=\sF_g(0)$]\genusstatus{AI-proposed, unchecked}
\label{prop:rig-flop}
Let $g\geq 2$ and let $\sF_g(\infty)$ denote the value of $\sF_g$ at $\sQ=0$ (which exists by \cref{lem:rig-inf}).
\begin{enumerate}
\item $\sF_g(\infty)=\sF_g(0)$ in $\KCY$.
\item At $\sQ=0$ all $\dd/\dd t$-derivatives vanish, so $\rr_{1,i}=\Lv_i\rr_{0,i}$ with $(\Lv_1,\Lv_2)=(\alpha_1,\alpha_2)$; the rows $\rZ_{0,i}=\hnw\rpt_{0,i}$ of $\RZ$ take the values
\[
\rZ_{0,1}(z;\infty)=h_{\wts'_-}(z),\qquad \rZ_{0,2}(z;\infty)=h_{\wts'_+}(z),
\]
\[
\wts'_-\coloneqq(\alpha_2-\alpha_1,\alpha_1,\wB,-\epsilon_1,-\epsilon_2),\qquad
\wts'_+\coloneqq(\alpha_1-\alpha_2,\alpha_2,\wA,-\epsilon_1,-\epsilon_2),
\]
and $\Delta_1(\infty)=\eul(\wts'_-)$, $\Delta_2(\infty)=\eul(\wts'_+)$. Here $\wts'_\pm$ are the tangent weights at the two fixed points of the flop $\Zfl=\Xfl\times\bbC^2$ of $Z$, $\Xfl$ the other small resolution of the conifold.
\item Consequently the constant-map terms of $Z$ and of $\Zfl$ agree:
\[
\sum_{p=\pm}\int_{\Mbar_g}\prod_{a=1}^{5}\frac{\Lamv_g(w^{(p)}_a)}{w^{(p)}_a}\quad\text{is the same for }(\wts_+,\wts_-)\text{ and for }(\wts'_+,\wts'_-),
\]
an identity among constant-map Hodge integrals for which no direct proof is given here.
\end{enumerate}
\end{prop}

\begin{proof}
(i) By \eqref{eq:rig-main} and \cref{lem:rig-li}, $\sF_g(\infty)=\sF_g(0)+\Pg\Li_{3-2g}(\infty)/(\epsilon_1\epsilon_2)=\sF_g(0)$.

(ii) By \cref{lem:rig-inf}, every factor is a power series in $\sQ$ and $\dd/\dd t=-\sQ\,\dd/\dd\sQ$ kills the value at $\sQ=0$; so at $\sQ=0$ the operators $\Dop_i=\dd/\dd t-\Delta_i'/(2\Delta_i)$ vanish, $\rr_{1,i}=\Lv_i\rr_{0,i}$ and $\Lv_i(0)=\alpha_1,\alpha_2$ by the proof of \cref{lem:rig-inf}. The values of the rows: the point-frame row $\rpt_{0,i}=\sum_mz^m\wick_m(\cW_2(U_i),\cW_3(U_i),\dots)$ (\cref{def:set-rows}) has $z$-coefficients that are polynomials in the $\cW_k(U_i)$, $k\geq 3$, and $1/\cW_2(U_i)$, each a power series in $\sQ$ by \cref{lem:rig-inf}, so its value at $\sQ=0$ is obtained by replacing each $\cW_k(U_i)$ by its value there. These values are the derivatives of a limiting phase. At colour $1$: writing $1-\cQ U=-\cQ U\,(1-\sQ/U)$, the superpotential $\cW(U)=-\alpha_2\log U+\wA\log(1-U)+\alpha_1\log(1-\cQ U)$ of \cref{lem:set-crit} becomes
\[
\cW(U)=(\alpha_1-\alpha_2)\log U+\wA\log(1-U)+\alpha_1\log(1-\sQ/U)+(\text{terms independent of }U),
\]
and $\thU\log(1-\sQ/U)=\sQ/(U-\sQ)$, so $\thU^k\log(1-\sQ/U)$ is $O(\sQ)$ for every $k\geq 1$ at $U=U_1(\sQ)$, since $U_1(0)=(\alpha_1-\alpha_2)/\wB\neq 0$. Hence $\cW_k(U_1)(\sQ=0)$, $k\geq 2$, are the derivatives $\thU^k$ of the phase $a\log U+b\log(1-U)$ with $(a,b)=(\alpha_1-\alpha_2,\wA)$, at $U=U_1(0)=(\alpha_1-\alpha_2)/\wB=a/(a+b)$, its critical point ($a+b=\wB$). At colour $2$: with $V=\cQ U$, $\thU=\theta_V$, $-\alpha_2\log U=-\alpha_2\log V+\alpha_2\log\cQ$, and $\wA\log(1-U)=\wA\log(1-\sQ V)$ with $\theta_V\log(1-\sQ V)=-\sQ V/(1-\sQ V)=O(\sQ)$, as are all higher $\theta_V$-derivatives; so
\[
\cW=-\alpha_2\log V+\alpha_1\log(1-V)+\wA\log(1-\sQ V)+(\text{terms independent of }V),
\]
and $\cW_k(U_2)(\sQ=0)$, $k\geq 2$, are the $\theta_V^k$-derivatives of $a\log V+b\log(1-V)$ with $(a,b)=(-\alpha_2,\alpha_1)$ at $V=V_2(0)=-\alpha_2/(\alpha_1-\alpha_2)=a/(a+b)$, again its critical point. In both cases $ab(a+b)\neq 0$ in $\KCY$ ($(\alpha_1-\alpha_2)\wA\wB\neq 0$ and $\alpha_2\alpha_1(\alpha_1-\alpha_2)\neq 0$) and the second derivative is $-a(a+b)/b\neq 0$ (the values computed in the proof of \cref{lem:rig-inf}). By \cref{lem:row-betaphase}, the normalised expansion of the phase $a\log U+b\log(1-U)$ in $y=\log U$ at its critical point is $h_{(-a,-b,a+b)}(z)$. Therefore
\[
\rpt_{0,1}(z;\infty)=h_{(\alpha_2-\alpha_1,\,-\wA,\,\wB)}(z),\qquad \rpt_{0,2}(z;\infty)=h_{(\alpha_2,\,-\alpha_1,\,\alpha_1-\alpha_2)}(z).
\]
Multiplying by $\hnw=h_{(\alpha_1,\wA,-\epsilon_1,-\epsilon_2)}$ and cancelling the pairs $v,-v$ ($h_{u+v}=h_uh_v$, $h_{-v}=h_v^{-1}$),
\begin{align*}
\rZ_{0,1}(z;\infty)&=h_{(\alpha_1,\wA,-\epsilon_1,-\epsilon_2,\,\alpha_2-\alpha_1,-\wA,\wB)}=h_{(\alpha_2-\alpha_1,\,\alpha_1,\,\wB,\,-\epsilon_1,-\epsilon_2)}=h_{\wts'_-},\\
\rZ_{0,2}(z;\infty)&=h_{(\alpha_1,\wA,-\epsilon_1,-\epsilon_2,\,\alpha_2,-\alpha_1,\alpha_1-\alpha_2)}=h_{(\alpha_1-\alpha_2,\,\alpha_2,\,\wA,\,-\epsilon_1,-\epsilon_2)}=h_{\wts'_+}.
\end{align*}
By \eqref{eq:rig-Dinf}, $\Delta_1(\infty)=-\alpha_1\wB(\alpha_1-\alpha_2)\epsilon_1\epsilon_2=(\alpha_2-\alpha_1)\alpha_1\wB\epsilon_1\epsilon_2=\eul(\wts'_-)$ and $\Delta_2(\infty)=(\alpha_1-\alpha_2)\alpha_2\wA\epsilon_1\epsilon_2=\eul(\wts'_+)$.

The geometry: the conifold $\{xy=zw\}$ has the four coordinate weights $\alpha_1,\wA,\alpha_2,\wB$ with $\alpha_1+\wA=\alpha_2+\wB=E$, the two monomials of the relation pairing $\{\alpha_1,\wA\}$ and $\{\alpha_2,\wB\}$. The resolution $X$ has the normal weights $(\alpha_1,\alpha_2)$ at $p_+$ and $(\wB,\wA)$ at $p_-$, with tangent weights $\pm\alpha_0$ along $\bbP^1$; the other small resolution $\Xfl$ has the normal weights $(\alpha_2,\wA)$ at one fixed point and $(\alpha_1,\wB)$ at the other, with tangent weights $\pm(\alpha_1-\alpha_2)$ along its exceptional curve. So $\wts'_\pm$ are the tangent weights of $\Zfl=\Xfl\times\bbC^2$ at its two fixed points, and the Calabi--Yau condition is preserved: the sum of the first three entries of $\wts'_\pm$ is $E=\epsilon_1+\epsilon_2$. \draftnote{The description of the flop and its weights is stated as in the source; the authors may wish to add a reference or a line of justification for the weights of the two small resolutions.}

(iii) At $\sQ=0$ the data of the graph sum of $Z$ are: constant rows $\rZ_{0,i}=h_{\wts'_\mp}$ with vanishing $t$-derivatives, $\rr_{1,i}=\Lv_i\rr_{0,i}$, and $\Delta_i=\eul(\wts'_\mp)$. These are the data of the graph sum of $\Zfl$ at its own point $\cQ'=0$ (\cref{lem:ide-norm} and \cref{prop:ide-const} read for $\Zfl$, whose tangent weights are $\wts'_\pm$, all non-zero in $\KCY$, with $\cQ'=0$ semisimple because $\alpha_1-\alpha_2\neq 0$), where the graph sum takes the constant-map value
\[
\sF^{\Zfl}_g(0)=\sum_{p=\pm}\eul(\wts'_p)^{g-1}\int_{\Mbar_g}\prod_{a=1}^{5}\frac{\Lamv_g(w'^{(p)}_a)}{(w'^{(p)}_a)^{g}}=\sum_{p=\pm}\int_{\Mbar_g}\prod_{a=1}^{5}\frac{\Lamv_g(w'^{(p)}_a)}{w'^{(p)}_a}.
\]
Hence $\sF_g(\infty)=\sF^{\Zfl}_g(0)$, and with (i) and \cref{rmk:rig-d0} this is the identity (iii). \draftnote{What is written in the source for this step is the sentence that the graph sum at $\sQ=0$ is the $\cQ'=0$ value of the frame-B sum of $\Zfl$, by the constants lemma read for $\Zfl$. The authors should check the identification of the factors: at $\sQ=0$ the eigenvalues of $Z$ are $(\alpha_1,\alpha_2)$ and the pairing is $\eta^{-1}$ of $Z$, while at $\cQ'=0$ the eigenvalues of $\Zfl$ are those of its own flat basis; the factors $\Dil_i$, $\Edge_{ij}$ depend on the eigenvalues and the pairing only through the rows $\rr_{0,i}$, the norms $\Delta_i$ and the orthogonality $\langle(1,\Lv_i),(1,\Lv_j)\rangle=\delta_{ij}\Delta_i$ of the eigen-covectors, which holds in both cases; this reduction is not spelled out in the source.}
\end{proof}

\begin{rmk}[the constant term]
\label{rmk:rig-F0}
Nothing in this section determines $\sF_g(0)=\GW_{g,0}$, the rational function of \cref{rmk:rig-d0}, of degree $2g-2$ in the weights. It is not part of the rigidity conjecture, and there is no reason for it to be independent of $\alpha$; the finite-order values recorded in \cref{sec:evi-flop} show that it is not. On the fibrewise torus $\TF$ it equals $-(\epsilon_1^{2g-2}+\epsilon_2^{2g-2})\,|B_{2g}B_{2g-2}|/\big(4g(2g-2)(2g-2)!\big)$, with $B_{2g}$ the Bernoulli numbers, from the localisation formula of \cref{rmk:rig-d0} and \cite[Theorem~4]{FP00}. \draftnote{The value on $\TF$ is quoted from the sources, where it is derived from [GP99] and [FP00, Theorem 4]; the derivation is not reproduced here.} \Cref{thm:rig-main} says nothing about $\sF_g(0)$ beyond \cref{prop:rig-flop}.
\end{rmk}

\subsection{The three-torus \texorpdfstring{$\alpha_0=0$}{alpha0 = 0}}

Let $\Tthree\coloneqq\{\alpha_0=0\}\cap\TCY$: the weights are $(0,\alpha_1,\alpha_2,\epsilon_1,\epsilon_2)$ with $\alpha_1+\alpha_2=\epsilon_1+\epsilon_2=E$, and $\alpha_1,\alpha_2,\epsilon_1$ are free; the function field is $\Kthree=\bbQ(\alpha_1,\alpha_2,\epsilon_1)$. Non-degeneracy (ND) reads $\alpha_1\alpha_2\epsilon_1\epsilon_2E\neq 0$ (here $\wA=\alpha_2$ and $\wB=\alpha_1$). $\Tthree$ contains the fibrewise torus $\TF$ ($\alpha_i=\epsilon_i$). On $\Tthree$ the point $\cQ=0$ is not semisimple (both roots of the critical equation degenerate, \cref{lem:set-limits}), so Theorem B is not available there directly; it is obtained by specialisation.

\begin{lem}[Theorem B on $\Tthree$]\genusstatus{AI-proposed, unchecked}
\label{lem:rig-3Tid}
Let $g\geq 2$. On $\Tthree$, with $\alpha_1\alpha_2\epsilon_1\epsilon_2E\neq 0$, the frame-B graph sum \eqref{eq:set-graphsum} formed with the rows $\rr_{0,i}=h_{(\alpha_1,\alpha_2,-E)}\rpt_{0,i}$, $\rr_{1,i}=(\Lv_i+z\Dop_i)\rr_{0,i}$ (i.e.\ with $\hnw$ specialised to $h_{(\alpha_1,\alpha_2,-\epsilon_1,-\epsilon_2)}$ and $\bw=(E,-\epsilon_1,-\epsilon_2)$) lies in $\Kthree(\cQ)$ and equals $\sF^{\Tthree}_g=\sum_{d\geq 0}\GW^{\Tthree}_{g,d}\cQ^d$, the series of the $\Tthree$-equivariant invariants in the localisation definition.
\end{lem}

\begin{proof}
The identity \eqref{eq:ide-main} holds in $\KCY(\cQ)$; we apply to it the ring map $\alpha_0\mapsto 0$ (on the subring of $\KCY$ of elements regular along $\alpha_0=0$, with values in $\Kthree$). Equivalently one may start from the identity at the five indeterminate weights (\cref{rmk:ide-alt}) and apply $\alpha_0\mapsto 0$ followed by $\epsilon_2\mapsto\alpha_1+\alpha_2-\epsilon_1$: the two composite specialisations agree. \draftnote{The source specialises the five-weight identity in two steps; the note's Theorem B is stated over $\KCY$, and the specialisation $\alpha_0\mapsto 0$ applied to it is the same ring map. Check that the wording of \cref{prop:ide-main} and \cref{rmk:ide-alt} supports this.} Three things are to be verified: the Gromov--Witten side is regular along $\alpha_0=0$ and specialises to the $\Tthree$-invariants; the graph side is regular along $\alpha_0=0$ at generic $\cQ$ and specialises to the graph sum of the statement; and the frame change persists.

\emph{The invariants, $d\geq 1$.} By \cref{lem:rig-poly}, $\GW_{g,d}\in(\epsilon_1\epsilon_2)^{-1}\bbQ[\alpha_0,\alpha_1,\alpha_2,\epsilon_1]$ is regular along $\alpha_0=0$, and its value there is the restriction to $\Tthree$ of the five-torus polynomial $\epsilon_1\epsilon_2\GW^{\Tmax}_{g,d}$, divided by $\epsilon_1\epsilon_2$. This restriction is $\GW^{\Tthree}_{g,d}$ by the contravariance of \cref{lem:rig-poly}: the $\Tmax$-equivariant integral over the proper $M=\Mbar_g(\bbP^1,d)$ restricts to the $\Tthree$-equivariant integral of the restricted class, and the latter is the localisation definition of the $\Tthree$-invariant, because the $\Tthree$-fixed locus of $\Mbar_g(Z,d)$ is $M$ (every compact curve lies in the zero section, which $\Tthree$ fixes pointwise since $\alpha_0=0$), with moving part $R\pi_*f^*\Nb$ of non-zero weights.
\draftnote{The source states the restriction of the invariants for $d\geq 1$ as a consequence of its polynomiality lemma without further comment, and spells out the argument only for $d=0$ (next paragraph); the sentence on the $\Tthree$-fixed locus for $d\geq 1$ follows that argument and should be checked.}

\emph{The constant-map term.} $\GW_{g,0}=\sum_p\int_{\Mbar_g}\prod_{a=1}^5\Lamv_g(w^{(p)}_a)/w^{(p)}_a$ (\cref{rmk:rig-d0}) is regular at $\alpha_0=0$: only the factors $\Lamv_g(\alpha_0)/\alpha_0$ (at $p_+$) and $\Lamv_g(-\alpha_0)/(-\alpha_0)$ (at $p_-$) have a pole there, a simple one, and the two residues cancel. Indeed $\Lamv_g(\alpha_0)/\alpha_0=\sum_i(-1)^i\lambda_i\alpha_0^{g-i-1}$ has the $\alpha_0^{-1}$ term $(-1)^g\lambda_g/\alpha_0$, and the other four factors at $p_+$ are regular at $\alpha_0=0$ (on $\TCY$ the weight $\epsilon_2=\alpha_0+\alpha_1+\alpha_2-\epsilon_1$ does involve $\alpha_0$; only its value at $\alpha_0=0$ enters the residue), so the residue of the $p_+$ term is
\[
(-1)^g\int_{\Mbar_g}\lambda_g\,\frac{\Lamv_g(\alpha_1)\Lamv_g(\alpha_2)\Lamv_g(-\epsilon_1)\Lamv_g(-\epsilon_2)}{\alpha_1\alpha_2\epsilon_1\epsilon_2};
\]
and $\Lamv_g(-\alpha_0)/(-\alpha_0)=\sum_i(-1)^i\lambda_i(-\alpha_0)^{g-i-1}$ has the $\alpha_0^{-1}$ term $(-1)^g\lambda_g/(-\alpha_0)=-(-1)^g\lambda_g/\alpha_0$, times the value at $\alpha_0=0$ of $\Lamv_g(\wB)\Lamv_g(\wA)\Lamv_g(-\epsilon_1)\Lamv_g(-\epsilon_2)/(\wA\wB\epsilon_1\epsilon_2)$, which is regular there with $\wA\to\alpha_2$, $\wB\to\alpha_1$: the residue of the $p_-$ term is minus that of the $p_+$ term. Moreover $\GW^{\Tthree}_{g,0}$ in the localisation definition equals this value at $\alpha_0=0$, by the same contravariance and localisation argument as for $d\geq 1$: the $\Tthree$-fixed locus of $\Mbar_g(Z,0)=\Mbar_g\times Z$ is the proper $\Mbar_g\times\bbP^1$ (the zero section is fixed pointwise by $\Tthree$; its normal weights $\alpha_1,\alpha_2,-\epsilon_1,-\epsilon_2$ are non-zero), with virtual class $e(\Hdg^\vee\boxtimes T_{\bbP^1})\cap[\Mbar_g\times\bbP^1]$ (the $\Tthree$-fixed part of the obstruction bundle $\Hdg^\vee\boxtimes T_Z$) and virtual normal bundle $\Nb\otimes(\cO-\Hdg^\vee)$; the same integral taken $\Tmax$-equivariantly over the proper $\Mbar_g\times\bbP^1$ lies in $\bbQ[\alpha_0,\alpha_1,\alpha_2,\epsilon_1,\epsilon_2][1/(\alpha_1\alpha_2\wA\wB\epsilon_1\epsilon_2)]$ (its denominators are the $\Tmax$-weights of $\Nb$ at $p_\pm$), restricts to $\Tthree$ by contravariance, and equals $\sum_p\int_{\Mbar_g}\prod_a\Lamv_g(w^{(p)}_a)/w^{(p)}_a$ by localisation at the two $\Tmax$-fixed points of $\bbP^1$ (\cite{GP99}; at $p$ the factor $\Lamv_g(\pm\alpha_0)/(\pm\alpha_0)$ is $e(\Hdg^\vee\otimes T_p\bbP^1)/e(T_p\bbP^1)$).

\emph{The graph side.} At generic $\cQ$ every factor is regular along $\alpha_0=0$: the two roots stay distinct ($\Disc|_{\alpha_0=0}=4\alpha_1\alpha_2\cQ+(\alpha_1-\alpha_2)^2\cQ^2\neq 0$ for generic $\cQ$); $U_i\notin\{0,1,1/\cQ\}$ (\cref{lem:set-crit}: $\Pq(1)=\wA(\cQ-1)$, $\Pq(1/\cQ)=\alpha_1(1-\cQ)/\cQ$, $\Pq(0)=-\alpha_2$, all non-zero at $\alpha_0=0$ for $\cQ\neq 1$); $\cW_2(U_i)$ and $\Delta_i$ are finite and non-zero; $\eta^{-1}$ is polynomial in the weights; and $\hnw\to h_{(\alpha_1,\alpha_2,-\epsilon_1,-\epsilon_2)}$ is finite. So the colour-summed total, written as a quotient of polynomials, has a denominator that does not vanish identically on $\alpha_0=0$, and its specialisation is the colour-summed point-frame sum formed with the rows $\hnw|_{\alpha_0=0}\rpt_{0,i}$. At $\alpha_0=0$, $\nw=(\alpha_1,\alpha_2,-\epsilon_1,-\epsilon_2)$ and $\bw=(E,-\epsilon_1,-\epsilon_2)$, so $\hnw\hw^{-1}=h_{(\alpha_1,\alpha_2,-\epsilon_1,-\epsilon_2,-E,\epsilon_1,\epsilon_2)}=h_{(\alpha_1,\alpha_2,-E)}$, independent of $\epsilon_1,\epsilon_2$ (pairs $v,-v$ cancel). The frame change $\RB=\RZ\hw^{-1}$ of \cref{prop:ide-frameB} applies verbatim ($\bw$ has non-zero entries and $\Hg_{\bw}=\hw.1$), which gives the frame-B form of the statement.
\end{proof}

\begin{cor}[Theorem C on $\Tthree$]\genusstatus{AI-proposed, unchecked}
\label{cor:rig-3T}
Let $g\geq 2$. On $\Tthree$, at the generic point and at every specialisation with $\alpha_1\alpha_2\epsilon_1\epsilon_2E\neq 0$, Theorem C (\cref{cor:reg-main}) holds for the frame-B graph sum of \cref{lem:rig-3Tid}: every coloured stable graph with an edge is regular at $t=0$, the single-vertex graph has polar part exactly $\gapg t^{2-2g}$ at colour $1$ and is regular at colour $2$, and the polar part of $\sF^{\Tthree}_g$ at $\cQ=1$ is exactly $\gapg t^{2-2g}=\beta_{2g-2}(\epsilon_1,\epsilon_2)t^{2-2g}$. No new row analysis is needed: the rows are the fibrewise rows rescaled, and the vanishing theorem is used at $\tauw=-\epsilon_1/E$.
\end{cor}

\begin{proof}
(1) \emph{The rows.} By \cref{cor:lea-a0}, at $\alpha_0=0$ with $\xc\coloneqq\alpha_1/\alpha_2$ and $\cc\coloneqq 1+\xc=E/\alpha_2$, the frame-B objects are the fibrewise frame-B objects at parameter $\xc$ read through the rescaling dictionary
\[
\rr_{0,i}(z;t)=\rr^{\TF}_{0,i}(z/\alpha_2;t),\quad \rr_{1,i}(z;t)=\alpha_2\,\rr^{\TF}_{1,i}(z/\alpha_2;t),\quad \Lv_i=\alpha_2\Lv^{\TF}_i,\quad \Delta_i=\kap\,\Delta^{\TF}_i,\quad \kap=\frac{\epsilon_1\epsilon_2\alpha_2^4}{\alpha_1},
\]
\[
\Dil_i(d)=\alpha_2^{1-d}\Dil^{\TF}_i(d),\quad \Edge_{ij}(m,n)=\kap\,\alpha_2^{-1-m-n}\Edge^{\TF}_{ij}(m,n),\quad \rho=\rho^{\TF},\quad \lf_j(\psi)=\alpha_2^{-1}\lf^{\TF}_j(\psi/\alpha_2),
\]
the pairing being $\kap$ times the fibrewise pairing at $(z/\alpha_2,w/\alpha_2)$. The frame-B constant is $\hnw\hw^{-1}=h_{(\alpha_1,\alpha_2,-E)}$, independent of $\epsilon_1,\epsilon_2$ (\cref{lem:rig-3Tid}).

(2) \emph{Theorem A on $\Tthree$.} Theorem A (\cref{thm:lea-main}) holds at $\alpha_0=0$ directly: its proof never uses the point $\cQ=0$, and the constant $\hnw$ enters it only through the multiset identity of \cref{lem:set-hmult} (remarks after \cref{cor:lea-a0}; the statements that need $\alpha_0\neq 0$ in their citations extend to $\alpha_0=0$ by rationality once the factors are defined there, which (1) does). Alternatively, by (1), each statement (F1)--(F3), (L1)--(L3) at $\alpha_0=0$ is the corresponding fibrewise statement at $\xc$ ($\xc\neq 0,-1$ because $\alpha_1E\neq 0$) read through the dictionary: under $z\mapsto z/\alpha_2$ the fibrewise leaves $1/(\cc-j\psi)$ become $\alpha_2/(E-j\psi)=\alpha_2\lf_j$, the spans $\Lsp_{k+1}$ are preserved, $\rho=(\Delta_1t^2/(-\epsilon_1\epsilon_2E^3))^{-1/2}$ equals the fibrewise $\rho$ (since $\Delta_1t^2/(-\epsilon_1\epsilon_2E^3)=\kap\Delta^{\TF}_1t^2/(-\epsilon_1\epsilon_2\alpha_2^3\cc^3)=\Delta^{\TF}_1t^2/(-\xc\cc^3)$), and the memberships hold with the same coefficients up to monomials in $\alpha_2$.

(3) \emph{Theorem C on $\Tthree$.} The proof of \cref{thm:reg-main} applies verbatim. The vertex class is $\Hg_{\bw}$ with $\bw=(E,-\epsilon_1,-\epsilon_2)$, which is not the fibrewise class (that would be $(E,-\alpha_1,-\alpha_2)$ rescaled) but is the class to which the vanishing theorem \cref{thm:reg-L} applies at $\tauw=-\epsilon_1/E$, $\st=-\epsilon_1\epsilon_2/E^2\neq 0$, and for which \cref{thm:reg-VI} gives the value $(-\epsilon_1\epsilon_2E^3)^{g-1}\ZZ[g,0]=\beta_{2g-2}(\epsilon_1,\epsilon_2)=(2g-3)!\Pg/(\epsilon_1\epsilon_2)$ (\cref{lem:reg-beta}). The decoupling makes this consistent: the rows and factors carry $(\alpha_1,\alpha_2)$ only, the vertex class carries $(\epsilon_1,\epsilon_2)$ only, and the proof of \cref{thm:reg-main} needs no relation between them, because $\Mirr(h)$ is a vanishing statement, $\Xiz_h(\msM)=0$ for each multiset separately, so the $\alpha$-dependent leaf coefficients need satisfy no relation. The identification of the graph sum with $\sF^{\Tthree}_g$, and hence the statement about the polar part of $\sF^{\Tthree}_g$, is \cref{lem:rig-3Tid}.
\end{proof}

\begin{rmk}
\label{rmk:rig-3Trig}
Rigidity on $\Tthree$ is not a separate statement: it is the restriction of \cref{thm:rig-main} to $\alpha_0=0$ by polynomiality (the last paragraph of its proof). The $\alpha_0$-adic argument of \cref{lem:rig-disc} is not available on $\Tthree$ itself, where $\alpha_0=0$; so \cref{cor:rig-3T} alone does not give \eqref{eq:rig-main}.
\end{rmk}


\appendix
% a_evidence. Appendix file; label prefix: evi.
\section{Finite-order checks}
\noindent\genusstatus{numerically supported} The finite-order checks recorded here are computations (exact arithmetic unless stated); they support the statements at the orders and genera given, and are not proofs. The code and outputs are in \texttt{code/} of this project.

\label{sec:evi}

\subsection{Scope}
\label{sec:evi-scope}

Nothing in this appendix is a proof step: the proofs of the note are complete without it, and no statement of \cref{sec:set,sec:row,sec:lea,sec:ide,sec:reg,sec:rig} refers to it. What follows records the finite-order computations that were made while the proofs were written, as evidence that the statements are correct at the orders and weight points tested, and as a record of what was tested and what was not. The arithmetic is exact throughout, either in $\bbQ$ (rational weights, truncated series) or in a prime field $\bbF_p$ with $p=2^{127}-1$ or $p=2^{89}-1$ (random residues), unless floating point is said; a zero test in $\bbF_p$ is probabilistic. Several independent implementations were used, written separately; where a statement was checked by more than one of them this is said. A control is a deliberately wrong input (a wrong class, a missing constant, a point off $\TCY$) which must make a check fail; the controls that failed as required are listed, because they show that the checks have the power to detect an error. Weight points are written $(\alpha_0,\alpha_1,\alpha_2,\epsilon_1,\epsilon_2)$, or $(\alpha_0,\alpha_1,\alpha_2,\epsilon_1)$ with $\epsilon_2=E-\epsilon_1$ understood when the point lies on $\TCY$.

\subsection{Identities and rows}
\label{sec:evi-rows}

\emph{Identities of \cref{sec:set,sec:row,sec:lea}.} At random rational points of $\TCY$ with $\alpha_0\neq 0$ and random rational $s$, $\tau$, $z$: the amplitude identity \cref{lem:row-amp} multiplied by $(1-\tau s)^3$ (30 points); the operator identity \cref{lem:lea-op} applied to $\tau^n$, $-3\leq n\leq 6$ (12 points); and, at 25 rational points where $\delta^2$ is a rational square, for both roots of the critical polynomial: $\Ns(s)=0$; $\Lv=\wA/(\tau s)-\alpha_0$ from the superpotential; $\Lv_1+\Lv_2$, $\Lv_1\Lv_2$ and $\tau(\Lv_1-\Lv_2)=\delta$ (\cref{lem:set-eig}); $\alpha_1\alpha_2-E\Lv_2=-\wB\tau s^\ast\Lv_2$ (\cref{lem:set-Delta2}); the two forms of \cref{lem:set-N}(ii); the Hessian $\phi_\tau''(s)=-\delta/(s(1-s)U)$ and $\cW_2(U)=U^2\phi_\tau''/\tau^2$ (\cref{lem:row-hess}); the three expressions $\Delta=\epsilon_1\epsilon_2(\Lv_1-\Lv_2)(\alpha_1\alpha_2-E\Lv_1)=\alpha_1\wA\epsilon_1\epsilon_2\cW_2=-\epsilon_1\epsilon_2E^3\alg^2/\tau^2$ (\cref{prop:set-idem}, \cref{prop:lea-key}); the eikonal identity $(1-\tau)\Eik=\wB\Ns(s)$ and $\Ptr=-\Ns'(s)$ (\cref{lem:row-eik}); $\langle\chiv_1,\chiv_2\rangle=0$ and $\langle\chiv_1,\chiv_1\rangle=\Delta_1$; the three forms of $\Thf$; $\Delta_2=\epsilon_1\epsilon_2\wB\delta s^\ast\Lv_2$; the self-adjointness of $\Aop$ (\cref{prop:lea-A}) on random covectors; $\Lv_1-\dphi'=E/\tau-\wA$ (\cref{lem:row-canon}); the derivatives $\dd s^\ast/\dd\tau=s^\ast\Lv_2/\delta$ and $(\log\alg^2)_\tau=2\Rtr/\delta$ (\cref{prop:row-transport}); the identities involving $\uq$; $\eta\eta^{-1}=1$; the multiset identity of \cref{lem:set-hmult} as an identity of power sums $\sum_av_a^{1-2k}$, $k=1,2,3$; and the first-order values of \cref{rmk:row-first} ($\lrem(0)$, $(s^\ast)'(0)$, $\Delta_2'(0)=0$, the first coefficient of $\alg$). All hold. A second set of exact tests, at about 50 random rational points, covered the amplitude identity, the operator identity, $\eta^{-1}$, the identities involving $\uq$, the discriminant of $\Disc$ (\cref{lem:set-disc}(ii)), the rescaling of the rows at $\alpha_0=0$ (\cref{cor:lea-a0}) and the power-sum form of the multiset identities; a third, at 50 and 30 random points, covered the identities of \cref{lem:set-disc}(i)--(iii) ($\Disc=\cQ^2[(E-\alpha_0\tau)^2-4\alpha_1\alpha_2\tau]$, its discriminant $16\alpha_1\alpha_2\wA\wB$, $\Disc(1)=E^2$), the product identity $(\alpha_1\alpha_2-E\Lv_1)(\alpha_1\alpha_2-E\Lv_2)=\alpha_1\alpha_2\wA\wB$, and, modulo the critical polynomial $\Pq(U)$: $U\cW_U(1-U)(1-\cQ U)=\Pq(U)$, $\mathrm{disc}_U\Pq=\Disc$, the quadratic equation satisfied by $\Lv(U)$, the equality of the three forms of $\Delta$, $\Pq(1)=\wA(\cQ-1)$, $\Pq(1/\cQ)=\alpha_1(1-\cQ)/\cQ$, and $\Pq|_{\alpha_0=0}=\alpha_2(\xc\cQ U^2+\cQ(1-\xc)U-1)$ (\cref{lem:set-crit}).

\emph{The rows at the singular idempotent.} At the two points $(5/3,-4/7,13/6,9/11)$ and $(-2/5,7/4,3,1/6)$ of $\TCY$, through orders $\tau^6z^9$ and again $\tau^8z^{11}$: the point-frame row in $y=\log U$, computed from the rescaled derivatives $\wrk_k(\tau)$ through the Wick polynomials (partitions of $2n$ into $n-m$ parts $\geq 3$), multiplied by $h_{(\alpha_1,\wA,-E)}(z)$, equals $(E/(E+z))\Fhyp\Ephi\alg$ coefficientwise (\cref{thm:row-main}), and the negative powers of $z$ on the right cancel; $\rpt_{0,1}(z;0)=h_{(E,-\alpha_1,-\wA)}E/(E+z)$ (\cref{prop:row-beta}); the $z$-degree test of (L1): $\coef{t^k}(z\Gk)\prod_{j\leq k+1}(E+jz)$ has $z$-degree exactly $k$ for $1\leq k\leq 6$, resp.\ $k\leq 8$ (\cref{thm:lea-L1}, \cref{cor:lea-top}); the closed form of the leaf coefficients $\Aleaf_j$ reproduces $\coef{t^k}(z\Gk)-E\delta_{k0}$ for $k\leq 5$, resp.\ $k\leq 7$ (\cref{thm:lea-Aj}), with $\coef{t^1}(z\Gk)=(\uq/2)(1/(E+z)-1/(E+2z))$ and $\rho_1=\uq/(2E^2)$; (F2): $\rho\rr_{1,1}=E/t+\alpha_1\alpha_2/(E+z)+O(t)$ for all $z$-orders tested, with $\rr_{1,1}$ built from the $H$-row and $\Delta_1=\alpha_1\wA\epsilon_1\epsilon_2\wrk_2/\tau^2$ (\cref{cor:row-F2}, \cref{prop:lea-sigma}); row 2 of the quantum differential equation at $U_1$ (\cref{prop:lea-row2}); $\Kone=\langle\vv_1(z),\vv_1(-z)\rangle/\Delta_1=1$ (\cref{prop:lea-K1}); $\Edge_{11}(0,0)$ and $\Edge_{11}(1,0)$ have no negative powers of $t$, with values $-\epsilon_1\epsilon_2\alpha_1\alpha_2\wA\wB/(2E^2)$ and $-\epsilon_1\epsilon_2\alpha_1\alpha_2\wA\wB/(2E^3)$ at $t=0$; $\Thf(0)=0$ and $\Thf_1=-\epsilon_1\epsilon_2\alpha_1\alpha_2\wA\wB/E^2$; the flow $\dd/\dd t\,(t^2\rho^2\Edge_{11})=\Thf\,\Dg\,\Dg$ for $m+n\leq 3$ (\cref{thm:lea-flow}); the bidegree bound of (L3) for $k\leq 2$ (\cref{thm:lea-L3}); and, at the point $(0,5/2,-3/7,4/9)$, the rescaling of \cref{cor:lea-a0}: the frame-B row at $\alpha_0=0$ equals the fibrewise frame-B row at $\xc=\alpha_1/\alpha_2$ with $z\mapsto z/\alpha_2$, and $h_{(\alpha_1,\alpha_2,-E)}(z)=h_{(\cc,-\xc,-1)}(z/\alpha_2)^{-1}$. All hold at both orders.

\emph{An independent computation to higher order.} A second implementation, in $\bbF_p$ with $p=2^{127}-1$ and $p=2^{89}-1$ and truncated Laurent series in $t$ to $t^{14}$ and $z^{23}$, tested at nine weight points of $\TCY$, including points with $\alpha_1=\alpha_2$, with $\uq=0$ and with $\epsilon_1=\epsilon_2$: the formula \eqref{eq:lea-key}; (F1)--(F3); (L1) for $k\leq 13$; (L2) for $k\leq 9$; (L3) for $k\leq 9$; the flow of \cref{thm:lea-flow}; the quantum differential equation and symplecticity at both idempotents (\cref{thm:lea-qde}); the values at the regular idempotent (\cref{prop:lea-regrows}); \cref{cor:lea-a0} at $\alpha_0=0$; and the statement of \cref{rmk:lea-cy} at points off $\TCY$. All passed, and four controls failed as required: the class $(E+1,-\epsilon_1,-\epsilon_2)$ in place of $\bw$; the omission of $\hnw$; $\wA\mapsto\alpha_2$ in the rows; and $\lf_j=1/(E+j\psi)$ in place of $1/(E-j\psi)$. The same computation observed the two instances of sharpness predicted by \cref{cor:lea-top}: at $\uq=0$ the top leaf $\lf_{k+1}$ is absent from $\coef{t^k}(\rho\qt)$ for odd $k$, and at $E=9\alpha_2/7$ the coefficient of $\lf_9$ in $\coef{t^8}(\rho\qt)$ vanishes (for odd $k$ and $E=2\alpha_2$, resp.\ $E=2\alpha_1$, the product $\Psh_k$ contains the factor $(k+1)\alpha_2-\tfrac{k+1}{2}E$, resp.\ $(k+1)\wA-\tfrac{k+1}{2}E$; $\Psh_8$ contains $9\alpha_2-7E$).

\emph{Further implementations.} Three further implementations, written for the exploration that preceded the proofs, tested at five points of $\TCY$, namely
\begin{gather*}
(1,2,4,3,4),\quad(3/5,-7/3,11/2,5/7,641/210),\\
(-2,9,1/3,13/4,49/12),\quad(2/3,5/7,-9/4,3/11,-1055/924),
\end{gather*}
and the fibrewise point $(0,5/3,2/7,5/3,2/7)$; at the points
\begin{gather*}
(3,5,3,5),\ (5,11,5,11),\ (11,5,5,11),\ (3,13,5,11),\ (8,3,5,11),\\
(3,4,5,11),\ (2,7,5,11),\ (7,13,2,29),\ (17,-5,19,-8),\ (3/7,-11/5,13/9,17/4),
\end{gather*}
written as $(\alpha_1,\alpha_2,\epsilon_1,\epsilon_2)$ with $\alpha_0=E-\alpha_1-\alpha_2$ (the fifth has $\uq=0$); and at random residues: the formula \eqref{eq:lea-key} through $t^6z^{12}$ and $t^7z^{12}$; $\rpt_{0,1}(z;0)=h_{(E,-\alpha_1,-\wA)}E/(E+z)$ and $\hnw\rpt_{0,1}(z;0)=\hw E/(E+z)$ (\cref{cor:row-F2}); $\rpt_{0,2}(z;0)=h_{(\alpha_2,-E,\wB)}$ and $\vv_2(z;0)$ proportional to $(1,\alpha_1\alpha_2/E)$ with the scalar $\hreg$ (\cref{prop:lea-regrows}); the closed form of $\Aleaf_j$ for $k\leq 6$ and the exact degree $k$ of $\coef{t^k}(z\Gk)\prod(E+jz)$ for $k\leq 7$, sharp except where $\Psh_k=0$; (L3) for $k\leq 5$ (78 equations, the anti-diagonal of \cref{cor:lea-anti} being needed); (L2) for $j=1,2$, $b\leq 3$, $k\leq 5$, and the exact degrees for $j=2$, $k\leq 3$, $b\leq 2$ (\cref{prop:lea-L2}); $\Edge_{21}(b,a)=\Edge_{12}(a,b)$ (\cref{prop:lea-transp}); all $\Edge_{ij}(m,n)$ with $m+n\leq 8$ and all $\Dil_i(d)$ of $t$-order $\geq 0$ (F3); $\langle\vv_i(z),\vv_j(-z)\rangle=\delta_{ij}\Delta_i$ through $z^{13}$; row 2 of the quantum differential equation at both colours through $z^{11}$; the flow for $m+n\leq 11$; the hand values $\rho_1=\uq/(2E^2)$, the subleading term of $\Delta_1$, $\coef{t^1}(\rho\qt)=(\uq/2)(\lf_1-\lf_2)$, the value $\alpha_1\alpha_2/(E+z)$ of the $H$-row, $\Edge_{11}|_{t=0}=-(\epsilon_1\epsilon_2\alpha_1\alpha_2\wA\wB/2)\lf_1(\psi)\lf_1(\psi')$ and the rank-one $\Edge_{12}|_{t=0}$ (\cref{rmk:lea-low}, \cref{rmk:lea-legs}); the constants at $\cQ=0$ of \cref{lem:ide-norm} at both roots through $z^9$ at three points; and the $\alpha_0=0$ rescaling to the fibrewise rows. Controls: the class $(E,-\alpha_1,-\wA)$ fails at every point with $\alpha_i\neq\epsilon_i$ (it agrees only on $\TF$), the scale $E+1$ gives the full degree $13$ in the degree test, and the point frame and the class $(E,-\alpha_1,-\wA)$ both give $\Edge_{11}$ of order $t^{-2}$. At the point $(2/9,11/4,-7/6,3/5)$ the degree test of (L1) was extended: $\coef{t^k}(\rho\qt)\prod_j(E-j\psi)$ has degree exactly $k$ for $k=1,\dots,10$ through $\psi^{19}$ ($k=0$ gives the constant $E$), and the bidegree bound of (L3) holds for $k\leq 6$ and $m+n\leq 14$ with no violation in 71 to 119 tests per $k$.

\subsection{Graph regularity}
\label{sec:evi-graphs}

The frame-B graph sum \eqref{eq:set-graphsum} was evaluated per coloured stable graph, summing over decorations, in $\bbF_p$ with Laurent precision in $t$, by three implementations. The counts of stable graphs of genus $2,3,4,5$ without legs are $7,42,379,4555$; the counts of coloured stable graphs with at least one edge are $20,350,8596,284806$.

\emph{Regularity of the coloured graphs with an edge (\cref{cor:reg-main}(a)).} Genus 2: $20$ of $20$ regular at $t=0$ at the point $(1,2,4,3,4)$ and on $\TF$. Genus 3: $350$ of $350$ at four points of $\TCY$. Genus 4: $8596$ of $8596$ at three points of $\TCY$. Genus 5: $284806$ of $284806$ at the single point $(3/5,-7/3,11/2,5/7,641/210)$. A second implementation found, at three generic points of $\TCY$, that of the $22$, $352$ and $8598$ coloured stable graphs of genus $2,3,4$ exactly one is singular, the single-vertex graph at colour $1$; the first implementation was re-run at the fresh points $(7/5,-3/8,13/6,11/9)$ and $(-5/4,19/7,2/3,17/5)$ with the same result. Special points of $\TCY$ were tested to exclude an artefact of a generic sample, all with the same outcome at genus $2,3,4$: $\uq=0$ at $(1,3,2,5/2)$; $\epsilon_1=\epsilon_2$ at $(3/2,-1/3,5/7,79/84)$; $\alpha_1=\alpha_2$ at $(2/3,5/4,5/4,1/7)$; $\alpha_1=\alpha_2$ with $\epsilon_1=\epsilon_2$ at $(2,3/5,3/5,8/5)$; $3\alpha_2=E$ at $(1,3,2,7/5)$; $\wA=\alpha_1$ at $(1,3,2,1/9)$; $\TF$ at $(0,5/3,2/7,5/3)$; a point with negative weights and small $E$, $(-3,1,5/2,7)$; $\alpha_1=-\alpha_2$ at $(5,7/3,-7/3,2)$; and $\alpha_0=\alpha_1$ at $(4/3,4/3,1/5,3)$.

\emph{The single-vertex graph (\cref{cor:reg-main}(b)).} In every run above, the colour-$1$ term has the single polar term $t^{2-2g}$ with coefficient $\gapg$, and the colour-$2$ term has no pole.

\emph{Tallies.} With the pole order of a decorated graph (a coloured graph with a choice of exponents and dilaton degrees) and the bound $\sum_v(2g_v-2)_+$ of \cref{thm:reg-main}, the tallies in frame B, written as (pole order, bound) followed by the number of decorated graphs, are, identically at all generic points tested:
\begin{align*}
\text{genus }2\colon&\quad\{(0,0)\colon 20,\ (2,2)\colon 7\};\\
\text{genus }3\colon&\quad\{(0,0)\colon 504,\ (0,4)\colon 1,\ (2,2)\colon 159,\ (4,4)\colon 29\};\\
\text{genus }4\colon&\quad\{(0,0)\colon 19912,\ (0,4)\colon 25,\ (0,6)\colon 1,\ (2,2)\colon 6069,\ (4,4)\colon 1069,\ (6,6)\colon 96\}.
\end{align*}
These are the fibrewise tallies. No decorated graph exceeds its bound and none has an odd leading order. The $159$ singular decorated graphs with an edge at genus 3 (and $7138$ at genus 4) lie in graphs with a vertex of genus $\geq 2$; decorated graphs of graphs with all $g_v\leq 1$ are individually regular, while at a vertex of genus $\geq 2$ only the decoration sum at fixed colouring is regular, as the proof of \cref{thm:reg-main} predicts. The polar coefficients of individual decorated graphs depend on $\alpha$: at fixed $(\epsilon_1,\epsilon_2)=(5,11)$ and twenty values of $\alpha$, the coefficient of $t^{2-2g}$ is independent of $\alpha$, that of $t^{3-2g}$ is $\uq$ times a constant, and at genus 3 the coefficients of $t^{-2}$ and $t^{-1}$ are polynomials in $(\uq,\alpha_0^2)$ of degrees $2$ and $3$, the latter divisible by $\uq$; no proof of these observations is given.

\emph{Totals.} The total minus $\Pg\Li_{3-2g}(\cQ)/(\epsilon_1\epsilon_2)$ has zero coefficients at all non-zero powers of $t$ from $t^{1-2g}$ through $t^5$ at genus $2,3,4$ and through $t^3$ at genus 5, and its $t^0$ coefficient equals the constant-map term of \cref{rmk:rig-d0} whenever $\alpha_0\neq 0$. At a generic point $\cQ_0\in\bbF_p$ away from $\cQ=1$, with $\cQ=\cQ_0\re^t$, the total equals $\Pg\Li_{3-2g}(\cQ)/(\epsilon_1\epsilon_2)+\sF_g(0)$ through $t^4$ at genus $2,3,4$ at $(7/5,-3/8,13/6,11/9)$ and at genus 3 at $(-5/4,19/7,2/3,17/5)$; the class $(E,-\alpha_1,-\wA)$ at genus 2 and the point frame at genus 3 give the same total, as a frame change must.

\emph{Controls.} Earlier stored data were reproduced with zero mismatches: the point-frame terms on $\TCY$ at two generic points ($17$ and $378$ decorated graphs at genus 2, 3) and the frame-B terms on $\TF$ at $\xc=3/5$ and $7/3$ ($27$ and $693$). Hodge integrals computed by Mumford's recursion reproduce known values ($\int\lambda_1\lambda_2\lambda_3=1/1451520$ and $\int\lambda_2^3=1/725760$ on $\Mbar_3$, among others). Frames that must fail: the class $(E,-\alpha_1,-\wA)$ at colour $1$ gives singular edge graphs in every genus ($6$ of $6$ at genus 2, $85$ of $350$ coloured graphs or $41$ of $41$ stable graphs at genus 3, $2522$ of $8596$ at genus 4) with the total unchanged; the point frame gives $6$ of $6$ singular at genus 2 and $41$ of $41$ at genus 3, and in it every decorated graph is singular with pole order $2g-2$, the poles cancelling only in the total; the classes $\wts_\pm$ themselves at the two colours, the swapped assignment, and $\nw$ at both colours give singular edge graphs ($6$ of $6$, $41$ of $41$); the classes $(-E,\epsilon_1,\epsilon_2)$ and $(E+1,-\epsilon_1,-\epsilon_2)$, used consistently, give the correct total but singular edge graphs. The colour-$2$ class is not constrained by these tests: with the class $\bw$ at colour $1$ and the point frame, $(E,-\alpha_1,\alpha_1-E)$, $\wts_+$ or $\wts_-$ at colour $2$, all edge graphs are regular ($0$ of $6$, $0$ of $41$). Controls on the constants: omitting the constant $\hnw$, or replacing $\epsilon_1$ by $\epsilon_1+1$ in it, gives a wrong total and singular edge graphs; twisting the rows by $\hw$ while keeping the class $(E+1,-\epsilon_1,-\epsilon_2)$ gives a wrong total.

\subsection{Direct localisation}
\label{sec:evi-loc}

Independently of the R-matrix, the invariants $\GW_{g,d}$ were computed by virtual localisation on $\Mbar_g(\bbP^1,d)$ (an independent enumeration of the fixed loci with its own vertex and edge factors; Hodge integrals by Mumford's recursion), in $\bbF_p$, at points of $\TCY$ not used elsewhere, and compared with $\Pg d^{2g-3}/(\epsilon_1\epsilon_2)$: $g\leq 6$ at $d=1$ at $(7/5,-3/8,13/6,11/9)$; $g\leq 5$ at $d\leq 3$ at $(-5/4,19/7,2/3,17/5)$; $g\leq 4$ at $d\leq 4$ at $(3/2,-1/3,5/7,8/3)$; $g\leq 7$ at $d=1$ at $(3/4,-2/5,9/7,13/3)$; $g\leq 4$ at $d\leq 2$ at the point $(1,2,-3,5,-5)$ with $E=0$ and at the point $(2/3,5/4,5/4,1/7)$ with $\alpha_1=\alpha_2$. All agree, including genus $0$ and genus $1$ (expected $1/(d^3\epsilon_1\epsilon_2)$ and $\Ph{1}/(d\epsilon_1\epsilon_2)$, \cref{prop:rig-g0}, \cref{lem:rig-g1}). The control off $\TCY$ at $(7/5,-3/8,13/6,11/9,2)$ fails for every $g\geq 1$. A separate direct-localisation computation at $d=1$, $g\leq 2$, confirmed $\epsilon_1\epsilon_2\GW^{\Tmax}_{1,1}$ of \cref{lem:rig-g1} on the five-torus and $\epsilon_1\epsilon_2\GW_{2,1}=\Ph{2}$ on $\TCY$; both agree with the genus $\leq 2$ part of the five-torus expansion displayed in \cite{BS} off $\TCY$. \draftnote{BS label: the displayed five-torus invariants of degree one in the resolved-conifold section, eq: res con GW 2.}

\subsection{Rational reconstruction in \texorpdfstring{$\cQ$}{Q}}
\label{sec:evi-rational}

From $60$ to $70$ evaluations at points $\cQ\in\bbF_p$, the frame-B graph terms were reconstructed as rational functions of $\cQ$ at a generic point of $\TCY$. Genus 2: the total is $N/(1-\cQ)^2$ with $\deg N=2$; each of the $6$ edge graphs is $N/\Disc^4$ with $\deg N=8$; the vertex graph is $N/(\Disc^4(1-\cQ)^2)$ with $\deg N=10$. Genus 3: the total is $N/(1-\cQ)^4$ with $\deg N=4$; each of the $41$ edge graphs is $N/\Disc^8$ with $\deg N=16$; the vertex graph is $N/(\Disc^8(1-\cQ)^4)$ with $\deg N=20$. So the individual graphs have poles on $\Disc=0$, of orders $4$ and $8$, which cancel in the total (\cref{lem:rig-disc}); the edge graphs are regular at $\cQ=1$ and all graphs are regular at $\cQ=\infty$ (\cref{lem:rig-inf}); and the total has the pole structure of $\Pg\Li_{3-2g}/(\epsilon_1\epsilon_2)+\sF_g(0)$ (\cref{thm:rig-main}). In the point frame at genus 2 every graph has a pole of order $2$ at $\cQ=1$.

\subsection{The rigidity argument}
\label{sec:evi-rig}

Exact rational arithmetic at random rational weights:
\begin{enumerate}
\item[(A)] $\Ssh(y)=y/(2\sinh(y/2))=1-y^2/24+7y^4/5760-31y^6/967680+127y^8/154828800-\dots$; $\Ph{2}$ and $\Ph{3}$ as polynomials; $(2g-3)!\Pg/(\epsilon_1\epsilon_2)=\beta_2,\beta_4$ as displayed in \cite{BS} (\cref{lem:reg-beta}); $\gapg$ at $(\epsilon_1,\epsilon_2)=(\xc,1)$, $g=2$, equals $(7+10\xc^2+7\xc^4)/(5760\xc)$. \draftnote{BS label: eq:betapol.}
\item[(B)] At 50 points: $\Disc=\cQ^2[(E-\alpha_0\tau)^2-4\alpha_1\alpha_2\tau]$; $\mathrm{disc}_{\cQ}\Disc=16\alpha_1\alpha_2\wA\wB$; $\Disc(1)=E^2$; $(\alpha_1\alpha_2-E\Lv_1)(\alpha_1\alpha_2-E\Lv_2)=\alpha_1\alpha_2\wA\wB$ (\cref{lem:set-disc}, \cref{prop:set-idem}).
\item[(C)] In floating point at $\cQ=10^6$ and $10^9$: $U_1\to(\alpha_1-\alpha_2)/\wB$, $\cQ U_2\to-\alpha_2/(\alpha_1-\alpha_2)$, $\Lv_i\to\alpha_1,\alpha_2$, $\Delta_1\to-\alpha_1\wB(\alpha_1-\alpha_2)\epsilon_1\epsilon_2$, $\Delta_2\to\alpha_2\wA(\alpha_1-\alpha_2)\epsilon_1\epsilon_2$ (\cref{lem:rig-inf}, \eqref{eq:rig-Dinf}).
\item[(D)] The simple-root lifting of \cref{lem:rig-root} iterated in $\Kred\fps{\alpha_0}$ to $O(\alpha_0^8)$ at five random $(\alpha_1,\alpha_2)$: $\Disc(\cQ_1)=0$ to that order, with the displayed coefficients of $\alpha_0^2,\alpha_0^3,\alpha_0^4$.
\item[(E)] $\Li_{-n}=\Aeul_n/(1-\cQ)^{n+1}$ with $\Aeul_n(0)=0$ and $\deg\Aeul_n=n$ ($\Aeul_1=\cQ$, $\Aeul_2=\cQ+\cQ^2$, $\Aeul_3=\cQ+4\cQ^2+\cQ^3$, $\Aeul_4=\cQ+11\cQ^2+11\cQ^3+\cQ^4$); $\Li_{-n}(1/\cQ)=(-1)^{n+1}\Li_{-n}(\cQ)$ for $n=1,\dots,7$ at ten random $\cQ$ each; agreement with the series at $\cQ=1/3$ to $10^{-30}$; $\Li_0(1/\cQ)=-1-\Li_0(\cQ)$; the expansion of $\Li_{-n}(\re^t)$ at $t=0$ has polar part exactly $(-1)^{n+1}n!\,t^{-n-1}$ for $n=1,\dots,6$ and no other negative power, with $\Li_{-1}(\re^t)=t^{-2}-1/12+O(t^2)$ and $\Li_{-3}(\re^t)=6t^{-4}+1/120+O(t^2)$ (\cref{lem:rig-li}).
\item[(F)] The multiset identities as power sums $\sum_av_a^{1-2k}$, $k=1,2,3$, at 20 points: $\nw+(E,-\alpha_1,-\wA)=(E,-\epsilon_1,-\epsilon_2)$ (\cref{lem:set-hmult}); at $\alpha_0=0$, $\nw-\bw=(\alpha_1,\alpha_2,-E)$ (\cref{lem:rig-3Tid}); at $\cQ=\infty$, $(\alpha_2-\alpha_1,-\wA,\wB)+\nw=\wts'_-$ and $(\alpha_2,-\alpha_1,\alpha_1-\alpha_2)+\nw=\wts'_+$ (\cref{prop:rig-flop}).
\end{enumerate}
The $\alpha_0$-adic step of \cref{lem:rig-disc} and the regularity at $\cQ=\infty$ of \cref{lem:rig-inf} were tested only through the totals and the rational reconstruction of \cref{sec:evi-rational}.

\subsection{The flop and the constant term}
\label{sec:evi-flop}

The identity $\sF_g(0)=\sF^{\Zfl}_g(0)$ of \cref{prop:rig-flop}(iii), with both sides computed as the constant-map Hodge integrals of \cref{rmk:rig-d0} at the weights $\wts_\pm$ and $\wts'_\pm$, holds for $g=2,\dots,6$ at two points of $\TCY$ and fails for $g=2,\dots,6$ at the point $(1,2,4,3,5)$ off $\TCY$. The critical points at $\sQ=0$, the limits of $\cW_2(U_i)$ and the weights $\wts'_\pm$ were also checked exactly by a second computation. The constant term $\sF_g(0)$ itself depends on $\alpha$: at fixed $(\epsilon_1,\epsilon_2)$ and three values of $\alpha$ it does not take the same value at the three for $g=2,3,4$ (\cref{rmk:rig-F0}).

\subsection{Off the Calabi--Yau torus}
\label{sec:evi-offCY}

With the class $(E_\alpha,-\epsilon_1,-\epsilon_2)$, $E_\alpha\coloneqq\alpha_0+\alpha_1+\alpha_2\neq\epsilon_1+\epsilon_2$, at the points $(2,3,4,5,11)$, $(3,3,4,5,11)$, $(1,7,2,5,3)$ and $(1,2,4,3,5)$: the singular row still freezes, $\rr_{0,1}(z;0)=E_\alpha/(E_\alpha+z)$, and $\Delta_1t^2\to-E_\alpha^3\epsilon_1\epsilon_2$ (\cref{rmk:lea-cy}); at genus 2 the edge graphs are regular but the vertex graph acquires a $t^{-1}$ term; at genus 3 four stable graphs are singular (those with a vertex of genus $\geq 2$; $4$ of $350$ coloured graphs), at genus 4 fifty-four of $8596$ coloured graphs are singular, the vertex graph has all polar orders, and the total is not $\Pg\Li_{3-2g}/(\epsilon_1\epsilon_2)+\sF_g(0)$. Direct localisation off $\TCY$ fails for every $g\geq 1$ (\cref{sec:evi-loc}), and the flop identity fails (\cref{sec:evi-flop}). The leaf structure survives off $\TCY$ in these tests; what is lost is the vanishing theorem, which needs the sum-zero class $\bw$ (\cref{thm:reg-L}, \cref{rmk:lea-cy}).

\subsection{What the checks do not cover}
\label{sec:evi-limits}

All checks are at finite genus ($g\leq 5$ for the graph sums, with a single point at genus 5; $g\leq 7$ for direct localisation at $d=1$; $g\leq 6$ for the flop identity), at finite orders in $t$, $z$, $\psi$, and at finitely many weight points; the zero tests in $\bbF_p$ are probabilistic. The exact agreement in $\cQ$ on $\TCY$ exists only through the rational reconstruction at genus $2,3$. Not tested at all: the stability argument (Fl6) and the descent of \cref{prop:ide-teleman}, the analytic input \cref{lem:row-betaphase} beyond the values it predicts, and the identification of \cref{prop:rig-flop}(iii) beyond the Hodge-integral identity itself. The two implementations of the graph sum agree with each other on the points compared, and with the direct localisation, which does not use the R-matrix at all.


\bibliographystyle{amsalpha}
\bibliography{refs}

\end{document}
